\documentclass[11pt]{article}
\usepackage[margin=1in]{geometry}
\usepackage{booktabs}
\usepackage{graphicx}
\usepackage{hyperref}
\usepackage{amsmath}
\usepackage{microtype}
\usepackage{authblk}
\usepackage{caption}
\captionsetup{font=small,labelfont=bf}
\usepackage{array}
\usepackage{pdflscape}

\title{Vocabulary flux in U.S.\ trademark filings:\\
       a Kullback--Leibler surprise analysis of the 16.5-million-filing backfile}
\author{Eric Silver}
\affil{}
\date{April 2026}

\begin{document}
\maketitle

\begin{abstract}
We propose a Kullback--Leibler surprise difference $\Delta KL =
KL_{\text{pros}} - KL_{\text{retr}}$ as an innovation measure for
trademark filings: $\Delta KL > 0$ when a filing's goods/services
vocabulary is prospectively novel (unusual against the recent past)
and retrospectively obvious (similar to the near future) --- the
directional resonance signal of Murdock, Allen \& DeDeo
(\emph{Cognition}, 2017) applied to commercially incentivised text.
We compute $\Delta KL$ on each of the $16.5$ million U.S.\ trademark
filings since 1985 against per-day exponentially decay-weighted
reference distributions of the same NICE class (half-life $H = 2$
years; chosen by signal-maximising sweep). The hypothesis going in
was that firms with positive $\Delta KL$ would outperform on
standard measures of firm success, and that $\Delta KL$ would be
additive to the standard patent-count innovation proxy. Both
predictions are confirmed: $\Delta KL$ predicts 4-year stock
excess returns and gross margin conditional on positive $\Delta KL$,
is essentially uncorrelated with patent counts ($r = +0.01$ on
$174{,}569$ firm-years), and is sub-additive with patents in the
$R^2$ decomposition --- it indexes a distinct innovation channel.

\textbf{Two response curves, two theories.} The innovation side of
$\Delta KL$ (positive) is itself a clean saturating-exponential
response: $P(\text{survived}\,5y) = a + b\,e^{-c\,\Delta KL}$ with
$a = 0.995$, $b = -0.55$, $c = 0.063$ on the $5$M-debut panel ($R^2 =
0.499$ across $20$ bins), concave-down, no inverted-$U$ inflection.
This is the empirical content of optimal distinctiveness
(Deephouse 1999, \emph{SMJ}; Zhao et al.\ 2017, \emph{SMJ}) translated
from \emph{level} of category-typicality to \emph{directional}
vocabulary flux: novelty pays with diminishing returns. The
\textbf{new finding} is that filings at the \emph{opposite} flux
extreme --- prospectively obvious and retrospectively novel,
$\Delta KL < 0$ --- have their own saturating-exponential response,
fitted much more tightly ($a = 0.573$, $b = -0.137$, $c = 3.33$,
$R^2 = 0.978$). On the harvest side, survival \emph{increases} with
$|\Delta KL|$ at a much faster rate and saturates closer to zero:
the deeper a debut filing sits into the receding legacy vocabulary
of its category, the more durable its mark. We read this as
\emph{increasing returns to late-stage imitation}, complementary to
the diminishing-returns innovation curve. Together the two
saturating exponentials produce the survival $U$ on $\Delta KL$ in
the pooled corpus. The pattern holds at the per-industry level in
$32$ of $44$ NICE classes; only $1$ industry shows a significant
inverted-$U$; $11$ are indeterminate. Token-level survival rates
corroborate the U at the corpus level: negative-stratum tokens
(\texttt{microfiche}, \texttt{dial-up}, \texttt{pager},
\texttt{lager}, \texttt{annuity}) and positive-stratum tokens
(\texttt{blockchain}, \texttt{cannabidiol}, \texttt{non-fungible})
both survive well above the corpus baseline; flux-neutral tokens
(\texttt{organisational}, \texttt{conformity}, \texttt{freighting})
survive at $11$--$26\%$ vs.\ a corpus baseline of $\sim 49\%$. In
our hand-curated panel of $41$ canonical ``winner'' firms,
$37/41$ have positive firm-mean $\Delta KL$. What the mechanism on
the harvest side is --- intelligent harvesting of vanishing
concepts, recognition of emerging ones, or selection on an
unobservable correlate --- is not separable in the trademark record;
what is unambiguous is that filings whose vocabulary places them
squarely \emph{on trend} in the centre of the distribution survive at
the lowest rate in the corpus.

\textbf{Novelty of language and novelty of claim are different signals.}
A pioneer-ladder analysis on five category-defining buzzwords
(\texttt{website}, \texttt{wireless}, \texttt{online retail},
\texttt{social media}, \texttt{blockchain}) ranks every firm in the
corpus by date-of-first-adoption; the very-first $10$--$50$ adopters
do badly, the fast-follower band at ranks $50$--$1{,}000$ has
$1.5$--$4\times$ the late-entrant IPO rate, and the mass late wave
does worst. Canonical winners are mostly absent from the very-early
ranks (Amazon $174$\textsuperscript{th} on \texttt{online retail},
Coinbase $96$\textsuperscript{th} on \texttt{blockchain}; Uber, Lyft,
Airbnb never used their category buzzword at all) but most have
positive first-filing $\Delta KL$, often strongly so (Oracle $+4.39$
nats, Intel $+3.26$, Microsoft $+2.99$, eBay $+1.73$, Apple $+1.68$,
Coinbase $+0.98$, Lyft $+0.76$, Twitter $+0.75$, Uber $+0.70$,
LinkedIn $+0.58$, Salesforce $+0.43$). The two ``novelties'' decouple:
the novelty of the \emph{claim} (which vocabulary the firm adopts) and
the novelty of the \emph{language} (the texture of their
goods/services prose) are independent levers.

$\Delta KL$ is essentially uncorrelated with the dominant alternative
innovation proxy (firm-year patent counts, $r = +0.010$ on $174{,}569$
firm-years), positively correlated with post-registration survival,
and positively correlated with profitability among publicly-traded
firms ($+3.2$ pp gross margin per $\sigma$, SE $0.6$ pp).
\end{abstract}

\section{Background}

Innovation is a central object of empirical economic and management
research, but the construct is hard to measure. The dominant quantitative
approaches are R\&D expenditure (an input proxy that rewards spend rather
than novelty), patent counts and patent citations (which require
strategic legal filing, and which Trajtenberg, Hagedoorn, and others have
shown to be at best a noisy proxy), and the market value premium of equity
over book value (which conflates innovation with every other investor
perception). Survey-based scales such as the Community Innovation Survey
and Gatignon et al.'s instrument capture richer constructs but suffer from
common-method bias and limited sample size. The lack of agreement on what
to measure is itself a finding: across measures, pairwise correlations are
typically $\rho \approx 0.3$.

We propose using trademark filings as a complementary novelty channel.
Trademarks have several useful properties that patents do not: applicants
file a single, brief, machine-readable goods-and-services description that
is sharply incentivized to be exact (broad enough to protect future use,
narrow enough to clear examination); the data are released into the
public domain by the USPTO without lag; and the universe of filers is
broader than the universe of patentees, including most consumer-facing
firms. The cost is that trademarks measure novelty in the
goods-and-services lexicon, not in the underlying technology --- a brand
that solves an old problem with a new identity (e.g.\ \textsc{Liquid Death}'s
canned water) is largely invisible.

Our broader motivation is to ask whether innovative firms are rewarded or
punished. Two pieces of folk wisdom are in tension. The first, conveyed
in the phrase ``you can tell the pioneers by the arrows in their backs,''
is that firms that depart from convention bear higher search and adoption
costs, get out-competed by fast-followers, and underperform on average.
The second, expressed in the older saying ``build a better mousetrap and
the world will beat a path to your door,'' is that meaningful novelty
attracts customers and rewards its inventor with durable advantage. Both
can be true simultaneously: novelty might raise variance --- some
pioneers fail spectacularly, others enjoy durable category leadership.
This paper offers an empirical test of that joint pattern using the full
USPTO trademark backfile.

\section{Data}

The trademark corpus is the full backfile published by the USPTO Open
Data Portal, current through end-of-year 2025. After filtering to
filings with non-empty goods/services text in NICE classes 1--45 and
a parseable filing date, the corpus contains
\textbf{$16{,}526{,}850$ filings} from \textbf{$5{,}673{,}680$ unique
registrant names} across all 45 NICE international classes (35 goods
classes plus 10 services classes), spanning 1985 to early 2026. Each
record carries the registrant's name as it appeared on the
application (this is the owner at the time of filing; subsequent
assignments are not reflected), the filing date, a status code that
records the current legal state of the mark, an optional registration
date, the NICE class assignments, and one or more goods-and-services
descriptive statements. Statements describing the visual mark itself
rather than the business (USPTO type-code prefix D) are excluded;
only goods-and-services statements (prefix GS) are kept. Two classes
are presented in narrative detail throughout the paper: NICE Class~9
(Software \& Electronics) and NICE Class~32 (Beer \& Soft Drinks).
The remaining 43 industries are summarized in
Section~\ref{sec:cross}, and per-class summary statistics for the two
focal classes appear in Appendix~\ref{app:summary}.

The financial-outcome data are the SEC EDGAR Financial Statement
Data Sets, 64 quarterly XBRL releases spanning fiscal years 2009
through 2024, covering $13{,}127$ unique CIKs that filed at least
one 10-K. We extract revenue, cost of revenue, gross profit, R\&D
expense, operating income, and net income.

\paragraph{Inclusion filter (\textsc{clean}).}
Almost every analysis in this paper is restricted to the same
inclusion filter, which we call \textsc{clean}. The filter has five
clauses, each addressing a separate failure mode: two outcome-window
clauses (\texttt{year} bounds) that ensure the survival outcome is
fully observed and that the prospective reference window is non-empty,
two density floors on the per-class reference distributions
(\texttt{n\_ref\_pros} and \texttt{n\_ref\_retr} $\geq 1{,}000$) that
ensure $KL$ scores reflect a real vocabulary distribution rather than
small-sample noise, and a per-filing token-count floor
($n_{\text{terms}} \geq 3$) that prevents 1- or 2-token descriptions
from producing degenerate $KL$ values. Of $16{,}530{,}448$ raw
filings with parsed surprise data, $\textbf{10{,}060{,}073}$
($60.9\%$) survive the full filter and form the analytic corpus.

\begin{table}[h]
\centering
\small
\begin{tabular}{lrr}
\toprule
Clause & dropped alone & added back if relaxed \\
       & (vs.\ raw, \%)& (vs.\ \textsc{clean}, \%) \\
\midrule
\texttt{year} in $[1990, 2021]$      & $4{,}490{,}092$ ($27.2\%$)  & $+2{,}985{,}245$ ($+29.7\%$) \\
\texttt{n\_ref\_pros} $\geq 1{,}000$ & $2{,}503{,}062$ ($15.1\%$)  & $+3{,}371$ ($+0.0\%$) \\
\texttt{n\_ref\_retr} $\geq 1{,}000$ & $2{,}543{,}635$ ($15.4\%$)  & $+0$ ($+0.0\%$) \\
$n_{\text{terms}} \geq 3$           & $3{,}440{,}391$ ($20.8\%$)  & $+745{,}431$ ($+7.4\%$) \\
\midrule
\textsc{clean} (all of the above)   & $6{,}470{,}375$ ($39.1\%$ excluded) & --- \\
\bottomrule
\end{tabular}
\caption{Per-clause filter effects on the $16{,}530{,}448$-filing raw
corpus. ``Dropped alone'' is the share that fails the clause when
applied in isolation to the raw corpus; ``added back'' is what would
re-enter the analytic corpus if we relaxed only that one clause while
keeping the others. The year window is the dominant binder; the two
$N_{\text{eff}}$ density floors are nearly redundant once the year
window is in place because they bind on the same burn-in years
(1985--1989 on the past side, 2022--2025 on the future side) that
the year clause already excludes.}
\label{tab:clean-filter}
\end{table}

The dominant binder is the year window, $[1990, 2021]$. The lower
bound is the literal arithmetic burn-in described in the Method
section (the corpus begins in 1984 and we leave a margin for the
decay reference). The upper bound is the cohort cap: with a 2026
freeze date and a 5-year post-registration survival window, filings
must be from 2021 or earlier for \texttt{survived\_5y} to be
observable. The token-count floor, $n_{\text{terms}} \geq 3$, is the
only other clause that uniquely excludes a meaningful share ($+745$k
filings rescued if dropped); when we look at the templated
short-description analysis in Section~\ref{sec:templated-survival} we
report results both with and without this clause, since the most
heavily templated G/S strings (\textit{e.g.}~``restaurant services'',
``legal services'') are 2-token strings and would otherwise be
excluded.

\paragraph{Registration is determined by status code, not registration
date.} We classify a mark as having ``completed registration'' from its
USPTO status code (6xx live-registered or 8xx post-registration-cancelled),
not from the \texttt{registration\_date} field, which is unreliable in
older records (mostly pre-2003): many records with a valid registered
status code carry a missing \texttt{registration\_date} in the
published XML. The pipeline that extracts and analyses these data is
published with the paper; technical details live in the companion
repository.

\section{Method}

For each filing $i$ at year $t$ with goods-and-services text $w_i$ we
construct a term-frequency distribution $P_i$ over a fixed dictionary $V$
of unigram and bigram tokens. The dictionary is built once per NICE class
with a custom analyzer that splits each description at semicolons before
tokenizing, so bigrams cannot cross goods/services-item boundaries. We
retain only unigrams and bigrams that appear in at least $50$ distinct
filings (documents) and in no more than $50\%$ of all filings, on the
view that very rare tokens are typos and very common tokens are USPTO
boilerplate.

For each filing $i$ at filing-date $t_i$ we form two reference
distributions by per-day exponential decay-weighting of every other
filing in the same NICE class:
\begin{align}
Q_i^{\text{pros}}(v) &\propto \alpha + \sum_{j:\, t_j < t_i}
  e^{-\lambda_{\text{pros}}\,(t_i - t_j)}\,\mathrm{count}(v, j), \\
Q_i^{\text{retr}}(v) &\propto \alpha + \sum_{j:\, t_j > t_i}
  e^{-\lambda_{\text{retr}}\,(t_j - t_i)}\,\mathrm{count}(v, j),
\end{align}
with $\lambda = \ln(2)/H$ and \textbf{half-life $H = 2$ years} on both
sides and Laplace smoothing $\alpha = 1.0$. The half-life is chosen by
a sweep over $H \in \{0.5, 1, 2, 3, 5\}$ years on Class~9 that
maximises corr$(\Delta KL, \text{survived\_5y})$ and the logit
pseudo-$R^2$ on the same outcome (Figure~\ref{fig:halflife-signal}).
Density on the reference distributions is reported as Kish effective
sample size $N_{\text{eff}} = (\sum_j w_j)^2 / \sum_j w_j^2$; we
restrict all filing-level analyses to filings with
$N_{\text{eff}} \geq 1{,}000$ on both sides.

The prospective and retrospective surprise of filing $i$ are the
support-restricted KL divergences
\[
KL_i^{\text{pros}} = \sum_{v: P_i(v)>0} P_i(v) \, \log\frac{P_i(v)}{Q_i^{\text{pros}}(v)},
\quad
KL_i^{\text{retr}} = \sum_{v: P_i(v)>0} P_i(v) \, \log\frac{P_i(v)}{Q_i^{\text{retr}}(v)}.
\]
With $\log$ in natural units, $KL_i$ is in nats. A useful intuition:
$KL_i^{\text{pros}}$ is the expected log-odds per token that a reader
who had seen only the filing's recent past would assign to the filing's
own term distribution versus the past's. $KL_i = 3$ nats means a typical
token in this filing is $e^3 \approx 20\times$ more probable under the
filing's own distribution than under the recent corpus. KL is sensitive
to filing length (short filings with one rare token can score high
mechanically); every regression in the paper controls for
$\log(n_\text{terms})$ and most report $z$-scored $\Delta KL$.

\begin{figure}[h]
\centering
\includegraphics[width=0.9\linewidth]{results/halflife_signal_class009.png}
\caption{Half-life sweep on Class~9. For each candidate $H \in
\{0.5,1,2,3,5\}$ years we recompute the decay-weighted reference
distribution and report four signal metrics on the $\Delta KL \to$
survival relationship. $H = 2$--$3$ years maximises corr$(\Delta KL,
\text{survived\_5y})$ and the logit pseudo-$R^2$; we adopt $H = 2$.}
\label{fig:halflife-signal}
\end{figure}

The directional flux score is $\Delta KL_i = KL_i^{\text{pros}} -
KL_i^{\text{retr}}$: the change between how unlike the recent past a
filing reads and how unlike the near future it reads. Positive
$\Delta KL$ means the filing's vocabulary is closer to the future than
to the past (\emph{flux-positive}; the field migrated toward this
filing); negative $\Delta KL$ means closer to the past than to the
future (\emph{flux-negative}; the field migrated away).

\paragraph{Magnitudes.}
To ground the reader: across the all-class analytic pool of $13.0$M
filings with $N_{\text{eff}} \geq 1{,}000$ on both sides, the 5th,
50th, and 95th percentiles are
$KL^{\text{pros}} = [2.99,\, 4.85,\, 7.20]$ nats (std $1.29$) and
$KL^{\text{retr}} = [2.98,\, 4.84,\, 7.22]$ nats (std $1.30$). The two
components are correlated at $\rho = +0.99$ pooled across all classes
(see Table~\ref{tab:components}); they measure the same
general-distinctiveness signal in two temporal directions, and their
per-filing levels track each other tightly. Because they share that
much distinctiveness against the per-class boilerplate, their
difference $\Delta KL$ has a markedly tighter range: 5th/50th/95th
percentiles $[-0.30,\, -0.00,\, +0.34]$ nats (std $0.22$, median
essentially zero by construction of the look-behind / look-ahead
window). A filing in the top 5\% of $\Delta KL$ is therefore $\sim
0.6$ nats more flux-positive than one in the bottom 5\%, on a scale
where each component routinely exceeds $5$ nats outright. The
substantive direction lives in this comparatively small flux
residual, not in the much larger common distinctiveness factor.

\paragraph{Identical-text filings are signal.} Trademark practice rewards
templated language. Many breweries file the boilerplate token
\texttt{Beer}; many app developers file \texttt{downloadable software}.
We do not deduplicate identical text: filings that use the same words
\emph{should} receive the same score, because deviations from the
boilerplate are themselves the signal.

\paragraph{Burn-in window.}
The decay reference distributions need filings on both sides of $t_i$
for the weighted sums to be well populated. The $N_{\text{eff}} \geq
1{,}000$ density filter above automatically excludes small NICE
classes' early-decade cohorts and the corpus right-edge filings whose
forward window is too sparse. For regressions on the panel we restrict
to filings from \textbf{1990 onward} (after the corpus's left-edge
digitisation boundary in 1984 plus a margin to populate the
look-behind) and to filings whose year has at least 5 years of forward
observation, currently \textbf{1990--2021} given a 2026 freeze date.
Figure~\ref{fig:burnin} shades the two burn-in windows.

\begin{figure}[t]
\centering
\includegraphics[width=0.85\linewidth]{results/burnin.png}
\caption{Filings per year with both 5-year reference windows
populated. Forward-hatched grey: literal 5-year past-side burn-in
(1985--1989) --- filings in these years cannot have a complete
look-behind because the digitised corpus begins in 1984.
Backward-hatched grey: literal 5-year future-side burn-in
(2021--2025) --- filings in these years cannot have a complete
look-ahead at the time of writing. The bands are chosen to print
distinguishably in black-and-white. A separate density filter
(\texttt{n\_ref\_prospective} $\geq 1000$) handles the additional
small-class exclusions visible in the long left tail of the plot.}
\label{fig:burnin}
\end{figure}

\subsection*{What $\Delta KL$ measures (and what it does not)}
\label{sec:flux-not-distinctiveness}

The two component KLs are very highly correlated: per-filing
$\rho(KL^{\text{pros}}, KL^{\text{retr}}) = 0.99$ on the pooled clean
panel. A reader plotting any outcome against either component
univariately sees nearly identical curves because the two component
levels are, to first order, the same number --- both measuring
\emph{distinctiveness} (how unusual the filing's vocabulary is against
the rest of the corpus, in either temporal direction).

The PCA decomposition of $(KL^{\text{pros}}, KL^{\text{retr}})$ makes
the geometry explicit:
\begin{align*}
\text{PC1} \;\approx\; \tfrac{1}{\sqrt{2}}(KL^{\text{pros}} + KL^{\text{retr}})
  &\;\;\;\text{= general distinctiveness} \\
\text{PC2} \;\approx\; \tfrac{1}{\sqrt{2}}(KL^{\text{pros}} - KL^{\text{retr}})
  &\;\;\;\propto \Delta KL = \text{vocabulary flux}
\end{align*}

$\Delta KL$ is the orthogonal residual after PC1 is removed; it is the
\emph{only} part of the joint signal with a forward/backward direction.
PC1 says ``this filing is unusual.'' PC2 says ``this filing is unusual
\emph{and the field's vocabulary subsequently moved toward it rather
than away}.'' We claim throughout that the substantive innovation
signal lives in PC2; PC1 is largely a length/rarity covariate. The
closest-adjacent management-theory framework, optimal distinctiveness
(Deephouse 1999, \emph{SMJ}; Zhao et al.\ 2017, \emph{SMJ}), predicts
an inverted-U on the \emph{level} of distinctiveness. The framework we
are extending is DeDeo's prospective/retrospective surprise scheme
(Murdock, Allen \& DeDeo 2017, \emph{Cognition}; Barron, Huang, Spang
\& DeDeo 2018, \emph{PNAS}), which decomposes a temporal corpus into
novelty, transience, and resonance and which we apply for the first
time to commercially incentivised text.

\section{What the components do separately}
\label{sec:components}

Before introducing $\Delta KL$ as a single scalar, it is worth asking
whether the prospective and retrospective KL components carry
independent information about a filing. A natural prior --- and one we
held coming in --- is that they should pull in opposite directions:
prospective surprise should make a filing harder to clear examination
(novel goods/services language draws office actions), while
retrospective surprise should reward the filer for ``being right'' (the
future ends up resembling them, suggesting they understood a market
others did not). Under that view, the value of $\Delta KL$ is that it
nets two opposing forces.

\textbf{The data do not support that decomposition.} Table~\ref{tab:components}
shows that the per-filing correlation between $KL_{\text{pros}}$ and
$KL_{\text{retr}}$ is $+0.97$ in software (Class 9), $+0.93$ in beverages
(Class 32), and $+0.97$ in the 11.4-million-filing all-class pool. The
two scores are computed against the same per-filing distribution
$P_i(v)$, with the only difference being whether the reference window
is the five preceding or the five following years; the bulk of variance
in either KL therefore comes from $P_i$ itself, and only a thin sliver
comes from the past-vs-future shift in the reference. As a consequence,
plotted against any per-filing outcome the two components produce
nearly identical curves.

\begin{table}[h]
\centering
\caption{Per-filing correlation of the two KL components and
single-component slopes against \emph{reaching registration} and
\emph{passing five-year survival} (per-filing linear fits on binned
means, 15 bins between the 2nd and 98th percentile of the predictor).
The two components carry essentially the same signal at the per-filing
level.}
\label{tab:components}
\begin{tabular}{lcccc}
\toprule
Scope & \shortstack{$\rho(KL_{\text{pros}},$\\ $KL_{\text{retr}})$} & \shortstack{slope on reg.\\ vs $KL_{\text{pros}}$} & \shortstack{slope on reg.\\ vs $KL_{\text{retr}}$} & \shortstack{slope on pass-5y\\ vs $KL_{\text{pros}}$} \\
\midrule
Software (Cl.\ 9)   & $+0.967$ & $+0.024$ & $+0.021$ & $+0.008$ \\
Beverages (Cl.\ 32) & $+0.933$ & $+0.018$ & $+0.017$ & $+0.009$ \\
All 45 (pooled)     & $+0.967$ & $+0.026$ & $+0.028$ & $+0.013$ \\
\bottomrule
\end{tabular}
\end{table}

\begin{figure}[h]
\centering
\includegraphics[width=\linewidth]{results/per_filing_pros_vs_retr_software.png}
\caption{Per-filing outcomes vs prospective and retrospective KL ---
software (Class 9). Top row: probability the filing reaches registration.
Bottom row: probability the filing reaches registration AND remains
live today. Left column conditions on prospective KL, right column on
retrospective KL. Each blue point is the bin mean (15 bins between the
2nd and 98th percentile of the per-filing predictor); shaded band is the
95\% Wilson interval around each bin's proportion; red dashed line is
the linear fit. The two columns are essentially the same curve --- per-filing
$\rho(KL_{\text{pros}}, KL_{\text{retr}}) = +0.97$.}
\label{fig:components-perfiling-sw}
\end{figure}

The curves themselves are not lines and not exponentials. They are
\textbf{monotone-saturating}: registration probability rises from about
$0.43$ at $KL = 3$ nats to about $0.53$ at $KL = 7.5$ nats in software,
and similarly $0.46 \to 0.59$ in beverages, with most of the rise
concentrated in the moderate range and the curve flattening at the
extremes (Figure~\ref{fig:components-perfiling-sw}; beverage and pooled
versions live in
\texttt{paper/results/per\_filing\_pros\_vs\_retr\_\{beverages,all\}.png}).
At the firm-mean level the same outcomes against the same components
trace inverted-U arcs in software (peak survival and peak public-listing
both at firm-mean $KL \approx 5$ nats; tails decline in both directions);
beverages produces a hump on survival and a monotone decline on
public-listing (high-novelty beverage firms are systematically less
likely to be public --- they are craft brewers, not Coca-Cola). The
firm-level views are in
\texttt{paper/results/pros\_vs\_retr\_curves\_\{software,beverages,all\}.png}.

\paragraph{The sharp implication.} Because the components are nearly
collinear, the signal in their \emph{difference}
$\Delta KL_i = KL_i^{\text{pros}} - KL_i^{\text{retr}}$ is a small
share of either component's variance --- but it is the only part of the
joint $(KL_{\text{pros}}, KL_{\text{retr}})$ score that does not cancel
out under the shared-$P_i$ structure. A filing whose
language was more unusual against the past than against the near future
is one whose language the near future ended up resembling. The shared
``how unusual was the language at all'' term cancels; what remains is
``did the unusualness fade because others adopted it.'' The rest of the
paper takes that residual --- the imitated-innovator score --- as the
primary axis. Software and beverages stay as our running examples
through Outcomes 1--3; the back matter establishes that the basic story
holds across the other 43 NICE classes.

\section{Face validity}
\label{sec:fv}

Figure~\ref{fig:quadrant} plots each filing in the prospective--retrospective
KL plane. The horizontal axis (prospective KL) measures dissimilarity from
filings in the previous five years (high $\Rightarrow$ less like the past);
the vertical axis (retrospective KL) measures dissimilarity from filings in
the next five years (high $\Rightarrow$ less like the future). Filings
\emph{below} the diagonal $KL_{\text{pros}} = KL_{\text{retr}}$ are
innovators by our criterion --- the past does a worse job of predicting
them than the future does. Filings above the diagonal are tired. The
labelled marks include expectation-confirming category-creators
(\textsc{OpenAI} 2016, \textsc{Claude} 2023, \textsc{Instagram} 2011,
\textsc{White Claw Seltzer Works} 2016, \textsc{Rockstar} 2002) and
expectation-disconfirming late-cycle marks (\textsc{Hard Mtn Dew} 2021).

\begin{figure}[t]
\centering
\includegraphics[width=\linewidth]{results/quadrant.png}
\caption{Per-filing prospective vs.\ retrospective KL surprise, zoomed
to the central [2, 10] nat range where the labelled examples sit. Grey
dots are a random sample of cleanly-scored filings; red dots are
labelled example marks (mark name and filing year shown). \textbf{The
KL scores compare each filing only to other filings in the same NICE
category and same five-year window.} An AI-related filing is therefore
not directly comparable to a beverage filing: \textsc{White Claw Seltzer
Works} is novel \emph{within Beer \& Soft Drinks 2011--2021}, not novel
in the absolute sense that an artificial-intelligence trademark is
novel against the history of human technology. Below-diagonal:
innovator (the future of the same category resembles this filing more
than the past did). Above-diagonal: tired. Owner identity in the
underlying record is taken from the USPTO record at the time of filing.}
\label{fig:quadrant}
\end{figure}

The bottled-water category illustrates a useful failure mode. The 2019
\textsc{Liquid Death} filing for canned water describes its goods as
\emph{still water; sparkling water} --- the same lexicon as
established bottled-water brands --- and so lands near the diagonal.
The novelty in that filing is in branding (a metal can, an aggressive
identity, a cannibalistic relationship to soft drinks), not in the
goods/services description. Comparable bottled-water filings
(\textsc{Saratoga}, \textsc{Smartwater}, \textsc{Essentia}) cluster in
the same region. The method is, by construction, blind to non-textual
novelty.

\section{Canonical winners and failed firms as KL profiles}
\label{sec:canonical-winners}

A direct face-validity check: take a curated list of widely-recognised
firms from the last fifty years, find their first post-founding
trademark filing, compute the firm-mean $\Delta KL$ across all their
filings, and compare against a control list of widely-recognised
\emph{failed} firms. If the construct is meaningful, the winners should
read as forward-flux (positive $\Delta KL$) and the failures should
spread out --- some failed despite being trend-setters, some failed
\emph{because} they were exploiting a receding vocabulary.

\paragraph{Caveat on entity matching.}
The trademark XML records each filing's owner only as a free-text
\texttt{owner\_name} string with no entity identifier, so any
``firm $\to$ all its filings'' join is necessarily heuristic. For the
table below we use hand-curated regular expressions that require the
brand string to co-occur with a corporate suffix (\texttt{Inc.},
\texttt{LLC}, \texttt{Corporation}, \texttt{Technologies}, etc.) and we
restrict each firm's panel to filings on or after its founding year.
The regexes are not perfect: we are confident that at least some hits
are entity collisions and we have tightened individual regexes where we
identified collisions. Amazon's regex, for example, matches only the
\texttt{Amazon.com} variants because the string ``Amazon Technologies,
Inc.'' was used as a registered owner by a different firm in
1994--1995, before Amazon's IP-holding subsidiary of that name became
active around 2001; we lose some post-2001 IP-holder filings but
eliminate the collision. Readers should treat the firm-mean numbers
below as indicative rather than authoritative; the qualitative pattern
(winners cluster on positive firm-mean $\Delta KL$; failed firms split
across both extremes) is robust to the residual collision rate but
specific decimal places are not.

\begin{table}[h]
\centering
\caption{Firm-mean and first-real-filing KL profile for 41 canonical
``winner'' firms and the 17 widely-recognised failed firms with
sufficient post-founding clean-panel data. Filings filtered to
filing-date $\geq$ founding-year; see entity-matching caveat above.
Corpus median pros $=4.85$, retr $=4.84$.}
\label{tab:big-firm-kl}
\tiny
\begin{tabular}{lllrrrrrrr}
\toprule
\textbf{cohort} & \textbf{brand} & \textbf{founded} & \textbf{n} & \textbf{mn\_pros} & \textbf{mn\_retr} & \textbf{mn\_$\Delta$KL} & \textbf{1st\_pros} & \textbf{1st\_retr} & \textbf{1st\_$\Delta$KL} \\
\midrule
\textsc{win} & Amazon & 1994 & 75 & 4.75 & 4.67 & +0.075 & 5.08 & 5.88 & -0.79 \\
\textsc{win} & eBay & 1995 & 286 & 4.64 & 4.60 & +0.037 & 7.45 & 5.72 & +1.73 \\
\textsc{win} & Etsy & 2005 & 58 & 3.72 & 3.73 & -0.013 & 5.42 & 5.48 & -0.05 \\
\textsc{win} & Shopify & 2006 & 284 & 3.40 & 3.32 & +0.080 & 5.98 & 5.88 & +0.10 \\
\textsc{win} & Wayfair & 2002 & 689 & 4.34 & 4.35 & -0.011 & 5.97 & 6.23 & -0.26 \\
\textsc{win} & Chewy & 2011 & 235 & 4.86 & 4.73 & +0.133 & 4.80 & 5.00 & -0.20 \\
\textsc{win} & Google & 1998 & 2,826 & 4.40 & 4.32 & +0.080 & 4.57 & 4.23 & +0.33 \\
\textsc{win} & Yahoo & 1994 & 826 & 4.33 & 4.26 & +0.068 & 5.92 & 5.86 & +0.06 \\
\textsc{win} & Netflix & 1997 & 179 & 3.83 & 3.81 & +0.016 & 4.40 & 4.17 & +0.23 \\
\textsc{win} & Dropbox & 2007 & 81 & 4.30 & 4.28 & +0.020 & 4.11 & 3.82 & +0.30 \\
\textsc{win} & Facebook/Meta & 2004 & 1,063 & 3.68 & 3.57 & +0.116 & 4.28 & 4.12 & +0.16 \\
\textsc{win} & Twitter/X & 2006 & 420 & 3.73 & 3.70 & +0.037 & 4.30 & 3.55 & +0.75 \\
\textsc{win} & LinkedIn & 2002 & 286 & 4.09 & 4.05 & +0.047 & 6.54 & 5.96 & +0.58 \\
\textsc{win} & Pinterest & 2010 & 44 & 3.44 & 3.35 & +0.089 & 5.05 & 4.74 & +0.32 \\
\textsc{win} & Snap & 2011 & 441 & 4.01 & 3.96 & +0.050 & 4.32 & 4.41 & -0.09 \\
\textsc{win} & Coinbase & 2012 & 147 & 4.33 & 4.16 & +0.161 & 6.48 & 5.50 & +0.98 \\
\textsc{win} & Block/Square & 2009 & 376 & 4.49 & 4.46 & +0.038 & 5.73 & 5.48 & +0.26 \\
\textsc{win} & Stripe & 2010 & 114 & 4.89 & 4.83 & +0.063 & 5.45 & 5.55 & -0.11 \\
\textsc{win} & Zoom & 2011 & 18 & 4.88 & 4.82 & +0.063 & 5.10 & 5.07 & +0.03 \\
\textsc{win} & Slack & 2009 & 1 & 5.54 & 5.30 & +0.243 & 5.54 & 5.30 & +0.24 \\
\textsc{win} & Salesforce & 1999 & 884 & 4.09 & 4.01 & +0.084 & 4.22 & 3.79 & +0.43 \\
\textsc{win} & Workday & 2005 & 77 & 4.15 & 4.02 & +0.133 & 5.73 & 5.53 & +0.20 \\
\textsc{win} & ServiceNow & 2004 & 148 & 4.16 & 4.06 & +0.098 & 4.26 & 4.31 & -0.05 \\
\textsc{win} & Snowflake & 2012 & 125 & 4.35 & 4.26 & +0.098 & 4.54 & 4.32 & +0.22 \\
\textsc{win} & Tesla & 2003 & 225 & 4.91 & 4.86 & +0.051 & 6.07 & 6.13 & -0.06 \\
\textsc{win} & Uber & 2009 & 348 & 4.23 & 4.10 & +0.139 & 6.45 & 5.75 & +0.70 \\
\textsc{win} & Airbnb & 2008 & 171 & 4.25 & 4.18 & +0.066 & 4.63 & 4.85 & -0.22 \\
\textsc{win} & Lyft & 2012 & 220 & 4.44 & 4.34 & +0.093 & 6.44 & 5.67 & +0.76 \\
\textsc{win} & Spotify & 2006 & 271 & 4.07 & 4.03 & +0.044 & 4.02 & 3.71 & +0.31 \\
\textsc{win} & Palantir & 2003 & 81 & 3.86 & 3.76 & +0.101 & 4.10 & 4.17 & -0.07 \\
\textsc{win} & Cisco & 1984 & 1,551 & 4.50 & 4.47 & +0.030 & 7.36 & 7.08 & +0.28 \\
\textsc{win} & Oracle & 1977 & 615 & 4.35 & 4.29 & +0.055 & 9.98 & 5.59 & +4.39 \\
\textsc{win} & SAP & 1972 & 342 & 4.03 & 3.98 & +0.051 & 5.30 & 5.24 & +0.07 \\
\textsc{win} & NVIDIA & 1993 & 519 & 3.85 & 3.79 & +0.060 & 6.27 & 6.18 & +0.09 \\
\textsc{win} & AMD & 1969 & 315 & 5.32 & 5.29 & +0.025 & 6.74 & 6.05 & +0.69 \\
\textsc{win} & Adobe & 1982 & 696 & 4.67 & 4.70 & -0.027 & 6.70 & 5.43 & +1.28 \\
\textsc{win} & Intuit & 1983 & 905 & 4.30 & 4.23 & +0.074 & 4.09 & 3.89 & +0.21 \\
\textsc{win} & Apple & 1976 & 2,460 & 4.48 & 4.47 & +0.017 & 7.19 & 5.51 & +1.68 \\
\textsc{win} & Microsoft & 1975 & 4,171 & 4.18 & 4.17 & +0.015 & 8.83 & 5.84 & +2.99 \\
\textsc{win} & Intel & 1968 & 690 & 4.34 & 4.33 & +0.003 & 7.95 & 4.70 & +3.26 \\
\textsc{win} & Qualcomm & 1985 & 1,302 & 4.92 & 4.94 & -0.013 & 5.17 & 4.93 & +0.24 \\
\textsc{fail} & Pets.com & 1998 & 17 & 5.36 & 5.39 & -0.032 & 4.54 & 4.55 & -0.00 \\
\textsc{fail} & Webvan & 1996 & 32 & 4.47 & 4.69 & -0.217 & 2.71 & 2.87 & -0.15 \\
\textsc{fail} & Theranos & 2003 & 305 & 3.15 & 3.10 & +0.045 & 3.86 & 3.78 & +0.08 \\
\textsc{fail} & MoviePass & 2011 & 4 & 7.03 & 7.30 & -0.266 & 7.03 & 7.30 & -0.27 \\
\textsc{fail} & WeWork & 2010 & 313 & 4.24 & 4.20 & +0.040 & 4.99 & 4.99 & +0.00 \\
\textsc{fail} & FTX & 2019 & 1 & 4.34 & 4.26 & +0.075 & 4.34 & 4.26 & +0.08 \\
\textsc{fail} & Solyndra & 2005 & 2 & 6.87 & 6.43 & +0.437 & 6.13 & 5.93 & +0.21 \\
\textsc{fail} & Bird & 2017 & 4 & 5.27 & 5.25 & +0.022 & 7.06 & 7.14 & -0.07 \\
\textsc{fail} & Casper & 2014 & 147 & 4.57 & 4.51 & +0.059 & 5.57 & 5.11 & +0.46 \\
\textsc{fail} & Better.com & 2014 & 4 & 4.36 & 4.48 & -0.119 & 4.60 & 4.81 & -0.21 \\
\textsc{fail} & Lordstown & 2018 & 1 & 5.96 & 5.97 & -0.016 & 5.96 & 5.97 & -0.02 \\
\textsc{fail} & Nikola & 2014 & 49 & 4.39 & 4.36 & +0.028 & 3.81 & 3.57 & +0.24 \\
\textsc{fail} & Beyond Meat & 2009 & 136 & 4.89 & 4.56 & +0.335 & 6.14 & 5.40 & +0.74 \\
\textsc{fail} & Peloton & 2012 & 125 & 4.56 & 4.49 & +0.066 & 3.89 & 3.83 & +0.06 \\
\textsc{fail} & Friendster & 2002 & 13 & 4.39 & 4.35 & +0.036 & 5.74 & 5.55 & +0.19 \\
\textsc{fail} & Myspace & 2003 & 176 & 3.51 & 3.44 & +0.071 & 5.32 & 5.25 & +0.08 \\
\textsc{fail} & Juicero & 2013 & 1 & 4.42 & 4.20 & +0.218 & 4.42 & 4.20 & +0.22 \\
\bottomrule
\end{tabular}
\end{table}

37 of 41 winners have \emph{positive} firm-mean $\Delta KL$ (the four
exceptions are Adobe ($-0.03$), Qualcomm ($-0.01$), Etsy ($-0.01$),
and Wayfair ($-0.01$)). The winners are not extreme outliers on
distinctiveness: their firm-mean prospective KL is mostly $3.4$--$4.8$
nats, around or just below the corpus median. They use mainstream
vocabulary, but the mainstream vocabulary they use is consistently
the kind the field migrates toward, not away from.

\paragraph{First-real-filing $\Delta KL$ is sharper on a subset.}
Thirteen canonical winners had first-filing $\Delta KL$ above $+0.4$
nats (roughly $2\sigma$): \textbf{Oracle} ($+4.39$, 1977),
\textbf{Intel} ($+3.26$, 1968), \textbf{Microsoft} ($+2.99$, 1975),
\textbf{eBay} ($+1.73$, 1997), \textbf{Apple} ($+1.68$, 1976),
\textbf{Adobe} ($+1.28$, 1982), \textbf{Coinbase} ($+0.98$, 2013),
\textbf{Lyft} ($+0.76$, 2012), \textbf{Twitter} ($+0.75$, 2007),
\textbf{Uber} ($+0.70$, 2009), \textbf{AMD} ($+0.69$, 1969),
\textbf{LinkedIn} ($+0.58$, 2003), and \textbf{Salesforce} ($+0.43$,
1999). Oracle's first filing ($+4.39$) is the largest first-filing
$\Delta KL$ in the curated set; the eBay filing's text ``on-line
trading services in which seller posts items to be auctioned and
bidding is done electronically'' ($+1.73$) introduced phrasing that
downstream filings extensively echoed.

\paragraph{The failed-firm panel does not collapse to negative
$\Delta KL$.} Among 17 widely-recognised failures, the split is
roughly even. \textbf{Webvan} ($-0.22$), \textbf{MoviePass} ($-0.27$),
\textbf{Better.com} ($-0.12$), and \textbf{Pets.com} ($-0.03$) failed
on the receding side. But \textbf{Theranos} ($+0.05$),
\textbf{Beyond Meat} ($+0.34$, first $+0.74$), \textbf{Casper}
($+0.06$, first $+0.46$), \textbf{Solyndra} ($+0.44$), \textbf{Nikola}
($+0.03$, first $+0.24$), and \textbf{Myspace} ($+0.07$) failed
despite trend-setting vocabulary. $\Delta KL$ is a vocabulary signal,
not a business-model signal. \emph{Negative $\Delta KL$ raises the
probability of being on the wrong side of an industry transition.
Positive $\Delta KL$ raises the probability of being on the right
side. Neither is destiny.}

\section{Pioneers, fast-followers, and the late wave}
\label{sec:pioneer-ladder}

To take ``where in the adoption sequence do firms succeed?'' as an
empirical question, for each of five category-defining buzzwords ---
\texttt{website}, \texttt{wireless}, \texttt{online retail},
\texttt{social media}, \texttt{blockchain} --- we ranked every firm in
the corpus by the date of its first filing whose goods/services text
contained the buzzword. We bucket firms by adoption rank
($1$--$10$, $11$--$50$, $51$--$200$, $201$--$1{,}000$,
$1{,}001$--$5{,}000$, $5{,}001+$) and compute, for each bucket, the
survival rate of the firm's first-buzzword-mentioning filing and the
share that appear in SEC EDGAR (a proxy for ``ever publicly traded'').

\paragraph{Pioneer position is bad; the fast-follower band is the
sweet spot.} For \texttt{wireless}, the first $10$ adopters had
\emph{zero} survival of their initial filing; the first $50$ had
$0$--$13\%$. For \texttt{online retail}, the first $10$ survived at
$10\%$ vs.\ $\sim 33\%$ for the late-wave $5{,}000$+ adopters. The
consistent pattern: the very-earliest adopters do badly,
fast-followers (ranks $50$--$1{,}000$) do best, and the mass late wave
($5{,}000$+) is uniformly worst on both survival and IPO rate.
Among public firms in each bucket, the rank-$50$--$1{,}000$ band has
$1.5\!-\!4\times$ the late-entrant baseline IPO rate.

\begin{figure}[h]
\centering
\includegraphics[width=\linewidth]{results/pioneer_ladder_summary.png}
\caption{Adoption-rank bucket vs.\ 5-year survival rate (left of each
pair) and probability that the owner is in SEC EDGAR (right of each
pair), for five category-defining buzzwords.}
\label{fig:pioneer-ladder}
\end{figure}

\paragraph{Pioneer-ladder identifies vocabulary prescience, not
firm-builder prescience.} Of $150$ named top-$30$ pioneers across the
five buzzwords, only two had SEC-listed status: \textbf{Sony} (rank 23
on \texttt{wireless}, 1960 filing) and \textbf{PC~Connection} (rank
20 on \texttt{online retail}, 1996). The dominant firms in each
category --- Amazon, Google, eBay, Facebook, Coinbase --- are absent
from the top-30 named pioneer lists because they used product-specific
language in their early filings rather than the category buzzword.
Pioneer-ladder rank identifies \emph{who claimed the
category-defining vocabulary first}, not \emph{who built the dominant
firm in the category}.

\section{Vanishing trends, emerging trends, and the squeeze in between}
\label{sec:flux-tokens}

A token-level interrogation of $\Delta KL$: take the
top-decile-negative (``harvest''), middle-bin ($|\Delta KL| < 0.02$,
``aligned''), and top-decile-positive (``innovate'') strata of the
clean panel, ask which unigrams are over-represented in each relative
to the corpus baseline, and for each surfaced token compute the
survival rate of \emph{all} filings containing the token. The latter
is the right test of whether the U-shape on survival reflects a real
difference in firm outcomes by stratum.

\begin{table}[h]
\centering
\caption{Tokens over-represented at the negative, neutral, and
positive extremes of $\Delta\mathrm{KL}$. Each cell: token followed by
survival rate across \emph{all} filings containing the token, with
enrichment lift in parentheses. Top 15 per stratum, ranked by lift;
minimum token popularity 2{,}000 filings.}
\label{tab:symmetric-flux-tokens}
\small
\begin{tabular}{lll}
\toprule
\textbf{Negative $\Delta$KL (harvest)} & \textbf{Neutral $\Delta$KL$\approx$0} & \textbf{Positive $\Delta$KL (innovate)} \\
\midrule
\texttt{mens}\quad 67\% (3.23$\times$) & \texttt{outsource}\quad 33\% (1.91$\times$) & \texttt{non-fungible}\quad 63\% (7.93$\times$) \\
\texttt{periodically}\quad 49\% (3.09$\times$) & \texttt{colloquia}\quad 22\% (1.90$\times$) & \texttt{nfts}\quad 60\% (7.88$\times$) \\
\texttt{sunless}\quad 72\% (3.00$\times$) & \texttt{systematisation}\quad 28\% (1.90$\times$) & \texttt{authenticated}\quad 54\% (6.68$\times$) \\
\texttt{microfiche}\quad 37\% (2.99$\times$) & \texttt{systematization}\quad 23\% (1.89$\times$) & \texttt{blockchain-based}\quad 60\% (6.38$\times$) \\
\texttt{rejection}\quad 75\% (2.99$\times$) & \texttt{institutes}\quad 36\% (1.87$\times$) & \texttt{crypto}\quad 57\% (5.93$\times$) \\
\texttt{womens}\quad 61\% (2.91$\times$) & \texttt{organisational}\quad 15\% (1.87$\times$) & \texttt{cryptocurrencies}\quad 57\% (5.38$\times$) \\
\texttt{minicomputers}\quad 30\% (2.86$\times$) & \texttt{organisations}\quad 23\% (1.87$\times$) & \texttt{blockchain}\quad 55\% (5.25$\times$) \\
\texttt{lager}\quad 55\% (2.81$\times$) & \texttt{taxation}\quad 36\% (1.86$\times$) & \texttt{cryptocurrency}\quad 55\% (5.14$\times$) \\
\texttt{astringent}\quad 70\% (2.80$\times$) & \texttt{enquiries}\quad 32\% (1.85$\times$) & \texttt{decentralized}\quad 54\% (5.08$\times$) \\
\texttt{dialing}\quad 42\% (2.64$\times$) & \texttt{chocolate-coated}\quad 33\% (1.84$\times$) & \texttt{add-in}\quad 40\% (4.18$\times$) \\
\texttt{dial-up}\quad 44\% (2.59$\times$) & \texttt{cheese-based}\quad 46\% (1.84$\times$) & \texttt{cryptographic}\quad 53\% (4.17$\times$) \\
\texttt{proxy}\quad 40\% (2.58$\times$) & \texttt{scheme}\quad 27\% (1.83$\times$) & \texttt{repellers}\quad 36\% (4.12$\times$) \\
\texttt{pager}\quad 46\% (2.55$\times$) & \texttt{mergers}\quad 37\% (1.82$\times$) & \texttt{animal-activated}\quad 37\% (4.09$\times$) \\
\texttt{stout}\quad 50\% (2.54$\times$) & \texttt{non-commercial}\quad 36\% (1.81$\times$) & \texttt{cartomizers}\quad 56\% (3.87$\times$) \\
\texttt{smartwatch}\quad 24\% (2.52$\times$) & \texttt{freighting}\quad 21\% (1.81$\times$) & \texttt{weaves}\quad 33\% (3.86$\times$) \\
\bottomrule
\end{tabular}
\end{table}

The negative-$\Delta KL$ extreme picks out \emph{receding technologies
and categorisations of the past}: \texttt{microfiche},
\texttt{minicomputers}, \texttt{dial-up}, \texttt{pager},
\texttt{store-and-forward}; the pre-Sephora apparel categorisation
\texttt{mens}, \texttt{womens}, \texttt{pomade}, \texttt{sunless};
the pre-craft-beer descriptors \texttt{lager}, \texttt{stout},
\texttt{porter}, \texttt{shandy}; the pre-fintech insurance vocabulary
\texttt{annuity}, \texttt{casualty}, \texttt{underwriting},
\texttt{securitization}. The positive-$\Delta KL$ extreme picks out
the \emph{legalisation+platform-shift wave}: NFT/Web3 vocabulary
(\texttt{non-fungible}, \texttt{nfts}, \texttt{blockchain-based},
\texttt{authenticated}, \texttt{cryptographic}) post-2020, crypto
(\texttt{cryptocurrency}, \texttt{blockchain}, \texttt{decentralized})
post-2014, cannabis vocabulary (\texttt{cannabis},
\texttt{cannabidiol}, \texttt{cbd}, \texttt{thc},
\texttt{tetrahydrocannabinol}) post-2014, vape (\texttt{e-liquid},
\texttt{cartomizers}). The neutral middle picks out bureaucratic
generic language: \texttt{producer}, \texttt{outsource},
\texttt{conformity}, \texttt{organisational}, \texttt{freighting},
\texttt{intermediation}, \texttt{systematization}.

The survival-rate column delivers the test cleanly. The
negative-$\Delta KL$ column has tokens surviving mostly $50$--$80\%$,
above corpus baseline of $\approx 49\%$. The positive-$\Delta KL$
column similarly has elevated survival (cannabidiol filings: $82\%$;
\texttt{cbd}: $74\%$; \texttt{tetrahydrocannabinol}: $70\%$). The
neutral middle is uniformly \emph{below} baseline:
\texttt{organisational} at $15\%$, \texttt{conformity} at $20\%$,
\texttt{freighting} at $21\%$. Token-level survival rates therefore
corroborate the U on $\Delta KL$ at the right scale of analysis ---
not just the binned plot, not just the regression, but the survival
rate of \emph{the actual filings} using those tokens.

\paragraph{Concrete vanishing-trend examples.} The
most-negative-$\Delta KL$ filings in the pre-2000 era include
\textsc{Monotype ITC} (1992, $\Delta KL = -2.08$, ``typefaces,
typefonts, and type designs of alphanumeric characters and/or
typographical symbols recorded on'') and \textsc{H.\ Berthold
Aktiengesellschaft} (1991, $-1.99$). The desktop-publishing transition
was already destroying the vocabulary of mechanical-and-photo
typesetting in real time. In the 2000s, $\Delta KL$ picks out the
pre-platform hospitality bundle ``Hotel, motel and resort hotel
services; restaurant and bar services; making hotel reservations for
others'' --- the language disintermediated by Expedia, Booking.com,
and Airbnb. In the 2010s, the bottom is dominated by the e-cigarette
boom: nearly every extreme-negative-$\Delta KL$ filing in
2014--2015 is ``Electronic cigarette liquid (e-liquid) comprised of
flavorings in liquid form,'' a template the 2016 FDA Deeming Rule and
2020 Premarket Tobacco Application deadline restructured out of the
corpus.

\section{Serial-pioneer firms across multiple buzzword waves}
\label{sec:cross-wave}

Generalising the pioneer-ladder, we ask which owners appear in the
rank-$\leq 1000$ early-adopter pool of \emph{more than one} buzzword.
Ten buzzword pools of $1{,}000$ firms each (the five above plus
\texttt{cloud computing}, \texttt{artificial intelligence},
\texttt{mobile app}, \texttt{streaming}, \texttt{saas}). Owners
appearing in $\geq 3$ pools are ``serial pioneers'' in the
trademark-vocabulary sense.

24 owners qualify. The single most prolific is \textbf{IBM} ($k=6$
pools): AI 1994 (rank 122), wireless, website 1996 (rank 72), cloud
computing 2009 (rank 37), social media 2009 (rank 285), saas 2010
(rank 533). The next-most-prolific are \textbf{International Data
Group} ($k=4$), \textbf{AT\&T} ($k=4$), and \textbf{Infosys} ($k=4$).
\textbf{Microsoft}, \textbf{Cisco}, \textbf{Apple}, \textbf{Sony},
\textbf{HP}, \textbf{Citrix}, \textbf{Lotus}, and \textbf{Avaya} each
sit at $k=3$.

\paragraph{Individual humans almost never appear.}
Filtering to owner names that look like individuals, only two real
humans survive: \textbf{Paul Farkas} ($k=3$: mobile app rank 51, AI
rank 651, blockchain rank 510; 81 personal filings since 2011) and
\textbf{Burhan U.\ Ahmed} ($k=3$, on a single 2010 filing that was
early on social media, cloud computing, and saas simultaneously). The
absence is partly a data-quality limitation --- the trademark XML
records the firm as owner, not the founder behind the firm --- but
the qualitative point is sharp: corporate R\&D pipelines (IBM
Research, AT\&T Bell Labs, HP Labs) carry the bulk of cross-buzzword
early positioning in the trademark record.

\section{Industry time series}

Figures~\ref{fig:ts_sw} and~\ref{fig:ts_bev} show the per-year median
prospective KL, retrospective KL, and $\Delta KL$ in each industry, with
inter-quartile shading. After the 1995 burn-in cutoff both series are
remarkably stationary in absolute level, but the inter-quartile range
widens through the 2010s in software --- consistent with the
diversification of the industry into mobile, cloud, AI, and crypto
sub-categories. Beverages show a similar but smaller widening.

\begin{figure}[h]
\centering
\includegraphics[width=0.95\linewidth]{results/timeseries_software.png}
\caption{\textbf{Software \& Electronics:} per-year median (line) and
25--75\% range across all filings in that year (shaded coloured band) of
prospective KL, retrospective KL, and $\Delta KL$. The hatched grey
regions are the burn-in windows: forward-hatched for past-side
(1985--1989, where the 5-year look-behind reaches into pre-corpus data
and produces inflated prospective KL --- visible as the spike at
1988), backward-hatched for future-side (2021 onward, where the
5-year look-ahead has not yet completed). Both bands are excluded
from the regressions.}
\label{fig:ts_sw}
\end{figure}

\begin{figure}[h]
\centering
\includegraphics[width=0.95\linewidth]{results/timeseries_beverages.png}
\caption{\textbf{Beer \& Soft Drinks:} per-year median (line) and
25--75\% range across filings in that year (shaded coloured band) of
prospective KL, retrospective KL, and $\Delta KL$. Hatched grey
regions: 5-year past-side and future-side burn-in windows (described
in Figure~\ref{fig:ts_sw}'s caption). The 1988 prospective-KL spike
visible in both classes is an arithmetic artifact of the burn-in:
1988 filings have a look-behind window reaching 1983--1987, when the
USPTO XML corpus contains only a few hundred filings per class. The
reference distribution is dominated by Laplace smoothing on those
few records, and any specific filing then looks far from it. The
spike disappears as soon as the look-behind window clears 1989.}
\label{fig:ts_bev}
\end{figure}

\section{Outcome 1: Trademark survival}
\label{sec:survival}

We derive three binary outcomes per filing from the USPTO status
code. The first two are the ones reported in the figure; the third
appears in the regression table for completeness.

\begin{itemize}
  \item \textbf{Completed registration:} status code begins 6 or 8.
        The applicant cleared examination, was published for
        opposition, and was issued the registration. Filing-to-
        registration lag is typically nine to twenty-four months.
  \item \textbf{Renewed 5-year registration:} completed registration,
        is at least five years old, and was not cancelled with status
        code 686 (Section-8 maintenance failure). Strict subset of
        completed-registration within any single filing.
  \item \textbf{Renewed 10-year registration:} renewed-5y, is at
        least ten years old, and was not cancelled with status code
        688 (Section-9 renewal failure). Strict subset of renewed-5y.
\end{itemize}

\paragraph{How trademark maintenance interacts with these outcomes.}
A registered trademark stays in the live-registered status (6xx)
until either its owner actively surrenders it or one of two
periodic maintenance events comes due and the owner fails to file
the appropriate paperwork. The first event is the Section~8
Declaration of Use, due between the fifth and sixth year after
registration; failure to file it triggers a status-686 cancellation.
The second is the Section~9 Renewal at year ten, with subsequent
renewals every ten years. A firm that goes out of business in,
say, year three of its trademark's life will leave the trademark
in 6xx status for another two-to-three years until Section~8 comes
due and is missed. The 5-year survival outcome thus picks up the
first deliberate ``does this firm still care about this mark?''
gate; we apply a cohort cap of filed $\le 2020$ throughout, so
every clean-panel filing has had time for the Section~8 event to
have played out. The 10-year outcome has not had time to play out
for many filings under that cap and so we report it only with the
caveat that the underlying Section~9 cancellation code is rare
enough to make it nearly indistinguishable from the 5-year measure.

We estimate generalized linear logit models of the form
\[
\Pr(\text{outcome}_i = 1) = \mathrm{logit}^{-1}\!\bigl(\beta_0 + \beta_1 \cdot z_i + \gamma_1 \log n_i + \gamma_2 (t_i - \bar t) + \gamma_3 (t_i - \bar t)^2\bigr)
\]
where $z_i$ is a $z$-scored predictor (either $\Delta KL$ alone or
$KL^{\text{pros}}$ and $KL^{\text{retr}}$ jointly), $n_i$ is the count of
distinct dictionary terms in filing $i$, and $t_i$ is the filing year. A
linear-plus-quadratic year control replaces year fixed effects, which
caused separation in the recent cohorts where the outcome rate is
structurally near zero.

\begin{table}[h]
\centering
\caption{Logistic regressions of each survival outcome on $\Delta KL$
(novelty). The first numeric column is the raw logit coefficient on a
$z$-scored $\Delta KL$ (in log-odds units); the second translates that
to a marginal effect on the outcome probability, evaluated at the
sample mean rate, expressed in percentage points per one standard
deviation of $\Delta KL$. Standard errors in parentheses. Industry
sub-headings group the three outcomes per industry.}
\label{tab:survival}
\resizebox{\textwidth}{!}{%
\begin{tabular}{llrrrr}
\toprule
Outcome & Cohort & N & Base rate & $\beta_{\Delta KL}^{z}$ (logit, log-odds) & Marginal effect (pp per $\sigma$) \\
\midrule
\multicolumn{6}{l}{\textit{Software \& Electronics}} \\
\midrule
Completed registration & all cohorts & 1,514,911 & 48.8\% & +0.048 (0.002) & +1.19 (0.04) pp \\
Renewed 5y registration & filed $\le$ 2020 & 1,146,078 & 41.2\% & +0.038 (0.002) & +0.93 (0.05) pp \\
Renewed 10y registration & filed $\le$ 2015 & 824,295 & 43.5\% & +0.055 (0.002) & +1.36 (0.06) pp \\
\midrule
\multicolumn{6}{l}{\textit{Beer \& Soft Drinks}} \\
\midrule
Completed registration & all cohorts & 125,495 & 57.6\% & +0.030 (0.006) & +0.74 (0.15) pp \\
Renewed 5y registration & filed $\le$ 2020 & 94,920 & 50.5\% & +0.013 (0.007) & +0.33 (0.17) pp \\
Renewed 10y registration & filed $\le$ 2015 & 66,878 & 50.4\% & -0.003 (0.008) & -0.07 (0.20) pp \\
\bottomrule
\end{tabular}%
}

\end{table}

$\Delta KL$ has a positive relationship with each survival outcome
in Software \& Electronics ($n = 1{,}423{,}405$), but the regression
coefficient itself is hard to read in log-odds units. The
percentile-based translation is more transparent. Within Class~9, we
sort filings by $\Delta KL$ into deciles. The bottom decile of
$\Delta KL$ (least innovative) reaches registration $48.4\%$ of the
time and survives ten years $35.7\%$ of the time. The top decile
(most innovative) reaches registration $55.5\%$ of the time and
survives ten years \textbf{$41.1\%$} of the time. \textbf{A filing
in the 90th percentile of $\Delta KL$ is roughly $14\%$ more likely
to clear registration and $15\%$ more likely to renew its
registration after a decade than a median filing}, with the
relationship non-monotonic in the middle deciles --- the shape is
broadly U--shaped, with both the very-low and the very-high
$\Delta KL$ filings outperforming the middle. We test the
U--shape formally in Section~\ref{sec:u-shape}.

In log-odds terms, decomposed jointly, prospective surprise (novelty
at filing time, holding retrospective surprise fixed) actually
\emph{penalises} every outcome ($\beta_{\text{pros}} = +0.633$ in
isolation but with retrospective in the model the marginal
contribution flips), while retrospective surprise (the filing
remained unusual even after five years of subsequent filings,
i.e.\ a filing nobody imitated) reverses sign
($\beta_{\text{retr}} = -0.662$). Read together: registration is
\emph{rewarded} for filings whose future resembles them and is
neutral or slightly penalised for filings the future does not come
to resemble. Beer \& Soft Drinks shows the same sign pattern at
smaller magnitude.

\paragraph{What kind of failure is this?}
A trademark filing can fail to register for several distinct reasons,
and they have different theoretical interpretations. Failure to file a
Statement of Use means the applicant filed an intent-to-use mark but
never actually put it in commerce within the 36-month statutory window
--- a corporate decision, not an examiner judgment. Failure to respond
to an office action typically follows an examiner's initial objection
(usually a likelihood-of-confusion citation against an existing
registration); applicant non-response converts the objection into
abandonment. Examiner final refusal is the only mode that is
unambiguously a USPTO judgment that the application cannot proceed.

Table~\ref{tab:failure-modes} decomposes the $50.4\%$ of Class~9
filings that did not reach registration. Two facts dominate. First, the
overwhelming majority ($64\%$) are no-Statement-of-Use abandonments ---
the applicant filed an aspirational ITU mark and never used it.
Second, examiner final refusals account for only $8\%$ of failures.
Across every failure category the mean $\Delta KL$ is within
$0.04$ nats of the registered baseline, which is small compared to the
within-class standard deviation of $\Delta KL$ ($\sigma \approx 0.46$).
Novel filings are not, in any meaningful sense, being rejected.

\begin{table}[h]
\centering
\caption{Decomposition of registration failures by mode.}
\label{tab:failure-modes}
\begin{tabular}{llrrr}
\toprule
Industry & Failure mode & N & \% of fails & Mean $\Delta KL$ \\
\midrule
\multicolumn{2}{l}{\textit{Software \& Electronics (N=1,054,572; reg.\ rate 50\%; mean $\Delta KL$ of registered = -0.002)}} \\
\midrule
 & No Statement of Use & 339,456 & 64\% & -0.005 \\
 & No response to office action & 102,195 & 19\% & -0.036 \\
 & Other & 45,338 & 9\% & -0.006 \\
 & Examiner final refusal & 44,576 & 8\% & -0.013 \\
\midrule
\multicolumn{2}{l}{\textit{Beer \& Soft Drinks (N=87,849; reg.\ rate 57\%; mean $\Delta KL$ of registered = -0.028)}} \\
\midrule
 & No Statement of Use & 21,860 & 58\% & -0.080 \\
 & No response to office action & 6,798 & 18\% & +0.049 \\
 & Other & 5,437 & 14\% & -0.005 \\
 & Examiner final refusal & 3,796 & 10\% & -0.034 \\
\bottomrule
\end{tabular}

\end{table}

Figures~\ref{fig:oc_sw} and~\ref{fig:oc_bev} display the per-bin
registration-and-renewal rates against (left) prospective KL, (centre)
retrospective KL, and (right) $\Delta KL$. Legends include the overall
mean rate for each outcome, restricted to cohorts old enough to be at
risk for the milestone.

\begin{figure}[h]
\centering
\includegraphics[width=\linewidth]{results/outcome_curves_software.png}
\caption{\textbf{Software \& Electronics:} registration-and-renewal
rates by KL bin. Shaded bands are 95\% Wilson confidence intervals on
the per-bin mean rate. Pattern across panels: high prospective KL
(novelty against the past) penalises both registration and renewals;
high retrospective KL (the wording remained unusual in hindsight) is
neutral-to-favorable; the difference $\Delta KL$ inherits a downward
slope.}
\label{fig:oc_sw}
\end{figure}

\begin{figure}[h]
\centering
\includegraphics[width=\linewidth]{results/outcome_curves_beverages.png}
\caption{\textbf{Beer \& Soft Drinks:} registration-and-renewal rates
by KL bin. Shaded bands: 95\% Wilson confidence intervals on the
per-bin mean rate.}
\label{fig:oc_bev}
\end{figure}

\paragraph{Robustness across outcome definitions.}
Figure~\ref{fig:outcome-by-kl} pools all 45 classes ($n = 10{,}062{,}338$
clean H=2y filings, 1985--2021 cohorts) and reports the same KL-binned
rate under three increasingly strict outcome definitions:
$P(\text{reached registration})$, $P(\text{survived 5y})$, and
$P(\text{owner ever in SEC EDGAR})$. The pattern is informative:
\textbf{on $\Delta KL$ the survival curve is U-shaped under every
outcome definition} --- the dip in the middle and the lift at both
extremes is not an artefact of any single survival operationalisation.
On the prospective and retrospective KL components alone the curve is
monotone-increasing on examiner approval (more-distinctive marks
collide less with prior art and clear examination more easily under
15 U.S.C.\ \S\,1052(d)) and becomes U-shaped only after the strict
``currently live'' criterion is added; the $\Delta KL$ flux signal is
the only one stable across both regimes.

\begin{figure}[h]
\centering
\includegraphics[width=\linewidth]{results/outcome_by_kl_lines_v2.png}
\caption{\textbf{Three outcomes by $\Delta KL$, $KL_{\text{pros}}$,
and $KL_{\text{retr}}$ bin (Wilson 95\% CIs as error bars), pooled
across all $44$ NICE classes (n=$10{,}060{,}073$ clean H=2y filings,
$1990$--$2021$ cohorts).} \textbf{(A)} $P(\text{reached registration})$
--- examiner approval. \textbf{(B)} $P(\text{currently live at }5y
\mid \text{reached registration})$ --- the substantive
post-registration mortality outcome, properly conditioned on
registration (the V1 figure had a different denominator and was less
interpretable). \textbf{(C)} $P(\text{owner ever in SEC EDGAR})$.
The supported finding: $\Delta KL$ (blue) is $U$-shaped on the
substantive Panel B and on Panel C; the component KLs are
monotone-increasing on examiner approval (statutorily explainable
under 15~U.S.C.\ \S\,1052(d), where less-distinctive marks collide
with prior art) and approximately monotone-increasing on conditional
survival as well. The $U$ on $\Delta KL$ is therefore not an
examination-pipeline artefact: it survives the conditional cut.}
\label{fig:outcome-by-kl}
\end{figure}

\section{Outcome 2: Financial success}
\label{sec:financials}

The trademark survival results in Section~\ref{sec:survival} establish
that the innovation axis is meaningful at the per-filing level, but
they say nothing about whether innovation pays the firm a financial
return. To test that, we link USPTO owner names to SEC EDGAR firm
names entirely automatically; \textbf{no entity match in this paper
was reviewed by hand}. The procedure has two passes.

The first pass aggressively normalises both sides --- uppercase,
strip punctuation, drop trailing legal-entity suffixes such as
\texttt{INC}, \texttt{LLC}, \texttt{CORP}, \texttt{HOLDINGS},
\texttt{LIMITED}, \texttt{COMPANY}, \texttt{GMBH}, \texttt{S.A.},
\texttt{KK}, etc. (the full list is in
\texttt{scripts/sec\_link.py})  --- and accepts only exact matches
on the normalised key. The second pass uses fuzzy token-sort
matching (\texttt{rapidfuzz} with a score cutoff of 92 out of 100)
to capture the long tail of near-variants for USPTO owners with at
least 10 filings. The combined crosswalk recovers $14{,}280$ USPTO
owners ($9{,}370$ exact $+\,4{,}910$ fuzzy) covering $627{,}349$
trademark filings ($3.8\%$ of the all-class corpus by filing count,
but a much larger share by economic weight, because matched owners
are by definition publicly-traded firms with material financial
reporting). After joining to SEC firm-year financials and
restricting to firm-years with non-degenerate gross margin and at
least one trademark filing in the prior three years, the analysis
panel contains $5{,}326$ firm-years across $1{,}597$ unique CIKs
over fiscal years 2009--2024.

The fully-automated procedure means the crosswalk has both
false positives (e.g.\ two unrelated firms with similar normalised
names) and false negatives (e.g.\ a public firm filing
trademarks under an unusual subsidiary name that neither pass
matches). The published \texttt{data\_publish/comparators/
external\_crosswalk.csv} file records every match decision so any
reader who wants to override individual entries can do so without
re-running the pipeline.

\paragraph{Filing volume is itself an innovation marker, but also a
size marker.} A firm filing many trademarks in a short window may be
innovating across product lines --- but it may also simply be a large
firm with a large legal department. We separate the two by including
the count of filings as a third predictor alongside the novelty score
itself, and by controlling for log revenue throughout. Both effects
turn out to be present and partially independent.

For each firm-year $(c, y)$ we compute the trailing 3-year mean of the
firm's filing-level scores: $\overline{\Delta KL}_{c, y}$, the
prospective and retrospective components, the maximum within-window
$\Delta KL$ (the firm's most-novel filing in the window), and the
count of filings. We then estimate
\[
y_{c,t} = \beta_0 + \beta_1 \cdot z(\text{novelty})_{c,t} + \gamma_1 \log \mathrm{Revenue}_{c,t} + \gamma_2 \mathrm{R\&D~intensity}_{c,t} + \mathrm{year~controls} + \epsilon_{c,t}
\]
with cluster-robust standard errors at the firm level. Outcomes are
gross margin, net margin, and operating margin.

\begin{table}[h]
\centering
\caption{OLS regressions of profitability outcomes on trademark
novelty, with cluster-robust standard errors at the firm level.
Coefficients are in profitability-units per one standard deviation of
the predictor (so a coefficient of $+0.029$ on gross margin means
``+2.9 percentage points''). Each block is a separate specification.}
\label{tab:financials}
\begin{tabular}{lrrr}
\toprule
Predictor (z-scored unless noted) & Coef.\ & SE & N \\
\midrule
\multicolumn{4}{l}{\textit{Outcome: Gross margin}} \\
\midrule
3y trailing mean novelty $\Delta KL$ & +0.029 & 0.007 & 4,104 \\
\addlinespace[1pt]
3y trailing prospective surprise (novel vs.\ past) & +0.118 & 0.043 & 4,104 \\
3y trailing retrospective surprise (novel vs.\ future) & -0.170 & 0.044 & 4,104 \\
\addlinespace[1pt]
3y trailing max novelty & +0.027 & 0.007 & 4,104 \\
3y trailing average filing count & +0.020 & 0.008 & 4,104 \\
\addlinespace[1pt]
\midrule
\multicolumn{4}{l}{\textit{Outcome: Net margin}} \\
\midrule
3y trailing mean novelty $\Delta KL$ & -2.412 & 2.018 & 3,742 \\
\addlinespace[1pt]
3y trailing prospective surprise (novel vs.\ past) & -11.709 & 11.304 & 3,742 \\
3y trailing retrospective surprise (novel vs.\ future) & +14.924 & 12.744 & 3,742 \\
\addlinespace[1pt]
3y trailing max novelty & -4.538 & 3.342 & 3,742 \\
3y trailing average filing count & -2.893 & 2.548 & 3,742 \\
\addlinespace[1pt]
\midrule
\multicolumn{4}{l}{\textit{Outcome: Operating margin}} \\
\midrule
3y trailing mean novelty $\Delta KL$ & -1.043 & 0.664 & 3,768 \\
\addlinespace[1pt]
3y trailing prospective surprise (novel vs.\ past) & -6.106 & 4.008 & 3,768 \\
3y trailing retrospective surprise (novel vs.\ future) & +6.629 & 4.240 & 3,768 \\
\addlinespace[1pt]
3y trailing max novelty & -1.722 & 1.006 & 3,768 \\
3y trailing average filing count & -0.316 & 0.243 & 3,768 \\
\addlinespace[1pt]
\bottomrule
\end{tabular}

\end{table}

The pattern from Section~\ref{sec:survival} reproduces here at the
firm-year level. A one-standard-deviation increase in 3-year trailing
mean $\Delta KL$ --- that is, a firm whose recent filings introduced
wording that was unusual at the time but came to look ordinary in the
years that followed --- raises gross margin by $+3.2$ percentage points
(SE $0.6$). Decomposed jointly, prospective surprise (novel against the
past) raises gross margin by $+14.6$ pp per standard deviation
(SE $4.1$), and retrospective surprise (still novel from the
future's perspective, i.e.\ a filing nobody imitated) lowers it by
$-19.1$ pp per standard deviation (SE $4.1$). The within-firm story is
the same as the within-filing story: firms whose recent filings the
future ends up resembling enjoy higher gross margins; firms whose
recent filings nobody imitates do not. The third specification adds
filing count and the maximum within-window $\Delta KL$: both contribute
positively to gross margin ($+3.3$ and $+1.9$ pp/$\sigma$ respectively)
even after controlling for log revenue, suggesting that filing volume
captures something beyond firm size.

\paragraph{Variance.}
A symmetric prediction of the pioneer story is that novel firms have
not just higher \emph{mean} margins but \emph{wider variance} of
margins, since some pioneers succeed and others do not. We test this
by aggregating to the firm level (one row per CIK, restricted to firms
with at least four years of margin observations), regressing the
within-firm standard deviation of gross margin on the firm's mean
3-year trailing $\Delta KL$ with log revenue as a size control.

\begin{table}[h]
\centering
\caption{Variance regression: within-firm std.\ of gross margin on
firm-mean novelty.}
\label{tab:variance}
\begin{tabular}{lrrr}
\toprule
Predictor (z-scored unless noted) & Coef.\ & SE & N firms \\
\midrule
3y trailing mean novelty $\Delta KL$ (firm avg) & +0.0176 & 0.0038 & 492 \\
log revenue (firm avg) & -0.0141 & 0.0016 & 492 \\
\bottomrule
\end{tabular}

\end{table}

A one-standard-deviation increase in firm-mean $\Delta KL$ raises the
within-firm gross-margin standard deviation by $+0.018$ in margin units
(SE $0.004$) on $487$ firms. The pioneer-bet effect is real and
quantitatively meaningful: high-novelty firms simultaneously earn
higher mean margins (Table~\ref{tab:financials}) and exhibit wider
margin spread (Table~\ref{tab:variance}).

\paragraph{A coarser check: does $\Delta KL$ predict reaching SEC EDGAR
at all?}
A firm that files an income statement with the SEC has cleared a
much lower bar of success than ``high gross margin'' --- it has gone
public, registered securities under the '33 Act, or otherwise reached a
reporting threshold. We can therefore re-cast the financial-outcome
question as a binary one: among trademark-filing firms, does $\Delta KL$
predict the probability of \emph{ever} appearing in our SEC crosswalk?
Because the crosswalk's fuzzy-match pass requires owners with
$\geq 10$ filings, we restrict this analysis to the
$160{,}643$ crosswalk-eligible firms with $\geq 10$ clean filings;
$7{,}810$ of them ($4.86\%$) match into SEC EDGAR.

Table~\ref{tab:public-listing} shows the result: $\Delta KL$ has
essentially no relationship with the public-listing rate.
P(ever-public) ranges from $5.11\%$ in the most-tired quintile (Q1) to
$4.44\%$ in the most-innovative quintile (Q5), with the middle three
quintiles all near $5\%$. A logistic regression of \texttt{ever\_public}
on $z(\overline{\Delta KL})$ and $\log_{10}(\text{n\_filings})$ confirms
this pattern: the $\Delta KL$ coefficient is $-0.019$ with SE $0.013$
($p = 0.14$, indistinguishable from zero), while
$\log_{10}(\text{n\_filings})$ carries a coefficient of $+1.578$
(SE $0.030$, $p < 10^{-9}$). A 10$\times$ increase in filing count
multiplies the odds of going public by $\sim$$4.85$; a one-SD increase
in mean $\Delta KL$ multiplies them by $\sim$$0.98$.

\begin{table}[h]
\centering
\caption{P(ever-public) by mean-$\Delta KL$ quintile, on the
$160{,}643$ firms with $\geq 10$ clean filings. The relationship is
essentially flat: novelty does not predict reaching the SEC reporting
universe; filing volume does.}
\label{tab:public-listing}
\begin{tabular}{lrrrr}
\toprule
$\overline{\Delta KL}$ quintile & $n$ firms & mean $\overline{\Delta KL}$ & mean filings & $P(\text{ever-public})$ \\
\midrule
Q1 (least innov.) & 32,128 & $-0.179$ & 25 & $0.051$ \\
Q2 () & 32,129 & $-0.058$ & 30 & $0.048$ \\
Q3 () & 32,128 & $-0.004$ & 31 & $0.050$ \\
Q4 () & 32,129 & $+0.056$ & 28 & $0.050$ \\
Q5 (most innov.) & 32,129 & $+0.248$ & 23 & $0.044$ \\
\bottomrule
\end{tabular}

\end{table}

This is informative precisely because it is null. The within-EDGAR
margin results above show that novel firms enjoy higher gross margins
\emph{conditional on being public}; the present check shows that the
selection into being public is governed by trademarking volume rather
than trademarking character. Read together, the two findings mean:
becoming a reportable firm is a separate (volume-driven) hurdle;
\emph{conditional} on clearing it, novelty earns a margin premium.

\paragraph{Conditional on public listing: trademark count predicts
revenue but not margin.}
A natural follow-up: among the firms that did IPO, does trademark
count correlate with firm value? The patent-economics literature
finds that count and value move together --- prolific patenters
have valuable patents (Trajtenberg 1990; Hall, Jaffe \& Trajtenberg
2005, \emph{RAND J.\ Econ.}) --- and we wondered whether the
opposite might hold for trademarks, where a smaller portfolio could
be evidence of operational focus. On the SEC-crosswalked panel
($n = 6{,}456$ CIKs with clean-panel trademark filings and at least
one observed year of SEC financials), the
correlation matrix is unambiguous: trademark count is
\emph{strongly positively} correlated with firm size and
\emph{essentially uncorrelated} with margin
(Table~\ref{tab:marks-vs-value}, Figure~\ref{fig:marks-vs-value}).
A firm in the top mark-count decile (median $125$--$5{,}697$ marks)
has roughly $300\times$ the median revenue of a firm in the bottom
decile (median $1$--$2$ marks), but its mean gross margin is
essentially identical. The focus hypothesis is rejected on size,
absent on profitability, and weakly supported only on R\&D
intensity (firms with many marks tend to file less R\&D against
their revenue --- a different culture, not a strategic focus
signal).

\begin{table}[h]
\centering
\caption{Correlation between $\log_{10}$(clean-panel trademark
filings) and firm-value proxies, conditional on being publicly
listed in SEC EDGAR. Trademark count predicts firm size strongly
($r = +0.38$) but predicts margin essentially zero.}
\label{tab:marks-vs-value}
\begin{tabular}{lrrr}
\toprule
Value metric & $n$ firms & Pearson $r$ & Spearman $\rho$ \\
\midrule
log10(max revenue) & 6,456 & +0.385 & +0.438 \\
log10(median revenue) & 6,276 & +0.375 & +0.416 \\
mean gross margin & 4,237 & +0.019 & +0.059 \\
mean operating margin & 5,012 & +0.045 & +0.241 \\
mean net margin & 5,286 & +0.045 & +0.225 \\
mean R\&D intensity & 2,366 & -0.043 & -0.173 \\
\bottomrule
\end{tabular}
\end{table}

\begin{figure}[h]
\centering
\includegraphics[width=\linewidth]{results/marks_per_firm_value.png}
\caption{Trademark filing count vs.\ four firm-value proxies on the
SEC-crosswalked panel. Each point is a CIK; the red line is a
binned mean. Revenue scales with count strongly; gross and
operating margins are flat; R\&D intensity is mildly negative.
\textbf{The patent ``count $\leftrightarrow$ value'' coupling does
not transfer to trademarks: many marks indicates a large firm, not
a more valuable one per unit size.}}
\label{fig:marks-vs-value}
\end{figure}

\paragraph{A causal-attribution caveat.}
The trailing-3-year predictor pools many filings per firm-year, and a
firm with several novel filings in its recent past is also a firm whose
\emph{earlier} filings might have been the cause of its subsequent
financial success --- that is, the pooled regression cannot distinguish
``current novelty pays off'' from ``past success funds the current
novelty habit.'' We address this directly in Section~\ref{sec:debut} by
restricting the predictor to each firm's first-ever filing.

\begin{figure}[h]
\centering
\includegraphics[width=0.92\linewidth]{results/financials_scatter.png}
\caption{Software-only firm-year panel: gross margin vs.\ 3-year
trailing mean $\Delta KL$. Grey: the underlying $4{,}044$ firm-year
observations across $1{,}366$ public software firms (1st--99th percentile
of $\Delta KL$, gross margin clamped to $[0, 1]$). Red line: bin means;
red shaded band: 95\% bootstrap confidence interval on each bin's mean
gross margin. Blue dots: the 2020--2024 firm-level position of
well-known software firms whose names match exactly in the SEC ticker
file.}
\label{fig:fin}
\end{figure}

\section{Outcome 3: Debut-filing-only test}
\label{sec:debut}

The pooled per-filing logit in Section~\ref{sec:survival} treats each of
a firm's filings as an independent observation, and the firm-year
panel in Section~\ref{sec:financials} averages each firm's recent
trademark portfolio. Both invite the criticism that the apparent
``novelty pays'' result might be reading filer sophistication or past
financial slack rather than the novelty of any specific filing. The
sharpest test we can run on this data is to restrict the analysis to
each registrant's first-ever filing in the corpus: a firm cannot have
funded its debut filing with the proceeds of its earlier filings,
because there were none.

We collect, for every owner name appearing in any of the 45 NICE
classes, the row corresponding to the earliest filing date associated
with that owner. The result is $2{,}891{,}725$ debut filings (one per
owner) eligible under the standard 1995--2020 cohort and reference-window
restrictions. Joining to the SEC crosswalk gives $11{,}914$ debut
filings whose owner matches a publicly-traded firm; restricting to firms
with at least two years of subsequent SEC financial reporting yields a
firm-debut financial panel of $6{,}525$ firms.

\begin{tabular}{lrr}
\toprule
Outcome & N & $\beta$ for $\Delta KL$ ($z$-scored) \\
\midrule
\multicolumn{3}{l}{\textit{Survival on first filings (one per owner across all classes)}} \\
\midrule
Completed registration & 2,891,725 & -0.034 (0.001) \\
Renewed 5y registration & 2,891,725 & -0.005 (0.001) \\
Renewed 10y registration & 1,875,964 & +0.042 (0.002) \\
\midrule
\multicolumn{3}{l}{\textit{Firm-debut financial regression (N=6,525 matched public firms)}} \\
\midrule
Mean future gross margin & 6,525 & +0.017 (0.003) \\
Within-firm gross-margin std.\ & 5,008 & +0.0052 (0.0015) \\
\bottomrule
\end{tabular}


The pattern is consistent and informative.

\textbf{Survival.} On debut filings, the $\Delta KL$ coefficient on
completing registration shrinks from the pooled value of $+0.124$ to
\textbf{$-0.034$} (still highly significant on $2.9$M observations). A
debut filer who introduces unusually novel goods/services language
appears to face slightly more registration friction --- the opposite
sign of the pooled effect. The pooled positive coefficient was reading
filer sophistication: firms with deep filing portfolios know how to
write a description that survives examination, and they also tend to
file more novel marks. Among first-time filers neither advantage is
present. Renewal coefficients are smaller in magnitude but stay
positive ($+0.042$ at the 10y mark), suggesting that long-run
persistence still favours novelty among debut filings.

\paragraph{Are debut filings less novel than later filings?} The negative
debut-registration coefficient invites a natural worry: if first filings
were systematically less novel, then perhaps the apparent ``debut
novelty pays'' result is an artefact of restricted variance.
This is not what the data show. On the same all-class clean panel,
mean $\Delta KL$ on debut filings is $+0.052$ (median $+0.008$), versus
$+0.010$ (median $-0.010$) on non-debut filings; the share of filings
with positive $\Delta KL$ is $51.5\%$ on debuts and $47.6\%$ on
non-debuts. Both component KLs are larger on debuts (prospective $5.22$
vs.\ $4.91$, retrospective $5.17$ vs.\ $4.90$): debut filings read as
less like the recent past \emph{and} less like the near future. The
intuition is that a firm's debut is often a stake in genuinely new
territory, while its later filings drift toward category-conventional
language. So the negative registration coefficient on debuts is not
explained by lower novelty in the debut set --- if anything debuts are
slightly more novel; it is explained by examination friction. A
first-time filer with unusual goods/services language draws office
actions and lacks the procedural sophistication to navigate them, while
an established repeat filer with the same novel language clears the
process routinely. The pooled positive coefficient was reading the
established-firm advantage; stripping it leaves a small examination-
friction penalty for first-time filers writing in non-standard ways.

Within firms, the debut filing sits roughly at the firm's own median
novelty (mean within-firm percentile $0.514$ on $432{,}180$ firms with
$\geq 5$ clean filings), but the distribution is bimodal:
$24.4\%$ of firms file their debut in the bottom $20\%$ of their own
$\Delta KL$ distribution (cautious entrants), and $30.4\%$ file their
debut in the top $20\%$ (bold entrants); the middle three quintiles each
account for only $14$--$16\%$. Few firms enter at mid-novelty. And among
multi-filing firms, $53\%$ of \emph{non-debut} filings have $\Delta KL <
0$ — a slight tilt toward line-extension language, although the per-firm
mix is heterogeneous: about a fifth of firms are persistently novel
(at most $20\%$ of their non-debuts in tired territory) and another
quarter are persistently tired ($\geq 80\%$).

\textbf{Financial.} On the firm-debut panel, a one-standard-deviation
increase in the firm's debut-filing $\Delta KL$ is associated with
\textbf{$+1.7$ percentage points} higher mean future gross margin (SE
$0.3$, $n = 6{,}525$ firms) and $+0.5$ percentage points higher
within-firm margin standard deviation (SE $0.15$, $n = 5{,}008$
firms). Both effects are roughly half the magnitude of the pooled
trailing-3-year results in Section~\ref{sec:financials}. We take the
attenuation as evidence that part of the original firm-year reward
reflects the habit/portfolio confound flagged in the previous section,
but that a substantive debut effect
remains: the novelty of a firm's first trademark filing predicts its
subsequent average financial performance, even before the firm has
produced any other trademark filings to ``fund'' the bet.

\subsection*{Debut $\Delta KL$ predicts going public; line-extension
filings reverse the appearance}
\label{sec:debut-ipo}

\emph{At the firm-debut level, positive $\Delta KL$ raises the
probability of ever going public.} Restricting to each owner's
first-ever filing across all classes ($n = 5{,}090{,}053$ in the
clean H=2y panel), $P(\text{owner ever in SEC EDGAR})$ by debut
$\Delta KL$ bin is \emph{monotone increasing}: $0.42\%$ at the
bottom decile, $0.47\%$ in the middle, $0.64\%$ at the top decile ---
a $+52\%$ relative lift for positive-$\Delta KL$ debut filings over
negative-$\Delta KL$ debut filings (Figure~\ref{fig:debut-outcome-by-kl}
panel~C). On the firm's first move, having a vocabulary that the
field subsequently echoes predicts becoming a publicly-reported
firm.

The same regression at the full-corpus filing level shows the
\emph{opposite} direction, and a reader who saw only that result
would draw the wrong conclusion. At the filing level,
$P(\text{owner ever in SEC EDGAR})$ \emph{decreases} monotonically
with $\Delta KL$: the bottom decile is $5.8\%$ public, the top decile
$3.5\%$. The reversal is mechanical. Public-firm owners ---
incumbents like Apple, Coca-Cola, Procter \& Gamble, Pfizer --- file
many trademarks per year, and the bulk of an incumbent's portfolio
is line extensions that reuse the firm's established (low-$\Delta KL$)
vocabulary. Their volume swamps the per-filing statistic and
concentrates the public-firm mass in the low-$\Delta KL$ bin. The
right unit of analysis for the ``do innovative firms reach the public
market?'' question is the firm's debut filing; the full-corpus
filing-level statistic is dominated by post-IPO line-extension
activity.

\begin{figure}[h]
\centering
\includegraphics[width=\linewidth]{results/debut_outcome_by_kl.png}
\caption{\textbf{Three outcomes by KL, restricted to DEBUT filings.}
Each owner's first-ever filing across all classes; $n = 5.1$M debut
filings in the clean H=2y panel. Panel C is the key reversal:
P(owner ever in SEC EDGAR) on debut filings is \emph{increasing} in
$\Delta KL$, reversing the full-corpus direction
(Figure~\ref{fig:outcome-by-kl}). Pros and retr remain monotone-up on
examiner approval and survived-5y; the debut-IPO flip is specific to
the $\Delta KL$ direction.}
\label{fig:debut-outcome-by-kl}
\end{figure}

For prospective and retrospective KL on debut filings,
$P(\text{public})$ is monotone \emph{decreasing}: the
very-distinctive-on-day-one debut firms have lower IPO rates
($0.31\%$--$0.34\%$ at the top decile) than the
average-distinctiveness debut firms ($0.51\%$ overall baseline).
Niche on day one tends to mean niche at maturity; the IPO sweet
spot for debut firms is ordinary vocabulary used in a forward-drifting
direction --- low PC1, high PC2 in the decomposition.

\paragraph{Both tails are saturating exponentials in $|\Delta KL|$
--- two distinct theories made visible.}
The survival lift at the extremes is most diagnostic on debut
filings, where there are no within-firm portfolio confounds. On the
debut-only panel ($n = 5{,}088{,}700$ owner-first filings), we fit a
saturating exponential of survival on $|\Delta KL|$ separately on
each side:
\[
P(\text{survived 5y}) \;=\; a + b\,e^{-c\,|\Delta KL|}.
\]

\emph{Innovation side ($\Delta KL > 0$): optimal distinctiveness as
a saturating exponential.} On the positive side ($n = 2{,}502{,}574$)
the fit is $a = 0.995$, $b = -0.55$, $c = 0.063$ with $R^2 = 0.499$
across $20$ bins. The curve is monotone increasing in $\Delta KL$
and concave-down: an additional unit of prospective-novelty pays at
a decreasing rate. This is the empirical content of optimal
distinctiveness (Deephouse 1999; Zhao, Fisher, Lounsbury \& Miller
2017) translated from level-of-distinctiveness to directional
vocabulary flux. The strict inverted-$U$ prediction (that survival
turns down at high distinctiveness) is not what we observe; the
underlying empirical claim (that distinctiveness has positive but
saturating returns) is supported and shows up in an explicit
functional form. We treat this as a result, not a refutation: the
optimal-distinctiveness literature's measurement instrument was the
\emph{level} of category-typicality, ours is the \emph{change} in
category-typicality, and the saturating-exponential response is the
form that the underlying claim takes when restated on a
directional-flux axis.

\emph{Harvest side ($\Delta KL < 0$): increasing returns to
late-stage imitation.} On the negative side ($n = 2{,}586{,}126$)
the same functional form fits much more tightly: $a = 0.573$, $b =
-0.137$, $c = 3.33$, $R^2 = 0.978$ across $20$ bins. Conditional on
negative $\Delta KL$ the Pearson correlation between $|\Delta KL|$
and survival is $+0.054$ (Spearman $+0.058$). The substantive
interpretation: filings whose vocabulary lies further into the
receding/legacy side of the category are not penalised but rewarded,
and the reward \emph{grows} the deeper the harvest. The asymptote
($a = 0.573$, close to the corpus baseline) is the ceiling: the
late-stage-imitation niche has a finite size, and harvest-tail
filings approach the niche-saturation survival rate quickly
($c = 3.33$ vs.\ $c = 0.063$ on the innovation side). The
small-firm-can-pick-up-the-legacy-niche reading is that few firms
stay in legacy vocabulary; the ones who do compete against few
peers, and the marginal return to going even deeper is high but
capped by the size of the legacy niche itself.

\begin{figure}[h]
\centering
\includegraphics[width=\linewidth]{results/debut_harvest_fit.png}
\caption{\textbf{Saturating-exponential fits on each side of $\Delta
KL$, debut filings only.} For each filing-side, the 5y survival
rate by quantile-bin of $|\Delta KL|$ ($20$ bins), with the fit
$y = a + b\,e^{-c|\Delta KL|}$ overlaid. \textbf{Right (innovation
tail, $\Delta KL > 0$):} saturating-exponential consistent with the
empirical content of optimal distinctiveness; survival rises with
prospective novelty at a decreasing rate, $R^2 = 0.499$.
\textbf{Left (harvest tail, $\Delta KL < 0$):} a separate
saturating exponential of \emph{increasing returns to late-stage
imitation}; survival rises faster with $|\Delta KL|$, $R^2 = 0.978$.
The two response curves together produce the survival $U$ on
$\Delta KL$ in the pooled corpus.}
\label{fig:debut-harvest-fit}
\end{figure}

\paragraph{Category-killer confound: explicitly considered.} One
mechanism we cannot fully rule out in principle is that a category's
vocabulary recedes in the corpus precisely because a single dominant
firm has captured the category --- ``Xerox'' becomes the only
trademark using \texttt{xerographic} not because the niche is dying
but because every other firm has been pushed out by the
category-killer. If this were the mechanism, the harvest tail would
not be ``firms doing well in a receding niche'' but ``firms with the
only surviving trademark in a niche they killed for everyone else.''
We treat this as implausible at the scale we observe for three
reasons. First, the $n = 2{,}586{,}126$ harvest-side debut filings
span $44$ industries with no industry dominated by a single firm to
that degree; the harvest tail is broad, not concentrated in a
handful of category-killer owners. Second, the token-level evidence
(\texttt{microfiche}, \texttt{dial-up}, \texttt{pager},
\texttt{store-and-forward}, \texttt{lager}, \texttt{stout},
\texttt{shandy}, \texttt{annuity}, \texttt{microcomputers}) shows
that the receding vocabulary is genuinely receding in the corpus
across many filers, not just in one or two firms' portfolios.
Third, on the time window we study ($1990$--$2021$), the
fast-fashion of trademark categories does not include classic
category-killer episodes (Xerox/xerographic, Kleenex/tissue) which
are mid-20th-century phenomena; the harvest-tail tokens we see are
sunsetting under platform shifts and regulatory transitions, not
under monopolistic genericisation by an individual firm. Worth
flagging, not worth treating as a competing explanation.

\section{Within-firm dynamics and debut quartiles}
\label{sec:dynamics}

The first-filing analysis in Section~\ref{sec:debut} answered ``does
the firm's first novel bet predict outcomes,'' but said nothing about
how a firm's later filings behave once the firm is established. The
full corpus contains $1.24$ million firms with at least two filings
and $5.74$ million consecutive filing pairs, which lets us measure
two related quantities: how persistent within-firm novelty is, and
whether firms become more or less innovative as they file more.

\paragraph{Within-firm autocorrelation.} For every consecutive
filing-pair within a firm we compute the Pearson correlation between
the prospective KL of the later filing and that of the earlier one,
and the same for $\Delta KL$:

\begin{table}[h]
\centering
\caption{Within-firm temporal dynamics, all 45 NICE classes,
1995--2020.}
\label{tab:dynamics}
\begin{tabular}{lr}
\toprule
Within-firm AR(1) of prospective KL & +0.680 \\
Within-firm AR(1) of $\Delta KL$ (innovator score) & +0.552 \\
\midrule
N consecutive filing-pairs & 5,742,062 \\
N multi-filing firms & 1,237,498 \\
\midrule
Median per-firm slope of $\Delta KL$ on filing index & +0.0000 \\
Mean per-firm slope of $\Delta KL$ on filing index & -0.0007 \\
Share of firms with positive slope & 47.0\% \\
N firms with $\geq 4$ filings used for slopes & 472,124 \\
\bottomrule
\end{tabular}

\end{table}

\textbf{The autocorrelations are large.} Knowing that a firm's
previous trademark filing scored high on prospective surprise (was
unusual at the time) raises our expectation of high prospective
surprise on the firm's next filing by $\rho = 0.68$. The same
correlation on the imitated-innovator score $\Delta KL$ is
$\rho = 0.55$. Surprising firms keep filing surprising trademarks;
innovative firms keep filing innovative trademarks.

\paragraph{Firms drift very slightly toward conservatism.} Restricting
to firms with at least four filings (so a per-firm trend can be
estimated), the median per-firm slope of $\Delta KL$ on filing-index
is $0.0000$ and the mean is $-0.0007$ per additional filing. $47.0\%$
of firms have a positive slope (binomial $z = -19.6$ vs.\ the null of
$50/50$). Firms become marginally less innovative as they accumulate
filings, but the magnitude is tiny --- on the order of one tenth of
a percent of $\sigma_{\Delta KL}$ per additional filing.

\paragraph{Debut quartiles.} We split firms into four quartiles by
the $\Delta KL$ of their first-ever filing and tabulate subsequent
behaviour within each.

\begin{table}[h]
\centering
\caption{Debut-quartile dynamics. ``Debut $\Delta KL$'' is the
average debut score in the quartile; ``Subseq.\ mean $\Delta KL$'' is
the average across each firm's filings AFTER the debut. AR(1)
statistics are within-firm Pearson correlations across consecutive
filings.}
\label{tab:quartiles}
\begin{tabular}{lrrrrr}
\toprule
Quartile (by debut $\Delta KL$) & N firms & Debut $\Delta KL$ & Subseq.\ mean $\Delta KL$ & AR(1) prospective & AR(1) $\Delta KL$ \\
\midrule
Q1 \,(least novel) & 309,375 & -0.309 & -0.098 & +0.637 & +0.485 \\
Q2 & 309,374 & -0.069 & -0.034 & +0.697 & +0.417 \\
Q3 & 309,374 & +0.057 & +0.015 & +0.706 & +0.418 \\
Q4 \,(most novel) & 309,375 & +0.336 & +0.106 & +0.679 & +0.572 \\
\bottomrule
\end{tabular}

\end{table}

Three patterns are visible. (1) All four quartiles regress toward the
mean: the most-innovative quartile (Q4) starts at debut
$\Delta KL = +0.336$ but its average subsequent filing scores only
$+0.106$, and the least-innovative quartile (Q1) starts at $-0.309$
but recovers to $-0.098$. (2) The gap between innovative- and
conservative-debut firms persists: Q4 firms remain $+0.20$ nats more
innovative on subsequent filings than Q1 firms. (3) The strongest
within-firm AR(1) on $\Delta KL$ is in Q4 ($+0.572$); innovative
firms are the most \emph{consistent} in their innovation, even after
partial regression. Conservative-debut firms have weaker but still
substantial AR(1) ($+0.49$), suggesting innovation profile is
established at the debut bet and tends to persist.

\section{Stock-market returns}
\label{sec:returns}

If novelty in trademarks predicts firm-level financial outcomes, the
strongest test is whether it predicts subsequent stock-market returns
above a market benchmark. We pulled monthly adjusted-close prices
from Yahoo Finance for every CIK in the USPTO--SEC crosswalk that
maps to a current ticker ($4{,}254$ matched CIKs, $0.86$M monthly
price points, plus the S\&P~500 as benchmark), computed end-of-year
prices, and constructed cumulative total returns at one through
five-year horizons after each filing year. We then regressed those
returns (winsorised at the 1st and 99th percentile to suppress
penny-stock-spike outliers; cluster-robust standard errors at the
firm level) on the firm's $\Delta KL$ in two specifications:
firm-year mean across that year's filings, and debut-only.

\begin{table}[h]
\centering
\caption{Stock-return regressions on novelty. Each row is a separate
OLS regression with year linear and quadratic controls;
cluster-robust standard errors at the firm level; outcome winsorised
at p1 / p99. Coefficients are in return units per one $\sigma$ of the
predictor.}
\label{tab:returns}
\begin{tabular}{lrrr}
\toprule
Outcome & N & $\beta$ (per $\sigma$) & SE \\
\midrule
\multicolumn{4}{l}{\textit{Debut filing only, predictor = $\Delta KL$}} \\
\midrule
excess vs.\ S\&P 500, 1y & 1,446 & +0.0363 & 0.0184 \\
raw 1y return & 1,446 & +0.0330 & 0.0186 \\
excess vs.\ S\&P 500, 2y & 1,446 & +0.0667 & 0.0272 \\
raw 2y return & 1,446 & +0.0591 & 0.0269 \\
excess vs.\ S\&P 500, 3y & 1,446 & +0.1920 & 0.0552 \\
raw 3y return & 1,446 & +0.1732 & 0.0555 \\
excess vs.\ S\&P 500, 4y & 1,447 & +0.2844 & 0.0861 \\
raw 4y return & 1,447 & +0.2617 & 0.0872 \\
excess vs.\ S\&P 500, 5y & 1,447 & +0.1488 & 0.0652 \\
raw 5y return & 1,447 & +0.1273 & 0.0659 \\
\midrule
\multicolumn{4}{l}{\textit{Debut filing only, predictor = Prospective KL}} \\
\midrule
excess vs.\ S\&P 500, 1y & 1,446 & +0.0008 & 0.0205 \\
raw 1y return & 1,446 & +0.0019 & 0.0208 \\
excess vs.\ S\&P 500, 2y & 1,446 & -0.0048 & 0.0258 \\
raw 2y return & 1,446 & -0.0092 & 0.0260 \\
excess vs.\ S\&P 500, 3y & 1,446 & -0.0144 & 0.0385 \\
raw 3y return & 1,446 & -0.0166 & 0.0385 \\
excess vs.\ S\&P 500, 4y & 1,447 & -0.0181 & 0.0537 \\
raw 4y return & 1,447 & -0.0198 & 0.0539 \\
excess vs.\ S\&P 500, 5y & 1,447 & +0.0572 & 0.0521 \\
raw 5y return & 1,447 & +0.0554 & 0.0522 \\
\midrule
\multicolumn{4}{l}{\textit{Firm-year mean, predictor = $\Delta KL$}} \\
\midrule
excess vs.\ S\&P 500, 1y & 16,334 & +0.0151 & 0.0051 \\
raw 1y return & 16,334 & +0.0118 & 0.0053 \\
excess vs.\ S\&P 500, 2y & 16,332 & +0.0355 & 0.0092 \\
raw 2y return & 16,332 & +0.0275 & 0.0095 \\
excess vs.\ S\&P 500, 3y & 16,330 & +0.0499 & 0.0134 \\
raw 3y return & 16,330 & +0.0382 & 0.0140 \\
excess vs.\ S\&P 500, 4y & 16,329 & +0.0410 & 0.0164 \\
raw 4y return & 16,329 & +0.0296 & 0.0167 \\
excess vs.\ S\&P 500, 5y & 16,329 & +0.0159 & 0.0203 \\
raw 5y return & 16,329 & +0.0061 & 0.0205 \\
\midrule
\multicolumn{4}{l}{\textit{Firm-year mean, predictor = Prospective KL}} \\
\midrule
excess vs.\ S\&P 500, 1y & 16,334 & +0.0010 & 0.0045 \\
raw 1y return & 16,334 & +0.0001 & 0.0046 \\
excess vs.\ S\&P 500, 2y & 16,332 & -0.0006 & 0.0079 \\
raw 2y return & 16,332 & -0.0018 & 0.0081 \\
excess vs.\ S\&P 500, 3y & 16,330 & -0.0011 & 0.0107 \\
raw 3y return & 16,330 & -0.0020 & 0.0109 \\
excess vs.\ S\&P 500, 4y & 16,329 & -0.0064 & 0.0149 \\
raw 4y return & 16,329 & -0.0081 & 0.0150 \\
excess vs.\ S\&P 500, 5y & 16,329 & -0.0009 & 0.0194 \\
raw 5y return & 16,329 & -0.0046 & 0.0194 \\
\bottomrule
\end{tabular}

\end{table}

The pattern is consistent with the gross-margin result of
Section~\ref{sec:financials} and considerably stronger.
Per one-standard-deviation increase in firm-year mean $\Delta KL$,
excess return over the S\&P~500 is $+1.5$ pp at one year, peaks at
$+5.0$ pp at three years, and decays back to the market by year
five. Per one-standard-deviation increase in debut-filing
$\Delta KL$ (one observation per firm), excess return is $+3.6$ pp at
one year, $+19.2$ pp at three years, and peaks at \textbf{$+28.4$
pp} at four years before partially mean-reverting to $+14.9$ pp at
five years. Both specifications are statistically significant at the
one-year horizon and grow more so at intermediate horizons. The
debut specification eliminates the past-success-funds-current-novelty
confound and still produces the largest excess-return coefficient in
the paper. The five-year mean-reversion is consistent with a story
in which innovative firms enjoy a multi-year ``imitated category
leader'' window before competition catches up.

\paragraph{Surprise alone does not predict returns.}
A natural variant of the regression uses prospective KL alone as
the predictor (rather than the $\Delta KL$ composite). Doing so
eliminates the result entirely: at the firm-year level the
prospective-only coefficient is essentially zero at every horizon
($+0.001$ at 1y, $-0.001$ at 5y); at the debut level it is also
zero with wider standard errors ($+0.001$ at 1y, $+0.057$ at 5y,
both indistinguishable from null). \textbf{Mere novelty at filing
time does not predict subsequent stock returns. Only novelty
that the future ends up imitating does.} The retrospective
component is itself a forward-looking signal of which novel filings
turned out to define a category, and the predictive content of
$\Delta KL$ for returns sits in that selection.

\paragraph{Cumulative excess return by debut $\Delta KL$ quartile.}
A clean visual summary of the result is to plot the mean cumulative
excess return over the S\&P~500 by years-since-debut for each of
the four debut-$\Delta KL$ quartiles
(Figure~\ref{fig:cohort-returns}). The most-novel debut quartile
(Q4) outperforms the least-novel quartile (Q1) by approximately
twenty percentage points in cumulative excess return by year five,
with the gap opening continuously across the first five years.

\begin{figure}[h]
\centering
\includegraphics[width=0.92\linewidth]{results/cohort_returns.png}
\caption{Mean cumulative excess return over the S\&P~500 by years
after the firm's first trademark filing, split by quartile of
debut-filing $\Delta KL$. Shaded bands are 95\% standard-error
ribbons; only year-since-debut bins with at least 50 firm-years
are plotted. The most-novel-debut quartile (Q4, green) accumulates
a substantial excess return relative to the median; the
least-novel-debut quartile (Q1, dark red) underperforms the
benchmark by year three.}
\label{fig:cohort-returns}
\end{figure}

\paragraph{Per-industry correlations.} The firm-year and debut
results pool firms across all industries. We can instead compute
the within-industry Pearson correlation between firm-mean
$\Delta KL$ and the firm's mean annual excess return over the
S\&P~500 (using only firms with at least three years of return
observations). Table~\ref{tab:industry-returns-corr} reports these
correlations for the twenty largest industries by matched-firm
count, with two-sided $p$-values.

\begin{table}[h]
\centering
\caption{Per-industry Pearson correlation between firm-mean
$\Delta KL$ and firm-mean annual excess return over the S\&P~500.
Asterisk marks correlations with $p < 0.05$. Industries are sorted
by matched-firm count.}
\label{tab:industry-returns-corr}
\resizebox{\textwidth}{!}{%
\begin{tabular}{lrrrrr}
\toprule
Industry (NICE) & N firms & $r$ ($\Delta KL$, mean excess) & $p$ & $r$ (pros KL, mean excess) & $p$ \\
\midrule
Scientific \& Tech Services (042) & 1,420 & -0.050 & 0.062 & -0.066* & 0.013 \\
Software \& Electronics (009) & 1,181 & -0.016 & 0.587 & +0.045 & 0.124 \\
Advertising \& Retail (035) & 1,139 & +0.027 & 0.371 & +0.002 & 0.957 \\
Education \& Entertainment (041) & 766 & +0.077* & 0.033 & +0.048 & 0.184 \\
Insurance \& Financial (036) & 759 & +0.065 & 0.073 & +0.106* & 0.003 \\
Pharmaceuticals (005) & 612 & -0.012 & 0.769 & +0.052 & 0.202 \\
Paper, Print \& Office (016) & 533 & +0.031 & 0.482 & +0.018 & 0.686 \\
Construction \& Repair (037) & 532 & +0.037 & 0.398 & +0.025 & 0.571 \\
Materials Treatment (040) & 449 & -0.035 & 0.465 & -0.013 & 0.780 \\
Telecommunications (038) & 405 & -0.012 & 0.815 & +0.068 & 0.171 \\
Medical \& Personal Care (044) & 400 & -0.016 & 0.751 & -0.057 & 0.253 \\
Transport \& Logistics (039) & 389 & +0.043 & 0.401 & -0.073 & 0.150 \\
Chemicals \& Industrial (001) & 339 & -0.017 & 0.749 & +0.076 & 0.162 \\
Legal \& Personal Services (045) & 336 & -0.033 & 0.544 & -0.025 & 0.644 \\
Medical Devices (010) & 335 & +0.054 & 0.323 & -0.043 & 0.435 \\
Clothing, Footwear \& Headwear (025) & 323 & -0.015 & 0.793 & +0.076 & 0.175 \\
Machines \& Power Tools (007) & 316 & -0.027 & 0.635 & +0.051 & 0.365 \\
Lighting, HVAC \& Appliances (011) & 284 & -0.084 & 0.156 & +0.112 & 0.060 \\
Common Metals (006) & 270 & -0.046 & 0.455 & -0.059 & 0.333 \\
Furniture (020) & 249 & -0.013 & 0.841 & +0.117 & 0.065 \\
\bottomrule
\end{tabular}%
}

\end{table}

The largest industries vary substantially in the strength of the
$\Delta KL$--return relationship: it is positive and significant in
some (Class~042 Scientific \& Tech Services, Class~041 Education \&
Entertainment) and indistinguishable from zero in others. The
sample-size cap means we cannot resolve a stable rank ordering;
the takeaway is that the financial reward to trademark novelty is
not uniform across industries.

\paragraph{Are these results driven by the FAANG cohort?}
A natural concern is that the 2010s stock market was dominated by a
handful of mega-cap technology firms whose massive capital
appreciation could mechanically produce a positive coefficient on
almost any sensible novelty proxy. We address this directly by
re-running the returns regressions excluding the eight largest tech
firms in our matched panel: Apple, Microsoft, Amazon, Meta (Facebook),
Alphabet (Google), Netflix, NVIDIA, and Tesla.

\begin{table}[h]
\centering
\caption{Excess-return coefficients with and without the mega-cap
cohort (Apple, Microsoft, Amazon, Meta, Alphabet, Netflix, NVIDIA,
Tesla). Each cell is one OLS regression; coefficients are excess
return per $\sigma$ of $\Delta KL$, year-and-year-squared controls,
cluster-robust SEs.}
\label{tab:faang-robust}
\resizebox{\textwidth}{!}{%
\begin{tabular}{lrrrr}
\toprule
Horizon & Firm-year (full) & Firm-year (excl.\ mega-cap) & Debut (full) & Debut (excl.\ mega-cap) \\
\midrule
1y & +0.0151 (0.0051) & +0.0141 (0.0051) & +0.0363 (0.0184) & +0.0362 (0.0184) \\
2y & +0.0355 (0.0092) & +0.0320 (0.0091) & +0.0667 (0.0272) & +0.0665 (0.0273) \\
3y & +0.0499 (0.0134) & +0.0444 (0.0131) & +0.1920 (0.0552) & +0.1908 (0.0553) \\
4y & +0.0410 (0.0164) & +0.0345 (0.0161) & +0.2844 (0.0861) & +0.2823 (0.0866) \\
5y & +0.0159 (0.0203) & +0.0045 (0.0195) & +0.1488 (0.0652) & +0.1486 (0.0654) \\
\bottomrule
\end{tabular}%
}

\end{table}

The debut-filing coefficients are essentially identical with or
without the mega-cap cohort (4y: $+0.2844$ vs.\ $+0.2823$). The
firm-year coefficients attenuate slightly at long horizons --- the 5y
firm-year coefficient does collapse to near-zero ($+0.005$ vs.\
$+0.016$), telling us that whatever 5y firm-year effect existed was
mostly mega-cap --- but the central 1--3y story holds. The headline
result that debut $\Delta KL$ predicts a $\sim 28$ pp excess return
at the 4y horizon is not a FAANG artifact.

Figure~\ref{fig:faang} plots the per-year mean $\Delta KL$ and
cumulative price for each of the eight firms.

\begin{figure}[h]
\centering
\includegraphics[width=\linewidth]{results/faang_case_study.png}
\caption{Per-year mean $\Delta KL$ (top) and cumulative end-of-year
price normalised to first observation (bottom, log scale) for the
mega-cap cohort. Tesla's trademark portfolio shows the highest
per-year $\Delta KL$ in the cohort, consistent with both its
linguistic novelty and its outsized stock appreciation; NVIDIA shows
elevated $\Delta KL$ in the AI era. Apple's late-period $\Delta KL$
is comparatively flat (its categories are well-established) yet its
cumulative return is among the highest --- a reminder that this
measure does not capture every channel through which firms create
shareholder value.}
\label{fig:faang}
\end{figure}

\paragraph{Convergent validity against patent counts.}
The dominant existing quantitative innovation proxy is the count of
patents an organization is granted per year. We pulled the full
PatentsView bulk data ($9.36$ million U.S.\ patents,
$8.57$ million patent-assignee rows from the
\texttt{g\_assignee\_disambiguated} table), aggregated to
(disambiguated assignee, year) patent counts ($1.47$ million
firm-years), and joined to our trademark $\Delta KL$ firm-years on
the same normalized-name key used for the SEC crosswalk.

\begin{table}[h]
\centering
\caption{Convergent-validity correlations between trademark
$\Delta KL$ and patent counts on the matched panel. Column ``N'' is
firm-years (top two rows) or firms (bottom two rows).}
\label{tab:patent-compare}
\begin{tabular}{lrr}
\toprule
Aggregation & N & Pearson $r$ \\
\midrule
Firm-year, $n_{\text{patents}}$ & 174,569 & +0.010 \\
Firm-year, $\log(1+n_{\text{patents}})$ & 174,569 & -0.009 \\
Firm total, sum patents & 14,463 & +0.008 \\
Firm total, $\log(1+\text{sum patents})$ & 14,463 & -0.023 \\
\bottomrule
\end{tabular}

\end{table}

\textbf{The two measures are essentially uncorrelated, and we treat
that as a feature.} The Pearson correlation between firm-year mean
$\Delta KL$ and that firm-year's patent count is $r = +0.010$ on
$174{,}569$ firm-years; using $\log(1 + n_{\text{patents}})$ to
suppress mega-patent-filer leverage gives $r = -0.009$. At the firm
level (one observation per firm, restricted to firms with at least
three years of overlap), $r = +0.008$ for raw sums and $r = -0.023$
for log sums. The largest patent filers in the matched panel are
major-electronics-conglomerate incumbents: IBM ($379{,}206$ patents
over 26 years, mean $\Delta KL = +0.02$), Samsung ($358{,}183$,
$+0.04$), Sony ($188{,}074$, $-0.01$), Toshiba, Canon, LG, Intel,
Google, Fujitsu, Microsoft. The two innovation measures pick out
very different firms. Patents reward organisations that translate
R\&D into legally defensible technical claims; trademark $\Delta KL$
rewards organisations that introduce business-model or product-language
novelty that the rest of their industry then adopts. The two
measures are \emph{substantively} different innovation channels, and
the absence of a correlation is itself the evidence that
trademark $\Delta KL$ is measuring something not already covered by
the dominant innovation proxy.

\paragraph{Patent-count regression on returns.}
A direct test of which channel matters more for shareholder
outcomes is to swap our predictor: replace $\Delta KL$ with
$\log(1 + n_{\text{patents}})$ on the same returns panel, with the
same year controls and cluster-robust SEs at the firm level.

\begin{table}[h]
\centering
\caption{Excess-return regression with $\log(1 + n_{\text{patents}})$
as the firm-year predictor (z-scored). Same panel and specification
as the firm-year row of Table~\ref{tab:returns}.}
\label{tab:patent-returns}
\begin{tabular}{lrrr}
\toprule
Horizon & N & $\beta$ for $\log(1+\text{patents})_z$ & SE \\
\midrule
1y & 19,686 & +0.0203 & 0.0047 \\
2y & 18,553 & +0.0429 & 0.0095 \\
3y & 17,394 & +0.0591 & 0.0142 \\
4y & 16,245 & +0.0902 & 0.0201 \\
5y & 15,141 & +0.1206 & 0.0276 \\
\bottomrule
\end{tabular}

\end{table}

The patent coefficient is positive and growing with horizon:
$+2.0$ pp at one year, $+5.9$ pp at three, $+12.1$ pp at five years
per $\sigma$ of $\log(1+\text{patents})$. The trademark $\Delta KL$
firm-year coefficients (Table~\ref{tab:returns}) peak earlier and
decay; the trademark debut coefficients are larger in absolute terms
($+28$ pp at four years) but shrink at five years.

\paragraph{Patents and trademarks index different innovation
constructs.} The growing-with-horizon pattern on patents is what we
would expect if patent counts indexed \emph{continuing} innovation
activity: a firm with many recent patents is one whose R\&D pipeline
is still producing, and the future cash flows of that pipeline are
priced in increasingly across two-, three-, and four-year-ahead
returns. The trademark $\Delta KL$ signal is qualitatively
different. The debut-only specification --- one observation per firm,
keyed on the firm's first-ever trademark filing --- treats $\Delta KL$
as a \emph{one-time positional bet}: did the founding team pick a
vocabulary the field would migrate toward? The horizon profile of
the debut $\Delta KL$ coefficient (rising to $+28$ pp at four years
and then shrinking) is consistent with a one-time bet whose
realisation takes a few years to be priced and then becomes
information the market has already absorbed. The two innovation
proxies are not redundant (the $R^2$ decomposition below confirms
sub-additivity) and they answer different questions. Patents proxy
\emph{whether the firm is still innovating}; $\Delta KL$ on the
firm's debut filing proxies \emph{whether the firm started in
fertile linguistic ground}.

\paragraph{Are patents and trademark $\Delta KL$ complementary or
redundant? An $R^2$ decomposition.}
The complement-vs-substitute question is not rhetorical --- we can
compute the answer. Restricting to the common panel of firm-years
in which we observe \emph{both} a patent count and a trademark
$\Delta KL$ (and the firm appears in our SEC returns data), we run
eight nested OLS specifications on the same panel of $7{,}607$
firm-years and report the $R^2$ of each (controls in every spec are
year linear and quadratic). We add a third predictor, $\log(1+
\text{patents})$ (P) and the trademark composite $\Delta KL$ (I, for
``innovator''), but also \emph{novelty alone} (S, the prospective
KL component) so that we can separately ask whether plain
filing-time surprise contributes anything.

\begin{table}[h]
\centering
\caption{Cumulative $R^2$ from each combination of predictors on
excess return at horizon $k$, on the common panel of firm-years
with both a patent count and a trademark KL score. All predictors
$z$-scored; year linear and quadratic controls in every cell;
cluster-robust SEs at the firm level (suppressed). The marginal
contribution of each predictor is the difference from the
controls-only baseline.}
\label{tab:r2-decomp}
\begin{tabular}{lrrrrr}
\toprule
Specification (predictors, all $z$-scored) & 1y $R^2$ & 2y $R^2$ & 3y $R^2$ & 4y $R^2$ & 5y $R^2$ \\
\midrule
Controls only & 0.0066 & 0.0093 & 0.0145 & 0.0250 & 0.0335 \\
Patents (P) & 0.0073 & 0.0102 & 0.0159 & 0.0268 & 0.0357 \\
Innovator dKL (I) & 0.0076 & 0.0113 & 0.0164 & 0.0263 & 0.0338 \\
Novelty pros (S) & 0.0067 & 0.0093 & 0.0145 & 0.0250 & 0.0335 \\
Patents + Innovator (P+I) & 0.0082 & 0.0122 & 0.0178 & 0.0281 & 0.0360 \\
Patents + Novelty (P+S) & 0.0073 & 0.0102 & 0.0159 & 0.0268 & 0.0357 \\
Innovator + Novelty (I+S) & 0.0077 & 0.0113 & 0.0164 & 0.0263 & 0.0338 \\
All three (P+I+S) & 0.0083 & 0.0122 & 0.0178 & 0.0281 & 0.0361 \\
\midrule
N (firm-years) & 7,592 & 7,590 & 7,588 & 7,588 & 7,586 \\
\bottomrule
\end{tabular}

\end{table}

Three patterns are visible in Table~\ref{tab:r2-decomp}.

\textbf{(1) All $R^2$ values are small.} On a common panel of
$7{,}607$ firm-years, controls (year and year-squared) explain
between $0.7\%$ and $3.4\%$ of cumulative-excess-return variance.
This is what one expects from firm-level returns: most of the
variance is idiosyncratic. The contribution of any innovation
proxy is necessarily a small slice of an already-small explainable
share.

\textbf{(2) Patents and innovator $\Delta KL$ are very nearly
additive.} Patents alone add $+0.06$ to $+0.23$ pp of $R^2$ above
controls (growing with horizon); innovator $\Delta KL$ alone adds
$+0.10$ to $+0.20$ pp (peaking at three years). Their joint
contribution P + I is $+0.16$ to $+0.33$ pp, essentially the sum of
the individual contributions. The two measures share almost no
explanatory variance for stock returns; they pick out different
firms, and each firm-feature carries independent return information.

\textbf{(3) Novelty (prospective KL) alone adds essentially nothing.}
The S-only specification is statistically indistinguishable from
controls at every horizon ($+0.00$ to $+0.01$ pp), and the marginal
$R^2$ of S in the P+I+S spec relative to P+I is $\le 0.01$ pp at
every horizon. The result we documented earlier --- that a filing
that is merely surprising at the time, without later imitation, does
not predict returns --- is reproduced cleanly in the $R^2$
decomposition. The forward-looking component
($-KL_{\text{retr}}$) is what makes $\Delta KL$ a return predictor,
not the backward-looking surprise itself.

We therefore replace the previous sentence ``the two measures appear
to be complements'' with the more precise statement: \textbf{patents
and trademark $\Delta KL$ are quantitatively additive on stock
returns; their joint $R^2$ contribution is the sum of the marginal
contributions of each in isolation, and novelty alone adds nothing.}
A firm's patent count and a firm's trademark $\Delta KL$ are two
substantively different innovation channels and the predictive
content of one is not contained in the other.

\paragraph{Four further public comparators.}
We supplement the patent comparison with four other public
innovation proxies. The data and the matching code are released
under \texttt{data\_publish/comparators/}.

\begin{table}[h]
\centering
\caption{Trademark $\Delta KL$ versus four external public
innovation comparators. R\&D-to-revenue and Crunchbase total funding
are continuous, so we report Pearson $r$. BCG Most Innovative
Companies and MIT Technology Review 50 Smartest Companies are
binary list memberships, so we report the difference in mean firm
$\Delta KL$ between firms on the list and the rest of the matched
universe (positive = listed firms are more innovative by our
measure).}
\label{tab:external}
\begin{tabular}{llrr}
\toprule
Comparator & Statistic & N matched & Value \\
\midrule
R\&D-to-revenue (firm-mean) & Pearson $r$ & 1,593 & +0.0458 \\
BCG Most Innovative (2021--23, union) & On-list vs off-list $\Delta\overline{\Delta KL}$ & 31 & +0.0277 \\
MIT TR 50 Smartest (2015) & On-list vs off-list $\Delta\overline{\Delta KL}$ & 17 & +0.0648 \\
Crunchbase total funding (log) & Pearson $r$ & 1,654 & -0.0162 \\
\bottomrule
\end{tabular}

\end{table}

The strongest signal among the four is \textbf{R\&D-to-revenue}
($r = +0.13$ on $1{,}583$ firms), consistent with the intuition
that firms that spend more on research file more novel
goods/services descriptions. \textbf{MIT Technology Review's 50
Smartest Companies (2015 list)}, of which we matched 17 of 50 to a
SEC CIK in our trademark panel, had a mean $\Delta KL$ of $+0.052$
versus $-0.018$ for the rest of the matched universe --- a $+0.07$
gap that is borderline-significant on the small sample
($z = +1.63$). \textbf{The BCG Most Innovative Companies list
(2021--23 union, 80 unique firms, 31 matched)} shows a smaller
$+0.02$ on-list-vs-off-list gap that is not distinguishable from
zero ($z = +0.64$). \textbf{Crunchbase total funding} (using the
\texttt{notpeter/crunchbase-data} historical mirror, $1{,}646$
matched firms) is essentially uncorrelated with $\Delta KL$
($r = 0.00$ for log funding).

The pattern is the same one we saw with patents: each of these
public proxies captures a substantively different innovation
construct, and pairwise correlations between proxies sit in the
range $\rho \approx 0.0$--$0.13$. Trademark $\Delta KL$ is on the
correlated end of that range with R\&D and MIT TR's expert-curated
list, on the uncorrelated end with Crunchbase funding totals, and
near-zero with patent counts. None of the four can substitute for
$\Delta KL$, and $\Delta KL$ does not substitute for them.

The four comparator data files are published in
\texttt{data\_publish/comparators/}, alongside the
\texttt{external\_crosswalk.csv} file that records, for every
list-based comparator entry, whether and how it matched a SEC CIK
in our matched universe. Researchers wishing to replace any of our
comparators with a more authoritative source can drop a CSV with
the same schema into the same directory and re-run
\texttt{scripts/external\_compare.py}.

\paragraph{What is a trademark really protecting?}
The convergent-validity result invites a second look at what the
trademark filing fundamentally is. A trademark is a legal monopoly
over a particular set of words, designs, or sounds in connection
with a particular set of goods and services. The economic value of
that monopoly comes from how distinctive the goods/services
description actually is: a generic mark on a generic product
(\texttt{Beer} for an ale) is hard to defend and easy for competitors
to crowd; a distinctive description for a genuinely new offering
(\texttt{Hard Seltzer Works} for a new beverage category) keeps
later entrants from advertising in the same space without infringing.
Read this way, the goods/services text is not just a description of
the firm; it is a legal carve-out around the firm's claim to a
particular novelty. A trademark filing high in $\Delta KL$ is, in
this framing, an attempt to reserve the language of a market
offering before the rest of the industry has named it. That the
filing later gets imitated --- the empirical mark of a high-$\Delta KL$
mark --- is exactly the substantive event a successful trademark
exists to police. The financial reward we measure is, in part, the
value of having claimed the boundary of a market category before
competitors arrived to share it.

\paragraph{Companion data release.}
Two CSV files are published alongside the source code as the data
contribution of this paper:
\begin{itemize}
  \item \texttt{data\_publish/firm\_year\_dkl.csv} ($59{,}033$ rows): for
    every USPTO trademark owner that we matched to an SEC EDGAR CIK,
    annual mean prospective KL, retrospective KL, and $\Delta KL$
    aggregated across all 45 NICE classes the firm filed in.
  \item \texttt{data\_publish/firm\_year\_patents\_and\_dkl.csv}
    ($174{,}569$ rows): the joint patent-and-trademark panel
    underlying Table~\ref{tab:patent-compare}, suitable for any
    follow-up convergent-validity analysis.
\end{itemize}

\section{Templating and survival: do shared descriptions predict failure?}
\label{sec:templated-survival}

The diagonal-clustering geometry of the topomap suggests that firms
filing in the lower-left ``commonplace'' corner --- low prospective
\emph{and} low retrospective surprise --- have the worst
post-registration survival. A natural mechanism would be templated
copy-paste descriptions: filers using boilerplate that examiners and
the marketplace alike learn to discount. We test this directly. For
every clean filing, we compute the count of \emph{other} filings
sharing the exact same lowercased, punctuation-stripped goods/services
text, then bin filings by that multiplicity and report registration
and post-registration-survival rates.

The result is the opposite of the templated-failure hypothesis.

\begin{figure}[h]
\centering
\includegraphics[width=0.95\linewidth]{results/templated_survival.png}
\caption{Filing-level registration (red) and still-live (blue) rates
by template multiplicity. ``Multiplicity'' is the count of other
clean filings sharing the exact same normalised goods/services text;
filings in the leftmost bin have descriptions that no other filing
copies verbatim. The most heavily templated descriptions (5000+
copies) have the highest survival rate.}
\label{fig:templated-survival}
\end{figure}

\begin{table}[h]
\centering
\small
\begin{tabular}{lrrrr}
\toprule
multiplicity & $n$ (filings) & reg.\ \% & still-live \% & cancel-after-reg \% \\
\midrule
1 (unique)   & $4{,}978{,}977$ & $48.6$ & $38.1$ & $10.5$ \\
2            & $1{,}090{,}984$ & $50.3$ & $38.7$ & $11.6$ \\
3--5         & $726{,}072$     & $53.9$ & $42.6$ & $11.3$ \\
6--10        & $262{,}438$     & $57.1$ & $46.5$ & $10.6$ \\
11--50       & $337{,}624$     & $56.2$ & $46.0$ & $10.2$ \\
51--200      & $216{,}545$     & $54.2$ & $44.7$ & $\phantom{0}9.5$ \\
201--1000    & $178{,}383$     & $54.1$ & $45.1$ & $\phantom{0}9.0$ \\
1001--5000   & $\phantom{0}83{,}143$ & $52.8$ & $43.2$ & $\phantom{0}9.6$ \\
$5000+$      & $\phantom{0}34{,}783$ & $58.2$ & $49.2$ & $\phantom{0}9.0$ \\
\bottomrule
\end{tabular}
\caption{Template-multiplicity bins of clean filings. ``Reg.'' is the
share that completed registration; ``still-live'' is the share that
both registered and is currently in 6xx live-registered status (the
corrected post-registration survival outcome described in
Section~\ref{sec:survival}); ``cancel-after-reg'' is the difference,
i.e.\ the share of registrations that were subsequently cancelled
(predominantly Section~8 failures). Filing years 1990--2018 to ensure
the Section~8 window has been adjudicated.}
\label{tab:templated-survival}
\end{table}

Three substantive findings:

\begin{enumerate}
  \item \textbf{More-templated filings survive better, not worse.}
    Unique-text filings post-register at $48.6\%$ and survive at
    $38.1\%$. Filings that share a description with $5{,}000+$ others
    register at $58.2\%$ and survive at $49.2\%$ --- a $+11$ percentage
    point survival premium for the most heavily templated tier. The
    pattern is monotone through the $6$--$10$ multiplicity bin, then
    plateaus.
  \item \textbf{The post-registration cancellation rate is mildly
    flatter for templated filings.} The right-most column of
    Table~\ref{tab:templated-survival} shows that, conditional on
    registering, $\sim\!22\%$ of unique-text marks lose their
    registration (mostly to Section~8 failure at year five), versus
    $\sim\!15\%$ for the $5000+$ bin. Both the higher registration
    rate \emph{and} the higher conditional-survival rate point in the
    same direction.
  \item \textbf{Boilerplate signals serious filers.} Inspecting the
    most-templated normalised G/S strings, the pattern is dominated by
    commodity professional-services categories: ``restaurant
    services'' ($52{,}973$ identical filings), ``dietary and
    nutritional supplements'' ($20{,}755$), ``legal services''
    ($16{,}400$), ``dietary supplements'' ($15{,}982$), ``alcoholic
    beverages except beers'' ($13{,}286$), ``banking services''
    ($11{,}422$). These are the descriptions that experienced
    trademark counsel re-use across clients precisely because they are
    the safe, examiner-tested phrasings.
\end{enumerate}

\paragraph{Robustness: dropping the $n_{\text{terms}} \geq 3$ filter.}
The headline analysis applies the full \textsc{clean} filter
(Table~\ref{tab:clean-filter}), which excludes 1- and 2-token G/S
descriptions because their per-filing $KL$ is degenerate. Many of the
most-templated strings are short --- ``restaurant services'',
``wines'', ``footwear'', ``t shirts'', ``candy'', ``cosmetics'',
``jewelry'' --- and so are dropped from the headline panel even when
the analysis at hand does not depend on $KL$ at all. To check that
the templating-survival pattern is not an artefact of this filter, we
re-run the analysis on \textsc{clean}--without--$n_{\text{terms}}$
($8{,}499{,}455$ filings, $+591$k vs.\ headline). The pattern is
nearly identical for the $1$- through $1001$--$5000$ bins ($\leq 1$
percentage point of drift on either rate). The only meaningful change
is the $5000+$ bin: still-live rate falls from $49.2\%$ to $44.2\%$,
because the short commodity templates that the headline filter
excludes survive at corpus-average rates rather than the elevated
rates of the longer structured templates that dominate the
$n_{\text{terms}}$-filtered $5000+$ bin (e.g.\ ``athletic apparel
namely shirts pants jackets footwear hats and caps athletic
uniforms,'' $13{,}047$ copies, $61.6\%$ still-live).

\begin{figure}[h]
\centering
\includegraphics[width=0.95\linewidth]{results/templated_survival_no_nterms.png}
\caption{Robustness re-run of Figure~\ref{fig:templated-survival} on
\textsc{clean}--without--$n_{\text{terms}}\geq 3$. The lower bins are
visually identical; the $5000+$ bin loses $\sim\!5$pp because the
short commodity templates that the $n_{\text{terms}}$ filter
excludes survive at roughly corpus-average rates, not at the elevated
rates of the structured longer templates.}
\label{fig:templated-survival-nonterms}
\end{figure}

\paragraph{So why is the lower-left of the topomap so bad?}
The diagonal-clustering reading from the topomap survives this audit:
the worst-survival cell is firms whose filings cluster on the
diagonal at low $KL$. But the mechanism is not ``templated text gets
penalised.'' Quite the reverse --- explicitly templated text helps.
What appears to drive the lower-left corner is one-off, generic G/S
text that nobody else copies verbatim: descriptions written in the
class's modal vocabulary but unique enough to count as zero-template,
filed by firms without the procedural sophistication to use the
known-good phrasings their lawyers would otherwise re-use. Boilerplate,
when it exists, signals professional repetition; the bottom-left
cluster is filled with filers who lack the boilerplate.

\section{Both extremes outperform: a U-shape across industries}
\label{sec:u-shape}

The decile-rate plot referenced in Section~\ref{sec:survival} shows
that the relationship between $\Delta KL$ and survival is not
linear; it is U--shaped, with both the most-conservative-debut
filings (lowest decile of $\Delta KL$) and the most-innovative
filings (highest decile) outperforming the middle deciles. The
intuition is that both directions of distance from the modal
filing are rewarded: very-novel filings benefit from category
creation and the highest decile of $\Delta KL$ is dominated by
firms that the future imitated; very-conservative filings benefit
from being unambiguously defensible and easy for the examiner to
parse. Filings in the middle, neither novel enough to claim a new
category nor templated enough to coast on legal-clarity grounds,
fare the worst.

\paragraph{Quadratic regression.}
We test the curvature directly. For every industry's clean
filings between 1990 and 2020, we fit a logistic model with both
$z(\Delta KL)$ and $z(\Delta KL)^2$ as predictors, with year and
year-squared controls. Pooled across all 45 NICE classes and
$9.3$ million filings, the aggregate quadratic coefficient is
positive but small ($\beta_{(\Delta KL)^2}^z = +0.0002$ on 5y
survival; $+0.0013$ on 10y survival). The aggregate pooled
quadratic is not the primary evidence for the U-shape; the
per-industry evidence is. At the per-industry level the quadratic
is positive in $27$ of the $44$ industries we can fit.

\begin{table}[h]
\centering
\caption{Pooled quadratic logistic regression of survival outcomes
on $z$-scored $\Delta KL$ and its square. Year and year-squared
controls suppressed. A positive quadratic coefficient confirms
concave-up (``U-shape'').}
\label{tab:u-shape-pooled}
\begin{tabular}{lrrr}
\toprule
Outcome & N & Linear $\beta_{\Delta KL}^z$ & Quadratic $\beta_{(\Delta KL)^2}^z$ \\
\midrule
Completed registration & 9,255,511 & -0.0258 (0.0009) & +0.0053 (0.0002) \\
Renewed 5y registration & 9,255,511 & -0.0133 (0.0009) & +0.0002 (0.0002) \\
Renewed 10y registration & 6,362,813 & +0.0078 (0.0011) & +0.0013 (0.0002) \\
\bottomrule
\end{tabular}

\end{table}

\paragraph{The U-shape is present in most --- but not all ---
industries.} Fitting the quadratic per industry on the 5y-survival
outcome \emph{pooled across all years} ($n = 44$ industries with at
least $5{,}000$ clean filings), the quadratic coefficient on
$z(\Delta KL)^2$ is positive in $27$ of $44$ industries ($61\%$) and
individually significant at $p < 0.05$ in $21$ of those $44$ ($48\%$).
The U-shape dominates in consumer-goods classes (Household Utensils,
Furniture, Leather, Textiles, Common Metals, etc.), where heavy
templated filing by large incumbents and distinct novelty waves
coexist. Several large information-economy classes show the opposite
on the pooled-years fit, an apparent \emph{inverted-U}: Software \&
Electronics ($t = -18.2$), Insurance \& Financial ($t = -13.8$),
Scientific \& Tech Services ($t = -12.9$), Advertising \& Retail
($t = -10.7$), Telecommunications ($t = -10.6$).

\paragraph{Per-(industry $\times$ year) refinement.} The pooled-years
fit hides whether an industry is consistently $+U$, consistently
inverted-$U$, or shifting between regimes. We re-fit the quadratic
logit separately for every (industry, year) cell ($44$ industries
$\times$ $32$ years, requiring $\geq 500$ filings and convergent
logit; $1{,}267$ cells where both conditions hold), and classify each
industry by the majority-of-years rule with a margin of $2$ cells
(Table~\ref{tab:u-shape-classification}). On this refinement, the
pooled-years inverted-$U$ assignment to Software \& Electronics,
Insurance \& Financial, Scientific \& Tech, and Advertising \&
Retail is \emph{not} replicated at the year level: those industries
have year-by-year quadratic coefficients that flip sign and are
mostly indeterminate. Only one industry has a clean inverted-$U$ at
the year level (Telecommunications, NICE 038). Eleven additional
industries are indeterminate (no significant year-level curvature
in either direction). The $+U$ holds in $32$ of $44$ industries.

\begin{table}[h]
\centering
\small
\begin{tabular}{lrr}
\toprule
Statistic & count & \% of total \\
\midrule
\textbf{Industries (of 44 with sufficient data)} & & \\
  with significant $+U$ on $\Delta KL$ & 32 & 73\% \\
  with significant inverted-$U$ on $\Delta KL$ & 1 & 2\% \\
  indeterminate & 11 & 25\% \\
\midrule
\textbf{Industry $\times$ year cells (of 1,267 total)} & & \\
  significant $+U$ on $\Delta KL$ ($z > +1.96$) & 338 & 27\% \\
  significant inverted-$U$ ($z < -1.96$) & 81 & 6\% \\
  indeterminate quadratic & 848 & 67\% \\
  linear advantage $+\Delta KL$ ($z > +1.96$) & 197 & 16\% \\
  linear advantage $-\Delta KL$ ($z < -1.96$) & 484 & 38\% \\
  neither (linear $|z| \leq 1.96$) & 586 & 46\% \\
\bottomrule
\end{tabular}
\caption{$U$-shape classification of $44$ NICE classes individually
(majority-of-years rule, margin $= 2$) and of the $1{,}267$
$($industry, year$)$ cells where a logit fit converged on $\geq 500$
filings. The pooled $+U$ holds in $32$ of $44$ industries; only $1$
has a significant inverted-$U$ (Telecommunications); the remaining
$11$ are indeterminate, not inverted-$U$. Cell-level: $27\%$ have a
significant $+U$ on $\Delta KL$, $6\%$ a significant inverted-$U$,
$67\%$ indeterminate. Linear-direction counting at the cell level
shows a striking asymmetry: $38\%$ of cells have a significant
linear advantage to \emph{negative} $\Delta KL$ (the harvest
direction), versus $16\%$ with a significant linear advantage to
positive $\Delta KL$.}
\label{tab:u-shape-classification}
\end{table}

Table~\ref{tab:u-shape-industries}
reports the pooled-years ranking of the 15 industries with the
strongest U-curvature and the 5 with the strongest pooled-years
inverted-$U$, for reference and comparison with the year-level
classification above.

\begin{table}[h]
\centering
\caption{Per-industry quadratic coefficient on $z(\Delta KL)^2$
for 5y registration-renewal. Top 15 (most pronounced U) and
bottom 5. Asterisk = $|t| > 1.96$.}
\label{tab:u-shape-industries}
\resizebox{\textwidth}{!}{%
\begin{tabular}{lrrr}
\toprule
Industry (NICE) & N & Linear $\beta_{\Delta KL}^z$ & Quadratic $\beta_{(\Delta KL)^2}^z$ \\
\midrule
Materials Treatment (040) & 83,213 & +0.0575 & +0.0527* \\
Common Metals (006) & 81,525 & -0.0939 & +0.0422* \\
Pharmaceuticals (005) & 314,898 & +0.0127 & +0.0405* \\
Machines \& Power Tools (007) & 135,331 & -0.0201 & +0.0403* \\
Chemicals \& Industrial (001) & 105,461 & +0.0420 & +0.0375* \\
Legal \& Personal Services (045) & 114,693 & -0.0256 & +0.0307* \\
Rubber \& Insulators (017) & 43,458 & -0.0456 & +0.0264* \\
Household Utensils (021) & 172,345 & -0.1541 & +0.0256* \\
Leather \& Luggage (018) & 126,295 & -0.0902 & +0.0247* \\
Textiles \& Bedding (024) & 78,947 & -0.1246 & +0.0223* \\
Medical Devices (010) & 160,575 & -0.0560 & +0.0208* \\
Construction \& Repair (037) & 157,226 & -0.0235 & +0.0186* \\
Vehicles (012) & 107,119 & -0.0356 & +0.0179* \\
Restaurants \& Hotels (043) & 122,371 & -0.0151 & +0.0166* \\
Medical \& Personal Care (044) & 167,607 & +0.0084 & +0.0158* \\
\midrule
\multicolumn{4}{l}{\textit{...bottom 5 (most negative quadratic = inverted-U):}} \\
\midrule
Wine \& Spirits (033) & 59,728 & +0.0085 & -0.0072* \\
Trim, Lace \& Notions (026) & 33,395 & -0.1202 & -0.0159* \\
Telecommunications (038) & 153,225 & +0.0204 & -0.0218* \\
Carpets \& Wall Hangings (027) & 20,959 & -0.0334 & -0.0403* \\
Musical Instruments (015) & 10,609 & -0.0760 & -0.0547* \\
\bottomrule
\end{tabular}%
}

\end{table}

Figure~\ref{fig:u-shape} overlays each industry's per-decile
5y-renewal rate (thin grey lines) on the cross-industry mean
(red), making the consistent U-shape visible.

\begin{figure}[h]
\centering
\includegraphics[width=0.92\linewidth]{results/u_shape_curves.png}
\caption{Per-decile 5y registration-renewal rate by $\Delta KL$
decile. Each grey line is one of the 45 NICE industries; red is
the cross-industry mean. The U-shape is visible in nearly every
industry.}
\label{fig:u-shape}
\end{figure}

\paragraph{Industries ranked by U-shape strength.}
Industries differ in how cleanly the U appears.
Table~\ref{tab:u-shape-ranked} ranks all 44 industries with at
least 5{,}000 clean filings by the $t$-statistic on the quadratic
$\Delta KL^2$ coefficient. Pharmaceuticals (Class 005) tops the
list, followed by Materials Treatment (040), Advertising \&
Retail (035), Chemicals \& Industrial (001), and Machines \&
Power Tools (007). At the bottom of the ranking are
Telecommunications (038, $t=-9.8$, statistically inverted-U),
Carpets \& Wall Hangings (027), Musical Instruments (015),
Trim, Lace \& Notions (026), Wine \& Spirits (033), and Tobacco
\& Smokers' Articles (034) --- in each of these the curvature is
absent or slightly inverted.

\begin{table}[h]
\centering
\caption{Per-industry quadratic on 5y-survival, sorted by $t$ on
the quadratic coefficient (positive = U-shape, negative =
inverted-U).}
\label{tab:u-shape-ranked}
\resizebox{\textwidth}{!}{%
\begin{tabular}{lrrrr}
\toprule
Industry (NICE) & N & $\beta_{dKL}^z$ (linear) & $\beta_{dKL^2}^z$ (quadratic) & $t$ on quadratic \\
\midrule
\multicolumn{5}{l}{\textit{Industries with the strongest U-shape (top 15 by $t$ on $\beta_{dKL^2}$)}} \\
\midrule
Household Utensils (021) & 183,800 & -0.1586 & +0.0258 & +14.6 \\
Legal \& Personal Services (045) & 114,693 & -0.0256 & +0.0307 & +14.5 \\
Materials Treatment (040) & 87,017 & +0.0204 & +0.0234 & +12.3 \\
Medical Devices (010) & 174,201 & -0.0694 & +0.0183 & +10.9 \\
Restaurants \& Hotels (043) & 122,371 & -0.0151 & +0.0166 & +9.7 \\
Medical \& Personal Care (044) & 167,607 & +0.0084 & +0.0158 & +9.4 \\
Furniture (020) & 144,864 & -0.1177 & +0.0200 & +9.4 \\
Common Metals (006) & 89,730 & -0.1260 & +0.0188 & +8.7 \\
Leather \& Luggage (018) & 133,408 & -0.1136 & +0.0136 & +8.5 \\
Toys, Games \& Sporting Goods (028) & 292,436 & -0.0697 & +0.0105 & +8.2 \\
Textiles \& Bedding (024) & 84,770 & -0.1196 & +0.0183 & +7.8 \\
Cosmetics \& Cleaning (003) & 300,326 & -0.0388 & +0.0074 & +7.0 \\
Lighting, HVAC \& Appliances (011) & 171,839 & -0.0528 & +0.0105 & +6.2 \\
Rubber \& Insulators (017) & 48,186 & -0.0901 & +0.0169 & +5.3 \\
Chemicals \& Industrial (001) & 118,548 & +0.0147 & +0.0094 & +4.4 \\
\midrule
\multicolumn{5}{l}{\textit{Industries with the weakest or inverted U (bottom 5)}} \\
\midrule
Telecommunications (038) & 162,634 & -0.0261 & -0.0130 & -10.6 \\
Advertising \& Retail (035) & 924,501 & +0.0143 & -0.0048 & -10.7 \\
Scientific \& Tech Services (042) & 692,938 & +0.0168 & -0.0097 & -12.9 \\
Insurance \& Financial (036) & 374,634 & +0.0541 & -0.0122 & -13.8 \\
Software \& Electronics (009) & 1,146,078 & +0.0715 & -0.0107 & -18.2 \\
\bottomrule
\end{tabular}%
}

\end{table}

\textbf{Industries that lack a U have a stronger linear effect.}
Among the 29 industries with $t_{\text{quad}} \geq 1.96$, the
mean $|\beta_{\text{linear}}|$ is $0.048$. Among the 15 industries
without a significant U, the mean $|\beta_{\text{linear}}|$ is
$0.063$ --- noticeably higher. These are industries in which
$\Delta KL$ acts more like a linear (rather than U-shaped)
predictor. Telecommunications, the strongest inverted-U, has a
positive linear coefficient ($\beta_{\text{lin}} = +0.020$) on
top of the negative quadratic; Trim \& Notions, Carpets, and
Musical Instruments have negative linear coefficients large in
magnitude ($-0.04$ to $-0.12$). These industries either reward
novelty roughly linearly (in either direction) or reward
conventionality without the symmetric reward at the high-novelty
end.

\paragraph{Is the U driven by extreme-KL filings?}
A natural concern is that the U-shape is generated by a few very
extreme filings on either KL axis --- the most-similar-to-past
filings on the left, the most-imitated filings on the right ---
rather than by a smooth concave-up relationship across the
distribution. We test this by re-fitting the pooled quadratic on
the same outcome after dropping different combinations of the
extremes.

\begin{table}[h]
\centering
\caption{Pooled quadratic logit of 5y survival on $z(\Delta KL)$
and $z(\Delta KL)^2$, with year and year-squared controls,
re-fit after dropping different parts of the sample. The
quadratic coefficient on $\Delta KL^2$ is essentially unchanged
across drop choices --- the U-shape is a property of the central
distribution, not of the extreme tails.}
\label{tab:u-shape-drilldown}
\begin{tabular}{lrrrr}
\toprule
Sample & N & $\beta_{dKL}^z$ & $\beta_{dKL^2}^z$ & $t$ on quad. \\
\midrule
Full sample (none dropped) & 8,633,791 & +0.0046 & +0.0087 & +41.9 \\
Drop bottom 5\% of prospective KL (< 3.03) & 8,202,101 & +0.0102 & +0.0082 & +38.6 \\
Drop bottom 5\% of retrospective KL (< 3.03) & 8,202,101 & +0.0087 & +0.0085 & +39.9 \\
Drop bottom 5\% of both axes & 8,138,272 & +0.0117 & +0.0082 & +38.2 \\
Drop bottom + top 10\% of $\Delta KL$ (test extreme tails) & 6,907,035 & -0.0186 & +0.0343 & +48.7 \\
\bottomrule
\end{tabular}

\end{table}

The result is unambiguous: \textbf{the U-shape does not depend on
extreme-KL filings.} Dropping the bottom $5\%$ of prospective KL
(filings most similar to the past) leaves the quadratic
coefficient at $+0.0082$ ($t = +38.6$), down from $+0.0087$
($t = +41.9$) on the full sample. Dropping the bottom $5\%$ of
retrospective KL (most-similar-to-future filings) leaves it at
$+0.0085$ ($t = +39.9$). Dropping both: $+0.0082$ ($t = +38.2$).
Most strikingly, dropping the bottom and top $10\%$ of $\Delta KL$
\emph{itself} --- i.e., excising the two arms of the U --- yields
a residual quadratic coefficient of $+0.0343$ ($t = +48.7$),
\emph{four times larger} than the full-sample coefficient. The
curvature is intrinsic to the central $80\%$ of the $\Delta KL$
distribution. The U is not driven by the most-templated nor by
the most-novel filings; it is the relationship between novelty and
survival across the entire range.

\paragraph{Two-sided exponential fit.}
A separate test of the same shape is to split the data at
$\Delta KL = 0$ and ask whether the rate rises exponentially with
distance from zero on each side: $\log(\text{rate}) = a + b \cdot
|\Delta KL|$ separately for the negative-$\Delta KL$ and
positive-$\Delta KL$ halves. The pooled aggregate fits, on the
5y-survival outcome, give slopes $b_{\text{neg}} = +0.04$ and
$b_{\text{pos}} = +0.05$ on the bin-aggregated data, both with
$R^2 > 0.7$ across the deciles within each side. Both sides slope
positively in $|\Delta KL|$: each one-nat increase in distance
from zero (in either direction) is associated with roughly a
$4$-to-$5\%$ relative increase in the per-bin renewal rate.

\paragraph{What is consistent with this story?} A natural reach
is to say that the most-imitated novel filings are category
leaders the industry follows and the most-conventional filings
are templated marks that examiners and competitors leave alone.
But that reading credits both tails with prescient strategy ---
the right-arm filers ``saw'' the future their wording would
define, and the left-arm filers ``saw'' that the safest path
was to use the existing canonical wording. Without behavioural
data we cannot license either claim. A more parsimonious story
is that we observe the empirical fact (both extremes outperform
the middle) without needing to ascribe foresight to either tail.
The structural mechanism we \emph{can} document is examiner
heuristics --- the office-action evidence below ---  which gives
both tails an advantage at the registration gate without
requiring any awareness of what the future would do.

\paragraph{An office-action mechanism for the U.}
A more mechanical reading of the same shape comes from
likelihood-of-confusion. The status code 700 ``abandoned, no
response to office action'' is the modal post-examination
failure, and its share by $\Delta KL$ decile is itself U-shaped
in reverse (D1: $5.1\%$; middle deciles 3--5: $13.2$--$14.1\%$;
D10: $9.1\%$). Filings whose wording uses common, well-defined
mid-frequency vocabulary that happens to overlap with prior
registrations get examiner objections; filings with very narrowly
common vocabulary (the bottom decile) and filings with vocabulary
nobody else has used (the top decile) avoid these conflicts.
This mechanism does not require firms in either tail to have
strategic awareness of their own position --- it is a property
of the examiner's prior-art search, mediated by how distinctive
or how generic the goods/services description is. \textbf{We
flag this as the cleanest available mechanism precisely because
it does not ask us to credit the negative-$\Delta KL$ tail with
prescience that we would not also extend to the positive tail.}

\paragraph{We do not have a mechanism story commensurate with the
magnitude of the effect.}
We want to flag this directly. The U-shape on $\Delta KL$ is
empirically robust --- larger than the linear $\Delta KL$ effect we
report on gross margin, present in $27$ of $44$ industries, present
at the filing level even within firms that demonstrably abandoned
some marks and kept others alive, and present at the token level for
both extremes. But \emph{the magnitude of the survival lift at both
flux extremes is greater than any single mechanism we can articulate
explains}. A 5--10 percentage-point gap between flux-neutral and
flux-extreme filings, sustained across millions of filings and
dozens of industries, is not what we expected to find, and the
candidate mechanisms we can name (likelihood-of-confusion examination
artefacts, real-options harvest by incumbents, Schumpeterian creation
at the trend-leading edge) each predict only part of the curve and
none predicts the depth. We document the empirical pattern and
present what falsification we can; we do not claim a unified
mechanism story.

\paragraph{Falsification battery (summary).}
The U-shape is unusual enough to warrant a battery of robustness
checks. Table~\ref{tab:alternatives} reports re-fits of the pooled
survival quadratic on the candidate confounds in summary form;
Appendix~\ref{app:falsification} narrates each check in full. The
quadratic coefficient $\beta_{\Delta KL^2}^z$ stays strongly
positive on every survival sub-sample, with a single exception:
the firm-year gross-margin OLS, where the relationship is roughly
linear in $\Delta KL$ rather than U-shaped. Survival is U; margin
is linear.

\begin{table}[h]
\centering
\caption{Summary of the falsification battery for the survival
U-shape. The quadratic coefficient $\beta_{\Delta KL^2}^z$ stays
strongly positive everywhere except in the financial outcome.
Appendix~\ref{app:falsification} narrates each test.}
\label{tab:alternatives}
\resizebox{\textwidth}{!}{%
\begin{tabular}{lrrrr}
\toprule
Sub-sample & N & Outcome rate & Linear $\beta_{\Delta KL}^z$ & Quadratic $\beta_{\Delta KL^2}^z$ \\
\midrule
\multicolumn{5}{l}{\textit{Foreign-extension hypothesis}} \\
\midrule
Full sample & 600,000 & 39.2\% & +0.0043 (0.0028) & +0.0089 (0.0008) \\
US-domiciled only (foreign proxy = False) & 577,392 & 39.9\% & +0.0043 (0.0028) & +0.0078 (0.0008) \\
Foreign-domiciled only (foreign proxy = True) & 22,608 & 21.3\% & +0.0741 (0.0165) & +0.0100 (0.0038) \\
\midrule
\multicolumn{5}{l}{\textit{First-time vs recurring}} \\
\midrule
First-time filings only & 429,815 & 39.0\% & +0.0019 (0.0033) & +0.0080 (0.0009) \\
Recurring (2nd+) filings only & 170,185 & 39.6\% & +0.0084 (0.0051) & +0.0140 (0.0016) \\
\midrule
\multicolumn{5}{l}{\textit{Financial U-shape: firm-year gross margin}} \\
\midrule
Public firms (gross margin OLS) & 8,600 & --- & +0.0020 (0.0047) & -0.0022 (0.0012) \\
\bottomrule
\end{tabular}%
}

\end{table}

The ten checks (Table~\ref{tab:alternatives}, with the
firm-persistence, temporal-stability, competition-density, and
external-corpus tests added) are:
(1) Foreign-extension hypothesis (Madrid-Protocol marks could inflate
the low-$\Delta KL$ tail): \textbf{rejected}; US-only and
foreign-only sub-samples give nearly identical quadratics
($+0.0078$ vs.\ $+0.0100$).
(2) First-time vs.\ recurring filings: \textbf{both have the U},
recurring slightly stronger.
(3) Financial outcomes: \textbf{the U does not appear} in firm-year
gross margin (the relationship is linear); the U lives in survival.
(4) Survivor bias from large incumbents: \textbf{partial} support
for the left arm --- the bottom $\Delta KL$ decile has larger firm
filing-volumes than the middle.
(5) ``Liquid Death'' hypothesis (low-$\Delta KL$ might be novel
products in old vocabulary): \textbf{rejected} on text-length
splits.
(6) Filing-level vs.\ firm-level: \textbf{filing-level} --- the U
survives within firms that abandoned some marks and kept others
alive (Figure~\ref{fig:firm-persistence}).
(7) The blockchain / cryptocurrency / NFT / web3 wave as a possible
imitated-but-failed counterexample: \textbf{not a counterexample}
--- the crypto cohort outperforms on every survival margin.
(8) Temporal stability: \textbf{the U is present in most individual
years, not just in pooled data}. Per-year fits 1990--2021 ($n = 32$
cohorts) show the pooled quadratic is significantly positive in
$23 / 32$ years (Figure~\ref{fig:u-shape-temporal}), independently of
macro cycle. The pooled-across-years result is the average of mostly
year-by-year positive coefficients, not an aggregation artefact.
(9) Competition-density mediation: \textbf{not a confound, but a
mediator}. On 10{,}000 stratified Class-9 focal filings we count
near-competitors (cosine $\geq 0.7$ on g/s TF-IDF) within an
asymmetric hard window $[t-2y,\, t+5y]$ of the focal's filing date
(no decay weighting; this window is for the cosine-neighbour search
only and is distinct from the H=2y exponential-decay reference used
for the KL surprise itself --- a separate object). The pre-window of
2y captures established competitors at filing-time; the post-window
of 5y matches the focal filing's 5y survival horizon (i.e.\ peers
that could plausibly have shaped or been shaped by the focal in its
formative window). We count how many such peers survived 5y from
their own filing date. Flux extremes are in clusters with \emph{more}
surviving peers, not fewer; the flux-neutral middle has the fewest
live peers. Adding $\log(1 + \text{live\_competitors})$ as a control
absorbs $84\%$ of the $\Delta KL$ quadratic on survival
(Figure~\ref{fig:competition-density}).
The U is therefore mediated by niche viability: filings at flux
extremes sit in vocabulary clusters where peers also survive, and the
$\Delta KL$ signal is partly a proxy for which clusters can support real
businesses. The linear $z(\Delta KL)$ effect remains positive and
significant after the mediator is included ($z = +2.65$), so flux is
not fully redundant with cluster viability.
(10) External-corpus confound (Google Books n-grams): \textbf{partial
shared signal, but not redundant}. For 75 extreme-flux tokens (25 per
stratum) we pulled English-language frequency time series from Google
Books 1985--2021 and compared each token's per-year trademark-corpus
frequency to its per-year Google Books frequency.
Per-token Pearson correlation between trademark log-lift and Google
Books log-lift is $+0.47$ (Spearman $+0.39$, $n = 53$ tokens). About
$22\%$ of trademark-flux variance per token is shared with general
English-language usage; the remaining $78\%$ is corpus-specific. The
trademark signal is partly but not wholly a reflection of broader
linguistic drift.

\begin{figure}[h]
\centering
\includegraphics[width=\linewidth]{results/firm_persistence_test.png}
\caption{Filing-level 5y survival by $\Delta KL$ decile, separately
for three firm-cohorts: \emph{all-survived} (every filing made it),
\emph{mixed} (some filings were abandoned by the firm and some kept
alive), and \emph{never-survived} (no filings made it). The U is
preserved in the mixed cohort. The signal is filing-level, not
firm-level.}
\label{fig:firm-persistence}
\end{figure}

\begin{figure}[h]
\centering
\includegraphics[width=\linewidth]{results/u_shape_temporal_stability.png}
\caption{\textbf{Per-year fits 1990--2021.} Top: the quadratic
coefficient $\beta_{(\Delta KL)^2}$ on 5y survival fit separately for
each filing-year cohort, pooled across all 44 classes (black) and for
five representative classes. Pink bands are NBER recessions (2001,
2008--09, 2020). The pooled U-shape is significantly positive in
$23/32$ years, including across all three recessions; the per-year
result does not track macro cycles in a clean way and is not driven
by a small number of cohorts. Bottom: per-year mean 5y survival rate
for the same series, for context.}
\label{fig:u-shape-temporal}
\end{figure}

\begin{figure}[h]
\centering
\includegraphics[width=\linewidth]{results/competition_density.png}
\caption{\textbf{Live-competitor density mediates the survival U.}
Per $\Delta KL$ decile on $10{,}000$ stratified Class-9 focal filings.
Left panel: mean count of total and live (survived-5y) near-competitors
within $[t-2y,\, t+5y]$, cosine $\geq 0.7$ on g/s TF-IDF. The
flux-neutral middle has the \emph{fewest} live peers; both flux
extremes are in clusters with $\sim 20$--$30$ live peers each. Right
panel: focal 5y survival rate and the share of neighbors alive at
5y, by decile. The U-shape on focal survival mirrors the U on
neighbor survival, which is what the mediation regression formalises:
$\log(1 + \text{live\_competitors})$ absorbs $84\%$ of the $\Delta
KL$ quadratic on focal survival, but leaves the linear-positive
$\Delta KL$ effect intact ($z = +2.65$). The interpretation:
flux extremes mark vocabulary clusters where peers can also support
real businesses; the flux-neutral middle marks dying niches.}
\label{fig:competition-density}
\end{figure}

\paragraph{What is sensitive to methodology: per-industry rankings.
What is robust: the per-year U.} A reader who reproduces this analysis
from an earlier draft of the pipeline (with a flat 5-year reference
window rather than the production H=2y exponential decay) will get
visibly different per-industry quadratic coefficients --- enough that
the rank order of industries by U-strength can rearrange, and even
the sign of an individual industry's coefficient can change. We
flag this explicitly. The \emph{per-industry rank} is a sensitive
statistic that depends on how the reference distribution is built; we
would not stake a finding on the precise sign of any single industry's
quadratic. The \emph{per-year} robustness check
(Figure~\ref{fig:u-shape-temporal}) is the appropriate falsification
test for the headline U: if the pooled U were a few-cohort artefact,
it would not survive year-by-year. It does. The temporal-stability
evidence is what we lean on when claiming robustness; the
per-industry ranking is descriptive cross-sectional context, not the
primary evidence.

\paragraph{$|\Delta KL|$ as a single-number test.}
Collapsing the U-shape to one number, we regress each survival
outcome on $|\Delta KL|$ alone (with the same year and
log-n-terms controls). On a $500{,}000$ filing sample:

\begin{table}[h]
\centering
\caption{Single-predictor regression on $|\Delta KL|$ (the
absolute distance from $\Delta KL = 0$).}
\label{tab:abs-dkl}
\resizebox{\textwidth}{!}{%
\begin{tabular}{lrrrrrr}
\toprule
Outcome & N & $\beta_{|\Delta KL|}^z$ (logit) & $z$ & $p$ & McFadden pseudo-$R^2$ & $\Delta R^2$ (linear, vs.\ controls) \\
\midrule
Completed registration & 500,000 & +0.0072 (0.0030) & +2.4 & 1.6e-02 & 0.0284 & 0.001\,pp \\
Renewed 5y & 500,000 & +0.0276 (0.0030) & +9.2 & 0.0e+00 & 0.0086 & 0.018\,pp \\
Renewed 10y & 332,073 & +0.0495 (0.0035) & +14.1 & 0.0e+00 & 0.0103 & 0.060\,pp \\
\bottomrule
\end{tabular}%
}

\end{table}

The $z$-statistics on $|\Delta KL|$ are $+2.4$ for completed
registration, $+9.2$ for 5y survival, $+14.1$ for 10y survival,
all positive and the latter two highly significant. The
incremental linear-$R^2$ from $|\Delta KL|$ over the controls is
small (under $0.06$ percentage points) but real, consistent with
the U-shape we documented in the quadratic.

\paragraph{What else could be going on?}
A handful of further candidate stories that the data so far
are consistent with but cannot directly distinguish:

\begin{itemize}
  \item \emph{Narrative vs.\ analytical decision-making.} Firms
        in the middle of the $\Delta KL$ distribution may be
        chasing trends without committing either to the trend's
        canonical vocabulary or to a clearly differentiated
        wording. The data are consistent with this --- the
        worst-performing filings are exactly those that ``follow
        the herd without standing out from it'' --- but we have
        no behavioural data to confirm decision style.
  \item \emph{Pay-phones-after-cell-phones.} The middle deciles
        may include filings whose entrepreneurs are starting
        businesses on templates the industry is in the process
        of abandoning --- not low-novelty enough to coast on
        legal clarity, not high-novelty enough to claim a new
        category. The 1995 Beer \& Soft Drinks trough
        (Section~\ref{sec:u-shape}) is exactly such a cohort.
  \item \emph{Filing volume as a hedge.} Large incumbents file
        many marks per year and accept that some will fail;
        their templated bottom-decile filings benefit from the
        firm's resources for renewal and maintenance even when
        the underlying product fails. The mean firm filing-volume
        by decile evidence above is consistent.
  \item \emph{Examiner heuristics.} Examiners may grant clear
        templates (boilerplate is easy to evaluate) and clear
        novelties (no prior conflict to cite) at higher rates
        than middle-of-the-distribution filings, regardless of
        any underlying business outcome. The status-700 share
        evidence above is consistent.
\end{itemize}

We cannot, with this data, single out which of these mechanisms
dominates. What we can say is that the survival U-shape is
robust to the obvious confounds, that it does not appear in the
firm-year financial outcome (which is roughly linear in $\Delta
KL$), and that there is direct mechanical evidence
(office-action incidence) that middle-$\Delta KL$ filings face
more friction than either tail.

\paragraph{Did Beer \& Soft Drinks go in a failed direction
around 1995?} The notable trough in the Class~32 time-series
near 1995 (Figure~\ref{fig:ts_bev}) is consistent with this
story. Inspecting the lowest-$\Delta KL$ Class~32 filings of
1995 reveals a heavy concentration of \emph{New Age beverage}
brands: Clearly Canadian (\textsc{Tropical Tangerine},
\textsc{Clearly Sparkling Tea}), Koala Springs International
(\textsc{Aussie Slammer}, \textsc{Aussie Fruit Sparkler}), Royal
Crown's instant-drink mixes (\textsc{Yupii},
\textsc{Fruitii Instant Drink}), Empire State Beverage's
flavoured lemonades and teas (\textsc{Sahara Pink Lemonade},
\textsc{Rosewater Lemonade}, \textsc{Moonlight Tea}), and
Nestl\'e's exotic-fruit ginger drinks. This is the
mid-1990s clear-cola / sparkling-tea / iced-tea / exotic-fruit
wave that crested in $1993$--$95$ and was supplanted by the
energy-drink wave that began with Red Bull's U.S.\ launch in
$1997$ and accelerated through Monster ($2002$). The 1995 cohort
of Class~32 filings was, in retrospect, exactly the cohort that
the next decade's Class~32 industry stopped looking like. The
$\Delta KL$ measure correctly identifies it as a tired cohort.

\section{Per-industry novelty vs.\ financial outcomes}
\label{sec:industry-vs-finance}

The firm-level financial regression averages across industries. To
see whether the industry composition of the matched panel is doing
the work, we collapse to one observation per NICE class. For each
industry $c$, we compute the mean prospective KL, retrospective KL,
and $\Delta KL$ across all clean filings (1995--2020 cohorts) in
$c$, and the mean gross margin across all SEC firm-years for the
firms matched to filings in $c$. Figure~\ref{fig:industry-vs-finance}
plots all three together.

\begin{figure}[h]
\centering
\includegraphics[width=\linewidth]{results/industry_vs_finance.png}
\caption{Each dot is one of the 45 NICE international classes; dot
size is proportional to the number of matched public firms in the
class. Left panel: industry-mean prospective KL (surprise at filing
time) on the $x$-axis. Centre: industry-mean retrospective KL
(remained-novel-in-hindsight). Right: industry-mean $\Delta KL$
(imitated-innovator score). All three are plotted against
industry-mean gross margin on the $y$-axis. Dashed lines are OLS
fits, $\rho$ is the Pearson correlation across the 45 industries.}
\label{fig:industry-vs-finance}
\end{figure}

\paragraph{What was happening in the labelled industries during the
analysis window?}
The cross-industry correlation in
Figure~\ref{fig:industry-vs-finance} is positive but modest, and the
specific labelled industries help interpret what the measure picks
up.

\begin{itemize}
  \item \textbf{Tobacco \& Smokers' Articles (034)} sits at the
        far-right (highest $\Delta KL$) of the third panel and at
        the top of the cross-industry ranking
        (Figure~\ref{fig:ranks}). The $1995$--$2020$ window covers
        the entire e-cigarette / vape revolution: the Ruyan
        cigarette in 2003, NJOY's U.S.\ launch in 2007, the Juul
        explosion of $2015$--$18$, and the cannabis-vape category
        of $2018$--$20$ (the category leader by debut $\Delta KL$
        in our table is \textsc{Audio Technology of New York Inc.}
        a vape-cartridge filer, $+2.75$). The category went from a
        few smokeless-tobacco brands to several thousand
        device-and-cartridge SKUs in fifteen years.

  \item \textbf{Trim, Lace \& Notions (026)} is second-highest on
        $\Delta KL$. The catalyst is the rise of Asian
        wholesale hair-extension and synthetic-wig brands selling
        directly into U.S.\ retail; the top firms by debut
        $\Delta KL$ in Table~\ref{tab:industry-summary} are
        \textsc{Xuchang Best Trading} ($+1.93$), \textsc{Chubby
        Trading} ($+1.74$), and \textsc{Runway Premier
        International Hair} ($+1.62$). Almost no public-firm
        category had higher pivot rate during this window.

  \item \textbf{Medical \& Personal Care (044)} is in the middle
        of the $\Delta KL$ distribution but with high gross
        margin. The window encompasses the telehealth wave
        (Teladoc 2002, MDLive 2009, the COVID-era surge in
        $2020$+), the medical-spa explosion, and CBD/cannabis
        personal-care entries.

  \item \textbf{Telecommunications (038)} sits at the bottom of
        the $\Delta KL$ distribution despite having a comparatively
        high industry margin in our matched panel. The
        goods/services lexicon was already settled by 1995
        (\texttt{providing access to telecommunications networks},
        \texttt{wireless voice communication}); 4G and 5G
        deployments through the 2010s used the same vocabulary.
        The category was financially dynamic but linguistically
        stable.

  \item \textbf{Building Materials (019)} sits even lower: a
        mature, vocabulary-stable category through the entire
        window, despite genuine technical change in
        engineered-wood, low-carbon-cement, and
        cross-laminated-timber sub-segments that the
        goods/services lexicon does not appear to have absorbed.
\end{itemize}

\section{Cross-industry comparison}
\label{sec:per-industry-finance}
\label{sec:cross}

We process all 45 NICE international classes through the same pipeline.
Table~\ref{tab:industry-summary} reports per-class summary statistics;
Figure~\ref{fig:ranks} ranks industries by mean filing-level
$\Delta KL$; Figure~\ref{fig:industry-betas} reports the per-industry
logit coefficient of $\Delta KL$ on each survival outcome;
Table~\ref{tab:industry-fv} lists the highest-$\Delta KL$ mark per
industry, picked from the 2010--2020 cohort window where every class
has dense filing volume.

\begin{table}[h]
\centering
\caption{Per-industry summary statistics, restricted to filings from
1995--2020 with both 5-year reference windows populated.}
\label{tab:industry-summary}
\scriptsize
\resizebox{\textwidth}{!}{%
\begin{tabular}{lrrrrrrrr}
\toprule
Industry (NICE) & Filings & Clean & Owners & Years & $\Delta KL$ med. & $\Delta KL$ mean & 5y reg.\ rate & Reg.\ rate \\
\midrule
Tobacco \& Smokers' Articles (034) & 82,520 & 31,547 & 16,087 & 1995--2020 & +0.15 & +0.28 & 43\% & 48\% \\
Trim, Lace \& Notions (026) & 84,034 & 33,395 & 23,121 & 1995--2020 & +0.10 & +0.16 & 57\% & 64\% \\
Ropes, Tents \& Textiles (raw) (022) & 47,287 & 19,895 & 13,300 & 1995--2020 & +0.08 & +0.13 & 58\% & 68\% \\
Yarns \& Threads (023) & 15,210 & 2,889 & 1,756 & 2006--2020 & +0.10 & +0.12 & 66\% & 72\% \\
Household Utensils (021) & 403,933 & 172,345 & 109,053 & 1995--2020 & +0.06 & +0.10 & 56\% & 64\% \\
Carpets \& Wall Hangings (027) & 56,986 & 20,959 & 12,190 & 1995--2020 & +0.07 & +0.09 & 56\% & 64\% \\
Hand Tools (008) & 128,745 & 57,965 & 35,798 & 1995--2020 & +0.06 & +0.09 & 57\% & 68\% \\
Furniture (020) & 326,621 & 132,147 & 80,029 & 1995--2020 & +0.04 & +0.09 & 55\% & 64\% \\
Firearms \& Pyrotechnics (013) & 37,403 & 11,844 & 6,016 & 1995--2020 & +0.08 & +0.09 & 53\% & 66\% \\
Lighting, HVAC \& Appliances (011) & 333,507 & 160,109 & 90,464 & 1995--2020 & +0.04 & +0.07 & 59\% & 68\% \\
Textiles \& Bedding (024) & 173,759 & 78,947 & 47,704 & 1995--2020 & +0.05 & +0.07 & 54\% & 62\% \\
Wine \& Spirits (033) & 222,138 & 59,728 & 28,902 & 1995--2020 & +0.03 & +0.05 & 44\% & 52\% \\
Jewelry \& Watches (014) & 234,908 & 91,669 & 59,847 & 1995--2020 & +0.02 & +0.04 & 54\% & 61\% \\
Restaurants \& Hotels (043) & 285,610 & 122,371 & 74,815 & 2003--2020 & +0.02 & +0.03 & 50\% & 60\% \\
Vehicles (012) & 242,485 & 107,119 & 56,375 & 1995--2020 & +0.00 & +0.03 & 54\% & 65\% \\
Leather \& Luggage (018) & 249,052 & 126,295 & 80,619 & 1995--2020 & -0.01 & +0.02 & 54\% & 61\% \\
Medical Devices (010) & 302,192 & 160,575 & 74,616 & 1995--2020 & -0.02 & +0.02 & 51\% & 61\% \\
Cosmetics \& Cleaning (003) & 564,014 & 278,140 & 130,889 & 1995--2020 & +0.02 & +0.01 & 49\% & 56\% \\
Medical \& Personal Care (044) & 312,270 & 167,607 & 97,110 & 2003--2020 & -0.00 & +0.01 & 53\% & 63\% \\
Common Metals (006) & 175,437 & 81,525 & 46,824 & 1995--2020 & +0.01 & +0.01 & 54\% & 70\% \\
Toys, Games \& Sporting Goods (028) & 583,108 & 258,260 & 131,810 & 1995--2020 & -0.02 & +0.01 & 51\% & 58\% \\
Legal \& Personal Services (045) & 203,476 & 114,693 & 74,406 & 2003--2020 & -0.01 & +0.00 & 54\% & 63\% \\
Education \& Entertainment (041) & 1,344,519 & 815,077 & 471,515 & 1995--2020 & -0.01 & -0.00 & 52\% & 62\% \\
Rubber \& Insulators (017) & 88,430 & 43,458 & 22,524 & 1995--2020 & -0.00 & -0.00 & 55\% & 71\% \\
Lubricants \& Fuels (004) & 78,345 & 30,267 & 16,546 & 1995--2020 & -0.02 & -0.00 & 49\% & 61\% \\
Materials Treatment (040) & 145,590 & 83,213 & 52,597 & 1995--2020 & -0.02 & -0.01 & 54\% & 66\% \\
Musical Instruments (015) & 31,792 & 10,609 & 6,531 & 1995--2020 & -0.00 & -0.01 & 54\% & 68\% \\
Software \& Electronics (009) & 1,807,636 & 1,054,572 & 521,819 & 1995--2020 & -0.02 & -0.01 & 50\% & 58\% \\
Construction \& Repair (037) & 284,361 & 157,226 & 93,653 & 1995--2020 & -0.00 & -0.01 & 56\% & 69\% \\
Machines \& Power Tools (007) & 286,930 & 135,331 & 69,563 & 1995--2020 & -0.03 & -0.01 & 57\% & 72\% \\
Advertising \& Retail (035) & 1,418,345 & 899,682 & 528,027 & 1995--2020 & -0.03 & -0.02 & 52\% & 61\% \\
Scientific \& Tech Services (042) & 1,148,416 & 635,128 & 352,645 & 1995--2020 & +0.00 & -0.02 & 52\% & 60\% \\
Pharmaceuticals (005) & 600,473 & 314,898 & 118,037 & 1995--2020 & -0.02 & -0.02 & 42\% & 49\% \\
Coffee, Tea, Bakery \& Spices (030) & 428,854 & 173,569 & 89,747 & 1995--2020 & -0.02 & -0.02 & 47\% & 56\% \\
Paper, Print \& Office (016) & 712,163 & 367,643 & 205,093 & 1995--2020 & -0.02 & -0.02 & 48\% & 58\% \\
Chemicals \& Industrial (001) & 218,420 & 105,461 & 44,697 & 1995--2020 & -0.02 & -0.02 & 51\% & 66\% \\
Clothing, Footwear \& Headwear (025) & 1,265,966 & 561,490 & 361,440 & 1995--2020 & -0.05 & -0.03 & 47\% & 53\% \\
Beer \& Soft Drinks (032) & 206,304 & 87,849 & 43,714 & 1995--2020 & -0.03 & -0.03 & 41\% & 48\% \\
Transport \& Logistics (039) & 187,535 & 105,370 & 59,203 & 1995--2020 & -0.02 & -0.03 & 53\% & 65\% \\
Meat, Dairy \& Prepared Foods (029) & 253,936 & 114,185 & 55,196 & 1995--2020 & -0.04 & -0.04 & 46\% & 56\% \\
Building Materials (019) & 122,596 & 60,289 & 30,109 & 1995--2020 & -0.02 & -0.04 & 52\% & 66\% \\
Agriculture \& Live Animals (031) & 136,656 & 52,623 & 26,983 & 1995--2020 & -0.04 & -0.04 & 48\% & 59\% \\
Paints \& Coatings (002) & 64,892 & 28,286 & 14,228 & 1995--2020 & -0.05 & -0.04 & 53\% & 69\% \\
Insurance \& Financial (036) & 602,901 & 354,316 & 174,543 & 1995--2020 & -0.05 & -0.05 & 51\% & 63\% \\
Telecommunications (038) & 230,267 & 153,225 & 78,422 & 1995--2020 & -0.04 & -0.06 & 47\% & 53\% \\
\bottomrule
\end{tabular}%
}

\end{table}

\begin{figure}[h]
\centering
\includegraphics[height=0.92\textheight]{results/industry_ranks.png}
\caption{Industries ranked by mean filing-level $\Delta KL$ across
1995--2020 cohorts. Green: industries whose typical filing was
prospectively novel and retrospectively obvious (the future imitated
the past). Red: industries whose typical filing was the opposite.
Tobacco \& Smokers' Articles tops the list, driven by the 2014--2020
e-cigarette / vape vocabulary explosion. Telecommunications, Insurance
\& Financial, and Paints \& Coatings sit at the bottom: their
goods-and-services lexicon was already settled before 1995 and has
been stable since.}
\label{fig:ranks}
\end{figure}

\begin{figure}[p]
\centering
\includegraphics[width=\linewidth]{results/industry_survival_betas.png}
\caption{Per-industry effect of $\Delta KL$ on survival outcomes
(logistic-regression $z$-score coefficient with 95\% CI). Industries
are ordered by mean $\Delta KL$ (the same order as
Figure~\ref{fig:ranks}). \textbf{The sign of the innovation reward
varies by industry.} Some industries (e.g.\ Common Metals, Hand Tools,
Paints \& Coatings) actually \emph{reward} novel filings at the
registration gate, presumably because in those industries a novel
goods/services description can describe a real new use the examiner
recognises as patentable activity. Most industries match the software
pattern: novelty penalised at registration, neutral or favourable at
renewal.}
\label{fig:industry-betas}
\end{figure}

\begin{table}[h]
\centering
\caption{Top-1 highest-$\Delta KL$ mark per industry, 2010--2020
cohort. The 2014--2020 cannabis/CBD/blockchain wave appears in nearly
every consumer-goods class.}
\label{tab:industry-fv}
\scriptsize
\begin{tabular}{>{\raggedright\arraybackslash}p{4.0cm}rr>{\raggedright\arraybackslash}p{4.5cm}>{\raggedright\arraybackslash}p{4.5cm}}
\toprule
Industry (NICE) & Year & $\Delta KL$ & Mark & Owner at filing \\
\midrule
Tobacco \& Smokers' Articles (034) & 2016 & +4.51 & DAISY & Green Tank Technologies Corp. \\
Agriculture \& Live Animals (031) & 2018 & +4.20 & UNITED AMERICAN HEMP & United American Hemp LLC \\
Cosmetics \& Cleaning (003) & 2018 & +4.05 & LEEF & PALEO PAW CORP. \\
Chemicals \& Industrial (001) & 2017 & +4.03 & REEFERTILIZER & Reefertilizer Inc. \\
Coffee, Tea, Bakery \& Spices (030) & 2012 & +3.93 & RED \& WHITE & Forefront Enterprises, LLC \\
Household Utensils (021) & 2017 & +3.59 & IREPELLER & 10093939 Canada Corp. \\
Beer \& Soft Drinks (032) & 2018 & +3.56 & KANNABIER & Advanced Creative Intelligent  \\
Wine \& Spirits (033) & 2014 & +3.49 & CHATEAU DE CAYX & Boutique du Château de Cayx \\
Ropes, Tents \& Textiles (raw) (022) & 2013 & +3.33 & SUNAMI & TENSILE SHADE PRODUCTS, LLC \\
Scientific \& Tech Services (042) & 2014 & +3.30 & ZENPOOL & Geniuses at Work Corporation \\
Toys, Games \& Sporting Goods (028) & 2012 & +3.20 & BETCLOUD & BALLY GAMING, INC. \\
Materials Treatment (040) & 2018 & +3.17 & VERITAS | FARMS & 271 Lake Davis Holdings LLC \\
Meat, Dairy \& Prepared Foods (029) & 2016 & +3.15 & PHIVIDA & PHIVIDA HOLDINGS INC. \\
Jewelry \& Watches (014) & 2013 & +3.11 & OFF THE CLOSET & Rahhal, Nour \\
Software \& Electronics (009) & 2018 & +3.09 &  & Unify Earth UEX \\
Furniture (020) & 2016 & +3.07 & GO PET CLUB & Pel Wholesale, Inc. \\
Paper, Print \& Office (016) & 2020 & +3.05 & KAYLEEMAY LLC & Fall, Cayle M \\
Clothing, Footwear \& Headwear (025) & 2010 & +3.03 & CHARMANT FEMME & 2BB UNLIMITED, INC. \\
Medical Devices (010) & 2020 & +3.03 & ELASTIMASK & Bedford Industries, Inc. \\
Education \& Entertainment (041) & 2017 & +3.01 & ACADEMY & Kingsland School of Blockchain \\
Trim, Lace \& Notions (026) & 2015 & +2.99 & WOWQUEEN & XUE YONG BO \\
Advertising \& Retail (035) & 2018 & +2.85 & OPENSEA & Ozone Networks, Inc. \\
Lubricants \& Fuels (004) & 2017 & +2.82 &  & Mary Jane's Medicinals, LLC \\
Pharmaceuticals (005) & 2017 & +2.74 & CANINES BEST DEFENSE CBD & Nobrega, Kathryn \\
Insurance \& Financial (036) & 2017 & +2.71 & MOSAIC TOKEN & MOSAIC RESEARCH LIMITED \\
Medical \& Personal Care (044) & 2018 & +2.64 & FARMING A HEALTHIER FUTURE & Southern Hemp, LLC \\
Carpets \& Wall Hangings (027) & 2016 & +2.63 & LOVEMAT & Wuhan Xuanyuanlong Science and \\
Transport \& Logistics (039) & 2017 & +2.63 & MARIJUANA EXPRESS & Marijuana Express LLC \\
Lighting, HVAC \& Appliances (011) & 2010 & +2.42 & KESSIL & DiCon Fiberoptics, Inc. \\
Vehicles (012) & 2013 & +2.34 & PRECISION DRONE WWW.PRECISIONDRO & Precision Drone, LLC \\
Leather \& Luggage (018) & 2018 & +2.31 & NOMADZ & Ambit Global \\
Textiles \& Bedding (024) & 2017 & +2.27 & BESTDROP & Zhiyong Xie \\
Legal \& Personal Services (045) & 2011 & +2.26 & DIAMONDS IN THE RUFF & CINDY ROSS \\
Restaurants \& Hotels (043) & 2013 & +2.20 & GREEN FEET BREWING & Green Feet Brewing \\
Hand Tools (008) & 2016 & +2.02 & THE NAZOR & Nazor, LLC \\
Firearms \& Pyrotechnics (013) & 2010 & +1.94 & ALPHA RAIL & TROY INDUSTRIES, INC. \\
Machines \& Power Tools (007) & 2017 & +1.86 & COUNTERMATE & IMPAK Corporation \\
Musical Instruments (015) & 2012 & +1.75 & ERGODOCK & Hickey, Bruce \\
Common Metals (006) & 2017 & +1.70 & UP YOUR CASTER & CASTER CONNECTION, LLC \\
Yarns \& Threads (023) & 2011 & +1.69 & BERROCO NANUK & Berroco, Inc. \\
Rubber \& Insulators (017) & 2014 & +1.69 & B3D & Brandon Bell \\
Construction \& Repair (037) & 2010 & +1.52 & WILKS MASONRY & Wilks Masonry Corporation \\
Building Materials (019) & 2012 & +1.51 & FREE FORM YOGA FLOOR & Isabelle J. Barger \\
Telecommunications (038) & 2017 & +1.44 & NCRYPT & NCHAIN HOLDINGS LTD. \\
Paints \& Coatings (002) & 2016 & +1.32 & TRIPLE BEST & Sike Wu \\
\bottomrule
\end{tabular}

\end{table}

\paragraph{Three patterns visible across the 45 classes.}

First, the cannabis/CBD wave of 2014--2020 dominates the top-by-class
face-validity table: cosmetics (\textsc{Leef Paleo Paw}), agriculture
(\textsc{United American Hemp}), beverages (\textsc{Kannabier}),
restaurants, materials, even insurance products are all topped by
hemp/CBD/cannabis-themed marks. The blockchain wave (2014--2018)
sweeps Class 35 (\textsc{OpenSea}), Class 41
(\textsc{Kingsland School of Blockchain}), and Class 36
(\textsc{Mosaic Token}). The vape/e-cigarette wave is concentrated in
Class 34 (\textsc{Daisy}), which leads the entire ranking.

Second, the most-and-least innovative industries are surprising.
Tobacco \& Smokers' Articles ($+0.28$ mean $\Delta KL$) and the
textile-input classes (Trim \& Notions, Yarns \& Threads, Ropes)
sit at the top --- the textile classes because they have low baseline
filing volume and any new fibre or smart-textile entrant disrupts the
distribution. The bottom of the list is dominated by mature regulated
industries: Telecommunications, Insurance \& Financial, Building
Materials, and Agriculture. Neither end is what one would have
predicted in advance.

Third, the per-industry survival betas
(Figure~\ref{fig:industry-betas}) show that the
``novelty-is-penalised-at-the-registration-gate'' result reported in
Section~\ref{sec:survival} is \emph{not} universal. In some industries
(notably Common Metals, Hand Tools, Paints \& Coatings) the
coefficient flips sign: novel filings in those industries are
\emph{more} likely to register, not less. A plausible reading is that
in industries where category-creation is rare and meaningful, an
examiner reading a novel goods/services description perceives a real
new use rather than vague language. Software (and pharmaceuticals,
chemicals) follow the opposite pattern: novel filings face slower
examination, more office actions, and lower registration rates,
consistent with examiners pushing back on speculative or boilerplate
broadenings of established categories.

\paragraph{Do industries in flux earn supernumerary returns?}
Across the $44$ NICE classes with adequate data, we computed four
flux measures per class (mean $KL_\text{pros}$, mean $\Delta KL$,
std $\Delta KL$, mean $|\Delta KL|$) and four financial outcomes
(median gross margin, median operating margin, median 4y excess return
vs.\ S\&P 500, median 5y excess return). Excess return is the
firm-mean over firm-years in our matched panel,
$\text{ret}_{4y} - \text{sp500\_ret}_{4y}$, aggregated as a
class-level median to suppress single-firm outliers. We use median
rather than mean because a handful of small public firms in
mechanical-tool industries showed outsized returns that distorted
class means by orders of magnitude.

\textbf{The answer is no.} The cross-industry Pearson correlation
between flux intensity (mean $|\Delta KL|$) and median 4y excess
return is $-0.09$ (Spearman $+0.01$, $n = 42$ industries with
returns data). For 5y excess return: Pearson $-0.05$, Spearman
$+0.02$. The mean prospective KL is similarly uncorrelated with
returns. Industries in flux do not generate excess returns over the
S\&P 500 in this corpus.

The cross-industry pattern that \emph{is} present is on margins, not
on returns: industries with higher flux have systematically
\emph{lower} median gross margins (Pearson $-0.22$, Spearman
$-0.20$, for mean $|\Delta KL|$). High-flux industries are less
profitable per dollar of revenue on average, consistent with such
industries being more competitive (entry-heavy, churn-heavy). The
market appears to price in the lower margins; investors are not
compensated with excess returns for taking on flux risk at the
industry level. Table~\ref{tab:industry-flux-vs-returns} reports
the cross-industry correlation matrix;
Figure~\ref{fig:industry-flux-vs-returns} plots the four key
scatter pairs.

\begin{table}[h]
\centering
\caption{Cross-industry correlations between flux measures (industry
mean) and financial outcomes (industry median across firms), $n = 44$
NICE classes (returns columns drop two classes with insufficient
SEC-listed firms).}
\label{tab:industry-flux-vs-returns}
\small
\begin{tabular}{lrrr}
\toprule
Flux measure (industry mean) & Outcome & Pearson $r$ & $n$ \\
\midrule
mean $KL_{\text{pros}}$ & median gross margin & -0.055 & 44 \\
mean $KL_{\text{pros}}$ & median operating margin & +0.138 & 44 \\
mean $KL_{\text{pros}}$ & median excess 4y return & +0.015 & 42 \\
mean $KL_{\text{pros}}$ & median excess 5y return & -0.060 & 42 \\
mean $\Delta KL$ & median gross margin & -0.214 & 44 \\
mean $\Delta KL$ & median excess 4y return & -0.022 & 42 \\
mean $\Delta KL$ & median excess 5y return & +0.005 & 42 \\
std $\Delta KL$ & median gross margin & -0.179 & 44 \\
std $\Delta KL$ & median excess 4y return & -0.111 & 42 \\
std $\Delta KL$ & median excess 5y return & -0.078 & 42 \\
mean $|\Delta KL|$ & median gross margin & -0.216 & 44 \\
mean $|\Delta KL|$ & median excess 4y return & -0.088 & 42 \\
mean $|\Delta KL|$ & median excess 5y return & -0.049 & 42 \\
\bottomrule
\end{tabular}
\end{table}

\begin{figure}[h]
\centering
\includegraphics[width=\linewidth]{results/industry_flux_vs_returns.png}
\caption{Per-industry scatter of flux measures (mean prospective KL
top row; mean $|\Delta KL|$ bottom row) versus median gross margin
(left) and median 4y excess return vs.\ S\&P 500 (right). Point size
proportional to number of SEC-listed firms with primary class in the
focal industry. The three- and two-digit labels are NICE class
numbers for the top-3 and bottom-3 industries on each $y$-axis.
\textbf{Flux intensity predicts lower margins, not higher returns.}}
\label{fig:industry-flux-vs-returns}
\end{figure}

\textbf{Reconciliation with the firm-level result.} The firm-level
U-shape on $\Delta KL$ (Section~\ref{sec:u-shape}) says that
\emph{within} an industry, firms with flux-extreme vocabulary
outperform the flux-neutral middle on survival. The industry-level
null on excess returns is consistent with this: a high-flux
industry contains both flux-extreme winners and flux-neutral losers,
and after the within-industry sorting plays out the industry-mean
return is approximately flat. The flux signal is a within-industry
selection statistic, not an across-industry valuation premium.
``Invest in the high-flux industry'' is not a strategy; ``invest in
the flux-extreme firms within an industry'' might be.

\section{What is the construct, really?}
\label{sec:construct}

We have called $\Delta KL$ an ``innovation'' score throughout, but
the empirical pattern this paper documents is consistent with a
plainer description.

The two-way pattern --- both extremes outperform the middle, both
on survival
(Section~\ref{sec:u-shape}) and on financial variance
(Section~\ref{sec:financials}) --- is consistent with the
construct measuring \emph{how unconventional the firm's bets are},
in either direction. Filings with very high $\Delta KL$ are
unconventional in the direction of category creation; filings with
very low $\Delta KL$ are unconventional in the direction of
templated boilerplate that examiners and competitors leave alone.
Filings with $\Delta KL \approx 0$ are doing what the rest of the
industry was already doing, and unsurprisingly that is what the
data predict for them: average outcomes, average failure rate,
average margins, average returns.

If we step back from the language of ``innovation'' and ask only
what the construct actually correlates with:
\begin{itemize}
  \item It is essentially uncorrelated with patent counts ($r =
        +0.01$ in the matched panel) and with Crunchbase total
        funding ($r = 0.00$); modestly correlated with R\&D-to-revenue
        ratio ($r = +0.13$).
  \item It picks out qualitatively different firms than those
        proxies pick out (Section~\ref{sec:financials}'s patent
        comparison).
  \item It predicts firm-year stock returns above the S\&P~500 on
        a horizon of two-to-five years
        (Section~\ref{sec:returns}); on the firm-debut variant the
        peak is at four years and the magnitude is approximately
        $+28$ percentage points per standard deviation. The peak
        is several years out, not immediate.
  \item It also predicts within-firm gross-margin variance
        (Section~\ref{sec:financials}). Firms whose recent
        $\Delta KL$ is high earn higher mean margins \emph{and}
        spend more years far from those means.
\end{itemize}

These properties are consistent with a simple summary: \textbf{the
construct measures the smartness, and the riskiness, of the bets a
firm makes when it names its products}. Average bets ($\Delta KL$
near zero) earn average outcomes. Bets that depart from convention
in either direction get rewarded or punished above-average; the
rewards take several years to arrive, and the variance of outcomes
is correspondingly larger. Firms with novel bets that the rest of
the industry then imitates --- the high-$\Delta KL$ filings ---
are the ones that capture the rewards in the right tail; firms
with novel bets nobody adopts pay the cost in the left tail.
Boilerplate filings ($\Delta KL$ very negative) are a different
kind of bet: a bet that the existing language is good enough,
that defensibility is more important than category creation. The
data say that bet also works, but on a flatter return distribution.

This is, at the construct level, an evidentiary claim that
\emph{ideas matter}, in a fairly literal sense: a firm whose
trademark filings depart from the industry-mean vocabulary in any
direction has measurably different financial dynamics than one
whose filings hew to the centre. The five-year half-life of the
return effect suggests that good ideas need time to ferment;
imitators have to recognise the new category and copy it before
the originator's claim to it becomes valuable. The sub-additivity
between this measure and patent counts in the $R^2$ decomposition
(Section~\ref{sec:financials}) means we are not just re-discovering
patents in a different guise. Whatever this construct is measuring,
it is measuring something separately predictive of firm
performance.

\section{Discussion}

\paragraph{Two response curves, two theories.}
The empirical pattern we document is best read as two
saturating-exponential response curves meeting at $\Delta KL \approx
0$. On the innovation side ($\Delta KL > 0$) the curve is the
empirical content of optimal distinctiveness (Deephouse 1999;
Zhao et al.\ 2017) translated from \emph{level} of distinctiveness to
\emph{directional} vocabulary flux: survival rises with prospective
novelty, concave-down, with diminishing returns to additional
$\Delta KL$ and no inverted-$U$ inflection at the extreme. What we
do not find is the strict Deephouse prediction that survival turns
down at high distinctiveness; what we do find is the underlying
empirical claim of the optimal-distinctiveness literature ---
distinctiveness pays with diminishing returns --- in an explicit
functional form ($a = 0.995$, $b = -0.55$, $c = 0.063$ on the debut
panel; $R^2 = 0.499$ across $20$ bins). On the harvest side
($\Delta KL < 0$) the same functional form fits much more tightly
($R^2 = 0.978$) and tells a complementary story: survival rises
\emph{faster} with $|\Delta KL|$, saturates closer to the origin,
and reaches an asymptote of $\sim 0.57$. The interpretation is
\emph{increasing returns to late-stage imitation}: the deeper a
filing sits into the receding vocabulary of its category, the more
durable its mark turns out to be. The two extractive strategies
are recognisable in the token data: the negative extreme picks out
clearly sunsetting industries (typesetting in the early 1990s as
desktop publishing arrived, the pre-platform hospitality bundle in
the early 2000s as Expedia disintermediated, e-cigarette
boilerplate of 2014--2015 just before the FDA Deeming Rule); the
positive extreme picks out the legalisation+platform-shift waves
(cannabis, crypto, NFTs, vape, plant-based, EV). \emph{Both
strategies pay}, but they pay through opposite functional forms:
diminishing returns to novelty on one side, increasing returns to
imitation on the other.

\paragraph{Mechanism is not separable.}
Whether the survival lift at negative $\Delta KL$ reflects
intelligent harvesting (firms recognising an industry transition and
extracting from the legacy side), passive late-occupation (firms whose
business naturally uses the receding vocabulary because they are
doing the legacy thing), or selection on an unobservable correlate
(firms that survive long enough to file legacy-vocabulary trademarks
have already proved durable) is not separable in the trademark
record. The unmistakable empirical pattern is that flux-neutral
filings die at the highest rate in the corpus, in $27$ of $44$
industries, at the token level, and \emph{even within the same firm}:
the firm-persistence test
(Section~\ref{sec:u-shape}, falsification \#6) confirms the U
survives within firms that demonstrably abandoned some trademarks
but not all. The signal is filing-level, not firm-level.

\paragraph{The ``arrows in the back'' resolution.}
At the trademark gate, novelty is \emph{rewarded}, not penalised.
The most prospectively surprising filings clear examination at higher
rates and persist longer through the renewal cycle. The mousetrap
pattern dominates the survival outcomes \emph{at the filing level}.
The pioneer-ladder analysis (Section~\ref{sec:pioneer-ladder}) gives
a more nuanced firm-level view: the very-first $10$--$50$ adopters
of category vocabulary do badly (the arrows), the fast-follower band
at ranks $50$--$1{,}000$ does best (the mousetrap-builders), and the
mass late wave is uniformly worst. Lieberman \& Montgomery (1988)
and Suarez \& Lanzolla (2007) anticipated this empirical pattern in
the management literature; the trademark corpus provides large-$n$
evidence at the level of business-description vocabulary.

The arrows-in-the-back signal does, however, show up financially. The
matched panel of $1{,}597$ public firms reveals a clear pioneer-bet
structure: high-novelty firms simultaneously earn higher mean gross
margins (Section~\ref{sec:financials}) \emph{and} display wider
within-firm gross-margin standard deviations (Table~\ref{tab:variance}).
Some pioneers succeed; others do not; the survivors carry the average.
At the trademark registrar's office novelty is a positive signal, but
in the marketplace it raises the variance of outcomes in addition to
the mean.

To make the underlying mechanics concrete,
Table~\ref{tab:tokens} reports a small set of marks together with the
top dictionary tokens that drove their prospective and retrospective
surprise. Filings that score high on prospective surprise share a common
property: they contain tokens (\texttt{generative artificial},
\texttt{seltzer}, \texttt{cannabis}) that are common in the future
corpus but rare in the look-behind window. Filings that score low ---
or even negative --- on $\Delta KL$ contain tokens whose temporal
distribution is the opposite.

\begin{landscape}
\begin{table}[p]
\centering
\caption{Selected filings with per-token KL contributions to prospective
and retrospective surprise. \textbf{Surprise is an interaction across
tokens, not a property of any single word.} A filing's KL score is
$\sum_v P(v) \log[P(v)/Q(v)]$, summed across every token $v$ that
appears in the filing's goods/services description. The token
\texttt{software} is by itself common in Class 9, but its
\emph{combination} with rare neighbors --- \texttt{generative
artificial}, \texttt{conversational agents}, \texttt{cryptographic
keys} --- can give a filing high prospective surprise relative to the
look-behind window's mix of tokens. Tokens shown below are the top five
by per-token contribution $P(v) \log[P(v)/Q(v)]$ to each direction's
KL; common tokens like \texttt{software} can appear in the prospective
list when their multinomial weight in the filing is much larger than
their weight in the look-behind reference.}
\label{tab:tokens}
\scriptsize
\begin{tabular}{>{\raggedright\arraybackslash}p{3.4cm}>{\raggedright\arraybackslash}p{3.8cm}rrr >{\raggedright\arraybackslash}p{6.0cm} >{\raggedright\arraybackslash}p{6.0cm}}
\toprule
Mark / Year & Owner at filing & $KL_{\text{pros}}$ & $KL_{\text{retr}}$ & $\Delta KL$ & Top tokens driving prospective surprise & Top tokens driving retrospective surprise \\
\midrule
OPENAI (2016) & OpenAI, Inc. & 7.65 & 5.56 & +2.09 & \texttt{software artificial}, \texttt{artificial intelligence}, \texttt{intelligence}, \texttt{artificial}, \texttt{software} & \texttt{software artificial}, \texttt{artificial intelligence}, \texttt{artificial}, \texttt{intelligence}, \texttt{software} \\
CLAUDE (2023) & Anthropic, PBC & 7.12 & 5.06 & +2.06 & \texttt{model performing}, \texttt{document question-answering}, \texttt{tasks natural}, \texttt{generative text}, \texttt{performing generative} & \texttt{tasks}, \texttt{natural}, \texttt{model performing}, \texttt{question-answering simulating}, \texttt{document question-answering} \\
INSTAGRAM (2011) & INSTAGRAM, LLC & 7.16 & 6.03 & +1.13 & \texttt{appearance enabling}, \texttt{modifying appearance}, \texttt{software modifying}, \texttt{enabling transmission}, \texttt{appearance} & \texttt{appearance enabling}, \texttt{modifying appearance}, \texttt{software modifying}, \texttt{enabling transmission}, \texttt{appearance} \\
PHOTOSHOP (1992) & ADOBE INC. & 4.81 & 4.83 & -0.01 & \texttt{manipulating graphic}, \texttt{creating manipulating}, \texttt{images computer}, \texttt{graphic images}, \texttt{manipulating} & \texttt{manipulating graphic}, \texttt{creating manipulating}, \texttt{images computer}, \texttt{graphic images}, \texttt{programs creating} \\
KUBERNETES (2014) & LF PROJECTS, LLC & 4.70 & 4.62 & +0.08 & \texttt{servers secure}, \texttt{secure private}, \texttt{private networks}, \texttt{networks}, \texttt{private} & \texttt{servers secure}, \texttt{secure private}, \texttt{private networks}, \texttt{networks}, \texttt{private} \\
WHITE CLAW SELTZER WORKS (2016) & MARK ANTHONY INTERNATIONAL S & 9.12 & 7.35 & +1.78 & \texttt{sugar-based beer}, \texttt{sugar-based}, \texttt{seltzer}, \texttt{beer} & \texttt{sugar-based beer}, \texttt{sugar-based}, \texttt{seltzer}, \texttt{beer} \\
LIQUID DEATH (2023) & Supplying Demand, Inc. & 3.97 & 3.80 & +0.17 & \texttt{drink mixes}, \texttt{water nutritional}, \texttt{mixes concentrates}, \texttt{water sparkling}, \texttt{concentrates} & \texttt{drink mixes}, \texttt{mixes concentrates}, \texttt{water nutritional}, \texttt{concentrates}, \texttt{flavor enhanced} \\
HARD MTN DEW (2021) & PEPSICO, INC. & 4.81 & 4.96 & -0.15 & \texttt{flavored brewed}, \texttt{brewed}, \texttt{brewed malt-based}, \texttt{malt-based alcoholic}, \texttt{nature beer} & \texttt{flavored brewed}, \texttt{brewed}, \texttt{brewed malt}, \texttt{nature beer}, \texttt{malt-based alcoholic} \\
ROCKSTAR (2002) & ROCKSTAR ENERGY LLC & 4.74 & 3.92 & +0.81 & \texttt{drinks}, \texttt{energy drinks}, \texttt{energy} & \texttt{drinks}, \texttt{energy drinks}, \texttt{energy} \\
\bottomrule
\end{tabular}

\end{table}
\end{landscape}

\subsection*{Limitations}
\begin{itemize}
  \item \textbf{Goods/services text is one of several novelty channels.}
        Brand novelty (\textsc{Liquid Death}), design novelty
        (\textsc{AirPods} casing), and pricing novelty are invisible to
        this corpus. Estimates here should be read as
        innovation-in-business-description novelty.
  \item \textbf{The crosswalk is exact-match only.} A fuzzy-match second
        pass and manual review of the top-by-volume USPTO owners would
        plausibly double or triple the matched panel size and reduce
        survivorship bias toward firms that did not change their
        registered name.
  \item \textbf{Survivorship bias in the financial panel.} Firms appear
        in SEC FSDS only while public. Pre-IPO innovators (Anthropic,
        OpenAI, Stripe at most timepoints) and dead-but-once-public
        firms drop out. The financials reward we estimate is conditional
        on already being a public concern.
  \item \textbf{The reference distributions are smoothed multinomials,
        not topic models.} An LDA replication is the obvious
        artifact-check.
\end{itemize}

\subsection*{Next steps}
The cross-industry universal sweep (all 45 NICE classes pooled) would
pin down which industries are innovative by this measure and would
allow new-class births (cannabis beverages 2018, generative AI tools
2022) to be scored against a cross-industry background. A
fuzzy-match crosswalk would expand the financial panel from $\sim 1{,}400$
firms toward the ${\sim}5{,}000$ ceiling of the SEC EDGAR universe. With
either or both extensions in place, the natural follow-up is a
\emph{variance} regression: do high-$\Delta KL$ firms exhibit higher
margin variance, the empirical fingerprint of pioneers either
spectacularly succeeding or spectacularly failing?

\appendix
\section{Per-class summary statistics}
\label{app:summary}

\begin{tabular}{lrrrrrrr}
\toprule
Industry (NICE) & Filings & Clean & Years & Owners & $\Delta KL$ med. & p5 & p95 \\
\midrule
Software \& Electronics (009) & 1,807,636 & 1,515,300 & 1958--2025 & 754,869 & -0.02 & -0.36 & +0.45 \\
Beer \& Soft Drinks (032) & 206,304 & 125,512 & 1980--2025 & 62,180 & +0.01 & -0.53 & +0.60 \\
\bottomrule
\end{tabular}


\section{Selected filings, full text}
\label{app:filings}

\begin{description}
\item[\textbf{OPENAI (2016)}] Owner at filing: \emph{OpenAI, Inc.}. $KL_{\text{pros}}=7.65$, $KL_{\text{retr}}=5.56$, $\Delta KL=+2.09$.\\
Goods/services text: \emph{Computer software for artificial intelligence}\\

\item[\textbf{CLAUDE (2023)}] Owner at filing: \emph{Anthropic, PBC}. $KL_{\text{pros}}=7.12$, $KL_{\text{retr}}=5.06$, $\Delta KL=+2.06$.\\
Goods/services text: \emph{Downloadable computer software in the nature of an artificial intelligence model for performing generative text AI tasks and natural language processing AI tasks and for writing content based on a theme, summarizing text, document question-answering and simulating natural conversation; downloadable software in the nature of a downloadable mobile application featuring an artificial intelligence model for performing generative text AI tasks and natural language processing AI tasks and for writing content based on a t…}\\

\item[\textbf{INSTAGRAM (2011)}] Owner at filing: \emph{INSTAGRAM, LLC}. $KL_{\text{pros}}=7.16$, $KL_{\text{retr}}=6.03$, $\Delta KL=+1.13$.\\
Goods/services text: \emph{Downloadable computer software for modifying the appearance and enabling transmission of photographs}\\

\item[\textbf{PHOTOSHOP (1992)}] Owner at filing: \emph{ADOBE INC.}. $KL_{\text{pros}}=4.81$, $KL_{\text{retr}}=4.83$, $\Delta KL=-0.01$.\\
Goods/services text: \emph{computer programs for creating and manipulating graphic images on a computer [ and manuals for use therewith, sold as a unit ]}\\

\item[\textbf{KUBERNETES (2014)}] Owner at filing: \emph{LF PROJECTS, LLC}. $KL_{\text{pros}}=4.70$, $KL_{\text{retr}}=4.62$, $\Delta KL=+0.08$.\\
Goods/services text: \emph{Computer software for accessing, browsing, sharing, defining, maintaining, virtualizing and communicating information over computer networks, servers and secure private networks; computer software for use in creating, managing, controlling, changing, replicating, deploying, naming and linking information over computer networks, servers and secure private networks; computer software for assisting users in navigating through computer networks, servers and secure private networks; computer software for running web app…}\\

\item[\textbf{WHITE CLAW SELTZER WORKS (2016)}] Owner at filing: \emph{MARK ANTHONY INTERNATIONAL SRL}. $KL_{\text{pros}}=9.12$, $KL_{\text{retr}}=7.35$, $\Delta KL=+1.78$.\\
Goods/services text: \emph{Hard seltzer Brewed sugar-based beer}\\

\item[\textbf{LIQUID DEATH (2023)}] Owner at filing: \emph{Supplying Demand, Inc.}. $KL_{\text{pros}}=3.97$, $KL_{\text{retr}}=3.80$, $\Delta KL=+0.17$.\\
Goods/services text: \emph{Powdered drink mixes for making flavored water, sports drinks, and energy drinks; powders and concentrates for making flavor enhanced water and sparkling water Powdered nutritional supplement drink mixes and concentrates}\\

\item[\textbf{HARD MTN DEW (2021)}] Owner at filing: \emph{PEPSICO, INC.}. $KL_{\text{pros}}=4.81$, $KL_{\text{retr}}=4.96$, $\Delta KL=-0.15$.\\
Goods/services text: \emph{flavored brewed malt-based alcoholic beverages in the nature of beer Alcoholic flavored brewed malt beverages, except beers}\\

\item[\textbf{ROCKSTAR (2002)}] Owner at filing: \emph{ROCKSTAR ENERGY LLC}. $KL_{\text{pros}}=4.74$, $KL_{\text{retr}}=3.92$, $\Delta KL=+0.81$.\\
Goods/services text: \emph{Sports drinks, namely, energy drinks}\\

\end{description}


\section{Falsification battery for the U-shape}
\label{app:falsification}

The compact summary in Section~\ref{sec:u-shape}
(Table~\ref{tab:alternatives}) listed seven candidate explanations
for the survival U-shape on $\Delta KL$ and the outcome of each
re-fit. This appendix narrates each check in full.

\paragraph{(1) Foreign-extension hypothesis is rejected.}
A common worry is that the very-stale (low-$\Delta KL$) end is
inflated by Madrid-Protocol or Section-44 international extensions
of marks already registered abroad --- such marks would use
vocabulary established years before they reached the US filing
record. Using a foreign-owner proxy (name suffixes \texttt{GMBH},
\texttt{S.A.}, \texttt{AB}, \texttt{KK}, \texttt{Pty}, \texttt{Pte},
\texttt{PLC}, \texttt{N.V.}, \texttt{B.V.}, etc.; $3.8\%$ of the
sample) and re-fitting on US-only and foreign-only sub-samples,
the quadratic coefficient is essentially identical ($+0.0078$ and
$+0.0100$ respectively). The U-shape is not a foreign-extension
artefact.

\paragraph{(2) The U is present in both first-time and recurring
filings} ($\beta_{\text{quad}} = +0.0080$ vs.\ $+0.0140$), with
recurring filings showing a slightly stronger curvature. The
fast-followers-do-best pattern from
Section~\ref{sec:pioneer-ladder} is consistent with stronger
curvature in the recurring subset (firms that are still filing
are firms that are still operating).

\paragraph{(3) The financial U-shape is essentially absent.}
A quadratic refit of the firm-year gross-margin OLS gives a linear
coefficient indistinguishable from zero ($\beta_{\text{lin}} =
+0.002$, SE $0.005$) and a marginally negative quadratic
($\beta_{\text{quad}} = -0.002$, SE $0.001$) on $8{,}600$
firm-years. \textbf{The U-shape lives in survival; gross margin
scales roughly linearly with $\Delta KL$.} The
$+3.2$ pp / $\sigma$ linear effect we report in
Section~\ref{sec:financials} is the right summary.

\paragraph{(4) Survivor bias from large incumbents partially supports
the left arm.} Mean firm filing-volume by $\Delta KL$ decile shows
the bottom decile is dominated by larger filers (mean firm-size
$103.7$ filings) than the middle deciles ($63$--$71$); the top
decile is back to $71$. Templated marks disproportionately come
from firms with deep portfolios who can afford to file --- and
maintain --- many similar boilerplate marks. This is consistent
with part of the left-arm uplift being \emph{firm capability}
rather than \emph{filing distinctiveness}.

\paragraph{(5) The ``Liquid Death'' hypothesis is rejected.}
A sceptic might worry that the bottom-decile arm is high-novelty
firms whose products are described in old vocabulary --- like
Liquid Death's canned water marketed as ``still water; sparkling
water.'' Splitting the bottom decile by goods/services text length
does not show this: the shortest quartile and the longest quartile
both renew at $44$--$46\%$ at five years. Templated marks are not
novel products in disguise; they really are templated, and they
really do survive.

\paragraph{(6) The U is filing-level, not firm-level.}
A further worry: the U-shape might be a property of the firms that
file rather than of any given filing --- ``on-trend'' firms might
simply fail more often, dragging their filings down, while
``harvest'' and ``innovate'' firms might survive and keep their
filings alive collectively. To test this we classified each firm
with $\geq 3$ clean filings by whether \emph{all} of its filings
survived 5y (``always-survives,'' $197{,}583$ firms), \emph{all}
failed (``never-survives,'' $253{,}064$ firms), or some survived
while others were let lapse (``mixed,'' $349{,}819$ firms). In the
mixed cohort --- firms that demonstrably persisted past at least
one Section-8 lapse --- the U-shape in filing-level survival by
$\Delta KL$ is preserved: bottom-decile rate $0.551$, middle
$0.479$, top $0.508$ (depth $\approx +0.07$ relative to middle).
Within the same firm, filings whose vocabulary places them at
either flux extreme are more likely to be the ones the firm keeps
alive than the flux-neutral filings the firm lets lapse. The U is
therefore a property of the specific filing, not just of the firm
behind it; trademarks on the flux-neutral ridge really are less
valuable to the firms that file them.
Figure~\ref{fig:firm-persistence} shows the three cohorts.

\paragraph{(7) Imitated-but-failed waves: the blockchain check.}
A natural counterexample to the ``imitated novelty pays'' claim is
to ask whether any recognisable hype wave was both widely imitated
(high $\Delta KL$) and disproportionately a failure. The
blockchain / cryptocurrency / NFT / web3 cohort is a clean test:
$87{,}463$ filings 1995--2025 mention at least one crypto-related
keyword in their goods/services text; $9{,}408$ are clean by our
usual restrictions. The cohort's mean $\Delta KL$ is $+0.412$ vs.\
$-0.002$ for the corpus average --- a genuinely high-novelty wave.
Survival-wise the wave is \emph{out-performing}, not failing:
registration rate $51.5\%$ vs.\ $48.3\%$ corpus mean; five-year
survival $51.4\%$ vs.\ $39.1\%$ corpus mean. Restricted to the top
decile of $\Delta KL$, the crypto cohort beats non-crypto top-decile
filings on both margins ($57.3\%$ vs.\ $47.5\%$ on registration;
$57.2\%$ vs.\ $40.9\%$ on five-year survival). Whatever else can
be said about the blockchain wave's commercial track record, the
trademark filings themselves are surviving well, consistent with
the same likelihood-of-confusion mechanic --- the crypto vocabulary
did not yet exist in the trademark register before the wave, so
applicants did not collide with prior marks.

\paragraph{(8) Temporal stability: the U is present year-by-year.}
The cleanest test of whether the pooled U is an aggregation artefact
is to refit the quadratic separately for each filing-year cohort and
ask whether the U appears in individual years or only when years
are pooled. We fit a logistic regression
$P(\text{survived\_5y}) = \mathrm{logit}^{-1}(\beta_0 + \beta_1
z(\Delta KL) + \beta_2 z(\Delta KL)^2 + \gamma \log(n_\text{terms}))$
separately for each year 1990--2021, both pooled across all 44 classes
and for five representative classes covering the range of cross-
sectional outcomes (Pharmaceuticals 005, Software \& Electronics 009,
Paper \& Print 016, Beer \& Soft Drinks 032, Telecommunications 038).
Figure~\ref{fig:u-shape-temporal} plots the per-year $\beta_2$ with
95\% CIs alongside NBER recession bars; the underlying per-year
counts are in \texttt{u\_shape\_temporal\_stability.csv}.

The pooled per-year coefficient is positive in $24$ of $32$ years and
\emph{significantly} positive ($t > 1.96$) in $23$ of $32$. It is
significantly inverted in only $4$ of $32$ years. The U is present in
the 1990s, in the 2000s, and in the 2010s; it strengthens slightly in
the 2020 cohort but does not depend on a single decade. Pharmaceuticals
is the cleanest case at the class level: $22 / 32$ years
significantly positive, $0$ significantly inverted. Software is
$19 / 32$ significantly positive, $2$ significantly inverted, with the
inverted years concentrated in the late-1990s pre-dot-com cohort.
Telecommunications is genuinely inverted in $10 / 32$ years; the
inverted-U cross-sectional result for that class is not a pooling
artefact but a real industry-level phenomenon.

The U does not track the macro cycle in a clean way. The 2001
dot-com recession is the one year in which the pooled coefficient
softens (it remains positive but only marginally significant); the
2008--09 financial crisis leaves the pooled coefficient roughly
unchanged; the 2020 recession shows a \emph{strengthening} of the
U. Recessions do not produce or destroy the U-shape; the U-shape
is more stable than the macro cycle is.

The temporal-stability test answers the question of whether the U is
a property of specific vintages of firms that happen to cluster at
flux-extreme positions. It is not: the pattern holds in nearly every
cohort independently.

\paragraph{(9) Competition-density mediation.}
The most-revealing falsification. We took $10{,}000$ stratified
Class-9 (Software \& Electronics) focal filings, sampled $1{,}000$
per $\Delta KL$ decile. For each focal filing $F$ at filing-date
$t_F$, we vectorised the entire Class-9 corpus with TF-IDF
(unigrams + bigrams, custom analyzer respecting semicolons,
$\text{min\_df} = 50$, $\text{max\_df} = 0.5$, L2-normalised),
restricted candidate neighbours to filings in $[t_F - 2y,\,
t_F + 5y]$ (the focal's own 5-year survival window plus a 2-year
pre-existing-competition lead-in), computed cosine similarity, and
kept neighbours with sim $\geq 0.7$. For each focal we then counted
\texttt{total\_competitors} (all such neighbours) and
\texttt{live\_competitors} (the subset whose own \texttt{survived\_5y}
was true).

The empirical result is in Table~\ref{tab:competition-density-bins}
and Figure~\ref{fig:competition-density}. By $\Delta KL$ decile, mean
live competitors goes:
$20.1, 20.0, 18.5, 12.9, 14.3, 11.9, 10.8, 12.8, 12.8, 31.4$.
The flux-neutral middle is in clusters with the fewest surviving
peers ($\sim 11$--$14$); both flux extremes are in clusters with $2$
to $3\times$ more surviving peers. The hypothesis we had pencilled in
--- ``flux-neutral firms are over-competed and the extremes are
blue-ocean'' --- is rejected in the direction of the count. The
extremes are not less-competed; they are in clusters where peers can
also survive.

\begin{table}[h]
\centering
\caption{Live-competitor density by $\Delta KL$ decile on $10{,}000$
stratified Class-9 focal filings. The flux-neutral middle has the
fewest live peers; both flux extremes have $2$--$3\times$ more.}
\label{tab:competition-density-bins}
\begin{tabular}{lrr}
\toprule
Coefficient & Model 1 & Model 2 \\
 & (no density control) & (+ $\log(1+\text{live competitors})$) \\
\midrule
$z(\Delta KL)$ (linear) & +0.0493 (0.0230) & +0.0616 (0.0233) \\
$z(\Delta KL)^2$ (quadratic) & +0.0110 (0.0076) & +0.0017 (0.0074) \\
$\log(1+\text{live competitors})$ & --- & +0.1929 (0.0163) \\
\midrule
$n$ focal filings (Class 9) & 10,000 & 10,000 \\
\bottomrule
\end{tabular}
\end{table}

The mediation regression is decisive on the second question
(``does competition density explain the survival U?''). Logit on
\texttt{survived\_5y} with $z(\Delta KL)$, $z(\Delta KL)^2$,
$\log(n_\text{terms})$, year and year$^2$ controls, with and without
$\log(1 + \text{live\_competitors})$ added:
\begin{itemize}
  \item Without the mediator: $\beta_{(\Delta KL)^2} = +0.0110$
        ($z = +1.43$).
  \item With the mediator: $\beta_{(\Delta KL)^2} = +0.0017$
        ($z = +0.24$) --- $84\%$ attenuation.
  \item Mediator coefficient: $+0.193$ ($z = +11.84$) --- strong,
        positive.
  \item Linear $z(\Delta KL)$ goes from $+0.049$ ($z = +2.15$) to
        $+0.062$ ($z = +2.65$) when the mediator is added: the
        linear flux effect is \emph{not} explained by competition
        density.
\end{itemize}

The interpretation is that the survival U on $\Delta KL$ is largely
mediated by the viability of the surrounding vocabulary cluster ---
filings at flux extremes survive better because their cluster
neighbours also survive --- but the linear ``positive flux helps''
effect is a residual signal beyond cluster viability. The U
\emph{shape} is a niche-viability proxy; the linear ``forward flux is
good'' effect is a flux-specific signal.

\paragraph{(10) External-corpus confound test (Google Books).}
A reader worried that $\Delta KL$ might just be tracking
general-corpus English-language usage would predict near-perfect
correlation between per-token trademark-corpus trajectories and
per-token Google Books n-gram trajectories over the same window. We
tested this on $75$ extreme-flux tokens (25 per stratum, drawn from
Table~\ref{tab:symmetric-flux-tokens}). For each token we pulled the
Google Books English n-gram frequency time series 1985--2021 via the
public viewer endpoint, and compared:
\begin{itemize}
  \item \texttt{tm\_lift} = mean trademark-corpus frequency in
        2017--2021 divided by mean in 1985--1989 (per-token).
  \item \texttt{google\_lift} = same computation on Google Books.
\end{itemize}
On the 53 tokens with usable data on both sides, the per-token
correlations are:
\begin{itemize}
  \item Pearson $r(\log \texttt{tm\_lift}, \log \texttt{google\_lift})
        = +0.466$ (Spearman $+0.388$).
  \item Pearson $r(\text{trademark log-slope}, \text{Google log-slope})
        = +0.443$ (Spearman $+0.344$).
\end{itemize}
The trademark-flux signal is materially correlated with general
English-language usage (about $22\%$ of trademark-flux variance per
token is shared with Google Books), but the bulk of variance ($78\%$)
is trademark-corpus-specific. The signal is partly but not wholly a
reflection of broader linguistic drift; the corpus-specific component
is what carries the survival predictive power.

A cleaner version of this test --- adding a Google-Books-derived
token-flux control to the filing-level survival regression --- is a
natural follow-up. We have not done it because (a) the per-token
correlation is moderate enough that we expect modest attenuation,
not nullification, and (b) the cluster-viability mediator in test
(9) is already the dominant mediator. We flag it as the most
informative extension if a reviewer requests further confound work.

\section{Per-industry deep dives}
\label{app:per-industry}

For each NICE class with at least $5{,}000$ clean H=2y filings in the
1990--2021 window, this appendix gives summary statistics, the top
three firms by firm-mean $\Delta KL$ (firms with $\geq 10$ filings),
and one worked example each from the positive-flux extreme, the
flux-neutral middle, and the negative-flux extreme. The brief
cross-industry table in Section~\ref{sec:per-industry-finance}
(\texttt{industry\_summary.tex}) gives an aggregated view; this
appendix is the per-class expansion for readers who want to drill
into a specific domain.

\subsection*{Class 001: Chemicals \& Industrial}
\label{app:industry-001}

\begin{tabular}{ll}
Clean filings: & 125,272 \\
Unique owners: & 52,278 \\
Year range: & 1990--2021 \\
Mean $KL_\text{pros}$: & 5.36 nats \\
Mean $KL_\text{retr}$: & 5.37 nats \\
Mean / median $\Delta KL$: & -0.009 / -0.010 \\
5y survival rate: & 47.1\% \\
Share with SEC-listed owner: & 6.47\% \\
\end{tabular}

\textbf{Top 3 firms by firm-mean $\Delta KL$ ($\geq 10$ filings):} HerbalScience Singapore Pte Ltd ($+1.20$); HERBALSCIENCE, INC. ($+1.04$); NICCA U.S.A., Inc. ($+1.03$)

\textbf{Worked examples}:

  \textbf{Positive flux:} \textsc{REEFERTILIZER} (2017, $\Delta KL=+3.49$):\\
  \textit{Plant nutrients, for use only with non-cannabis plants and cannabis with a delta-9 tetrahydrocannabinol (THC) concentration of not more than 0.3 percent on a dry weight basis}\\
  \textbf{Flux-neutral:} \textsc{COSMCPREP} (2003, $\Delta KL=+0.00$):\\
  \textit{reagents and magnetic beads used to purify plasmid, cosmid, fosmid, and BAC DNA for scientific research}\\
  \textbf{Negative flux:} \textsc{CLOUD CANDY} (2014, $\Delta KL=-1.83$):\\
  \textit{Electronic cigarette liquid (e-liquid) comprised of propylene glycol; Electronic cigarette liquid (e-liquid) comprised of vegetable glycerin}\\

\subsection*{Class 002: Paints \& Coatings}
\label{app:industry-002}

\begin{tabular}{ll}
Clean filings: & 33,653 \\
Unique owners: & 16,695 \\
Year range: & 1990--2021 \\
Mean $KL_\text{pros}$: & 4.78 nats \\
Mean $KL_\text{retr}$: & 4.79 nats \\
Mean / median $\Delta KL$: & -0.012 / -0.023 \\
5y survival rate: & 44.6\% \\
Share with SEC-listed owner: & 4.40\% \\
\end{tabular}

\textbf{Top 3 firms by firm-mean $\Delta KL$ ($\geq 10$ filings):} J. M. Huber Corporation ($+0.68$); ASCONN PRINT TECHNOLOGY CO.,LTD ($+0.57$); CLOVER TECHNOLOGIES GROUP, LLC ($+0.53$)

\textbf{Worked examples}:

  \textbf{Positive flux:} \textsc{BLOXX-IT} (1990, $\Delta KL=+2.38$):\\
  \textit{paint primer, bonder and sealer coating for interior and exterior use}\\
  \textbf{Flux-neutral:} \textsc{OVP} (1997, $\Delta KL=-0.00$):\\
  \textit{light interference pigments for use in such goods as inks, dyes, paints, plastics, and textiles}\\
  \textbf{Negative flux:} \textsc{OPTICS BY SOLA} (1999, $\Delta KL=-1.50$):\\
  \textit{Coatings for ophthalmic lenses, anti-reflective coatings for ophthalmic lenses, abrasion resistant coatings for ophthalmic lenses; impact resistant coatings for ophthalmic lenses, anti-fog coatings fo}\\

\subsection*{Class 003: Cosmetics \& Cleaning}
\label{app:industry-003}

\begin{tabular}{ll}
Clean filings: & 326,430 \\
Unique owners: & 154,551 \\
Year range: & 1990--2021 \\
Mean $KL_\text{pros}$: & 4.32 nats \\
Mean $KL_\text{retr}$: & 4.31 nats \\
Mean / median $\Delta KL$: & +0.005 / +0.005 \\
5y survival rate: & 50.0\% \\
Share with SEC-listed owner: & 2.94\% \\
\end{tabular}

\textbf{Top 3 firms by firm-mean $\Delta KL$ ($\geq 10$ filings):} MNG 2005 INC ($+1.46$); Revlon-Realistic Professional Products, Inc. ($+1.19$); Infused MFG, LLC ($+0.94$)

\textbf{Worked examples}:

  \textbf{Positive flux:} \textsc{ADARA} (2008, $\Delta KL=+2.85$):\\
  \textit{Hair care preparations consisting of organic coconut virgin oil and coconut virgin oil; Non-medicated cosmetic skin care preparations consisting of organic coconut virgin oil and coconut virgin oil}\\
  \textbf{Flux-neutral:} \textsc{ELEANOR MAE COLLECTION} (2002, $\Delta KL=+0.00$):\\
  \textit{non-medicated foot cream; skin lotion; milk bath; lip balm; goat milk soap}\\
  \textbf{Negative flux:} \textsc{PICKLED ESSENCE FAMILY} (2009, $\Delta KL=-1.55$):\\
  \textit{Aromatherapy pillows that may be used to freshen clothing or lay over the eyes}\\

\subsection*{Class 004: Lubricants \& Fuels}
\label{app:industry-004}

\begin{tabular}{ll}
Clean filings: & 36,160 \\
Unique owners: & 19,926 \\
Year range: & 1990--2021 \\
Mean $KL_\text{pros}$: & 4.91 nats \\
Mean $KL_\text{retr}$: & 4.91 nats \\
Mean / median $\Delta KL$: & +0.004 / -0.017 \\
5y survival rate: & 47.8\% \\
Share with SEC-listed owner: & 5.58\% \\
\end{tabular}

\textbf{Top 3 firms by firm-mean $\Delta KL$ ($\geq 10$ filings):} Frostbeard Studio, LLC ($+0.81$); hep global GmbH ($+0.59$); Victoria's Secret Stores, Inc. ($+0.55$)

\textbf{Worked examples}:

  \textbf{Positive flux:} \textsc{HEMPTHOLOGY} (2018, $\Delta KL=+2.82$):\\
  \textit{Candles; Wax melts; all of the foregoing goods containing CBD derived solely from hemp with a delta-9 tetrahydrocannabinol (THC) concentration of not more than 0.3 percent on a dry weight basis}\\
  \textbf{Flux-neutral:} \textsc{TE HACE FALTA PRENDER MAS CARBON!} (2017, $\Delta KL=+0.00$):\\
  \textit{Charcoal briquettes; Coal briquettes; Firewood; Lighter fluid; Lighter fluid for charcoal; Wood chips for use as fuel; Wood chips of mesquite for smoking and grilling; Solid fire starters}\\
  \textbf{Negative flux:} \textsc{HARNESSING THE POWER OF HEMP, TO POWER Y} (2021, $\Delta KL=-1.92$):\\
  \textit{Candles containing hemp with a delta-9 tetrahydrocannabinol concentration of not more than 0.3 percent on a dry weight basis}\\

\subsection*{Class 005: Pharmaceuticals}
\label{app:industry-005}

\begin{tabular}{ll}
Clean filings: & 359,560 \\
Unique owners: & 134,184 \\
Year range: & 1990--2021 \\
Mean $KL_\text{pros}$: & 4.98 nats \\
Mean $KL_\text{retr}$: & 4.99 nats \\
Mean / median $\Delta KL$: & -0.011 / -0.011 \\
5y survival rate: & 56.4\% \\
Share with SEC-listed owner: & 10.60\% \\
\end{tabular}

\textbf{Top 3 firms by firm-mean $\Delta KL$ ($\geq 10$ filings):} Silver Spur Corporation ($+2.28$); SILVER SPUR CORPORATION ($+1.54$); Earthmade Natural Vitamins, Inc. ($+1.44$)

\textbf{Worked examples}:

  \textbf{Positive flux:} \textsc{COLON MAGIC} (2009, $\Delta KL=+2.72$):\\
  \textit{Dietary supplement for eliminating toxins from the intestinal tract; Nutritional supplement for eliminating toxins from the intestinal tract}\\
  \textbf{Flux-neutral:} \textsc{NEXT GENERATION CURCUMIN} (2019, $\Delta KL=+0.00$):\\
  \textit{Dietary supplements containing curcumin Dietary ingredient, namely, curcumin, for use as an ingredient in the manufacture of dietary supplements}\\
  \textbf{Negative flux:} \textsc{HEART SMART} (2011, $\Delta KL=-1.47$):\\
  \textit{Foods and snacks which are free of Wheat to accommodate special medical and health conditions}\\

\subsection*{Class 006: Common Metals}
\label{app:industry-006}

\begin{tabular}{ll}
Clean filings: & 97,003 \\
Unique owners: & 55,943 \\
Year range: & 1990--2021 \\
Mean $KL_\text{pros}$: & 5.48 nats \\
Mean $KL_\text{retr}$: & 5.46 nats \\
Mean / median $\Delta KL$: & +0.020 / +0.013 \\
5y survival rate: & 43.4\% \\
Share with SEC-listed owner: & 4.04\% \\
\end{tabular}

\textbf{Top 3 firms by firm-mean $\Delta KL$ ($\geq 10$ filings):} Quad Performance Products, Inc. ($+0.42$); Kason Industries, Inc. ($+0.36$); General Electric Company ($+0.29$)

\textbf{Worked examples}:

  \textbf{Positive flux:} \textsc{GOLD DIAMOND} (1990, $\Delta KL=+2.73$):\\
  \textit{key chains, key tags and fobs, not of precious metal}\\
  \textbf{Flux-neutral:} \textsc{GVS-K} (2015, $\Delta KL=-0.00$):\\
  \textit{Junctions of metal for pipes; valves of metal, other than parts of machines; reinforcing materials of metal for pipes; elbows of metal for pipes; clips of metal for pipes; drain traps valves of metal;}\\
  \textbf{Negative flux:} \textsc{VALUE PRESS} (1995, $\Delta KL=-1.10$):\\
  \textit{metal closure caps for containers and bottles}\\

\subsection*{Class 007: Machines \& Power Tools}
\label{app:industry-007}

\begin{tabular}{ll}
Clean filings: & 160,666 \\
Unique owners: & 83,246 \\
Year range: & 1990--2021 \\
Mean $KL_\text{pros}$: & 5.34 nats \\
Mean $KL_\text{retr}$: & 5.34 nats \\
Mean / median $\Delta KL$: & +0.003 / -0.016 \\
5y survival rate: & 40.3\% \\
Share with SEC-listed owner: & 4.35\% \\
\end{tabular}

\textbf{Top 3 firms by firm-mean $\Delta KL$ ($\geq 10$ filings):} Goodnature Products, Inc. ($+0.50$); Arovast Corporation ($+0.40$); ANC Technology INC ($+0.38$)

\textbf{Worked examples}:

  \textbf{Positive flux:} \textsc{FAFNIR TDC} (1990, $\Delta KL=+2.34$):\\
  \textit{bearings for axles, transmissions, drive shafts and controls for land vehicles}\\
  \textbf{Flux-neutral:} \textsc{AXIS AUTOMATION} (2015, $\Delta KL=-0.00$):\\
  \textit{Design of standard and custom machinery for others Food processing machines, namely, tray denesters for separating food trays, and conveyor lines in the nature of vacuum conveyors Dispenser for dispen}\\
  \textbf{Negative flux:} \textsc{STANNAH STAIRLIFTS} (1992, $\Delta KL=-0.94$):\\
  \textit{mobility aids for elderly or disabled persons; namely, stairlifts and parts thereof}\\

\subsection*{Class 008: Hand Tools}
\label{app:industry-008}

\begin{tabular}{ll}
Clean filings: & 71,252 \\
Unique owners: & 44,508 \\
Year range: & 1990--2021 \\
Mean $KL_\text{pros}$: & 4.65 nats \\
Mean $KL_\text{retr}$: & 4.58 nats \\
Mean / median $\Delta KL$: & +0.070 / +0.032 \\
5y survival rate: & 40.3\% \\
Share with SEC-listed owner: & 2.66\% \\
\end{tabular}

\textbf{Top 3 firms by firm-mean $\Delta KL$ ($\geq 10$ filings):} BORBOLETA BEAUTY INC. ($+0.68$); STILETTO TOOLS, INC. ($+0.66$); Marut Enterprises LLC ($+0.52$)

\textbf{Worked examples}:

  \textbf{Positive flux:} \textsc{GARD'N'GRIP} (1990, $\Delta KL=+2.70$):\\
  \textit{hand tools for lawn and garden use; namely, rakes, shovels, cutters, hoes, and forks}\\
  \textbf{Flux-neutral:} \textsc{RIPLEY} (1996, $\Delta KL=+0.00$):\\
  \textit{wire, cable, and optical fibre preparation tools, namely, cutters, strippers, rippers, crimpers, trimmers, and slitters for electrical, electronic, CATV, telecommunication, telephone, and industrial a}\\
  \textbf{Negative flux:} \textsc{COMFORT REACH} (1999, $\Delta KL=-0.89$):\\
  \textit{hand pruners, weeders, bulb planters and loppers}\\

\subsection*{Class 009: Software \& Electronics}
\label{app:industry-009}

\begin{tabular}{ll}
Clean filings: & 1,226,806 \\
Unique owners: & 602,439 \\
Year range: & 1990--2021 \\
Mean $KL_\text{pros}$: & 5.42 nats \\
Mean $KL_\text{retr}$: & 5.43 nats \\
Mean / median $\Delta KL$: & -0.011 / -0.023 \\
5y survival rate: & 48.6\% \\
Share with SEC-listed owner: & 5.52\% \\
\end{tabular}

\textbf{Top 3 firms by firm-mean $\Delta KL$ ($\geq 10$ filings):} Leading Technology, Inc. ($+2.50$); Redmond Productions, Inc. ($+1.37$); Geniuses at Work Corporation ($+1.18$)

\textbf{Worked examples}:

  \textbf{Positive flux:} \textsc{MEDINET} (1990, $\Delta KL=+3.40$):\\
  \textit{a computer-based network including computer hardware and software}\\
  \textbf{Flux-neutral:} \textsc{PRIME PROTECTION} (2002, $\Delta KL=-0.00$):\\
  \textit{Computerized ear tags or chips for the identification of cattle and related computer software for animal health record keeping Plastic ear tags for the identification of cattle}\\
  \textbf{Negative flux:} \textsc{UNDONE} (2011, $\Delta KL=-1.44$):\\
  \textit{Electric hair curlers; Electric hair curling irons; Electric hair straightener; Electric hair straightening irons; Electric hair-curlers}\\

\subsection*{Class 010: Medical Devices}
\label{app:industry-010}

\begin{tabular}{ll}
Clean filings: & 190,121 \\
Unique owners: & 88,206 \\
Year range: & 1990--2021 \\
Mean $KL_\text{pros}$: & 5.35 nats \\
Mean $KL_\text{retr}$: & 5.35 nats \\
Mean / median $\Delta KL$: & -0.001 / -0.017 \\
5y survival rate: & 47.4\% \\
Share with SEC-listed owner: & 7.80\% \\
\end{tabular}

\textbf{Top 3 firms by firm-mean $\Delta KL$ ($\geq 10$ filings):} Societe Vygon ($+0.67$); NuGeneration Technologies, LLC ($+0.60$); CALIFORNIA EXOTIC NOVELTIES, LLC ($+0.58$)

\textbf{Worked examples}:

  \textbf{Positive flux:} \textsc{COVAFLU} (2020, $\Delta KL=+3.20$):\\
  \textit{Face covers being sanitary masks for protection against viral infection; Face shields for medical use}\\
  \textbf{Flux-neutral:} \textsc{TRI-CHEW} (2010, $\Delta KL=-0.00$):\\
  \textit{Oral therapy apparatus, namely, a tool used by a professional therapist or trained caregiver for jaw movement, tongue movement, lip movement, oral exploration, or for the development of feeding skills}\\
  \textbf{Negative flux:} \textsc{XEO} (2019, $\Delta KL=-1.29$):\\
  \textit{Catheters; catheters for percutaneous transluminal coronary angioplasty (PTCA catheters) and catheters for percutaneous transluminal angioplasty (PTA catheters), namely, catheters for balloon dilatati}\\

\subsection*{Class 011: Lighting, HVAC \& Appliances}
\label{app:industry-011}

\begin{tabular}{ll}
Clean filings: & 194,994 \\
Unique owners: & 112,215 \\
Year range: & 1990--2021 \\
Mean $KL_\text{pros}$: & 4.84 nats \\
Mean $KL_\text{retr}$: & 4.79 nats \\
Mean / median $\Delta KL$: & +0.045 / +0.014 \\
5y survival rate: & 38.8\% \\
Share with SEC-listed owner: & 2.80\% \\
\end{tabular}

\textbf{Top 3 firms by firm-mean $\Delta KL$ ($\geq 10$ filings):} G.K.L. Corp. ($+0.59$); ALPHARERRY INC ($+0.59$); VANJO LLC ($+0.54$)

\textbf{Worked examples}:

  \textbf{Positive flux:} \textsc{VENTEX} (1990, $\Delta KL=+2.50$):\\
  \textit{ventilation (air conditioning) installations and apparatus, fans, air conditioning apparatus}\\
  \textbf{Flux-neutral:} \textsc{MASTER VALET} (2012, $\Delta KL=+0.00$):\\
  \textit{[ Electric laundry irons and parts therefor ] Garment steamers [ ; steam generators for household purposes and parts therefore ]}\\
  \textbf{Negative flux:} \textsc{SOOPSPOON} (2019, $\Delta KL=-1.74$):\\
  \textit{Electrically-powered handheld vaporizers for vaporizing vaporizable constituents of herbal and plant matter for creating an aroma for household use; Gas powered handheld vaporizers for vaporizing vapo}\\

\subsection*{Class 012: Vehicles}
\label{app:industry-012}

\begin{tabular}{ll}
Clean filings: & 128,694 \\
Unique owners: & 68,694 \\
Year range: & 1990--2021 \\
Mean $KL_\text{pros}$: & 4.93 nats \\
Mean $KL_\text{retr}$: & 4.90 nats \\
Mean / median $\Delta KL$: & +0.034 / +0.013 \\
5y survival rate: & 43.7\% \\
Share with SEC-listed owner: & 5.01\% \\
\end{tabular}

\textbf{Top 3 firms by firm-mean $\Delta KL$ ($\geq 10$ filings):} Legendary Auto Interiors, Ltd. ($+0.99$); Drop Stop, LLC ($+0.70$); Hooey, LLC ($+0.62$)

\textbf{Worked examples}:

  \textbf{Positive flux:} \textsc{PF 815} (1990, $\Delta KL=+2.36$):\\
  \textit{vehicle tires for automobiles, trucks and vans}\\
  \textbf{Flux-neutral:} \textsc{I WHEEL} (2021, $\Delta KL=+0.00$):\\
  \textit{Wheelchair spoke covers; Wheelchair spoke guards; Wheelchair tire covers; Handrims and pushrims for wheelchairs}\\
  \textbf{Negative flux:} \textsc{ZF-POWERLINE} (2011, $\Delta KL=-1.05$):\\
  \textit{Transmissions, manual transmissions, automated and automatic transmissions, dual clutch transmissions, hybrid transmissions, synchronized transmissions, powershift transmissions, continuously variable}\\

\subsection*{Class 013: Firearms \& Pyrotechnics}
\label{app:industry-013}

\begin{tabular}{ll}
Clean filings: & 13,777 \\
Unique owners: & 7,027 \\
Year range: & 1991--2021 \\
Mean $KL_\text{pros}$: & 4.53 nats \\
Mean $KL_\text{retr}$: & 4.47 nats \\
Mean / median $\Delta KL$: & +0.063 / +0.059 \\
5y survival rate: & 45.9\% \\
Share with SEC-listed owner: & 3.94\% \\
\end{tabular}

\textbf{Top 3 firms by firm-mean $\Delta KL$ ($\geq 10$ filings):} BATTLEARMS IP, LLC ($+0.63$); Ghost Inc. ($+0.56$); Battle Arms Development, Inc. ($+0.50$)

\textbf{Worked examples}:

  \textbf{Positive flux:} \textsc{STS} (2009, $\Delta KL=+2.14$):\\
  \textit{(Based on Use in Commerce) Firearms(Based on Intent to Use) Ammunition for firearms; Firearm attachments, namely, modular external rail systems for attaching accessories to firearms}\\
  \textbf{Flux-neutral:} \textsc{SHAMAN} (2021, $\Delta KL=+0.00$):\\
  \textit{Hand gun accessories, namely, belt clips for securing a gun without the use of a holster}\\
  \textbf{Negative flux:} \textsc{} (1994, $\Delta KL=-1.60$):\\
  \textit{non-explosive defensive weapon in the nature of an organic irritant packaged in a pressurized spray canister}\\

\subsection*{Class 014: Jewelry \& Watches}
\label{app:industry-014}

\begin{tabular}{ll}
Clean filings: & 109,375 \\
Unique owners: & 72,078 \\
Year range: & 1990--2021 \\
Mean $KL_\text{pros}$: & 4.52 nats \\
Mean $KL_\text{retr}$: & 4.50 nats \\
Mean / median $\Delta KL$: & +0.019 / +0.004 \\
5y survival rate: & 43.2\% \\
Share with SEC-listed owner: & 1.79\% \\
\end{tabular}

\textbf{Top 3 firms by firm-mean $\Delta KL$ ($\geq 10$ filings):} Miki Corporation ($+0.61$); MIKI CORPORATION ($+0.60$); Spence Diamonds Ltd. ($+0.49$)

\textbf{Worked examples}:

  \textbf{Positive flux:} \textsc{OFF THE CLOSET} (2013, $\Delta KL=+3.11$):\\
  \textit{Women's jewelry; Earrings; Necklace and earring combinations that can be worn separately or as one piece; Necklaces as jewelry; Rings as jewelry; Bracelets}\\
  \textbf{Flux-neutral:} \textsc{ANCHOR2ANCHOR} (2021, $\Delta KL=+0.00$):\\
  \textit{Hats; Scarfs; Shoes; Shorts; Sneakers; Socks; Sweaters; Sweatshirts; Underwear; Bomber jackets; Denim jackets; Denim jeans; Denim shorts; Denim skirts; Knit shirts; Knit tops; Knitted caps; Knitted un}\\
  \textbf{Negative flux:} \textsc{SJ} (1992, $\Delta KL=-1.29$):\\
  \textit{gold or platinum jewelry with or without diamonds precious and/or semi-precious stones}\\

\subsection*{Class 015: Musical Instruments}
\label{app:industry-015}

\begin{tabular}{ll}
Clean filings: & 12,496 \\
Unique owners: & 7,733 \\
Year range: & 1991--2021 \\
Mean $KL_\text{pros}$: & 4.47 nats \\
Mean $KL_\text{retr}$: & 4.48 nats \\
Mean / median $\Delta KL$: & -0.008 / -0.015 \\
5y survival rate: & 43.2\% \\
Share with SEC-listed owner: & 0.74\% \\
\end{tabular}

\textbf{Top 3 firms by firm-mean $\Delta KL$ ($\geq 10$ filings):} Lone Wolf Blues Company ($+0.38$); Zdenka Infeld Asset Management GmbH ($+0.32$); Godlyke, Inc. ($+0.29$)

\textbf{Worked examples}:

  \textbf{Positive flux:} \textsc{EZ CHORD} (2006, $\Delta KL=+1.27$):\\
  \textit{Guitar playing assistance devices, namely, a sliding clamp-like structure positioned onto the guitar neck to form and sound guitar chords and notes}\\
  \textbf{Flux-neutral:} \textsc{O-PORT} (2009, $\Delta KL=-0.00$):\\
  \textit{Guitar playing assistance devices, namely, a sliding clamp-like structure positioned onto the guitar neck to form and sound guitar chords and notes}\\
  \textbf{Negative flux:} \textsc{RHYTHM TOOLS} (2003, $\Delta KL=-1.04$):\\
  \textit{Musical instruments, namely drums; drumheads; a click clack, namely a set of jingles, usually metal, mounted within a generally u-shaped frame; a timpano, namely a tensioned membrane mounted in a rect}\\

\subsection*{Class 016: Paper, Print \& Office}
\label{app:industry-016}

\begin{tabular}{ll}
Clean filings: & 444,018 \\
Unique owners: & 245,687 \\
Year range: & 1990--2021 \\
Mean $KL_\text{pros}$: & 5.39 nats \\
Mean $KL_\text{retr}$: & 5.40 nats \\
Mean / median $\Delta KL$: & -0.008 / -0.012 \\
5y survival rate: & 50.2\% \\
Share with SEC-listed owner: & 3.87\% \\
\end{tabular}

\textbf{Top 3 firms by firm-mean $\Delta KL$ ($\geq 10$ filings):} Llumar Star Kites, Inc. ($+1.08$); Brown \& Bigelow, Inc. ($+0.90$); Atkins Daniel D ($+0.82$)

\textbf{Worked examples}:

  \textbf{Positive flux:} \textsc{IMPRINT JIM DIESO} (2010, $\Delta KL=+2.53$):\\
  \textit{Art prints comprised of digital illustrations originating from photographs}\\
  \textbf{Flux-neutral:} \textsc{STAND \& WITHSTAND INTEGRITY GROUP} (2021, $\Delta KL=+0.00$):\\
  \textit{Educational services, namely, developing curriculum for others in the field of how to get an education; Training services in the field of how to get an education; Arranging and conducting of seminars }\\
  \textbf{Negative flux:} \textsc{ITC MENDOZA ROMAN} (1992, $\Delta KL=-2.08$):\\
  \textit{typefaces, typefonts, and type designs of alphanumeric characters and/or typographical symbols typefaces, typefonts and type designs of alphanumeric characters and/or typographical symbols recorded on}\\

\subsection*{Class 017: Rubber \& Insulators}
\label{app:industry-017}

\begin{tabular}{ll}
Clean filings: & 51,254 \\
Unique owners: & 26,584 \\
Year range: & 1990--2021 \\
Mean $KL_\text{pros}$: & 4.97 nats \\
Mean $KL_\text{retr}$: & 4.96 nats \\
Mean / median $\Delta KL$: & +0.018 / +0.010 \\
5y survival rate: & 42.6\% \\
Share with SEC-listed owner: & 4.28\% \\
\end{tabular}

\textbf{Top 3 firms by firm-mean $\Delta KL$ ($\geq 10$ filings):} Manco, Inc. ($+0.44$); MadeSolid, Inc. ($+0.42$); Insta-Foam Products, Inc. ($+0.33$)

\textbf{Worked examples}:

  \textbf{Positive flux:} \textsc{ADHESIVE PLUS} (1990, $\Delta KL=+2.22$):\\
  \textit{chemical adhesive compound for general, construction and household use in sealing, caulking and insulating}\\
  \textbf{Flux-neutral:} \textsc{RENEWLENE} (2008, $\Delta KL=+0.00$):\\
  \textit{Extruded plastic in the form of bars, blocks, pellets, rods, sheets and tubes for use in manufacturing; Plastic in bars, blocks, pellets, rods, sheets and tubes for general industrial use; Plastic in }\\
  \textbf{Negative flux:} \textsc{DURRA BUILDING STRONG SOLUTIONS} (2004, $\Delta KL=-0.90$):\\
  \textit{acoustical wall and roof panels for buildings}\\

\subsection*{Class 018: Leather \& Luggage}
\label{app:industry-018}

\begin{tabular}{ll}
Clean filings: & 147,807 \\
Unique owners: & 95,097 \\
Year range: & 1990--2021 \\
Mean $KL_\text{pros}$: & 4.32 nats \\
Mean $KL_\text{retr}$: & 4.31 nats \\
Mean / median $\Delta KL$: & +0.008 / -0.015 \\
5y survival rate: & 43.4\% \\
Share with SEC-listed owner: & 1.98\% \\
\end{tabular}

\textbf{Top 3 firms by firm-mean $\Delta KL$ ($\geq 10$ filings):} Piper, Diane E. ($+0.51$); Romar International Corp. ($+0.50$); Bresnick, Tiffani ($+0.38$)

\textbf{Worked examples}:

  \textbf{Positive flux:} \textsc{STEPHANO SABETTI} (1990, $\Delta KL=+2.63$):\\
  \textit{men's, women's, and children's clothing - namely, pants, shirts, t-shirts, sweatshirts, jackets, coats, sweaters, scarves, gloves, socks, shoes, ties, and hats leather articles, namely, carry on bags,}\\
  \textbf{Flux-neutral:} \textsc{ROTUNDA} (1996, $\Delta KL=+0.00$):\\
  \textit{business card cases, shoe bags for travel, tote bags}\\
  \textbf{Negative flux:} \textsc{AXIOM EQUINE SPORT BOOT} (2012, $\Delta KL=-0.98$):\\
  \textit{Horse flanks and shanks, namely, horse protective padding, braces, and leg boots for the flanks and legs}\\

\subsection*{Class 019: Building Materials}
\label{app:industry-019}

\begin{tabular}{ll}
Clean filings: & 70,015 \\
Unique owners: & 34,729 \\
Year range: & 1990--2021 \\
Mean $KL_\text{pros}$: & 5.08 nats \\
Mean $KL_\text{retr}$: & 5.10 nats \\
Mean / median $\Delta KL$: & -0.015 / -0.014 \\
5y survival rate: & 45.6\% \\
Share with SEC-listed owner: & 3.45\% \\
\end{tabular}

\textbf{Top 3 firms by firm-mean $\Delta KL$ ($\geq 10$ filings):} Saroyan Lumber Company ($+0.82$); SCHLUTER-SYSTEMS KG ($+0.38$); Holyfield Management Inc. ($+0.35$)

\textbf{Worked examples}:

  \textbf{Positive flux:} \textsc{ACRAJOINT} (1990, $\Delta KL=+2.94$):\\
  \textit{building materials; namely, a system comprised of glass blocks and adhesive used to form a wall, sold as a unit}\\
  \textbf{Flux-neutral:} \textsc{MANTA} (1999, $\Delta KL=+0.00$):\\
  \textit{material in the form of bricks, tiles, sheets or mortar for screening electromagnetic interference; material comprising a blend of carbon and iron oxides (ferrite) in the form of bricks, tiles, sheets}\\
  \textbf{Negative flux:} \textsc{HOMESEAL} (1998, $\Delta KL=-1.03$):\\
  \textit{Asphalt, modified asphalt and rubberized asphalt based peel and stick membranes for use on roofing}\\

\subsection*{Class 020: Furniture}
\label{app:industry-020}

\begin{tabular}{ll}
Clean filings: & 165,306 \\
Unique owners: & 102,530 \\
Year range: & 1990--2021 \\
Mean $KL_\text{pros}$: & 5.16 nats \\
Mean $KL_\text{retr}$: & 5.09 nats \\
Mean / median $\Delta KL$: & +0.075 / +0.042 \\
5y survival rate: & 42.3\% \\
Share with SEC-listed owner: & 2.76\% \\
\end{tabular}

\textbf{Top 3 firms by firm-mean $\Delta KL$ ($\geq 10$ filings):} Xiamen Youjing E-commerce Co., Ltd. ($+0.66$); Everglades College, Inc. ($+0.63$); Shanxi Huacai Zhongxing E-commerce Co., Ltd. ($+0.62$)

\textbf{Worked examples}:

  \textbf{Positive flux:} \textsc{THE SIMPSONS} (1990, $\Delta KL=+2.05$):\\
  \textit{tables, chairs, desks, beds, children's furniture, mattresses, pillows, mirrors, picture frames, plastic key ring tags}\\
  \textbf{Flux-neutral:} \textsc{TILANG} (2020, $\Delta KL=+0.00$):\\
  \textit{Pet toys Planters for flowers and plants Shoe bags for travel Curtain rods; Mirrors}\\
  \textbf{Negative flux:} \textsc{ULTRA VIVID} (2016, $\Delta KL=-0.83$):\\
  \textit{Non-metal and non-leather key chains Treatment of materials by laser beam}\\

\subsection*{Class 021: Household Utensils}
\label{app:industry-021}

\begin{tabular}{ll}
Clean filings: & 219,996 \\
Unique owners: & 143,153 \\
Year range: & 1990--2021 \\
Mean $KL_\text{pros}$: & 4.85 nats \\
Mean $KL_\text{retr}$: & 4.78 nats \\
Mean / median $\Delta KL$: & +0.067 / +0.037 \\
5y survival rate: & 40.8\% \\
Share with SEC-listed owner: & 2.23\% \\
\end{tabular}

\textbf{Top 3 firms by firm-mean $\Delta KL$ ($\geq 10$ filings):} ALPHARERRY INC ($+0.81$); InnovatioNxt, LLC ($+0.72$); FLYCAMP TRADING LLC ($+0.67$)

\textbf{Worked examples}:

  \textbf{Positive flux:} \textsc{IREPELLER} (2017, $\Delta KL=+1.96$):\\
  \textit{Plug-in diffusers for mosquito repellents; Ultrasonic pest repellers}\\
  \textbf{Flux-neutral:} \textsc{HALO} (2017, $\Delta KL=+0.00$):\\
  \textit{Wine preservation devices, namely, bottle stoppers specifically adapted for use with wine bottles}\\
  \textbf{Negative flux:} \textsc{BEER RINGS} (2016, $\Delta KL=-1.03$):\\
  \textit{Charms for attachment to beverage glassware for identification purposes}\\

\subsection*{Class 022: Ropes, Tents \& Textiles (raw)}
\label{app:industry-022}

\begin{tabular}{ll}
Clean filings: & 24,515 \\
Unique owners: & 16,565 \\
Year range: & 1990--2021 \\
Mean $KL_\text{pros}$: & 4.70 nats \\
Mean $KL_\text{retr}$: & 4.59 nats \\
Mean / median $\Delta KL$: & +0.111 / +0.072 \\
5y survival rate: & 40.1\% \\
Share with SEC-listed owner: & 2.31\% \\
\end{tabular}

\textbf{Top 3 firms by firm-mean $\Delta KL$ ($\geq 10$ filings):} Caravan Canopy Int'l, Inc. ($+0.69$); Xiamen Youjing E-commerce Co., Ltd. ($+0.58$); High Caliper Growing, Inc. ($+0.54$)

\textbf{Worked examples}:

  \textbf{Positive flux:} \textsc{} (2013, $\Delta KL=+2.41$):\\
  \textit{Awnings; Awnings comprised primarily of tensile fabric membranes; Awnings for vehicles}\\
  \textbf{Flux-neutral:} \textsc{FAYONLINE} (2018, $\Delta KL=-0.00$):\\
  \textit{Sleeping bag pads; Futons Hammocks Blanket throws; Baby blankets; Children's blankets; Liners adapted to sleeping bags for camping}\\
  \textbf{Negative flux:} \textsc{STONE CONTAINER CORPORATION} (1990, $\Delta KL=-0.95$):\\
  \textit{[ logs, trees, wood chips and bark ] agricultural and forestry products; namely, market pulp sawdust}\\

\subsection*{Class 024: Textiles \& Bedding}
\label{app:industry-024}

\begin{tabular}{ll}
Clean filings: & 97,945 \\
Unique owners: & 59,850 \\
Year range: & 1990--2021 \\
Mean $KL_\text{pros}$: & 4.25 nats \\
Mean $KL_\text{retr}$: & 4.21 nats \\
Mean / median $\Delta KL$: & +0.042 / +0.016 \\
5y survival rate: & 42.9\% \\
Share with SEC-listed owner: & 2.23\% \\
\end{tabular}

\textbf{Top 3 firms by firm-mean $\Delta KL$ ($\geq 10$ filings):} FISHER TEXTILES, LLC ($+0.61$); An Le ($+0.51$); Zhu Tianding ($+0.42$)

\textbf{Worked examples}:

  \textbf{Positive flux:} \textsc{RECORD BREAKERS WORLD OF SPEED} (1990, $\Delta KL=+2.91$):\\
  \textit{bed sheets, pillow cases, bedspreads, comforters, draperies, curtains, blankets}\\
  \textbf{Flux-neutral:} \textsc{ARM \& HAMMER} (2010, $\Delta KL=+0.00$):\\
  \textit{box spring covers, bed blankets}\\
  \textbf{Negative flux:} \textsc{APPLETON LANE} (1996, $\Delta KL=-0.94$):\\
  \textit{vinyl wallcoverings and wallcovering borders textile fabrics for use in draperies and upholstery sample books of wallcoverings}\\

\subsection*{Class 025: Clothing, Footwear \& Headwear}
\label{app:industry-025}

\begin{tabular}{ll}
Clean filings: & 668,593 \\
Unique owners: & 426,635 \\
Year range: & 1990--2021 \\
Mean $KL_\text{pros}$: & 4.36 nats \\
Mean $KL_\text{retr}$: & 4.40 nats \\
Mean / median $\Delta KL$: & -0.041 / -0.051 \\
5y survival rate: & 51.6\% \\
Share with SEC-listed owner: & 1.39\% \\
\end{tabular}

\textbf{Top 3 firms by firm-mean $\Delta KL$ ($\geq 10$ filings):} CEO Lux, LLC ($+0.72$); CHILD SAFETY PRODUCTS, INC. ($+0.66$); Warner Bros. Inc. ($+0.57$)

\textbf{Worked examples}:

  \textbf{Positive flux:} \textsc{ÜBERION} (2007, $\Delta KL=+3.14$):\\
  \textit{Triathlon clothing, namely, triathlon tights, triathlon shorts, triathlon singlets, triathlon shirts and triathlon suits}\\
  \textbf{Flux-neutral:} \textsc{COBALT} (2019, $\Delta KL=-0.00$):\\
  \textit{wetsuits; wetsuit gloves; water socks; rash guards; aqua shoes; flip-flops for use as footwear}\\
  \textbf{Negative flux:} \textsc{GEFFEN BABY} (2011, $\Delta KL=-1.53$):\\
  \textit{Cloth diapers; Infant cloth diapers; Inserts specially adapted for cloth diapers made of hemp; Textile diapers}\\

\subsection*{Class 026: Trim, Lace \& Notions}
\label{app:industry-026}

\begin{tabular}{ll}
Clean filings: & 42,323 \\
Unique owners: & 29,873 \\
Year range: & 1990--2021 \\
Mean $KL_\text{pros}$: & 4.47 nats \\
Mean $KL_\text{retr}$: & 4.34 nats \\
Mean / median $\Delta KL$: & +0.124 / +0.075 \\
5y survival rate: & 40.4\% \\
Share with SEC-listed owner: & 1.76\% \\
\end{tabular}

\textbf{Top 3 firms by firm-mean $\Delta KL$ ($\geq 10$ filings):} Celebrity Signatures International, Inc. ($+1.07$); RED SLATE BRANDS LLC ($+0.82$); KIM, HYUN JOON ($+0.77$)

\textbf{Worked examples}:

  \textbf{Positive flux:} \textsc{BEAUTY BUN} (2011, $\Delta KL=+2.08$):\\
  \textit{Hair accessories, namely, plastic styling inserts that create height and volume on the crown of the head}\\
  \textbf{Flux-neutral:} \textsc{STOGEDEST} (2021, $\Delta KL=+0.00$):\\
  \textit{Wigs; Clown wigs; False hair; Hair extensions; Hair pieces; Hair pieces and wigs; Hair weaves; Human braiding hair; Wig caps}\\
  \textbf{Negative flux:} \textsc{TEXTURED LOCS} (2020, $\Delta KL=-2.66$):\\
  \textit{Wigs, hairpieces, and add-in and add-on hair accessories constructed primarily of synthetic and/or human hair}\\

\subsection*{Class 027: Carpets \& Wall Hangings}
\label{app:industry-027}

\begin{tabular}{ll}
Clean filings: & 26,727 \\
Unique owners: & 15,621 \\
Year range: & 1990--2021 \\
Mean $KL_\text{pros}$: & 4.55 nats \\
Mean $KL_\text{retr}$: & 4.49 nats \\
Mean / median $\Delta KL$: & +0.053 / +0.019 \\
5y survival rate: & 40.9\% \\
Share with SEC-listed owner: & 3.14\% \\
\end{tabular}

\textbf{Top 3 firms by firm-mean $\Delta KL$ ($\geq 10$ filings):} Marketing and Sales Management Corp. ($+0.75$); J. Josephson, Inc. ($+0.59$); Watershed Geosynthetics LLC ($+0.54$)

\textbf{Worked examples}:

  \textbf{Positive flux:} \textsc{THE SIMPSONS} (1990, $\Delta KL=+2.63$):\\
  \textit{carpets, rugs, vinyl wall coverings, wall hangings, wall paper}\\
  \textbf{Flux-neutral:} \textsc{MILLICARE} (1996, $\Delta KL=+0.00$):\\
  \textit{cleaning preparations consisting of porous, absorbent particles in powder form and a liquid concentrate for use in the cleaning of commercial carpets}\\
  \textbf{Negative flux:} \textsc{GENCORP} (1997, $\Delta KL=-2.06$):\\
  \textit{wallcovering having an outer surface of plastic}\\

\subsection*{Class 028: Toys, Games \& Sporting Goods}
\label{app:industry-028}

\begin{tabular}{ll}
Clean filings: & 325,521 \\
Unique owners: & 167,749 \\
Year range: & 1990--2021 \\
Mean $KL_\text{pros}$: & 4.94 nats \\
Mean $KL_\text{retr}$: & 4.94 nats \\
Mean / median $\Delta KL$: & -0.006 / -0.030 \\
5y survival rate: & 47.1\% \\
Share with SEC-listed owner: & 3.83\% \\
\end{tabular}

\textbf{Top 3 firms by firm-mean $\Delta KL$ ($\geq 10$ filings):} Aruze Gaming America, Inc. ($+1.17$); KING SHOW GAMES, INC. ($+0.73$); Nabisco Brands Company ($+0.73$)

\textbf{Worked examples}:

  \textbf{Positive flux:} \textsc{WASHINGTON SENATORS} (1990, $\Delta KL=+3.02$):\\
  \textit{toys and sporting goods; namely, Christmas tree ornaments}\\
  \textbf{Flux-neutral:} \textsc{CORTEZ} (2001, $\Delta KL=-0.00$):\\
  \textit{fishing products, namely, rods, reels, lures, fishing line, terminal tackle, tackle bags and tackle boxes}\\
  \textbf{Negative flux:} \textsc{SWAY BUSTER} (2013, $\Delta KL=-2.26$):\\
  \textit{Golf Training Apparatus, namely, devices used to help groove a repetitive swing and stance}\\

\subsection*{Class 029: Meat, Dairy \& Prepared Foods}
\label{app:industry-029}

\begin{tabular}{ll}
Clean filings: & 134,683 \\
Unique owners: & 64,279 \\
Year range: & 1990--2021 \\
Mean $KL_\text{pros}$: & 4.63 nats \\
Mean $KL_\text{retr}$: & 4.65 nats \\
Mean / median $\Delta KL$: & -0.020 / -0.023 \\
5y survival rate: & 51.8\% \\
Share with SEC-listed owner: & 3.42\% \\
\end{tabular}

\textbf{Top 3 firms by firm-mean $\Delta KL$ ($\geq 10$ filings):} BEYOND MEAT, INC. ($+0.94$); Aphria Inc. ($+0.78$); Crunchy Munchies LLC ($+0.62$)

\textbf{Worked examples}:

  \textbf{Positive flux:} \textsc{PHIVIDA} (2016, $\Delta KL=+2.68$):\\
  \textit{Healthy base beverages, namely, drinking water with vitamins containing hemp derived from industrial hemp with a delta-9 tetrahydrocannabinol (THC) concentration of not more than 0.3 percent on a dry }\\
  \textbf{Flux-neutral:} \textsc{MUSCLE BROTH} (2016, $\Delta KL=+0.00$):\\
  \textit{liquid or powdered meat extracts containing amino acids, vitamins and minerals}\\
  \textbf{Negative flux:} \textsc{EXCITE MY BITE} (2021, $\Delta KL=-1.31$):\\
  \textit{bouillon not containing or derived from cannabis with a delta-9 THC concentration of more than 0.3 percent on a dry weight basis or containing CBD}\\

\subsection*{Class 030: Coffee, Tea, Bakery \& Spices}
\label{app:industry-030}

\begin{tabular}{ll}
Clean filings: & 205,442 \\
Unique owners: & 104,841 \\
Year range: & 1990--2021 \\
Mean $KL_\text{pros}$: & 4.90 nats \\
Mean $KL_\text{retr}$: & 4.91 nats \\
Mean / median $\Delta KL$: & -0.011 / -0.010 \\
5y survival rate: & 51.3\% \\
Share with SEC-listed owner: & 2.72\% \\
\end{tabular}

\textbf{Top 3 firms by firm-mean $\Delta KL$ ($\geq 10$ filings):} ECBlend LLC ($+1.15$); VMR Products, LLC ($+0.97$); New Generation Foods, Inc. ($+0.85$)

\textbf{Worked examples}:

  \textbf{Positive flux:} \textsc{SWITCH} (2012, $\Delta KL=+2.67$):\\
  \textit{Electronic cigarette batteries; Electronic cigarette battery chargers Electronic cigarette liquid (e-liquid) comprised of flavorings in liquid form used to refill electronic cigarettes Electronic ciga}\\
  \textbf{Flux-neutral:} \textsc{FRUITABLES} (1995, $\Delta KL=-0.00$):\\
  \textit{fruit rolls, fruit leathers, dried fruit, fruit flavored confections, namely candied fruit, fruit flavored snacks}\\
  \textbf{Negative flux:} \textsc{EMERALD CITY VAPE PREMIUM VAPOR LIQUID S} (2015, $\Delta KL=-2.01$):\\
  \textit{Electronic cigarette liquid (e-liquid) comprised of flavorings in liquid form used to refill electronic cigarette cartridges}\\

\subsection*{Class 031: Agriculture \& Live Animals}
\label{app:industry-031}

\begin{tabular}{ll}
Clean filings: & 64,135 \\
Unique owners: & 32,411 \\
Year range: & 1990--2021 \\
Mean $KL_\text{pros}$: & 4.92 nats \\
Mean $KL_\text{retr}$: & 4.92 nats \\
Mean / median $\Delta KL$: & -0.004 / -0.013 \\
5y survival rate: & 49.8\% \\
Share with SEC-listed owner: & 2.21\% \\
\end{tabular}

\textbf{Top 3 firms by firm-mean $\Delta KL$ ($\geq 10$ filings):} TRIM READY INC ($+1.16$); DAVID AUSTIN ROSES LIMITED ($+1.09$); CMG PARTNERS, INC ($+0.76$)

\textbf{Worked examples}:

  \textbf{Positive flux:} \textsc{BARK AVENUE} (2017, $\Delta KL=+3.20$):\\
  \textit{Edible dog treats not containing CBD and not containing any ingredients with a delta-9 tetrahydrocannabinol concentration of more than 0.3 percent on a dry weight basis}\\
  \textbf{Flux-neutral:} \textsc{DOC'CHOPS} (2004, $\Delta KL=-0.00$):\\
  \textit{Animal snack food; Snack food for animals designed to contain medications}\\
  \textbf{Negative flux:} \textsc{12 TWENTY} (2019, $\Delta KL=-1.24$):\\
  \textit{Pet food Transdermal patches for therapeutic use in the treatment of addictions, for relief of pain, nausea, and discomfort, for relaxation, for reducing stress, nervousness and fatigue, for maintaini}\\

\subsection*{Class 032: Beer \& Soft Drinks}
\label{app:industry-032}

\begin{tabular}{ll}
Clean filings: & 100,951 \\
Unique owners: & 49,681 \\
Year range: & 1990--2021 \\
Mean $KL_\text{pros}$: & 4.18 nats \\
Mean $KL_\text{retr}$: & 4.20 nats \\
Mean / median $\Delta KL$: & -0.014 / -0.017 \\
5y survival rate: & 56.6\% \\
Share with SEC-listed owner: & 3.11\% \\
\end{tabular}

\textbf{Top 3 firms by firm-mean $\Delta KL$ ($\geq 10$ filings):} HONEST INC. ($+1.98$); Duclaw Brewery, LLC ($+0.96$); All Market Inc. ($+0.89$)

\textbf{Worked examples}:

  \textbf{Positive flux:} \textsc{PHIVIDA} (2016, $\Delta KL=+2.77$):\\
  \textit{Healthy base beverages, namely, drinking water with vitamins containing hemp derived from industrial hemp with a delta-9 tetrahydrocannabinol (THC) concentration of not more than 0.3 percent on a dry }\\
  \textbf{Flux-neutral:} \textsc{TRUE CAL-ORGANIC FARMS} (2013, $\Delta KL=-0.00$):\\
  \textit{Organic fresh fruits; organic fresh vegetables Organic fruit juices and beverages; organic vegetable juices and beverages}\\
  \textbf{Negative flux:} \textsc{NORTHERN EXPOSURE} (1997, $\Delta KL=-1.15$):\\
  \textit{fermented alcoholic beverages made with malt, namely, assorted beers}\\

\subsection*{Class 033: Wine \& Spirits}
\label{app:industry-033}

\begin{tabular}{ll}
Clean filings: & 68,232 \\
Unique owners: & 32,766 \\
Year range: & 1990--2021 \\
Mean $KL_\text{pros}$: & 4.09 nats \\
Mean $KL_\text{retr}$: & 4.06 nats \\
Mean / median $\Delta KL$: & +0.038 / +0.020 \\
5y survival rate: & 53.1\% \\
Share with SEC-listed owner: & 0.90\% \\
\end{tabular}

\textbf{Top 3 firms by firm-mean $\Delta KL$ ($\geq 10$ filings):} Stranger \& Stranger Ltd ($+0.76$); MAESTRO TEQUILERO, S.A. DE C.V. ($+0.72$); Destiladora Del Valle de Tequila, S.A. de C.V. ($+0.70$)

\textbf{Worked examples}:

  \textbf{Positive flux:} \textsc{MARYGIN} (2018, $\Delta KL=+2.60$):\\
  \textit{Liquor; all of the foregoing containing ingredients solely derived from hemp with a delta-9 tetrahydrocannabinol (THC) concentration of not more than 0.3 percent on a dry weight basis and not containi}\\
  \textbf{Flux-neutral:} \textsc{UNITED STARS} (2018, $\Delta KL=+0.00$):\\
  \textit{Alcoholic beverages, namely, vodka, prepared alcoholic cocktails, peppermint liqueurs, distilled spirits, alcoholic beverages containing fruit}\\
  \textbf{Negative flux:} \textsc{CHATEAU SOCIANDO-MALLET} (2012, $\Delta KL=-1.04$):\\
  \textit{Wines of guaranteed label of origin Haut Medoc from a particular wine-making estate}\\

\subsection*{Class 034: Tobacco \& Smokers' Articles}
\label{app:industry-034}

\begin{tabular}{ll}
Clean filings: & 36,678 \\
Unique owners: & 18,795 \\
Year range: & 1990--2021 \\
Mean $KL_\text{pros}$: & 3.66 nats \\
Mean $KL_\text{retr}$: & 3.50 nats \\
Mean / median $\Delta KL$: & +0.163 / +0.115 \\
5y survival rate: & 55.4\% \\
Share with SEC-listed owner: & 1.42\% \\
\end{tabular}

\textbf{Top 3 firms by firm-mean $\Delta KL$ ($\geq 10$ filings):} Freedom Smokeless LLC ($+1.65$); FRESH BAKED CO. ($+1.37$); Vaporfuel, LLC ($+1.26$)

\textbf{Worked examples}:

  \textbf{Positive flux:} \textsc{DAISY} (2016, $\Delta KL=+3.90$):\\
  \textit{Oral vaporizers for smoking purposes, namely, for smoking cannabis solely derived from hemp with a delta-9 tetrahydrocannabinol (THC) concentration of not more than 0.3 percent on a dry weight basis}\\
  \textbf{Flux-neutral:} \textsc{SATIVAH} (2010, $\Delta KL=+0.00$):\\
  \textit{Herbal smoking blend for cigarettes, cigars and pipes [ and smoking articles used therewith, namely, cigarette papers, cigar or cigarette lighters not made of precious metals and matches ]}\\
  \textbf{Negative flux:} \textsc{SITRA} (2021, $\Delta KL=-1.39$):\\
  \textit{Hookahs; hookah parts, namely, containers for housing charcoal separate from bowls while permitting heat transfer to tobacco placed in the bowls; smokers articles, namely, containers for housing hooka}\\

\subsection*{Class 035: Advertising \& Retail}
\label{app:industry-035}

\begin{tabular}{ll}
Clean filings: & 996,559 \\
Unique owners: & 584,964 \\
Year range: & 1990--2021 \\
Mean $KL_\text{pros}$: & 5.38 nats \\
Mean $KL_\text{retr}$: & 5.40 nats \\
Mean / median $\Delta KL$: & -0.019 / -0.027 \\
5y survival rate: & 47.0\% \\
Share with SEC-listed owner: & 3.65\% \\
\end{tabular}

\textbf{Top 3 firms by firm-mean $\Delta KL$ ($\geq 10$ filings):} Steven S. Lazar ($+1.17$); Beverly Hills Internet ($+1.08$); FINANZIARIA D'INVESTIMENTO - FININVEST SPA ($+0.75$)

\textbf{Worked examples}:

  \textbf{Positive flux:} \textsc{ACI AUDIO COMMUNICATIONS INCORPORATED} (1990, $\Delta KL=+3.44$):\\
  \textit{promoting the sale of the goods and/or services of others by recording telephone commercials and providing access thereto by the members of the purchasing public}\\
  \textbf{Flux-neutral:} \textsc{SAVE IT MAKES CENTS} (2014, $\Delta KL=-0.00$):\\
  \textit{Retail store services featuring convenience store items and gasoline}\\
  \textbf{Negative flux:} \textsc{RECLAIM YOUR SPACE} (2011, $\Delta KL=-1.51$):\\
  \textit{Providing organizing services for belongings and spaces in a residence or office such as papers and time}\\

\subsection*{Class 036: Insurance \& Financial}
\label{app:industry-036}

\begin{tabular}{ll}
Clean filings: & 399,896 \\
Unique owners: & 196,121 \\
Year range: & 1990--2021 \\
Mean $KL_\text{pros}$: & 4.92 nats \\
Mean $KL_\text{retr}$: & 4.96 nats \\
Mean / median $\Delta KL$: & -0.042 / -0.055 \\
5y survival rate: & 47.3\% \\
Share with SEC-listed owner: & 6.69\% \\
\end{tabular}

\textbf{Top 3 firms by firm-mean $\Delta KL$ ($\geq 10$ filings):} The Upper Deck Company ($+0.70$); Philip Lilien ($+0.69$); AMERICAN TELEPHONE AND TELEGRAPH COMPANY ($+0.67$)

\textbf{Worked examples}:

  \textbf{Positive flux:} \textsc{WE SPEAK YOUR LANGUAGE} (1990, $\Delta KL=+3.23$):\\
  \textit{insurance administration services in the field of health care}\\
  \textbf{Flux-neutral:} \textsc{SPACE WE'VE GOT SPACE TO FILL} (1999, $\Delta KL=+0.00$):\\
  \textit{Commercial real estate services in the nature of property leasing Commercial real estate services in the nature of commercial property development}\\
  \textbf{Negative flux:} \textsc{MOST-WANTED} (1992, $\Delta KL=-1.31$):\\
  \textit{fund-raising services for basic and applied research into neuromuscular disorders, and for funding of hospital services for persons afflicted by such disorders}\\

\subsection*{Class 037: Construction \& Repair}
\label{app:industry-037}

\begin{tabular}{ll}
Clean filings: & 178,756 \\
Unique owners: & 106,312 \\
Year range: & 1990--2021 \\
Mean $KL_\text{pros}$: & 5.18 nats \\
Mean $KL_\text{retr}$: & 5.18 nats \\
Mean / median $\Delta KL$: & +0.001 / +0.002 \\
5y survival rate: & 42.5\% \\
Share with SEC-listed owner: & 4.84\% \\
\end{tabular}

\textbf{Top 3 firms by firm-mean $\Delta KL$ ($\geq 10$ filings):} Aerus Concepts, L.P. ($+0.39$); Irvine Company, The ($+0.38$); Coastal Corporation, The ($+0.32$)

\textbf{Worked examples}:

  \textbf{Positive flux:} \textsc{MODEM EXPRESS} (1990, $\Delta KL=+2.53$):\\
  \textit{REPAIR SERVICES IN THE FIELD OF [COMPUTERS] AND COMPUTER PERIPHERALS [retail store services in the field of computers and computer peripherals]}\\
  \textbf{Flux-neutral:} \textsc{SONERI} (2015, $\Delta KL=-0.00$):\\
  \textit{Hotel services; Restaurant services Retail store services featuring groceries Financial advisory and consultancy services; Financial investment brokerage; Real estate agencies; Real estate brokerage B}\\
  \textbf{Negative flux:} \textsc{THE MAJORS} (2007, $\Delta KL=-1.21$):\\
  \textit{Homebuilding services; development of real property for single-family and multi-family residential and master-planned community purposes}\\

\subsection*{Class 038: Telecommunications}
\label{app:industry-038}

\begin{tabular}{ll}
Clean filings: & 170,826 \\
Unique owners: & 86,334 \\
Year range: & 1990--2021 \\
Mean $KL_\text{pros}$: & 4.32 nats \\
Mean $KL_\text{retr}$: & 4.37 nats \\
Mean / median $\Delta KL$: & -0.050 / -0.040 \\
5y survival rate: & 51.5\% \\
Share with SEC-listed owner: & 7.36\% \\
\end{tabular}

\textbf{Top 3 firms by firm-mean $\Delta KL$ ($\geq 10$ filings):} Los Angeles Cellular Telephone Company ($+0.90$); Iowa Network Services, Inc. ($+0.81$); Graphnet, Inc. ($+0.76$)

\textbf{Worked examples}:

  \textbf{Positive flux:} \textsc{CORPORATE TELEMANAGEMENT} (1990, $\Delta KL=+3.43$):\\
  \textit{telecommunication services, namely, providing long distance telephone services}\\
  \textbf{Flux-neutral:} \textsc{VZCREW} (2002, $\Delta KL=-0.00$):\\
  \textit{Telecommunication services, namely wireless telephone communication services and wireless messaging services providing online information in the field of entertainment}\\
  \textbf{Negative flux:} \textsc{SIMPLIFY THE CONNECTION} (1998, $\Delta KL=-1.07$):\\
  \textit{Long distance telecommunications services, namely, long-distance telephone calling services}\\

\subsection*{Class 039: Transport \& Logistics}
\label{app:industry-039}

\begin{tabular}{ll}
Clean filings: & 120,417 \\
Unique owners: & 67,340 \\
Year range: & 1990--2021 \\
Mean $KL_\text{pros}$: & 5.03 nats \\
Mean $KL_\text{retr}$: & 5.05 nats \\
Mean / median $\Delta KL$: & -0.020 / -0.017 \\
5y survival rate: & 45.4\% \\
Share with SEC-listed owner: & 6.90\% \\
\end{tabular}

\textbf{Top 3 firms by firm-mean $\Delta KL$ ($\geq 10$ filings):} Go America Tours, Inc. ($+0.67$); Coastal Corporation, The ($+0.55$); Sun Run Generation LLC ($+0.45$)

\textbf{Worked examples}:

  \textbf{Positive flux:} \textsc{JAL JAPAN AIR LINES A WORLD OF DIFFERENC} (1990, $\Delta KL=+2.68$):\\
  \textit{transportation of passengers and cargo by air; and storage services, namely the storage of cargo and luggage which has been or is to be transported by air}\\
  \textbf{Flux-neutral:} \textsc{VICTORIA CLIPPER} (2017, $\Delta KL=+0.00$):\\
  \textit{transportation of passengers by ship without on-board overnight accommodations}\\
  \textbf{Negative flux:} \textsc{YOUR ENERGY EDGE} (1999, $\Delta KL=-1.24$):\\
  \textit{public utility services, namely supplying natural gas to others (}\\

\subsection*{Class 040: Materials Treatment}
\label{app:industry-040}

\begin{tabular}{ll}
Clean filings: & 94,015 \\
Unique owners: & 59,227 \\
Year range: & 1990--2021 \\
Mean $KL_\text{pros}$: & 5.37 nats \\
Mean $KL_\text{retr}$: & 5.37 nats \\
Mean / median $\Delta KL$: & -0.004 / -0.011 \\
5y survival rate: & 44.6\% \\
Share with SEC-listed owner: & 5.75\% \\
\end{tabular}

\textbf{Top 3 firms by firm-mean $\Delta KL$ ($\geq 10$ filings):} Intervet Inc. ($+0.32$); United States Department of Energy ($+0.31$); Cherry Street Energy, LLC ($+0.30$)

\textbf{Worked examples}:

  \textbf{Positive flux:} \textsc{ETERNITY HANES HEAVYWEIGHT} (2007, $\Delta KL=+3.02$):\\
  \textit{Imprinting messages on wearing apparel, accessories and mugs}\\
  \textbf{Flux-neutral:} \textsc{MAGSHIRE BEL} (2020, $\Delta KL=-0.00$):\\
  \textit{Cryogenic preservation; Cryopreservation of human tissue; Cryopreservation services; Vitrification drying and dehydration preservation services for biologically active materials}\\
  \textbf{Negative flux:} \textsc{SUPERME!} (2010, $\Delta KL=-1.47$):\\
  \textit{Creation of custom portraits and paintings based on photographs, snapshots or images provided by the customer}\\

\subsection*{Class 041: Education \& Entertainment}
\label{app:industry-041}

\begin{tabular}{ll}
Clean filings: & 920,449 \\
Unique owners: & 531,807 \\
Year range: & 1990--2021 \\
Mean $KL_\text{pros}$: & 5.18 nats \\
Mean $KL_\text{retr}$: & 5.19 nats \\
Mean / median $\Delta KL$: & -0.006 / -0.016 \\
5y survival rate: & 47.0\% \\
Share with SEC-listed owner: & 3.06\% \\
\end{tabular}

\textbf{Top 3 firms by firm-mean $\Delta KL$ ($\geq 10$ filings):} CPI Church Street, Inc. ($+1.11$); Digital Gaming Corporation Limited ($+0.93$); Dapper Labs Inc. ($+0.86$)

\textbf{Worked examples}:

  \textbf{Positive flux:} \textsc{THE COASTERS} (1990, $\Delta KL=+3.88$):\\
  \textit{entertainment services in the nature of a musical group}\\
  \textbf{Flux-neutral:} \textsc{BEN MIYARES' PACKAGING MANAGEMENT UPDATE} (2005, $\Delta KL=+0.00$):\\
  \textit{Providing online electronic newsletters and trade journals in the field of packaging machinery, the packaging machinery industry and association membership affairs; providing electronic newsletters an}\\
  \textbf{Negative flux:} \textsc{MLE MAJOR LEAGUE EMPIRE} (2019, $\Delta KL=-1.29$):\\
  \textit{Entertainment services, namely, an on-line activity where you create your own music videos}\\

\subsection*{Class 042: Scientific \& Tech Services}
\label{app:industry-042}

\begin{tabular}{ll}
Clean filings: & 743,414 \\
Unique owners: & 411,644 \\
Year range: & 1990--2021 \\
Mean $KL_\text{pros}$: & 5.17 nats \\
Mean $KL_\text{retr}$: & 5.20 nats \\
Mean / median $\Delta KL$: & -0.022 / -0.006 \\
5y survival rate: & 46.1\% \\
Share with SEC-listed owner: & 5.61\% \\
\end{tabular}

\textbf{Top 3 firms by firm-mean $\Delta KL$ ($\geq 10$ filings):} Geniuses at Work Corporation ($+1.28$); Advance Publications, Inc. ($+0.98$); Albin, James R. ($+0.92$)

\textbf{Worked examples}:

  \textbf{Positive flux:} \textsc{PIPELINE} (1993, $\Delta KL=+3.47$):\\
  \textit{computer software for use in providing multiple-user access to a global computer information network for the transfer and dissemination of a wide range of information providing multiple-user access to}\\
  \textbf{Flux-neutral:} \textsc{IEEE: ADVANCING TECHNOLOGY FOR HUMANITY} (2009, $\Delta KL=-0.00$):\\
  \textit{Providing news and information in the field of science, engineering and engineering technology via the Internet; Accreditation services, namely, developing, evaluating, and testing standards with rega}\\
  \textbf{Negative flux:} \textsc{TAILGATE TEAM} (1996, $\Delta KL=-2.13$):\\
  \textit{retail grocery store and supermarket services}\\

\subsection*{Class 043: Restaurants \& Hotels}
\label{app:industry-043}

\begin{tabular}{ll}
Clean filings: & 136,699 \\
Unique owners: & 83,211 \\
Year range: & 2002--2021 \\
Mean $KL_\text{pros}$: & 4.87 nats \\
Mean $KL_\text{retr}$: & 4.85 nats \\
Mean / median $\Delta KL$: & +0.014 / +0.012 \\
5y survival rate: & 48.4\% \\
Share with SEC-listed owner: & 2.80\% \\
\end{tabular}

\textbf{Top 3 firms by firm-mean $\Delta KL$ ($\geq 10$ filings):} EX3 ($+0.42$); BRINKER INTERNATIONAL PAYROLL COMPANY, L.P. ($+0.37$); Dunhill 1530 Main, L.P. ($+0.34$)

\textbf{Worked examples}:

  \textbf{Positive flux:} \textsc{MARYLAND FRIED CHICKEN COMPLETE DINNERS } (2008, $\Delta KL=+2.39$):\\
  \textit{Restaurant services, including sit-down service of food and take-out restaurant services}\\
  \textbf{Flux-neutral:} \textsc{PANNA NEW LATINO FOOD} (2015, $\Delta KL=+0.00$):\\
  \textit{Café and restaurant services; Restaurant services featuring tequenos, cachitos, breads, sandwiches, hamburger, hotdogs, arepas, pastry, cachapas, empanadas, cheese sticks, salads, omelets, natural jui}\\
  \textbf{Negative flux:} \textsc{KIVO} (2007, $\Delta KL=-0.90$):\\
  \textit{Contract food services; preparing and serving foods in institutional dining facilities}\\

\subsection*{Class 044: Medical \& Personal Care}
\label{app:industry-044}

\begin{tabular}{ll}
Clean filings: & 188,794 \\
Unique owners: & 109,333 \\
Year range: & 2002--2021 \\
Mean $KL_\text{pros}$: & 4.96 nats \\
Mean $KL_\text{retr}$: & 4.95 nats \\
Mean / median $\Delta KL$: & +0.005 / -0.000 \\
5y survival rate: & 45.8\% \\
Share with SEC-listed owner: & 5.56\% \\
\end{tabular}

\textbf{Top 3 firms by firm-mean $\Delta KL$ ($\geq 10$ filings):} ClearChoice Holdings LLC ($+0.78$); Aspiration Partners, Inc. ($+0.53$); Bittner, Craig A. ($+0.43$)

\textbf{Worked examples}:

  \textbf{Positive flux:} \textsc{AMERICAN COLLECTIVE} (2009, $\Delta KL=+2.30$):\\
  \textit{Medical cannabis collective and dispensary services, namely, providing medical cannabis to medical cannabis patients}\\
  \textbf{Flux-neutral:} \textsc{NO BARRIERS. NO LIMITS. INDEPENDENCE FOR} (2019, $\Delta KL=+0.00$):\\
  \textit{Health care consulting in the field of occupational therapy; Occupational therapy services}\\
  \textbf{Negative flux:} \textsc{ROMAN HEMP CO.} (2021, $\Delta KL=-1.05$):\\
  \textit{Farming services in the field of hemp and hemp flower derived from hemp with a delta-9 tetrahydrocannabinol or a delta-8 tetrahydrocannabinol concentration of not more than 0.3 percent on a dry-weight}\\

\subsection*{Class 045: Legal \& Personal Services}
\label{app:industry-045}

\begin{tabular}{ll}
Clean filings: & 126,420 \\
Unique owners: & 81,811 \\
Year range: & 2002--2021 \\
Mean $KL_\text{pros}$: & 4.96 nats \\
Mean $KL_\text{retr}$: & 4.99 nats \\
Mean / median $\Delta KL$: & -0.022 / -0.022 \\
5y survival rate: & 45.2\% \\
Share with SEC-listed owner: & 3.47\% \\
\end{tabular}

\textbf{Top 3 firms by firm-mean $\Delta KL$ ($\geq 10$ filings):} YUGA LABS, INC. ($+0.56$); FamFriend, LLC ($+0.50$); Blend Images, LLC ($+0.39$)

\textbf{Worked examples}:

  \textbf{Positive flux:} \textsc{A COURSE IN MIRACLES MINISTER} (2014, $\Delta KL=+2.06$):\\
  \textit{Ecclesiastical services, namely, the licensing of ministers, namely, ordaining ministers to perform religious ceremonies}\\
  \textbf{Flux-neutral:} \textsc{ARTHUR MILLER JR ADNF FOUNDATION A DAUGH} (2017, $\Delta KL=+0.00$):\\
  \textit{Charitable outreach services, namely, providing counseling services in the field of interpersonal communication between law enforcement officers and community members who have faced tragedies or other}\\
  \textbf{Negative flux:} \textsc{NEFA FOUNDATION} (2008, $\Delta KL=-1.50$):\\
  \textit{Providing information regarding understanding and preventing terrorism}\\


\section{Per-industry deep dives: trending, ending, and public debuts}
\label{app:per-industry-deep}

This appendix complements Appendix~\ref{app:per-industry} with a
richer per-class snapshot. For each NICE class with at least
$5{,}000$ clean H=2y filings, we present:
\begin{itemize}
  \item Summary statistics and the SEC-listed share of filings.
  \item A scatter of all clean filings in the $(KL_\text{pros},
        KL_\text{retr})$ plane, with the $\Delta KL = 0$ diagonal
        and the top-10 SEC-listed firms whose first trademark filing
        occurred in 1990--2021 marked as blue stars and labeled.
  \item Top-10 \emph{trending} tokens (over-represented in the
        top-decile-$\Delta KL$ stratum within this class) with their
        intra-class 5y survival rates.
  \item Top-10 \emph{ending} tokens (over-represented in the
        bottom-decile-$\Delta KL$ stratum) with their intra-class
        survival rates.
  \item Top-10 flux-neutral tokens by frequency in the
        $|\Delta KL| \leq 0.02$ stratum, with survival rates.
  \item A table of the top-10 public firms with their debut filing's
        mark, year, $KL_\text{pros}$, $\Delta KL$, and a g/s excerpt.
\end{itemize}
``Top-10 public firms'' are ranked by maximum observed annual revenue
in the SEC EDGAR financial-statement panel. Firms whose primary
trademark class is the focal class are eligible; firms whose first
filing predates 1990 (i.e., they predated our observation window) are
excluded.

\subsection*{Class 001: Chemicals \& Industrial}
\label{app:deep-001}

\begin{tabular}{ll}
Clean filings (1990--2021): & 125,272 \\
Mean $KL_\text{pros}$ / $KL_\text{retr}$ / $\Delta KL$: & 5.36 / 5.37 / -0.009 nats \\
5y survival rate: & 47.1\% \\
Share with SEC-listed owner: & 6.47\% \\
\end{tabular}

\begin{center}
\includegraphics[width=0.7\linewidth]{results/per_industry_deep/class001_scatter.png}
\end{center}

\textbf{Trending tokens (top-decile $\Delta KL$)} --- with intra-class 5y survival rate.

\begin{tabular}{lrrr}
\toprule
token & lift & $n$ filings & intra-class survival \\
\midrule
\texttt{ampule} & 4.31$\times$ & 144 & 97\% \\
\texttt{delta-} & 4.12$\times$ & 391 & 63\% \\
\texttt{cannabis} & 4.04$\times$ & 223 & 62\% \\
\texttt{tetrahydrocannabinol} & 3.93$\times$ & 351 & 67\% \\
\texttt{percent} & 3.91$\times$ & 391 & 63\% \\
\texttt{thc} & 3.81$\times$ & 420 & 63\% \\
\texttt{tincture} & 3.71$\times$ & 170 & 88\% \\
\texttt{hemp} & 3.36$\times$ & 711 & 64\% \\
\texttt{cbd} & 3.35$\times$ & 379 & 70\% \\
\texttt{biomaterials} & 3.33$\times$ & 144 & 43\% \\
\bottomrule
\end{tabular}

\textbf{Ending tokens (bottom-decile $\Delta KL$)} --- intra-class 5y survival rate.

\begin{tabular}{lrrr}
\toprule
token & lift & $n$ filings & intra-class survival \\
\midrule
\texttt{e-liquid} & 6.03$\times$ & 375 & 54\% \\
\texttt{non-synthetic} & 6.02$\times$ & 211 & 16\% \\
\texttt{odorants} & 5.74$\times$ & 237 & 19\% \\
\texttt{bleaches} & 4.91$\times$ & 277 & 24\% \\
\texttt{brighteners} & 4.82$\times$ & 274 & 24\% \\
\texttt{colognes} & 4.24$\times$ & 330 & 34\% \\
\texttt{glycerin} & 3.91$\times$ & 486 & 47\% \\
\texttt{refill} & 3.51$\times$ & 276 & 57\% \\
\texttt{cigarette} & 3.39$\times$ & 814 & 51\% \\
\texttt{decalcifying} & 3.35$\times$ & 203 & 61\% \\
\bottomrule
\end{tabular}

\textbf{Flux-neutral most-frequent tokens ($|\Delta KL|\leq 0.02$)}.

\begin{tabular}{lrrr}
\toprule
token & $n$ in stratum & $n$ in class & intra-class survival \\
\midrule
\texttt{use} & 8,388 & 79,491 & 44\% \\
\texttt{chemical} & 5,379 & 46,269 & 41\% \\
\texttt{chemicals} & 4,406 & 35,839 & 39\% \\
\texttt{preparations} & 3,746 & 31,701 & 38\% \\
\texttt{manufacture} & 3,552 & 31,197 & 42\% \\
\texttt{industrial} & 3,456 & 26,413 & 36\% \\
\texttt{products} & 3,125 & 25,459 & 40\% \\
\texttt{purposes} & 3,068 & 24,195 & 35\% \\
\texttt{industry} & 2,672 & 21,184 & 32\% \\
\texttt{materials} & 2,251 & 15,521 & 35\% \\
\bottomrule
\end{tabular}

\textbf{Top-10 SEC-listed firms with debut in 1990--2021} (ranked by max observed revenue, debut filing g/s shown).

\begin{tabular}{p{4.0cm}p{2.2cm}rrrp{6.5cm}}
\toprule
Firm & Mark & year & $KL_\text{pros}$ & $\Delta KL$ & First-filing g/s excerpt \\
\midrule
Cellular Concrete Technologies, Inc. & FOAMCON & 1994 & 5.98 & +0.19 & foam concentrate for use in making cellular concrete, cement-based insulation and concrete building products \\
Compass Therapeutics LLC & CURING CANCER THROUGH COMBINAT & 2015 & 4.76 & +0.25 & Custom design and development of antibody therapeutic preparations Antibody therapeutic preparations for medical purposes Antibodies for scientific and research purposes \\
Global Enterprises & STEEL SEAL & 1993 & 6.93 & -0.25 & tire puncture sealing compounds \\
KX Technologies LLC & FINITY & 2013 & 4.29 & -0.11 & chemical and mineral fiber pulp preparations in the nature of filtering materials for the manufacturing of domestic, industrial, and commercial water filters, and for the manufacture of commercial and \\
Syntonix Pharmaceuticals, Inc. & TRANSCEPTOR & 2002 & 4.01 & -0.05 & scientific research and development services for others in the field of pharmaceutical preparations biological agents used in the manufacture of pharmaceutical preparations \\
Z Corporation & Z-CURE & 2008 & 6.71 & +0.17 & Chemical in liquid form used as an infiltrant or sealant agent in the manufacture of three-dimensional models; and Chemical preparations in liquid form for use as an infiltrant or sealant agent in the \\
WAL-MART STORES, INC. & SAM'S AMERICAN CHOICE & 1996 & 4.91 & -0.04 & pretzels and pretzels with peanut butter filling; pasta; macaroni and cheese; cheese flavored snack foods, namely, puffed corn snacks, popped popcorn, and corn chips; cookies; popped popcorn; chocolat \\
Wal-Mart Stores, Inc & GREAT VALUE & 2016 & 5.46 & -0.15 & algaecides Chlorine; shock treatments, namely, water treatment chemicals for use in swimming pools or spas \\
Wal-Mart Stores, Inc. & GARDEN BASICS & 1997 & 5.72 & -0.01 & flower pot holders and planters fertilizer, potting soil, and soil conditioners for domestic use bagged and potted live plants, flower bulbs, and gardening seeds gardening tools, namely, spreaders, sh \\
EXXON MOBIL CORPORATION & OPTEMA & 1990 & 5.32 & +0.28 & plastic resins for use in the packaging, rubber and textile industries \\
\bottomrule
\end{tabular}
\clearpage

\subsection*{Class 002: Paints \& Coatings}
\label{app:deep-002}

\begin{tabular}{ll}
Clean filings (1990--2021): & 33,653 \\
Mean $KL_\text{pros}$ / $KL_\text{retr}$ / $\Delta KL$: & 4.78 / 4.79 / -0.012 nats \\
5y survival rate: & 44.6\% \\
Share with SEC-listed owner: & 4.40\% \\
\end{tabular}

\begin{center}
\includegraphics[width=0.7\linewidth]{results/per_industry_deep/class002_scatter.png}
\end{center}

\textbf{Trending tokens (top-decile $\Delta KL$)} --- with intra-class 5y survival rate.

\begin{tabular}{lrrr}
\toprule
token & lift & $n$ filings & intra-class survival \\
\midrule
\texttt{photocopiers} & 5.42$\times$ & 1,398 & 21\% \\
\texttt{photocopier} & 5.19$\times$ & 801 & 27\% \\
\texttt{filled} & 4.81$\times$ & 2,270 & 26\% \\
\texttt{copiers} & 4.71$\times$ & 753 & 27\% \\
\texttt{inkjet} & 4.57$\times$ & 937 & 24\% \\
\texttt{toner} & 4.26$\times$ & 2,707 & 32\% \\
\texttt{thermochromic} & 4.18$\times$ & 268 & 15\% \\
\texttt{capture} & 4.17$\times$ & 429 & 18\% \\
\texttt{multifunction} & 4.15$\times$ & 414 & 20\% \\
\texttt{skin-dressing} & 4.10$\times$ & 139 & 27\% \\
\bottomrule
\end{tabular}

\textbf{Ending tokens (bottom-decile $\Delta KL$)} --- intra-class 5y survival rate.

\begin{tabular}{lrrr}
\toprule
token & lift & $n$ filings & intra-class survival \\
\midrule
\texttt{surfacers} & 3.11$\times$ & 61 & 36\% \\
\texttt{interior} & 2.41$\times$ & 4,299 & 47\% \\
\texttt{exterior} & 2.14$\times$ & 4,828 & 48\% \\
\texttt{latex} & 1.97$\times$ & 487 & 42\% \\
\texttt{character} & 1.96$\times$ & 148 & 32\% \\
\texttt{pigmented} & 1.95$\times$ & 1,250 & 48\% \\
\texttt{labelling} & 1.94$\times$ & 62 & 21\% \\
\texttt{courts} & 1.91$\times$ & 68 & 43\% \\
\texttt{equipped} & 1.90$\times$ & 58 & 19\% \\
\texttt{nozzle} & 1.88$\times$ & 69 & 32\% \\
\bottomrule
\end{tabular}

\textbf{Flux-neutral most-frequent tokens ($|\Delta KL|\leq 0.02$)}.

\begin{tabular}{lrrr}
\toprule
token & $n$ in stratum & $n$ in class & intra-class survival \\
\midrule
\texttt{use} & 1,396 & 15,797 & 41\% \\
\texttt{coatings} & 1,284 & 14,341 & 45\% \\
\texttt{nature} & 1,041 & 9,985 & 40\% \\
\texttt{paints} & 1,008 & 11,058 & 38\% \\
\texttt{coating} & 848 & 8,607 & 42\% \\
\texttt{paint} & 746 & 8,935 & 44\% \\
\texttt{industrial} & 718 & 6,496 & 37\% \\
\texttt{metal} & 577 & 5,160 & 39\% \\
\texttt{preparations} & 573 & 4,896 & 32\% \\
\texttt{wood} & 549 & 6,076 & 38\% \\
\bottomrule
\end{tabular}

\textbf{Top-10 SEC-listed firms with debut in 1990--2021} (ranked by max observed revenue, debut filing g/s shown).

\begin{tabular}{p{4.0cm}p{2.2cm}rrrp{6.5cm}}
\toprule
Firm & Mark & year & $KL_\text{pros}$ & $\Delta KL$ & First-filing g/s excerpt \\
\midrule
SOLA INTERNATIONAL INC. & INTERIOR & 1998 & 6.07 & +1.25 & coatings for ophthalmic lenses, anti-reflective coatings for ophthalmic lenses, abrasion resistant coatings for ophthalmic lenses optical goods and apparatus, namely, spectacles, sunglasses, safety go \\
Sola International Inc. & GLARE OUT & 1998 & 4.49 & -1.17 & coatings for ophthalmic lenses, anti-reflective coatings for ophthalmic lenses, abrasion resistant coatings for ophthalmic lenses; impact resistant coatings for ophthalmic lenses; anti-fog coatings fo \\
WAL-MART STORES, INC. & CURIOUSLY J & 2017 & 2.08 & -0.08 & Dyes for use in the manufacture of textiles and fabrics; food coloring; paints; finger paint; inks, toners, and dyes for use with or on textiles and fabrics; edible food coloring Alcoholic cocktail mi \\
Wal-Mart Stores, Inc & DAILY'S & 2000 & 3.29 & -0.14 & Cooking appliances, namely, electric Dutch ovens, food steamers, electric skillets, fryers, microwaves, toasters, breadmakers, popcorn poppers, waffle irons, electric coffee pots, coffee makers; elect \\
Wal-Mart Stores, Inc. & COLOR PLACE & 1994 & 4.01 & -0.30 & house paint, interior and exterior \\
Wal-mart Stores, Inc. &  & 2011 & 3.89 & -0.20 & On-line retail store services featuring a wide variety of consumer goods of others; Retail department store services; Retail store services featuring a wide variety of consumer goods of others Air bru \\
EXXON MOBIL CORPORATION &  & 2006 & 4.89 & +0.11 & Anti-rust greases; Anti-rust oils; Rust preventatives in the nature of a coating \\
GENERAL MOTORS LLC & SATURN & 1992 & 4.89 & -0.13 & paint for motor vehicles \\
General Motors Corporation & WHEN THE RIGHT WAY IS THE ONLY & 2000 & 5.79 & +0.37 & Paint for motor land vehicles \\
General Motors LLC & ENVIRONMINT & 2009 & 6.46 & -0.19 & Paint for use in the manufacture of automobiles \\
\bottomrule
\end{tabular}
\clearpage

\subsection*{Class 003: Cosmetics \& Cleaning}
\label{app:deep-003}

\begin{tabular}{ll}
Clean filings (1990--2021): & 326,430 \\
Mean $KL_\text{pros}$ / $KL_\text{retr}$ / $\Delta KL$: & 4.32 / 4.31 / +0.005 nats \\
5y survival rate: & 50.0\% \\
Share with SEC-listed owner: & 2.94\% \\
\end{tabular}

\begin{center}
\includegraphics[width=0.7\linewidth]{results/per_industry_deep/class003_scatter.png}
\end{center}

\textbf{Trending tokens (top-decile $\Delta KL$)} --- with intra-class 5y survival rate.

\begin{tabular}{lrrr}
\toprule
token & lift & $n$ filings & intra-class survival \\
\midrule
\texttt{e-liquid} & 4.04$\times$ & 800 & 76\% \\
\texttt{vaporizer} & 3.83$\times$ & 399 & 80\% \\
\texttt{attaching} & 3.20$\times$ & 1,286 & 47\% \\
\texttt{cannabis} & 3.07$\times$ & 1,158 & 64\% \\
\texttt{refill} & 3.05$\times$ & 780 & 64\% \\
\texttt{affixing} & 2.96$\times$ & 5,356 & 33\% \\
\texttt{non-veterinary} & 2.92$\times$ & 930 & 40\% \\
\texttt{trims} & 2.91$\times$ & 502 & 24\% \\
\texttt{anxiety} & 2.80$\times$ & 521 & 63\% \\
\texttt{cigarettes} & 2.75$\times$ & 1,463 & 62\% \\
\bottomrule
\end{tabular}

\textbf{Ending tokens (bottom-decile $\Delta KL$)} --- intra-class 5y survival rate.

\begin{tabular}{lrrr}
\toprule
token & lift & $n$ filings & intra-class survival \\
\midrule
\texttt{post-suntanning} & 6.28$\times$ & 492 & 90\% \\
\texttt{pre-suntanning} & 5.81$\times$ & 530 & 89\% \\
\texttt{degreaser} & 4.64$\times$ & 474 & 57\% \\
\texttt{cuticles} & 4.09$\times$ & 575 & 76\% \\
\texttt{straightener} & 3.80$\times$ & 532 & 76\% \\
\texttt{lightener} & 3.79$\times$ & 755 & 79\% \\
\texttt{sunless} & 3.78$\times$ & 1,352 & 72\% \\
\texttt{astringent} & 3.67$\times$ & 1,562 & 68\% \\
\texttt{shimmer} & 3.47$\times$ & 525 & 56\% \\
\texttt{soak} & 3.43$\times$ & 937 & 63\% \\
\bottomrule
\end{tabular}

\textbf{Flux-neutral most-frequent tokens ($|\Delta KL|\leq 0.02$)}.

\begin{tabular}{lrrr}
\toprule
token & $n$ in stratum & $n$ in class & intra-class survival \\
\midrule
\texttt{preparations} & 22,153 & 174,132 & 47\% \\
\texttt{care} & 19,078 & 148,035 & 51\% \\
\texttt{skin} & 17,196 & 138,329 & 50\% \\
\texttt{hair} & 15,788 & 117,462 & 48\% \\
\texttt{body} & 14,323 & 117,200 & 51\% \\
\texttt{lotions} & 14,120 & 108,293 & 49\% \\
\texttt{non-medicated} & 13,803 & 111,213 & 48\% \\
\texttt{creams} & 13,345 & 101,302 & 49\% \\
\texttt{cosmetics} & 13,266 & 110,457 & 46\% \\
\texttt{cosmetic} & 12,597 & 99,440 & 45\% \\
\bottomrule
\end{tabular}

\textbf{Top-10 SEC-listed firms with debut in 1990--2021} (ranked by max observed revenue, debut filing g/s shown).

\begin{tabular}{p{4.0cm}p{2.2cm}rrrp{6.5cm}}
\toprule
Firm & Mark & year & $KL_\text{pros}$ & $\Delta KL$ & First-filing g/s excerpt \\
\midrule
ACM & VITICOLOR & 2008 & 4.33 & -0.09 & Cosmetics; cosmetic products, namely, self-tanning lotions; cosmetic preparations for sun-tanning; skin creams; foundation for evening out skin tone; cosmetics, namely, pigmentation regulators \\
ATK Entertainment, Inc. & AALAM TUCK KHALED & 2020 & 2.78 & -0.07 & Clothing, namely, shirts, t-shirts, pants, dresses, skirts, jeans, slacks, belts, hats, caps being headwear, gloves, ear muffs, headwear, neckwear, sweaters, fleece pullovers, jogging suits, warm-up s \\
BETTERLIFE PHARMA INC. & SLEEP BETTERX & 2019 & 2.61 & -0.17 & Printed publications, namely, newsletters, brochures, reports and guides in the field of cannabis bath additives, namely, bath beads, bath bombs, bath crystals and bubble bath; bath herbs; bath oils;  \\
BetterLife Pharma Inc. & LIVE BETTERX & 2020 & 2.47 & -0.19 & Downloadable electronic publications, namely, newsletters, brochures, reports and guides in the field of cannabis Retail store and wholesale store services featuring industrial hemp food and body care \\
ECO PLANET INC. & ORGANIC OSCAR & 2009 & 5.00 & +0.12 & Organic pet care kits comprising shampoo, conditioner, body spray sold as a unit \\
Global Cosmetic Company & SEA ELEMENTS & 2003 & 5.27 & -0.14 & Hair Shampoo, Hair conditioners ,Hair  styling gels \\
Global Cosmetics Group Pty Ltd & COUTURE SILK & 2012 & 3.20 & -0.06 & Hair coloring preparations; hair color; hair dyes; hair care preparations; [ hair shampoo; hair conditioner; ] all of the aforementioned containing a proportion of silk extract; bleaching preparations \\
MY CO., LTD. & GOEE & 2020 & 3.18 & -0.04 & Under-eye creams; shampoos for the care of scalp; non-medicated toners for scalp; lip balm; cosmetic preparations for the hair and scalp; skin moisturizers; skin moisturizer masks; skin cleansing loti \\
Plantation Holdings Group, Inc. & THE PLANTATION SHOP & 2004 & 5.00 & +0.07 & Gourmet cookies, coffee, wafers, biscuits, snack nut mix consisting primarily of chocolate covered nuts and pralines, candy, toffee and cheese straws Jewelry Retail store services and retail shop at h \\
SM Global Enterprises, Inc. & ORIGINAL SHAVE COMPANY & 2017 & 3.01 & -0.12 & On-line wholesale and retail store services featuring Grooming supplies for Men and Women After shave lotions; after-shave; after-shave balms; after-shave creams; after-shave emulsions; after-shave ge \\
\bottomrule
\end{tabular}
\clearpage

\subsection*{Class 004: Lubricants \& Fuels}
\label{app:deep-004}

\begin{tabular}{ll}
Clean filings (1990--2021): & 36,160 \\
Mean $KL_\text{pros}$ / $KL_\text{retr}$ / $\Delta KL$: & 4.91 / 4.91 / +0.004 nats \\
5y survival rate: & 47.8\% \\
Share with SEC-listed owner: & 5.58\% \\
\end{tabular}

\begin{center}
\includegraphics[width=0.7\linewidth]{results/per_industry_deep/class004_scatter.png}
\end{center}

\textbf{Trending tokens (top-decile $\Delta KL$)} --- with intra-class 5y survival rate.

\begin{tabular}{lrrr}
\toprule
token & lift & $n$ filings & intra-class survival \\
\midrule
\texttt{clamshell} & 5.07$\times$ & 284 & 43\% \\
\texttt{bougies} & 4.33$\times$ & 240 & 43\% \\
\texttt{mini-pumpkins} & 4.28$\times$ & 339 & 40\% \\
\texttt{carved} & 4.24$\times$ & 354 & 38\% \\
\texttt{pumpkins} & 4.10$\times$ & 356 & 41\% \\
\texttt{gourds} & 4.04$\times$ & 359 & 38\% \\
\texttt{unity} & 3.89$\times$ & 311 & 40\% \\
\texttt{contained} & 3.70$\times$ & 424 & 46\% \\
\texttt{votive} & 3.48$\times$ & 1,160 & 43\% \\
\texttt{tealight} & 3.36$\times$ & 776 & 36\% \\
\bottomrule
\end{tabular}

\textbf{Ending tokens (bottom-decile $\Delta KL$)} --- intra-class 5y survival rate.

\begin{tabular}{lrrr}
\toprule
token & lift & $n$ filings & intra-class survival \\
\midrule
\texttt{briquets} & 7.00$\times$ & 60 & 43\% \\
\texttt{reconstructor} & 6.79$\times$ & 53 & 89\% \\
\texttt{mattifier} & 6.60$\times$ & 53 & 87\% \\
\texttt{luster} & 6.42$\times$ & 53 & 89\% \\
\texttt{exfoliator} & 5.38$\times$ & 80 & 82\% \\
\texttt{tint} & 5.33$\times$ & 75 & 77\% \\
\texttt{corrector} & 5.30$\times$ & 66 & 79\% \\
\texttt{straightener} & 5.16$\times$ & 62 & 84\% \\
\texttt{detangler} & 5.00$\times$ & 78 & 79\% \\
\texttt{polisher} & 4.89$\times$ & 139 & 86\% \\
\bottomrule
\end{tabular}

\textbf{Flux-neutral most-frequent tokens ($|\Delta KL|\leq 0.02$)}.

\begin{tabular}{lrrr}
\toprule
token & $n$ in stratum & $n$ in class & intra-class survival \\
\midrule
\texttt{use} & 1,378 & 14,053 & 43\% \\
\texttt{candles} & 1,353 & 15,757 & 46\% \\
\texttt{oils} & 1,204 & 12,471 & 42\% \\
\texttt{oil} & 923 & 10,398 & 46\% \\
\texttt{industrial} & 854 & 8,707 & 38\% \\
\texttt{fuel} & 799 & 8,302 & 44\% \\
\texttt{lubricants} & 751 & 8,599 & 41\% \\
\texttt{preparations} & 709 & 6,566 & 40\% \\
\texttt{purposes} & 706 & 5,686 & 34\% \\
\texttt{products} & 688 & 6,630 & 45\% \\
\bottomrule
\end{tabular}

\textbf{Top-10 SEC-listed firms with debut in 1990--2021} (ranked by max observed revenue, debut filing g/s shown).

\begin{tabular}{p{4.0cm}p{2.2cm}rrrp{6.5cm}}
\toprule
Firm & Mark & year & $KL_\text{pros}$ & $\Delta KL$ & First-filing g/s excerpt \\
\midrule
AGE International, Inc. & BLANTON'S & 2021 & 4.32 & -0.05 & wall mounted coat rack; wooden display stands for bottle stoppers, bottles and tasting glasses; wall art made of wood reclaimed from bourbon barrels face coverings being masks for protection against v \\
TGI SOLAR POWER GROUP INC. & ADVENT & 2021 & 4.71 & +0.18 & Batteries; Batteries and battery chargers; Battery boxes; Solar panels for the production of electricity; Grids for batteries; Off-grid power and deep cycle battery storage systems comprised of batter \\
WAL-MART STORES, INC. & CURIOUSLY J & 2017 & 2.02 & -0.06 & Dyes for use in the manufacture of textiles and fabrics; food coloring; paints; finger paint; inks, toners, and dyes for use with or on textiles and fabrics; edible food coloring Alcoholic cocktail mi \\
Wal-Mart Stores, Inc & DAILY'S & 2000 & 2.98 & -0.01 & Cooking appliances, namely, electric Dutch ovens, food steamers, electric skillets, fryers, microwaves, toasters, breadmakers, popcorn poppers, waffle irons, electric coffee pots, coffee makers; elect \\
Wal-Mart Stores, Inc. & TECH 2000 & 1995 & 5.01 & -0.73 & automotive oils and lubricants \\
Wal-mart Stores, Inc. & BACKYARD GR LL & 2011 & 4.85 & +0.06 & Barbecue smokers; Barbecues and grills; Ceramic briquettes for use in barbecue grills; Charcoal grills; Gas grills; Lighters primarily for lighting grills, fireplaces and candles; Roasters powered by  \\
EXXON MOBIL CORPORATION &  & 1991 & 5.97 & +0.20 & hydraulic oils \\
Exxon Mobil Corporation & E+2 & 1998 & 4.55 & -0.41 & fuel for motor vehicles, namely, gasoline \\
CHEVRON CORPORATION & PRO-GARD & 1990 & 6.97 & +1.28 & all purpose aerosol lubricant \\
Chevron Corporation & MECHANISM LPS & 1991 & 6.04 & +0.35 & hydraulic oil \\
\bottomrule
\end{tabular}
\clearpage

\subsection*{Class 005: Pharmaceuticals}
\label{app:deep-005}

\begin{tabular}{ll}
Clean filings (1990--2021): & 359,560 \\
Mean $KL_\text{pros}$ / $KL_\text{retr}$ / $\Delta KL$: & 4.98 / 4.99 / -0.011 nats \\
5y survival rate: & 56.4\% \\
Share with SEC-listed owner: & 10.60\% \\
\end{tabular}

\begin{center}
\includegraphics[width=0.7\linewidth]{results/per_industry_deep/class005_scatter.png}
\end{center}

\textbf{Trending tokens (top-decile $\Delta KL$)} --- with intra-class 5y survival rate.

\begin{tabular}{lrrr}
\toprule
token & lift & $n$ filings & intra-class survival \\
\midrule
\texttt{cannabis} & 3.40$\times$ & 1,393 & 78\% \\
\texttt{marijuana} & 3.39$\times$ & 366 & 84\% \\
\texttt{e-liquid} & 3.18$\times$ & 449 & 83\% \\
\texttt{vaporizer} & 3.08$\times$ & 419 & 82\% \\
\texttt{opioid} & 2.94$\times$ & 405 & 76\% \\
\texttt{cannabidiol} & 2.90$\times$ & 1,139 & 89\% \\
\texttt{delta-} & 2.88$\times$ & 2,174 & 84\% \\
\texttt{hand-sanitizing} & 2.85$\times$ & 3,404 & 54\% \\
\texttt{occurring} & 2.81$\times$ & 544 & 78\% \\
\texttt{treats} & 2.67$\times$ & 2,685 & 60\% \\
\bottomrule
\end{tabular}

\textbf{Ending tokens (bottom-decile $\Delta KL$)} --- intra-class 5y survival rate.

\begin{tabular}{lrrr}
\toprule
token & lift & $n$ filings & intra-class survival \\
\midrule
\texttt{antihistamine} & 4.89$\times$ & 380 & 61\% \\
\texttt{herbicide} & 4.16$\times$ & 418 & 57\% \\
\texttt{anti-virals} & 3.77$\times$ & 1,714 & 96\% \\
\texttt{percent} & 3.67$\times$ & 1,829 & 84\% \\
\texttt{tetrahydrocannabinol} & 3.66$\times$ & 1,681 & 85\% \\
\texttt{immunosuppressants} & 3.54$\times$ & 1,555 & 95\% \\
\texttt{anti-fungals} & 3.50$\times$ & 1,748 & 91\% \\
\texttt{delta-} & 3.41$\times$ & 2,174 & 84\% \\
\texttt{fungicide} & 3.31$\times$ & 680 & 59\% \\
\texttt{freshener} & 3.28$\times$ & 1,422 & 61\% \\
\bottomrule
\end{tabular}

\textbf{Flux-neutral most-frequent tokens ($|\Delta KL|\leq 0.02$)}.

\begin{tabular}{lrrr}
\toprule
token & $n$ in stratum & $n$ in class & intra-class survival \\
\midrule
\texttt{supplements} & 19,715 & 137,035 & 53\% \\
\texttt{dietary} & 16,936 & 110,853 & 53\% \\
\texttt{nutritional} & 16,700 & 102,677 & 54\% \\
\texttt{preparations} & 15,581 & 142,819 & 55\% \\
\texttt{use} & 13,562 & 120,130 & 53\% \\
\texttt{pharmaceutical} & 10,172 & 97,351 & 57\% \\
\texttt{medical} & 9,812 & 75,283 & 47\% \\
\texttt{treatment} & 9,176 & 87,149 & 57\% \\
\texttt{purposes} & 8,697 & 62,259 & 45\% \\
\texttt{food} & 7,472 & 55,012 & 50\% \\
\bottomrule
\end{tabular}

\textbf{Top-10 SEC-listed firms with debut in 1990--2021} (ranked by max observed revenue, debut filing g/s shown).

\begin{tabular}{p{4.0cm}p{2.2cm}rrrp{6.5cm}}
\toprule
Firm & Mark & year & $KL_\text{pros}$ & $\Delta KL$ & First-filing g/s excerpt \\
\midrule
1PM Industries, Inc. & AQUATRIM & 2015 & 4.27 & -0.13 & Dietary and nutritional supplements used for weight loss \\
ACM & MOLUTREX & 2010 & 5.15 & -0.11 & Pharmaceutical products for the treatment of verrucae \\
ATK Entertainment, Inc. & AALAM KHALED & 2020 & 3.36 & -0.06 & Clothing, namely, shirts, t-shirts, pants, dresses, skirts, jeans, slacks, belts, hats, caps being headwear, gloves, ear muffs, headwear, neckwear, sweaters, fleece pullovers, jogging suits, warm-up s \\
BETTERLIFE PHARMA INC. & THINK BETTERX & 2019 & 2.64 & -0.27 & Medicated oils, salves, and concentrated pastes, namely, medicated skin care preparations; Medicated oils, salves, and concentrated pastes, namely, medicated skin care preparations, all of the foregoi \\
Better For You Wellness, Inc. & BETTER FOR YOU WELLNESS & 2021 & 4.43 & -0.05 & Protein bars, namely, protein supplements formed and packaged as bars Advisory and consultation services in the fields of agronomy, agriculture, and horticulture; agronomic consulting services; agrono \\
CleanGoal Energy, Corp. & ASTALABS & 2018 & 3.73 & +0.01 & Dietary and nutritional supplements \\
CohBar, Inc. & COHBAR & 2014 & 4.34 & +0.24 & Pharmaceutical and medical preparations for the prevention and treatment of major diseases, namely, diabetes, cancer, Alzheimer's disease and atherosclerosis; peptide based medical therapeutic prepara \\
Compass Therapeutics LLC & CURING CANCER THROUGH COMBINAT & 2015 & 4.60 & +0.23 & Custom design and development of antibody therapeutic preparations Antibody therapeutic preparations for medical purposes Antibodies for scientific and research purposes \\
IPR Pharmaceuticals, Inc. &  & 2001 & 4.97 & +0.25 & Pharmaceutical and medicinal preparations and substances for the treatment of migraines \\
MIRA Pharmaceuticals, Inc. & MIRA-1A POT PILL & 2021 & 6.27 & -0.19 & Therapeutic pharmaceutical for the treatment of anxiety \\
\bottomrule
\end{tabular}
\clearpage

\subsection*{Class 006: Common Metals}
\label{app:deep-006}

\begin{tabular}{ll}
Clean filings (1990--2021): & 97,003 \\
Mean $KL_\text{pros}$ / $KL_\text{retr}$ / $\Delta KL$: & 5.48 / 5.46 / +0.020 nats \\
5y survival rate: & 43.4\% \\
Share with SEC-listed owner: & 4.04\% \\
\end{tabular}

\begin{center}
\includegraphics[width=0.7\linewidth]{results/per_industry_deep/class006_scatter.png}
\end{center}

\textbf{Trending tokens (top-decile $\Delta KL$)} --- with intra-class 5y survival rate.

\begin{tabular}{lrrr}
\toprule
token & lift & $n$ filings & intra-class survival \\
\midrule
\texttt{kept} & 4.96$\times$ & 270 & 26\% \\
\texttt{shutterboxes} & 4.72$\times$ & 572 & 23\% \\
\texttt{add} & 4.47$\times$ & 170 & 31\% \\
\texttt{pistol} & 4.36$\times$ & 307 & 28\% \\
\texttt{reach} & 4.30$\times$ & 316 & 28\% \\
\texttt{hat-hanging} & 4.00$\times$ & 562 & 19\% \\
\texttt{flashlight} & 3.82$\times$ & 361 & 32\% \\
\texttt{multi-point} & 3.82$\times$ & 144 & 36\% \\
\texttt{lockboxes} & 3.78$\times$ & 164 & 35\% \\
\texttt{slats} & 3.74$\times$ & 752 & 26\% \\
\bottomrule
\end{tabular}

\textbf{Ending tokens (bottom-decile $\Delta KL$)} --- intra-class 5y survival rate.

\begin{tabular}{lrrr}
\toprule
token & lift & $n$ filings & intra-class survival \\
\midrule
\texttt{non-leather} & 4.37$\times$ & 288 & 74\% \\
\texttt{yo-yos} & 4.20$\times$ & 188 & 67\% \\
\texttt{replica} & 4.10$\times$ & 144 & 55\% \\
\texttt{stationery-type} & 3.93$\times$ & 163 & 74\% \\
\texttt{participating} & 3.77$\times$ & 106 & 48\% \\
\texttt{steins} & 3.77$\times$ & 199 & 67\% \\
\texttt{sweatpants} & 3.76$\times$ & 473 & 60\% \\
\texttt{billfolds} & 3.75$\times$ & 224 & 77\% \\
\texttt{basketballs} & 3.67$\times$ & 354 & 65\% \\
\texttt{afghans} & 3.61$\times$ & 97 & 71\% \\
\bottomrule
\end{tabular}

\textbf{Flux-neutral most-frequent tokens ($|\Delta KL|\leq 0.02$)}.

\begin{tabular}{lrrr}
\toprule
token & $n$ in stratum & $n$ in class & intra-class survival \\
\midrule
\texttt{metal} & 9,062 & 84,324 & 42\% \\
\texttt{use} & 3,826 & 28,315 & 37\% \\
\texttt{parts} & 2,612 & 16,917 & 33\% \\
\texttt{materials} & 2,390 & 15,503 & 31\% \\
\texttt{fittings} & 1,977 & 15,163 & 29\% \\
\texttt{pipes} & 1,942 & 12,554 & 28\% \\
\texttt{nature} & 1,922 & 13,088 & 30\% \\
\texttt{building} & 1,915 & 13,124 & 32\% \\
\texttt{plastic} & 1,907 & 13,269 & 40\% \\
\texttt{hardware} & 1,886 & 15,693 & 32\% \\
\bottomrule
\end{tabular}

\textbf{Top-10 SEC-listed firms with debut in 1990--2021} (ranked by max observed revenue, debut filing g/s shown).

\begin{tabular}{p{4.0cm}p{2.2cm}rrrp{6.5cm}}
\toprule
Firm & Mark & year & $KL_\text{pros}$ & $\Delta KL$ & First-filing g/s excerpt \\
\midrule
AGE International, Inc. & BLANTON'S & 2021 & 4.79 & -0.09 & wall mounted coat rack; wooden display stands for bottle stoppers, bottles and tasting glasses; wall art made of wood reclaimed from bourbon barrels face coverings being masks for protection against v \\
American Retail Corporation & CABANA COVE & 2010 & 7.32 & +0.10 & Metal swimming pools \\
MIM Construction AB & YELLOW LINE & 2016 & 4.72 & +0.06 & storage containers adapted for use in vehicle boots; vehicle partitions for pets; fitted liners for the cargo area of vehicles; boot tidies, namely, storage bins especially adapted to fit in cars; dog \\
SECURE-IT, INC. & FLAVOR SAVER & 1999 & 6.48 & -0.05 & Cable and lock fasteners of metal for computer equipment \\
Wolverine Holding Company, LLC & PRECISION ENGINEERED PRODUCTS & 2002 & 5.00 & -0.00 & Metal hose stations for the food, dairy, beverage, pharmaceutical, health and beauty, and processing industries; metal tubing; metal pipe fittings; metal gaskets for the food, dairy, beverage, pharmac \\
WAL-MART STORES, INC. & PALM VALLEY & 2004 & 5.65 & -0.03 & Metal fence panels, posts and stays; metal garden fence panels; metal patio fences; hand-operated metal garden hose reels and cabinets Hand-operated non-metal garden hose reels and cabinets Non-metal  \\
Wal-Mart Stores, Inc & SLATE MOUNTAIN & 2016 & 4.10 & -0.03 & Barbecue mitts; bed linen; liners adapted to sleeping bags for camping Hand tools, namely, hand-operated pumps; mallets; pocket knives; utility knives Brushes for basting meat; cleaning brushes for ho \\
Wal-Mart Stores, Inc. & CREATE A CRAFT & 1994 & 5.29 & +0.01 & craft supplies, namely metal key rings and [ sleigh bells ] [ doll parts ] wooden, plastic, and metal embroidery hoops; sewing notions, namely [ moving eyes; ] beads; [ wooden craft sticks; chenille s \\
Exxon Mobil Corporation & MAZEFLO & 2010 & 5.20 & +0.07 & Metal sand control screens used in oil and gas well construction \\
Apple Inc. &  & 2008 & 5.32 & -0.15 & Keyrings of common metal; Metal rings and chains for keys; Metal key chains \\
\bottomrule
\end{tabular}
\clearpage

\subsection*{Class 007: Machines \& Power Tools}
\label{app:deep-007}

\begin{tabular}{ll}
Clean filings (1990--2021): & 160,666 \\
Mean $KL_\text{pros}$ / $KL_\text{retr}$ / $\Delta KL$: & 5.34 / 5.34 / +0.003 nats \\
5y survival rate: & 40.3\% \\
Share with SEC-listed owner: & 4.35\% \\
\end{tabular}

\begin{center}
\includegraphics[width=0.7\linewidth]{results/per_industry_deep/class007_scatter.png}
\end{center}

\textbf{Trending tokens (top-decile $\Delta KL$)} --- with intra-class 5y survival rate.

\begin{tabular}{lrrr}
\toprule
token & lift & $n$ filings & intra-class survival \\
\midrule
\texttt{disseminating} & 4.58$\times$ & 662 & 24\% \\
\texttt{spiralizers} & 4.25$\times$ & 252 & 28\% \\
\texttt{wet-dry} & 4.20$\times$ & 374 & 24\% \\
\texttt{drycleaning} & 4.19$\times$ & 215 & 23\% \\
\texttt{squeezers} & 4.07$\times$ & 1,386 & 21\% \\
\texttt{dicing} & 4.01$\times$ & 979 & 28\% \\
\texttt{perfumes} & 3.95$\times$ & 782 & 25\% \\
\texttt{electric-powered} & 3.94$\times$ & 363 & 31\% \\
\texttt{shampooing} & 3.87$\times$ & 1,121 & 23\% \\
\texttt{exhausting} & 3.85$\times$ & 1,773 & 22\% \\
\bottomrule
\end{tabular}

\textbf{Ending tokens (bottom-decile $\Delta KL$)} --- intra-class 5y survival rate.

\begin{tabular}{lrrr}
\toprule
token & lift & $n$ filings & intra-class survival \\
\midrule
\texttt{walk-behind} & 3.11$\times$ & 257 & 47\% \\
\texttt{subs} & 2.79$\times$ & 161 & 40\% \\
\texttt{riding} & 2.75$\times$ & 745 & 47\% \\
\texttt{oscillation} & 2.54$\times$ & 185 & 17\% \\
\texttt{completion} & 2.44$\times$ & 553 & 42\% \\
\texttt{sprocket} & 2.42$\times$ & 186 & 50\% \\
\texttt{surround} & 2.22$\times$ & 171 & 16\% \\
\texttt{flip} & 2.22$\times$ & 185 & 26\% \\
\texttt{downhole} & 2.22$\times$ & 1,101 & 48\% \\
\texttt{wristwatches} & 2.22$\times$ & 167 & 15\% \\
\bottomrule
\end{tabular}

\textbf{Flux-neutral most-frequent tokens ($|\Delta KL|\leq 0.02$)}.

\begin{tabular}{lrrr}
\toprule
token & $n$ in stratum & $n$ in class & intra-class survival \\
\midrule
\texttt{machines} & 11,129 & 85,739 & 32\% \\
\texttt{parts} & 8,297 & 58,796 & 34\% \\
\texttt{use} & 6,894 & 49,168 & 34\% \\
\texttt{electric} & 5,800 & 45,342 & 31\% \\
\texttt{machine} & 5,595 & 37,204 & 31\% \\
\texttt{industrial} & 4,751 & 31,003 & 30\% \\
\texttt{apparatus} & 4,728 & 28,485 & 26\% \\
\texttt{tools} & 4,569 & 32,200 & 33\% \\
\texttt{equipment} & 4,253 & 28,326 & 38\% \\
\texttt{power} & 4,207 & 32,213 & 36\% \\
\bottomrule
\end{tabular}

\textbf{Top-10 SEC-listed firms with debut in 1990--2021} (ranked by max observed revenue, debut filing g/s shown).

\begin{tabular}{p{4.0cm}p{2.2cm}rrrp{6.5cm}}
\toprule
Firm & Mark & year & $KL_\text{pros}$ & $\Delta KL$ & First-filing g/s excerpt \\
\midrule
INA ACQUISITION CORP. & CHIP & 1993 & 4.84 & +0.00 & mechanical devices for applying tubular bags or liners to the internal surfaces of pipes or passageways, such devices consisting of an interconnected arrangement of pressure chambers, valves, cylinder \\
INA Acquisition Corp. &  & 2007 & 6.79 & -0.00 & Pipe and conduit repair and rehabilitation machines, namely, insertion device for inserting flexible liners into existing pipes and conduits for the repair and rehabilitation of existing pipes and con \\
New World Technologies Inc. & RAD & 2015 & 7.93 & -0.11 & Pneumatic torque wrenches \\
One World Technologies Limited & EML EXTENDED MOTOR LIFE & 2003 & 5.88 & +0.02 & Power tools, namely, grinders, hammer drills, drills, and jig saws \\
One World Technologies, Ltd. & LASER-TECH & 2003 & 5.84 & -0.13 & Stud Sensors for sensing studs in walls; Laser Line Generators for projecting a laser line illumination; and Laser Torpedo Levels for measuring a horizontal level or a vertical level of an item and ha \\
REGI U.S., INC. & RAND CAM & 1994 & 4.23 & +0.16 & motors and engines for vehicles engines not for land vehicles; rotary pumps; rotary compressors for air conditioners and refrigerators; rotary air compressors for providing compressed air to shops and \\
SOLA INTERNATIONAL INC. & MATRIX & 1992 & 7.82 & +0.37 & machines used in the manufacture of eyeglass lenses \\
Sola International Inc. & S SOLA & 1995 & 7.06 & -0.30 & laminating machines used in the manufacture of eyeglass lenses \\
Wolverine Holding Company, LLC & PRECISION ENGINEERED PRODUCTS & 2002 & 4.82 & -0.02 & Metal hose stations for the food, dairy, beverage, pharmaceutical, health and beauty, and processing industries; metal tubing; metal pipe fittings; metal gaskets for the food, dairy, beverage, pharmac \\
WAL-MART STORES, INC. & EXPERT GARDENER & 2003 & 5.23 & -0.24 & Wheel barrows Garden apparel, namely, aprons, hats, boots Lawn sprinklers; gardening gloves Plant food; lawn and plant fertilizer for domestic use; potting soil; compost; spike-shaped fertilizer for h \\
\bottomrule
\end{tabular}
\clearpage

\subsection*{Class 008: Hand Tools}
\label{app:deep-008}

\begin{tabular}{ll}
Clean filings (1990--2021): & 71,252 \\
Mean $KL_\text{pros}$ / $KL_\text{retr}$ / $\Delta KL$: & 4.65 / 4.58 / +0.070 nats \\
5y survival rate: & 40.3\% \\
Share with SEC-listed owner: & 2.66\% \\
\end{tabular}

\begin{center}
\includegraphics[width=0.7\linewidth]{results/per_industry_deep/class008_scatter.png}
\end{center}

\textbf{Trending tokens (top-decile $\Delta KL$)} --- with intra-class 5y survival rate.

\begin{tabular}{lrrr}
\toprule
token & lift & $n$ filings & intra-class survival \\
\midrule
\texttt{magneto} & 5.15$\times$ & 388 & 23\% \\
\texttt{shim} & 4.87$\times$ & 412 & 25\% \\
\texttt{tappet} & 4.72$\times$ & 425 & 26\% \\
\texttt{flywheel} & 4.64$\times$ & 432 & 26\% \\
\texttt{absorber} & 4.45$\times$ & 449 & 25\% \\
\texttt{feeler} & 4.22$\times$ & 483 & 27\% \\
\texttt{carburetor} & 4.13$\times$ & 491 & 28\% \\
\texttt{tuning} & 3.79$\times$ & 529 & 26\% \\
\texttt{beneath} & 3.73$\times$ & 131 & 27\% \\
\texttt{adjustment} & 3.64$\times$ & 556 & 28\% \\
\bottomrule
\end{tabular}

\textbf{Ending tokens (bottom-decile $\Delta KL$)} --- intra-class 5y survival rate.

\begin{tabular}{lrrr}
\toprule
token & lift & $n$ filings & intra-class survival \\
\midrule
\texttt{wood-splitting} & 5.82$\times$ & 103 & 57\% \\
\texttt{sidewalk} & 5.69$\times$ & 109 & 57\% \\
\texttt{posthole} & 5.50$\times$ & 140 & 59\% \\
\texttt{vegetation} & 5.37$\times$ & 162 & 55\% \\
\texttt{barn} & 5.13$\times$ & 113 & 57\% \\
\texttt{tamping} & 4.41$\times$ & 145 & 50\% \\
\texttt{weeders} & 3.28$\times$ & 396 & 59\% \\
\texttt{floral} & 3.19$\times$ & 166 & 52\% \\
\texttt{sledge} & 3.14$\times$ & 226 & 52\% \\
\texttt{sod} & 3.00$\times$ & 300 & 60\% \\
\bottomrule
\end{tabular}

\textbf{Flux-neutral most-frequent tokens ($|\Delta KL|\leq 0.02$)}.

\begin{tabular}{lrrr}
\toprule
token & $n$ in stratum & $n$ in class & intra-class survival \\
\midrule
\texttt{tools} & 3,429 & 36,891 & 35\% \\
\texttt{hand} & 3,417 & 35,656 & 36\% \\
\texttt{knives} & 2,485 & 29,268 & 35\% \\
\texttt{use} & 2,232 & 20,225 & 34\% \\
\texttt{hand-operated} & 1,706 & 19,419 & 26\% \\
\texttt{electric} & 1,601 & 20,662 & 30\% \\
\texttt{forks} & 1,586 & 19,227 & 33\% \\
\texttt{spoons} & 1,407 & 17,205 & 33\% \\
\texttt{metal} & 1,391 & 11,300 & 39\% \\
\texttt{operated} & 1,329 & 13,142 & 34\% \\
\bottomrule
\end{tabular}

\textbf{Top-10 SEC-listed firms with debut in 1990--2021} (ranked by max observed revenue, debut filing g/s shown).

\begin{tabular}{p{4.0cm}p{2.2cm}rrrp{6.5cm}}
\toprule
Firm & Mark & year & $KL_\text{pros}$ & $\Delta KL$ & First-filing g/s excerpt \\
\midrule
ATK Entertainment, Inc. & AALAM TUCK KHALED & 2020 & 2.61 & -0.11 & Clothing, namely, shirts, t-shirts, pants, dresses, skirts, jeans, slacks, belts, hats, caps being headwear, gloves, ear muffs, headwear, neckwear, sweaters, fleece pullovers, jogging suits, warm-up s \\
New World Technologies Inc. &  & 2018 & 6.18 & -0.02 & Torque wrenches Torque measuring and calibrating devices for torque wrenches \\
WAL-MART STORES, INC. & EXPERT GARDENER & 2003 & 4.52 & -0.12 & Wheel barrows Garden apparel, namely, aprons, hats, boots Lawn sprinklers; gardening gloves Plant food; lawn and plant fertilizer for domestic use; potting soil; compost; spike-shaped fertilizer for h \\
Wal-Mart Stores, Inc & DAILY'S & 2000 & 2.79 & -0.01 & Cooking appliances, namely, electric Dutch ovens, food steamers, electric skillets, fryers, microwaves, toasters, breadmakers, popcorn poppers, waffle irons, electric coffee pots, coffee makers; elect \\
Wal-Mart Stores, Inc. & GADGETS 'N' GIZMOS & 1994 & 4.54 & +0.02 & timers, thermometers, scales household utensils, namely spoons, knives, graters, kitchen knives rolling pins, salt and pepper shakers \\
Wal-mart Stores, Inc. &  & 2011 & 4.17 & -0.01 & On-line retail store services featuring a wide variety of consumer goods of others; Retail department store services; Retail store services featuring a wide variety of consumer goods of others Air bru \\
MCKESSON CORPORATION & BRIGHT IDEA & 1990 & 3.51 & +0.57 & salad oil cans, sold empty ice picks, ice tongs, cutlery, namely, cleavers, cheese planes, grapefruit knives, spoons, forks, paring knives, filleting knives, boning knives, butcher knives, slicing kni \\
McKesson Corporation & SUNMARK & 2003 & 3.51 & -0.06 & humidifiers and vaporizers House mark for a line of personal care products, namely, tweezers, nail clippers, nail and cuticle scissors, nail files, razors and razor blades and eyelash curlers House ma \\
Microsoft Corporation & MICROSOFT & 1991 & 4.80 & +0.06 & knives; namely, pocket knives \\
ConocoPhillips Company & 76 & 2003 & 4.61 & -0.09 & Knives; hand held tools, namely, pliers, tool sets composed of hand held tools in the nature of a compact folding knife, pliers, multiple blades and leather and fabric knife sheaths \\
\bottomrule
\end{tabular}
\clearpage

\subsection*{Class 009: Software \& Electronics}
\label{app:deep-009}

\begin{tabular}{ll}
Clean filings (1990--2021): & 1,226,806 \\
Mean $KL_\text{pros}$ / $KL_\text{retr}$ / $\Delta KL$: & 5.42 / 5.43 / -0.011 nats \\
5y survival rate: & 48.6\% \\
Share with SEC-listed owner: & 5.52\% \\
\end{tabular}

\begin{center}
\includegraphics[width=0.7\linewidth]{results/per_industry_deep/class009_scatter.png}
\end{center}

\textbf{Trending tokens (top-decile $\Delta KL$)} --- with intra-class 5y survival rate.

\begin{tabular}{lrrr}
\toprule
token & lift & $n$ filings & intra-class survival \\
\midrule
\texttt{non-fungible} & 8.62$\times$ & 1,273 & 64\% \\
\texttt{blockchain-based} & 7.42$\times$ & 1,349 & 63\% \\
\texttt{blockchain} & 6.45$\times$ & 3,427 & 59\% \\
\texttt{cryptocurrency} & 6.33$\times$ & 3,255 & 58\% \\
\texttt{cryptographic} & 4.44$\times$ & 1,738 & 55\% \\
\texttt{currencies} & 4.35$\times$ & 1,602 & 51\% \\
\texttt{collectibles} & 4.20$\times$ & 1,293 & 56\% \\
\texttt{tokens} & 4.02$\times$ & 5,129 & 47\% \\
\texttt{wallet} & 3.90$\times$ & 2,428 & 54\% \\
\texttt{ledger} & 3.83$\times$ & 1,541 & 51\% \\
\bottomrule
\end{tabular}

\textbf{Ending tokens (bottom-decile $\Delta KL$)} --- intra-class 5y survival rate.

\begin{tabular}{lrrr}
\toprule
token & lift & $n$ filings & intra-class survival \\
\midrule
\texttt{typeface} & 3.53$\times$ & 1,333 & 51\% \\
\texttt{soundbar} & 3.17$\times$ & 1,629 & 23\% \\
\texttt{videorecorders} & 3.05$\times$ & 1,303 & 17\% \\
\texttt{selfie} & 3.02$\times$ & 6,129 & 20\% \\
\texttt{generates} & 2.89$\times$ & 3,501 & 49\% \\
\texttt{fonts} & 2.77$\times$ & 2,519 & 45\% \\
\texttt{wager} & 2.74$\times$ & 6,243 & 52\% \\
\texttt{smartwatch} & 2.72$\times$ & 3,370 & 22\% \\
\texttt{webcams} & 2.70$\times$ & 5,647 & 20\% \\
\texttt{outcomes} & 2.64$\times$ & 4,227 & 49\% \\
\bottomrule
\end{tabular}

\textbf{Flux-neutral most-frequent tokens ($|\Delta KL|\leq 0.02$)}.

\begin{tabular}{lrrr}
\toprule
token & $n$ in stratum & $n$ in class & intra-class survival \\
\midrule
\texttt{computer} & 71,265 & 654,366 & 48\% \\
\texttt{software} & 68,779 & 602,415 & 50\% \\
\texttt{use} & 43,630 & 388,588 & 45\% \\
\texttt{electronic} & 42,973 & 362,929 & 44\% \\
\texttt{downloadable} & 35,839 & 275,741 & 44\% \\
\texttt{data} & 32,816 & 261,388 & 43\% \\
\texttt{information} & 30,922 & 227,341 & 48\% \\
\texttt{video} & 30,379 & 265,478 & 48\% \\
\texttt{devices} & 29,782 & 258,931 & 43\% \\
\texttt{field} & 28,728 & 213,618 & 44\% \\
\bottomrule
\end{tabular}

\textbf{Top-10 SEC-listed firms with debut in 1990--2021} (ranked by max observed revenue, debut filing g/s shown).

\begin{tabular}{p{4.0cm}p{2.2cm}rrrp{6.5cm}}
\toprule
Firm & Mark & year & $KL_\text{pros}$ & $\Delta KL$ & First-filing g/s excerpt \\
\midrule
1X Technologies LLC & 1X TECHNOLOGIES & 2017 & 4.97 & -0.19 & Cable and wiring plates in the nature of protective sheaths for electrical cable and wiring; Cable junctions for electric cables; Cables for electrical and optical signal transmission systems; Cables  \\
ALTO CO., LTD. & ALTO & 2011 & 4.74 & -0.06 & [ Optical condensers for use in microscopes, enlargers, slide projectors, video projectors, and telescopes; magic lanterns; photometers; actinometers; signalling devices, namely, traffic light apparat \\
ATK Entertainment, Inc. & AALAM TUCK KHALED & 2020 & 3.14 & -0.07 & Clothing, namely, shirts, t-shirts, pants, dresses, skirts, jeans, slacks, belts, hats, caps being headwear, gloves, ear muffs, headwear, neckwear, sweaters, fleece pullovers, jogging suits, warm-up s \\
Air Unlimited & SLAM & 2021 & 6.71 & -0.10 & Suspension systems for automobiles and component parts therefor; air springs and air valves, all for suspension systems for automobiles Electronic controllers and electric wiring harnesses for suspens \\
Artix Entertainment, LLC & ADVENTUREQUEST & 2008 & 4.37 & -0.21 & Motion picture films featuring animation; audio and video recordings featuring animation; prerecorded audio and video tapes, laser discs, DVDs, and CDs featuring animation; downloadable audio and vide \\
BETTERLIFE PHARMA INC. & MOVE BETTERX & 2019 & 5.13 & +0.16 & Printed publications, namely, newsletters, brochures, reports and guides in the field of cannabis bath additives, namely, bath beads, bath bombs, bath crystals and bubble bath; bath herbs; bath oils;  \\
BLH ACQUISITION CO., LLC &  & 2012 & 6.19 & -0.35 & Magnetic coded gift cards Non-magnetically encoded prepaid purchase cards for allowing users to transfer financial value at the point of sale and on-line via a computuer network \\
BetterLife Pharma Inc. & LIVE BETTERX & 2020 & 4.85 & +0.03 & Downloadable electronic publications, namely, newsletters, brochures, reports and guides in the field of cannabis Retail store and wholesale store services featuring industrial hemp food and body care \\
Biosolar, Inc. & BIOSOLAR & 2006 & 6.48 & +0.27 & Apparatus for converting electronic radiation to electrical energy, namely, photovoltaic solar modules; Apparatus for converting electronic radiation to electrical energy, namely, photovoltaic solar-t \\
Blackstar Enterprise Group, Inc. & B. E. T. & 2021 & 4.81 & +0.18 & Financial services, namely, providing trading in securities services for trading SEC regulated securities on an immutable blockchain, accessible via a trading platform; transaction processing, namely, \\
\bottomrule
\end{tabular}
\clearpage

\subsection*{Class 010: Medical Devices}
\label{app:deep-010}

\begin{tabular}{ll}
Clean filings (1990--2021): & 190,121 \\
Mean $KL_\text{pros}$ / $KL_\text{retr}$ / $\Delta KL$: & 5.35 / 5.35 / -0.001 nats \\
5y survival rate: & 47.4\% \\
Share with SEC-listed owner: & 7.80\% \\
\end{tabular}

\begin{center}
\includegraphics[width=0.7\linewidth]{results/per_industry_deep/class010_scatter.png}
\end{center}

\textbf{Trending tokens (top-decile $\Delta KL$)} --- with intra-class 5y survival rate.

\begin{tabular}{lrrr}
\toprule
token & lift & $n$ filings & intra-class survival \\
\midrule
\texttt{rigidity} & 4.79$\times$ & 801 & 20\% \\
\texttt{wavelengths} & 4.78$\times$ & 561 & 29\% \\
\texttt{constriction} & 4.74$\times$ & 791 & 19\% \\
\texttt{non-ablative} & 4.68$\times$ & 711 & 25\% \\
\texttt{green} & 4.52$\times$ & 588 & 30\% \\
\texttt{life-sized} & 4.50$\times$ & 1,230 & 20\% \\
\texttt{orange} & 4.49$\times$ & 579 & 29\% \\
\texttt{yellow} & 4.46$\times$ & 596 & 30\% \\
\texttt{workbook} & 4.19$\times$ & 1,457 & 18\% \\
\texttt{penile} & 4.17$\times$ & 967 & 25\% \\
\bottomrule
\end{tabular}

\textbf{Ending tokens (bottom-decile $\Delta KL$)} --- intra-class 5y survival rate.

\begin{tabular}{lrrr}
\toprule
token & lift & $n$ filings & intra-class survival \\
\midrule
\texttt{percent} & 5.29$\times$ & 238 & 56\% \\
\texttt{dilatation} & 4.93$\times$ & 209 & 42\% \\
\texttt{atherectomy} & 4.81$\times$ & 210 & 31\% \\
\texttt{interbody} & 4.77$\times$ & 199 & 66\% \\
\texttt{pacing} & 4.46$\times$ & 280 & 59\% \\
\texttt{angiographic} & 4.27$\times$ & 227 & 27\% \\
\texttt{non-active} & 4.20$\times$ & 195 & 31\% \\
\texttt{energizing} & 4.17$\times$ & 199 & 31\% \\
\texttt{transluminal} & 4.12$\times$ & 216 & 32\% \\
\texttt{pedicle} & 4.08$\times$ & 233 & 45\% \\
\bottomrule
\end{tabular}

\textbf{Flux-neutral most-frequent tokens ($|\Delta KL|\leq 0.02$)}.

\begin{tabular}{lrrr}
\toprule
token & $n$ in stratum & $n$ in class & intra-class survival \\
\midrule
\texttt{medical} & 11,263 & 119,396 & 45\% \\
\texttt{use} & 8,005 & 82,136 & 41\% \\
\texttt{apparatus} & 7,708 & 74,909 & 38\% \\
\texttt{instruments} & 5,599 & 54,584 & 38\% \\
\texttt{purposes} & 5,390 & 52,059 & 34\% \\
\texttt{devices} & 5,313 & 55,395 & 47\% \\
\texttt{surgical} & 3,806 & 41,541 & 43\% \\
\texttt{treatment} & 2,928 & 26,564 & 41\% \\
\texttt{diagnostic} & 2,588 & 20,431 & 41\% \\
\texttt{equipment} & 2,403 & 21,187 & 40\% \\
\bottomrule
\end{tabular}

\textbf{Top-10 SEC-listed firms with debut in 1990--2021} (ranked by max observed revenue, debut filing g/s shown).

\begin{tabular}{p{4.0cm}p{2.2cm}rrrp{6.5cm}}
\toprule
Firm & Mark & year & $KL_\text{pros}$ & $\Delta KL$ & First-filing g/s excerpt \\
\midrule
AGE International, Inc. & BLANTON'S & 2021 & 5.09 & -0.04 & wall mounted coat rack; wooden display stands for bottle stoppers, bottles and tasting glasses; wall art made of wood reclaimed from bourbon barrels face coverings being masks for protection against v \\
ATK Entertainment, Inc. & AALAM & 2020 & 3.05 & -0.05 & Clothing, namely, shirts, t-shirts, pants, dresses, skirts, jeans, slacks, belts, hats, caps being headwear, gloves, ear muffs, headwear, neckwear, sweaters, fleece pullovers, jogging suits, warm-up s \\
Bone Biologics, Inc. & BIOBONE & 2006 & 6.80 & -0.28 & bone graft substitute, orthobiologic, bone growth factor, synthetic bone graft, demineralized bone graft \\
Event Cardio Group & NOW CARDIO & 2016 & 5.21 & +0.01 & Medical device, namely, a cardiac and cardiac event monitor having a sensor and capability to transmit cardiac and other physiological data of patients remotely Medical Services, namely, remote monito \\
HR PHARMACEUTICALS, INC. & MTG INSTANT CATH & 2007 & 5.44 & -0.18 & Urinary catheters; Urinary catheter products, namely, catheters and parts and fittings therefor; and Disposable intermittent catheters \\
R2 Technologies, Inc. & CRYOAESTHETICS & 2020 & 4.89 & +0.38 & Cryotherapy apparatus for medical and cosmetic use for performing aesthetic skin treatment procedures \\
SM Global Enterprises, Inc. & INTIMATENEED & 2017 & 3.74 & -0.02 & Personal sexual lubricants, creams for sexual aid, lotions for sexual aid and topical gels for use as body lubricants; topical gels for use as male and female genital sensitizers and desensitizers Adu \\
Sharps Technology, Inc. & PROMESIA & 2018 & 5.55 & -0.01 & Online retail store and online retail distribution services featuring syringes and needles for medical use and for injection, disposable syringes and needles for medical use and for injection, disposa \\
Sunrise & DREEM & 2019 & 2.71 & +0.07 & Scientific, weighing, measuring, signaling, checking, supervision, life-saving apparatus and instruments, namely, electric and electronic apparatus and instruments for the measuring, storage, encoding \\
Z Corporation & ZPHOTO & 2007 & 7.64 & -0.16 & medical images; Two-and-a-half dimensional medical images providing opaque and translucent contrast resolution \\
\bottomrule
\end{tabular}
\clearpage

\subsection*{Class 011: Lighting, HVAC \& Appliances}
\label{app:deep-011}

\begin{tabular}{ll}
Clean filings (1990--2021): & 194,994 \\
Mean $KL_\text{pros}$ / $KL_\text{retr}$ / $\Delta KL$: & 4.84 / 4.79 / +0.045 nats \\
5y survival rate: & 38.8\% \\
Share with SEC-listed owner: & 2.80\% \\
\end{tabular}

\begin{center}
\includegraphics[width=0.7\linewidth]{results/per_industry_deep/class011_scatter.png}
\end{center}

\textbf{Trending tokens (top-decile $\Delta KL$)} --- with intra-class 5y survival rate.

\begin{tabular}{lrrr}
\toprule
token & lift & $n$ filings & intra-class survival \\
\midrule
\texttt{party-themed} & 2.97$\times$ & 1,594 & 26\% \\
\texttt{chinese} & 2.88$\times$ & 2,682 & 19\% \\
\texttt{diving} & 2.81$\times$ & 5,317 & 19\% \\
\texttt{fairy} & 2.79$\times$ & 9,357 & 15\% \\
\texttt{christmas-themed} & 2.75$\times$ & 1,412 & 19\% \\
\texttt{underground} & 2.74$\times$ & 2,806 & 19\% \\
\texttt{all-weather} & 2.66$\times$ & 1,698 & 26\% \\
\texttt{turned} & 2.61$\times$ & 662 & 33\% \\
\texttt{searchlights} & 2.60$\times$ & 8,194 & 18\% \\
\texttt{branches} & 2.58$\times$ & 3,105 & 22\% \\
\bottomrule
\end{tabular}

\textbf{Ending tokens (bottom-decile $\Delta KL$)} --- intra-class 5y survival rate.

\begin{tabular}{lrrr}
\toprule
token & lift & $n$ filings & intra-class survival \\
\midrule
\texttt{vaporizable} & 3.57$\times$ & 255 & 49\% \\
\texttt{constituents} & 3.06$\times$ & 294 & 46\% \\
\texttt{robe} & 3.02$\times$ & 255 & 64\% \\
\texttt{thawing} & 2.99$\times$ & 294 & 48\% \\
\texttt{scenting} & 2.99$\times$ & 274 & 61\% \\
\texttt{vaporizing} & 2.72$\times$ & 364 & 45\% \\
\texttt{refill} & 2.42$\times$ & 265 & 46\% \\
\texttt{refrigerator-freezers} & 2.32$\times$ & 228 & 43\% \\
\texttt{glitter} & 2.28$\times$ & 254 & 54\% \\
\texttt{gloss} & 2.25$\times$ & 231 & 55\% \\
\bottomrule
\end{tabular}

\textbf{Flux-neutral most-frequent tokens ($|\Delta KL|\leq 0.02$)}.

\begin{tabular}{lrrr}
\toprule
token & $n$ in stratum & $n$ in class & intra-class survival \\
\midrule
\texttt{electric} & 7,491 & 89,973 & 31\% \\
\texttt{use} & 7,102 & 71,438 & 34\% \\
\texttt{apparatus} & 6,335 & 66,232 & 27\% \\
\texttt{water} & 5,831 & 56,555 & 35\% \\
\texttt{air} & 5,445 & 55,451 & 32\% \\
\texttt{lighting} & 4,566 & 63,268 & 29\% \\
\texttt{installations} & 4,140 & 44,286 & 24\% \\
\texttt{machines} & 4,097 & 39,740 & 28\% \\
\texttt{purposes} & 4,032 & 39,558 & 26\% \\
\texttt{lamps} & 3,863 & 55,787 & 25\% \\
\bottomrule
\end{tabular}

\textbf{Top-10 SEC-listed firms with debut in 1990--2021} (ranked by max observed revenue, debut filing g/s shown).

\begin{tabular}{p{4.0cm}p{2.2cm}rrrp{6.5cm}}
\toprule
Firm & Mark & year & $KL_\text{pros}$ & $\Delta KL$ & First-filing g/s excerpt \\
\midrule
ALTO CO., LTD. & ALTO & 2011 & 3.78 & +0.04 & [ Optical condensers for use in microscopes, enlargers, slide projectors, video projectors, and telescopes; magic lanterns; photometers; actinometers; signalling devices, namely, traffic light apparat \\
American Retail Corporation & AVENTURA & 2009 & 6.60 & -0.23 & Spas in the nature of heated pools \\
ECO Integrated Technologies, Inc. & AQUAVAP & 2021 & 4.77 & +0.03 & Water purification apparatus; Water purification, sedimentation and filtration apparatus for clarification, thickening and filtration of water, waste water and industrial water; Waste water purificati \\
KX TECHNOLOGIES LLC & MATRIKX + CTO & 1993 & 4.86 & -0.08 & components of water filtration units for producing potable water for domestic use; namely, activated carbon water filters \\
KX Technologies LLC & FINITY & 2013 & 3.96 & -0.11 & chemical and mineral fiber pulp preparations in the nature of filtering materials for the manufacturing of domestic, industrial, and commercial water filters, and for the manufacture of commercial and \\
MY CO., LTD. & WATER PLUS & 2020 & 6.04 & -0.13 & Bath and shower gels; beauty bath additives, namely, bubble bath, bath beads, bath foam, bath crystals, bath pearls, and non-medicated bath salts; non-medicated bath salts; bath and shower gels and sa \\
R4 TECHNOLOGIES INC. & R4 TECHNOLOGIES & 2013 & 4.49 & +0.20 & [ Electronic commerce payment services, namely, establishing funded accounts used to purchase goods and services via mobile devices; Financial services, namely, providing on-line stored value accounts \\
R4 TECHNOLOGIES, LLC. & SMART COOLER & 2012 & 5.23 & +0.16 & Refrigerated beverage dispensing units; Refrigerated dispensing units for beverages; Refrigerated self-service cases, refrigerated self-service deli cases and custom fabricated refrigeration systems u \\
SAVANT TECHNOLOGIES LLC & GENURA & 1995 & 4.85 & -0.21 & electric lamps \\
SPA ACQUISITION CORP. & L.A. SPAS & 1990 & 7.56 & +0.39 & acryllic hot tubs with fiberglass backing \\
\bottomrule
\end{tabular}
\clearpage

\subsection*{Class 012: Vehicles}
\label{app:deep-012}

\begin{tabular}{ll}
Clean filings (1990--2021): & 128,694 \\
Mean $KL_\text{pros}$ / $KL_\text{retr}$ / $\Delta KL$: & 4.93 / 4.90 / +0.034 nats \\
5y survival rate: & 43.7\% \\
Share with SEC-listed owner: & 5.01\% \\
\end{tabular}

\begin{center}
\includegraphics[width=0.7\linewidth]{results/per_industry_deep/class012_scatter.png}
\end{center}

\textbf{Trending tokens (top-decile $\Delta KL$)} --- with intra-class 5y survival rate.

\begin{tabular}{lrrr}
\toprule
token & lift & $n$ filings & intra-class survival \\
\midrule
\texttt{contracted} & 5.81$\times$ & 736 & 23\% \\
\texttt{dropping} & 5.53$\times$ & 776 & 22\% \\
\texttt{gap} & 5.42$\times$ & 799 & 24\% \\
\texttt{expanded} & 5.30$\times$ & 806 & 23\% \\
\texttt{droppings} & 5.26$\times$ & 247 & 11\% \\
\texttt{wedge} & 5.21$\times$ & 826 & 25\% \\
\texttt{durable} & 5.14$\times$ & 830 & 26\% \\
\texttt{paws} & 4.85$\times$ & 165 & 32\% \\
\texttt{fill} & 4.81$\times$ & 892 & 26\% \\
\texttt{anti-slipping} & 4.63$\times$ & 479 & 18\% \\
\bottomrule
\end{tabular}

\textbf{Ending tokens (bottom-decile $\Delta KL$)} --- intra-class 5y survival rate.

\begin{tabular}{lrrr}
\toprule
token & lift & $n$ filings & intra-class survival \\
\midrule
\texttt{sport-utility} & 4.68$\times$ & 314 & 79\% \\
\texttt{pillar} & 2.89$\times$ & 204 & 47\% \\
\texttt{motorhomes} & 2.64$\times$ & 216 & 40\% \\
\texttt{bob-sleighs} & 2.61$\times$ & 142 & 15\% \\
\texttt{fifth} & 2.52$\times$ & 892 & 56\% \\
\texttt{elderly} & 2.42$\times$ & 240 & 46\% \\
\texttt{recapping} & 2.42$\times$ & 236 & 60\% \\
\texttt{radial} & 2.35$\times$ & 200 & 52\% \\
\texttt{non-monetary} & 2.25$\times$ & 231 & 35\% \\
\texttt{softball} & 2.18$\times$ & 238 & 45\% \\
\bottomrule
\end{tabular}

\textbf{Flux-neutral most-frequent tokens ($|\Delta KL|\leq 0.02$)}.

\begin{tabular}{lrrr}
\toprule
token & $n$ in stratum & $n$ in class & intra-class survival \\
\midrule
\texttt{vehicles} & 7,431 & 67,870 & 39\% \\
\texttt{parts} & 6,427 & 56,049 & 40\% \\
\texttt{vehicle} & 4,596 & 41,986 & 36\% \\
\texttt{land} & 4,347 & 33,773 & 36\% \\
\texttt{structural} & 3,304 & 27,975 & 40\% \\
\texttt{automobiles} & 3,023 & 28,376 & 35\% \\
\texttt{motor} & 3,019 & 27,315 & 36\% \\
\texttt{use} & 2,925 & 25,248 & 35\% \\
\texttt{wheels} & 2,698 & 22,834 & 35\% \\
\texttt{electric} & 2,491 & 20,707 & 34\% \\
\bottomrule
\end{tabular}

\textbf{Top-10 SEC-listed firms with debut in 1990--2021} (ranked by max observed revenue, debut filing g/s shown).

\begin{tabular}{p{4.0cm}p{2.2cm}rrrp{6.5cm}}
\toprule
Firm & Mark & year & $KL_\text{pros}$ & $\Delta KL$ & First-filing g/s excerpt \\
\midrule
ATK Entertainment, Inc. & AALAM & 2020 & 2.93 & -0.15 & Clothing, namely, shirts, t-shirts, pants, dresses, skirts, jeans, slacks, belts, hats, caps being headwear, gloves, ear muffs, headwear, neckwear, sweaters, fleece pullovers, jogging suits, warm-up s \\
Air Unlimited & SLAM SPECIALTIES & 2021 & 4.73 & +0.05 & Suspension systems for automobiles and component parts therefor; air springs and air valves, all for suspension systems for automobiles Electronic controllers and electric wiring harnesses for suspens \\
Atlis Motor Vehicles, Inc. &  & 2021 & 4.95 & +0.06 & Electric vehicles, namely, passenger vehicles; pickup trucks; minivans; cargo vans; utility trucks; light commercial trucks; structural repair parts and replacement parts therefor \\
Irwin International, Inc. & AIRCRAFT SPRUCE & 2000 & 3.97 & -0.32 & Airplanes (structural parts for) \\
MIM Construction AB & CAREY & 2016 & 4.72 & +0.31 & storage containers adapted for use in vehicle boots; vehicle partitions for pets; fitted liners for the cargo area of vehicles; boot tidies, namely, storage bins especially adapted to fit in cars; dog \\
Manthey Redmond Corporation & MULTI-FIRE TECHNOLOGY & 2012 & 4.09 & -0.08 & Internal combustion engines for land vehicles \\
REGI U.S., INC. & RAND CAM & 1994 & 4.60 & +0.15 & motors and engines for vehicles engines not for land vehicles; rotary pumps; rotary compressors for air conditioners and refrigerators; rotary air compressors for providing compressed air to shops and \\
TRITON ACQUISITION, LLC & TRITON & 1997 & 5.36 & -0.25 & trailers, namely, motorcycle, personal watercraft, utility, snowmobile, pontoon, all terrain vehicle, and boat trailers \\
WAL-MART STORES, INC. & EXPERT GARDENER & 2003 & 5.61 & +0.04 & Wheel barrows Garden apparel, namely, aprons, hats, boots Lawn sprinklers; gardening gloves Plant food; lawn and plant fertilizer for domestic use; potting soil; compost; spike-shaped fertilizer for h \\
Wal-Mart Stores, Inc & SLATE MOUNTAIN & 2016 & 4.10 & -0.09 & Barbecue mitts; bed linen; liners adapted to sleeping bags for camping Hand tools, namely, hand-operated pumps; mallets; pocket knives; utility knives Brushes for basting meat; cleaning brushes for ho \\
\bottomrule
\end{tabular}
\clearpage

\subsection*{Class 013: Firearms \& Pyrotechnics}
\label{app:deep-013}

\begin{tabular}{ll}
Clean filings (1990--2021): & 13,777 \\
Mean $KL_\text{pros}$ / $KL_\text{retr}$ / $\Delta KL$: & 4.53 / 4.47 / +0.063 nats \\
5y survival rate: & 45.9\% \\
Share with SEC-listed owner: & 3.94\% \\
\end{tabular}

\begin{center}
\includegraphics[width=0.7\linewidth]{results/per_industry_deep/class013_scatter.png}
\end{center}

\textbf{Trending tokens (top-decile $\Delta KL$)} --- with intra-class 5y survival rate.

\begin{tabular}{lrrr}
\toprule
token & lift & $n$ filings & intra-class survival \\
\midrule
\texttt{breaks} & 4.14$\times$ & 70 & 43\% \\
\texttt{gunsights} & 3.16$\times$ & 98 & 64\% \\
\texttt{roman} & 2.42$\times$ & 62 & 53\% \\
\texttt{ancillary} & 2.34$\times$ & 350 & 55\% \\
\texttt{attaching} & 1.99$\times$ & 1,282 & 46\% \\
\texttt{pouch} & 1.95$\times$ & 174 & 36\% \\
\texttt{attachments} & 1.95$\times$ & 1,351 & 45\% \\
\texttt{sky} & 1.91$\times$ & 68 & 49\% \\
\texttt{pointing} & 1.79$\times$ & 390 & 46\% \\
\texttt{external} & 1.79$\times$ & 542 & 47\% \\
\bottomrule
\end{tabular}

\textbf{Ending tokens (bottom-decile $\Delta KL$)} --- intra-class 5y survival rate.

\begin{tabular}{lrrr}
\toprule
token & lift & $n$ filings & intra-class survival \\
\midrule
\texttt{cores} & 2.98$\times$ & 57 & 58\% \\
\texttt{colored} & 2.91$\times$ & 55 & 40\% \\
\texttt{effects} & 2.86$\times$ & 63 & 43\% \\
\texttt{non-explosive} & 2.68$\times$ & 56 & 59\% \\
\texttt{roman} & 2.58$\times$ & 62 & 53\% \\
\texttt{shotgun} & 2.56$\times$ & 398 & 52\% \\
\texttt{firecrackers} & 2.34$\times$ & 252 & 38\% \\
\texttt{smokeless} & 2.31$\times$ & 52 & 71\% \\
\texttt{defensive} & 2.30$\times$ & 74 & 57\% \\
\texttt{aerosol} & 2.29$\times$ & 118 & 40\% \\
\bottomrule
\end{tabular}

\textbf{Flux-neutral most-frequent tokens ($|\Delta KL|\leq 0.02$)}.

\begin{tabular}{lrrr}
\toprule
token & $n$ in stratum & $n$ in class & intra-class survival \\
\midrule
\texttt{firearms} & 482 & 6,870 & 43\% \\
\texttt{gun} & 367 & 4,493 & 43\% \\
\texttt{ammunition} & 333 & 4,390 & 43\% \\
\texttt{guns} & 260 & 3,735 & 40\% \\
\texttt{firearm} & 256 & 3,224 & 44\% \\
\texttt{parts} & 241 & 3,692 & 42\% \\
\texttt{cases} & 233 & 3,062 & 41\% \\
\texttt{use} & 232 & 2,798 & 43\% \\
\texttt{accessories} & 231 & 3,131 & 45\% \\
\texttt{rifle} & 187 & 2,498 & 40\% \\
\bottomrule
\end{tabular}

\textbf{Top-10 SEC-listed firms with debut in 1990--2021} (ranked by max observed revenue, debut filing g/s shown).

\begin{tabular}{p{4.0cm}p{2.2cm}rrrp{6.5cm}}
\toprule
Firm & Mark & year & $KL_\text{pros}$ & $\Delta KL$ & First-filing g/s excerpt \\
\midrule
Wal-Mart Stores, Inc. & BLACK CANYON CAMP GEAR & 1996 & 4.31 & -0.24 & hunting knives, pocket knives, utility knives, shovels gun cases dishes gas camping stoves, gas lanterns tents and hammocks sleeping bags \\
FORD MOTOR COMPANY & BUILT TO LAST & 1998 & 4.99 & +0.37 & Motor vehicles, namely, automobiles, trucks, vans, sport-utility vehicles, and their structural parts and engines \\
LOCKHEED MARTIN CORPORATION & MLRS & 1999 & 6.37 & +0.15 & Mobile projectile launch platforms mounted on a vehicle \\
Lockheed Martin Corporation &  & 1999 & 5.37 & +0.35 & guns, guided missiles, military rifles \\
RTX CORPORATION &  & 2020 & 1.79 & -0.16 & Training in servicing and repair of aircraft and systems and components, in use of electronic equipment/systems, pumps, blowers, and compressors; training on how to operate unmanned aircraft; training \\
SCHLUMBERGER LIMITED & S & 2000 & 2.66 & -0.05 & Drilling machines and parts therefor; boring machines; mechanical mixing machines; rotary, suction, sump and vacuum pumps Printed instructional, educational, and teaching materials about well drilling \\
Schlumberger Limited & SCHLUMBERGER & 2010 & 2.40 & -0.39 & Construction project management services for others in the fields of oil and gas exploration and production; oil and gas drilling; consulting services in the fields of oil and gas well drilling; sand  \\
Safeway Inc. & SIGNATURE HOME & 2015 & 3.80 & -0.25 & Fire igniters Paper, plastic and foam cups, dishes and bowls, scouring sponges, cleaning sponges, household rubber/latex gloves for general use, hypo vinyl household gloves for general use, disposable \\
General Dynamics Corporation & STRENGTH ON YOUR SIDE & 1998 & 4.01 & +0.01 & Weapons systems, namely guns, artillery guns, and anti-aircraft guns and parts therefor; ammunition handling systems comprising wagons, cases, magazines, loaders and unloaders and parts therefor; ammu \\
RAYTHEON COMPANY & TERM & 2001 & 6.32 & -0.36 & Munitions Research and development activities in the area of munitions \\
\bottomrule
\end{tabular}
\clearpage

\subsection*{Class 014: Jewelry \& Watches}
\label{app:deep-014}

\begin{tabular}{ll}
Clean filings (1990--2021): & 109,375 \\
Mean $KL_\text{pros}$ / $KL_\text{retr}$ / $\Delta KL$: & 4.52 / 4.50 / +0.019 nats \\
5y survival rate: & 43.2\% \\
Share with SEC-listed owner: & 1.79\% \\
\end{tabular}

\begin{center}
\includegraphics[width=0.7\linewidth]{results/per_industry_deep/class014_scatter.png}
\end{center}

\textbf{Trending tokens (top-decile $\Delta KL$)} --- with intra-class 5y survival rate.

\begin{tabular}{lrrr}
\toprule
token & lift & $n$ filings & intra-class survival \\
\midrule
\texttt{separately} & 4.02$\times$ & 1,559 & 28\% \\
\texttt{earring} & 3.67$\times$ & 1,781 & 30\% \\
\texttt{telemeters} & 3.38$\times$ & 136 & 18\% \\
\texttt{necklace} & 3.33$\times$ & 1,969 & 31\% \\
\texttt{reduce} & 3.05$\times$ & 118 & 19\% \\
\texttt{wristwatch} & 2.99$\times$ & 431 & 38\% \\
\texttt{pace} & 2.98$\times$ & 114 & 32\% \\
\texttt{multi} & 2.94$\times$ & 136 & 12\% \\
\texttt{hoop} & 2.84$\times$ & 1,596 & 25\% \\
\texttt{split} & 2.79$\times$ & 2,936 & 22\% \\
\bottomrule
\end{tabular}

\textbf{Ending tokens (bottom-decile $\Delta KL$)} --- intra-class 5y survival rate.

\begin{tabular}{lrrr}
\toprule
token & lift & $n$ filings & intra-class survival \\
\midrule
\texttt{headset} & 5.19$\times$ & 160 & 19\% \\
\texttt{karat} & 4.27$\times$ & 225 & 56\% \\
\texttt{cellphones} & 4.19$\times$ & 203 & 24\% \\
\texttt{telemeters} & 3.90$\times$ & 136 & 18\% \\
\texttt{surround} & 3.68$\times$ & 247 & 15\% \\
\texttt{decal} & 3.53$\times$ & 136 & 85\% \\
\texttt{hvac} & 3.44$\times$ & 125 & 10\% \\
\texttt{diode} & 3.27$\times$ & 355 & 19\% \\
\texttt{equalizers} & 3.24$\times$ & 247 & 25\% \\
\texttt{reduce} & 3.22$\times$ & 118 & 19\% \\
\bottomrule
\end{tabular}

\textbf{Flux-neutral most-frequent tokens ($|\Delta KL|\leq 0.02$)}.

\begin{tabular}{lrrr}
\toprule
token & $n$ in stratum & $n$ in class & intra-class survival \\
\midrule
\texttt{jewelry} & 7,396 & 79,271 & 43\% \\
\texttt{precious} & 4,969 & 46,170 & 36\% \\
\texttt{watches} & 4,953 & 47,812 & 42\% \\
\texttt{rings} & 4,487 & 45,597 & 36\% \\
\texttt{bracelets} & 4,445 & 46,380 & 37\% \\
\texttt{metal} & 3,870 & 33,694 & 36\% \\
\texttt{cases} & 3,845 & 30,446 & 37\% \\
\texttt{necklaces} & 3,629 & 38,823 & 35\% \\
\texttt{chains} & 3,568 & 34,325 & 35\% \\
\texttt{clocks} & 3,562 & 32,323 & 39\% \\
\bottomrule
\end{tabular}

\textbf{Top-10 SEC-listed firms with debut in 1990--2021} (ranked by max observed revenue, debut filing g/s shown).

\begin{tabular}{p{4.0cm}p{2.2cm}rrrp{6.5cm}}
\toprule
Firm & Mark & year & $KL_\text{pros}$ & $\Delta KL$ & First-filing g/s excerpt \\
\midrule
AGE International, Inc. & BLANTON'S & 2021 & 4.70 & +0.01 & wall mounted coat rack; wooden display stands for bottle stoppers, bottles and tasting glasses; wall art made of wood reclaimed from bourbon barrels face coverings being masks for protection against v \\
ATK Entertainment, Inc. & AALAM KHALED & 2020 & 2.04 & +0.03 & Clothing, namely, shirts, t-shirts, pants, dresses, skirts, jeans, slacks, belts, hats, caps being headwear, gloves, ear muffs, headwear, neckwear, sweaters, fleece pullovers, jogging suits, warm-up s \\
Lollipop AB & LOLLIPOP & 1996 & 3.78 & -0.04 & stationery, paper materials for packing paper and cardboard boxes, paper and plastic bags for merchandise packaging, photo and memory albums and scrapbooks, calendars and diaries, bookmarkers braid, h \\
ODIMO INCORPORATED & THE POWER TO TRUST & 2000 & 5.61 & +0.37 & Precious gemstones; fine jewelry; watches Computerized online ordering services featuring jewelry and gemstones; online retail store services featuring jewelry and gemstones \\
Plantation Holdings Group, Inc. & THE PLANTATION SHOP & 2004 & 4.67 & +0.02 & Gourmet cookies, coffee, wafers, biscuits, snack nut mix consisting primarily of chocolate covered nuts and pralines, candy, toffee and cheese straws Jewelry Retail store services and retail shop at h \\
TREND INNOVATIONS, LLC & NECKLET & 2021 & 3.90 & +0.21 & Online retail store services featuring jewelry and accessories Jewelry; Clasps for jewelry; Jewelry cases \\
WAL-MART STORES, INC. & FADED GLORY & 1996 & 6.39 & -0.72 & sunglasses watches \\
Wal-Mart Stores, Inc & DAILY'S & 2000 & 2.92 & -0.10 & Cooking appliances, namely, electric Dutch ovens, food steamers, electric skillets, fryers, microwaves, toasters, breadmakers, popcorn poppers, waffle irons, electric coffee pots, coffee makers; elect \\
Wal-Mart Stores, Inc. & COLOR TREND & 1994 & 6.37 & -0.33 & sunglasses handbags earrings \\
APPLE INC. &  & 2002 & 3.16 & +0.16 & [Jewelry; watches and clocks; cufflinks;] key rings of precious metal [ ; chronometers; brooches, charms; tie pins; tie clips;] badges [, bracelets, necklaces; medallions, belt buckles, pins; accessor \\
\bottomrule
\end{tabular}
\clearpage

\subsection*{Class 015: Musical Instruments}
\label{app:deep-015}

\begin{tabular}{ll}
Clean filings (1990--2021): & 12,496 \\
Mean $KL_\text{pros}$ / $KL_\text{retr}$ / $\Delta KL$: & 4.47 / 4.48 / -0.008 nats \\
5y survival rate: & 43.2\% \\
Share with SEC-listed owner: & 0.74\% \\
\end{tabular}

\begin{center}
\includegraphics[width=0.7\linewidth]{results/per_industry_deep/class015_scatter.png}
\end{center}

\textbf{Trending tokens (top-decile $\Delta KL$)} --- with intra-class 5y survival rate.

\begin{tabular}{lrrr}
\toprule
token & lift & $n$ filings & intra-class survival \\
\midrule
\texttt{xiao} & 3.33$\times$ & 81 & 14\% \\
\texttt{clamp-like} & 2.92$\times$ & 168 & 27\% \\
\texttt{balalaikas} & 2.89$\times$ & 76 & 28\% \\
\texttt{chords} & 2.85$\times$ & 179 & 27\% \\
\texttt{theaters} & 2.83$\times$ & 53 & 64\% \\
\texttt{positioned} & 2.77$\times$ & 177 & 28\% \\
\texttt{vertical} & 2.72$\times$ & 114 & 18\% \\
\texttt{bowed} & 2.59$\times$ & 58 & 45\% \\
\texttt{structure} & 2.56$\times$ & 191 & 28\% \\
\texttt{harpsichords} & 2.50$\times$ & 60 & 32\% \\
\bottomrule
\end{tabular}

\textbf{Ending tokens (bottom-decile $\Delta KL$)} --- intra-class 5y survival rate.

\begin{tabular}{lrrr}
\toprule
token & lift & $n$ filings & intra-class survival \\
\midrule
\texttt{fretted} & 2.94$\times$ & 51 & 61\% \\
\texttt{viola} & 2.21$\times$ & 95 & 54\% \\
\texttt{cello} & 2.19$\times$ & 64 & 52\% \\
\texttt{trumpet} & 2.16$\times$ & 51 & 61\% \\
\texttt{grand} & 2.14$\times$ & 84 & 21\% \\
\texttt{upright} & 2.11$\times$ & 90 & 29\% \\
\texttt{saxophone} & 1.87$\times$ & 75 & 47\% \\
\texttt{xiao} & 1.73$\times$ & 81 & 14\% \\
\texttt{theaters} & 1.70$\times$ & 53 & 64\% \\
\texttt{winders} & 1.64$\times$ & 61 & 52\% \\
\bottomrule
\end{tabular}

\textbf{Flux-neutral most-frequent tokens ($|\Delta KL|\leq 0.02$)}.

\begin{tabular}{lrrr}
\toprule
token & $n$ in stratum & $n$ in class & intra-class survival \\
\midrule
\texttt{musical} & 594 & 9,775 & 41\% \\
\texttt{instruments} & 541 & 9,256 & 40\% \\
\texttt{electronic} & 281 & 4,189 & 37\% \\
\texttt{music} & 276 & 3,804 & 37\% \\
\texttt{cases} & 259 & 3,732 & 38\% \\
\texttt{guitars} & 231 & 4,440 & 40\% \\
\texttt{stands} & 217 & 3,259 & 32\% \\
\texttt{accessories} & 201 & 2,902 & 40\% \\
\texttt{instrument} & 200 & 2,978 & 40\% \\
\texttt{use} & 193 & 2,540 & 39\% \\
\bottomrule
\end{tabular}

\textbf{Top-10 SEC-listed firms with debut in 1990--2021} (ranked by max observed revenue, debut filing g/s shown).

\begin{tabular}{p{4.0cm}p{2.2cm}rrrp{6.5cm}}
\toprule
Firm & Mark & year & $KL_\text{pros}$ & $\Delta KL$ & First-filing g/s excerpt \\
\midrule
Wal-Mart Stores, Inc. & MISTLETOE MAGIC & 1994 & 4.32 & -0.10 & music boxes \\
APPLE INC. &  & 2004 & 3.24 & -0.25 & [ automobiles ] electronic apparatus and software therefor for downloading, transmitting, receiving, playing, storing and organizing audio, video and navigation data for use in automobiles [, railway  \\
Apple Inc. & OSSIA & 2018 & 1.79 & -0.04 & Musical entertainment services, namely, live performances by musical bands; post-production music editing; music publishing services; consulting services in the field of music editing and production;  \\
WELLPOINT, INC. & C-6000 & 1993 & 5.62 & -0.08 & cymbals and percussion instruments \\
FORD MOTOR COMPANY & OFFICIAL LICENSED PRODUCT FORD & 1995 & 4.37 & -0.11 & music boxes \\
Meta Platforms, Inc. & LF & 2019 & 2.59 & -0.15 & Computer peripheral equipment; computer software, namely, adapted to the operation of computers, downloadable from the internet for operation of computers, operating system software for computers, for \\
INTERNATIONAL BUSINESS MACHINES CORPORATION & M & 1993 & 4.01 & +0.45 & music synthesizers computer hardware, and computer components, namely, communications adapters, data modems, telephone answering machines, sound generators, music synthesizers, printed circuit boards, \\
Coca-Cola Company, The & COCA-COLA & 1991 & 5.55 & +0.31 & glass holders and wastebaskets kitchen linens; namely, table runners, dishcloths, napkins and placemats toy banks metal signs [ musical diaramas ] \\
CBS Inc. & DR. QUINN MEDICINE WOMAN & 1995 & 2.63 & -0.04 & aprons, belts, suspenders, Halloween costumes and masks sold in connection therewith, role-playing costumes, girls' dresses, boots, rubbers, slippers, gloves, mittens, knit caps, painters caps, athlet \\
Coupang Corp. & COUPANG & 2020 & 1.28 & -0.16 & Leather trimmings for furniture; all purpose carrying bags; book bags; duffel bags; leather; leather thongs for fastening or securing items; boxes made of leather; collars for animals; nose bags being \\
\bottomrule
\end{tabular}
\clearpage

\subsection*{Class 016: Paper, Print \& Office}
\label{app:deep-016}

\begin{tabular}{ll}
Clean filings (1990--2021): & 444,018 \\
Mean $KL_\text{pros}$ / $KL_\text{retr}$ / $\Delta KL$: & 5.39 / 5.40 / -0.008 nats \\
5y survival rate: & 50.2\% \\
Share with SEC-listed owner: & 3.87\% \\
\end{tabular}

\begin{center}
\includegraphics[width=0.7\linewidth]{results/per_industry_deep/class016_scatter.png}
\end{center}

\textbf{Trending tokens (top-decile $\Delta KL$)} --- with intra-class 5y survival rate.

\begin{tabular}{lrrr}
\toprule
token & lift & $n$ filings & intra-class survival \\
\midrule
\texttt{originating} & 5.05$\times$ & 1,118 & 44\% \\
\texttt{stand-alone} & 4.26$\times$ & 1,631 & 83\% \\
\texttt{original} & 3.92$\times$ & 2,215 & 45\% \\
\texttt{giclee} & 3.82$\times$ & 521 & 39\% \\
\texttt{scoring} & 3.73$\times$ & 2,401 & 77\% \\
\texttt{watercolour} & 3.26$\times$ & 1,495 & 29\% \\
\texttt{murals} & 3.10$\times$ & 2,131 & 39\% \\
\texttt{tarot} & 3.08$\times$ & 503 & 44\% \\
\texttt{disposing} & 3.06$\times$ & 1,208 & 46\% \\
\texttt{calligraphic} & 3.05$\times$ & 1,250 & 23\% \\
\bottomrule
\end{tabular}

\textbf{Ending tokens (bottom-decile $\Delta KL$)} --- intra-class 5y survival rate.

\begin{tabular}{lrrr}
\toprule
token & lift & $n$ filings & intra-class survival \\
\midrule
\texttt{periodically} & 4.58$\times$ & 1,884 & 47\% \\
\texttt{published} & 3.87$\times$ & 3,147 & 51\% \\
\texttt{stamped} & 3.79$\times$ & 531 & 36\% \\
\texttt{column} & 3.72$\times$ & 1,289 & 44\% \\
\texttt{section} & 3.52$\times$ & 1,206 & 39\% \\
\texttt{pre-paid} & 3.49$\times$ & 2,321 & 54\% \\
\texttt{issued} & 3.30$\times$ & 512 & 46\% \\
\texttt{devoted} & 3.23$\times$ & 644 & 51\% \\
\texttt{circulation} & 3.20$\times$ & 2,213 & 55\% \\
\texttt{quarterly} & 3.15$\times$ & 653 & 46\% \\
\bottomrule
\end{tabular}

\textbf{Flux-neutral most-frequent tokens ($|\Delta KL|\leq 0.02$)}.

\begin{tabular}{lrrr}
\toprule
token & $n$ in stratum & $n$ in class & intra-class survival \\
\midrule
\texttt{printed} & 21,670 & 171,737 & 48\% \\
\texttt{books} & 19,531 & 161,759 & 53\% \\
\texttt{field} & 18,988 & 134,349 & 47\% \\
\texttt{paper} & 15,296 & 124,983 & 46\% \\
\texttt{materials} & 13,557 & 93,439 & 47\% \\
\texttt{featuring} & 13,256 & 92,664 & 49\% \\
\texttt{educational} & 11,559 & 76,630 & 48\% \\
\texttt{information} & 11,218 & 74,050 & 48\% \\
\texttt{use} & 11,132 & 80,387 & 44\% \\
\texttt{magazines} & 11,115 & 83,541 & 52\% \\
\bottomrule
\end{tabular}

\textbf{Top-10 SEC-listed firms with debut in 1990--2021} (ranked by max observed revenue, debut filing g/s shown).

\begin{tabular}{p{4.0cm}p{2.2cm}rrrp{6.5cm}}
\toprule
Firm & Mark & year & $KL_\text{pros}$ & $\Delta KL$ & First-filing g/s excerpt \\
\midrule
AGE International, Inc. & BLANTON'S & 2021 & 5.14 & -0.02 & wall mounted coat rack; wooden display stands for bottle stoppers, bottles and tasting glasses; wall art made of wood reclaimed from bourbon barrels face coverings being masks for protection against v \\
ATK Entertainment, Inc. & AALAM KHALED & 2020 & 2.39 & +0.02 & Clothing, namely, shirts, t-shirts, pants, dresses, skirts, jeans, slacks, belts, hats, caps being headwear, gloves, ear muffs, headwear, neckwear, sweaters, fleece pullovers, jogging suits, warm-up s \\
Artix Entertainment, LLC & ADVENTUREQUEST & 2008 & 4.86 & -0.25 & paper goods and printed matter, namely, children's books and novels \\
BETTERLIFE PHARMA INC. & THINK BETTERX & 2019 & 4.55 & +0.18 & Medicated oils, salves, and concentrated pastes, namely, medicated skin care preparations; Medicated oils, salves, and concentrated pastes, namely, medicated skin care preparations, all of the foregoi \\
BLH ACQUISITION CO., LLC &  & 2012 & 5.76 & -0.47 & Magnetic coded gift cards Non-magnetically encoded prepaid purchase cards for allowing users to transfer financial value at the point of sale and on-line via a computuer network \\
BetterLife Pharma Inc. & LIVE BETTERX & 2020 & 4.01 & -0.02 & Downloadable electronic publications, namely, newsletters, brochures, reports and guides in the field of cannabis Retail store and wholesale store services featuring industrial hemp food and body care \\
Entertainment Art Company, The & SUNDAY COMICS STORE FEATURING  & 1993 & 5.22 & -0.15 & printed sales catalogues issued periodically relating to goods involving comic strip characters retail store services in the field of memorabilia, collectibles and books involving comic strip characte \\
IPR Pharmaceuticals, Inc. & THE STATIN WITH ALL THE RIGHT  & 2002 & 5.21 & +0.45 & Audiotapes, videotapes, films in the nature of motion picture films and cinematographic film, CD's and DVD's on the subject of treatment of cardiovascular diseases and cardiovascular health; audiotape \\
Information Architects, Incorporated & SGML POWER TOOLS & 1996 & 4.93 & -0.10 & computer software printed material for utilizing SGMI, data modeling language \\
KX TECHNOLOGIES LLC & KX KX INDUSTRIES, L.P. & 1997 & 6.94 & +0.10 & [ newsletters on the subject of water filters and filtration ] water filters for domestic, commercial, and industrial use \\
\bottomrule
\end{tabular}
\clearpage

\subsection*{Class 017: Rubber \& Insulators}
\label{app:deep-017}

\begin{tabular}{ll}
Clean filings (1990--2021): & 51,254 \\
Mean $KL_\text{pros}$ / $KL_\text{retr}$ / $\Delta KL$: & 4.97 / 4.96 / +0.018 nats \\
5y survival rate: & 42.6\% \\
Share with SEC-listed owner: & 4.28\% \\
\end{tabular}

\begin{center}
\includegraphics[width=0.7\linewidth]{results/per_industry_deep/class017_scatter.png}
\end{center}

\textbf{Trending tokens (top-decile $\Delta KL$)} --- with intra-class 5y survival rate.

\begin{tabular}{lrrr}
\toprule
token & lift & $n$ filings & intra-class survival \\
\midrule
\texttt{pla} & 6.03$\times$ & 338 & 22\% \\
\texttt{riders} & 5.79$\times$ & 57 & 16\% \\
\texttt{polylactic} & 5.70$\times$ & 256 & 21\% \\
\texttt{slipping} & 5.23$\times$ & 220 & 25\% \\
\texttt{acrylonitrile} & 4.64$\times$ & 265 & 23\% \\
\texttt{abs} & 4.61$\times$ & 425 & 27\% \\
\texttt{sunlight} & 4.47$\times$ & 179 & 20\% \\
\texttt{filaments} & 3.79$\times$ & 955 & 25\% \\
\texttt{anti-slip} & 3.66$\times$ & 391 & 26\% \\
\texttt{potted} & 3.65$\times$ & 52 & 23\% \\
\bottomrule
\end{tabular}

\textbf{Ending tokens (bottom-decile $\Delta KL$)} --- intra-class 5y survival rate.

\begin{tabular}{lrrr}
\toprule
token & lift & $n$ filings & intra-class survival \\
\midrule
\texttt{major} & 2.78$\times$ & 72 & 51\% \\
\texttt{incorporation} & 2.39$\times$ & 67 & 61\% \\
\texttt{banknotes} & 2.37$\times$ & 59 & 24\% \\
\texttt{pry} & 2.28$\times$ & 57 & 42\% \\
\texttt{assistant} & 2.14$\times$ & 56 & 61\% \\
\texttt{dual} & 2.13$\times$ & 89 & 43\% \\
\texttt{caulks} & 2.13$\times$ & 122 & 61\% \\
\texttt{disconnects} & 2.12$\times$ & 52 & 56\% \\
\texttt{packagings} & 2.10$\times$ & 62 & 52\% \\
\texttt{self-adhering} & 2.07$\times$ & 58 & 33\% \\
\bottomrule
\end{tabular}

\textbf{Flux-neutral most-frequent tokens ($|\Delta KL|\leq 0.02$)}.

\begin{tabular}{lrrr}
\toprule
token & $n$ in stratum & $n$ in class & intra-class survival \\
\midrule
\texttt{use} & 4,259 & 31,292 & 39\% \\
\texttt{plastic} & 3,060 & 21,992 & 37\% \\
\texttt{materials} & 2,491 & 15,496 & 31\% \\
\texttt{industrial} & 2,272 & 14,228 & 35\% \\
\texttt{metal} & 1,953 & 11,077 & 32\% \\
\texttt{rubber} & 1,891 & 12,161 & 33\% \\
\texttt{form} & 1,736 & 11,352 & 31\% \\
\texttt{manufacture} & 1,540 & 10,280 & 34\% \\
\texttt{insulating} & 1,516 & 9,706 & 28\% \\
\texttt{purposes} & 1,483 & 8,690 & 27\% \\
\bottomrule
\end{tabular}

\textbf{Top-10 SEC-listed firms with debut in 1990--2021} (ranked by max observed revenue, debut filing g/s shown).

\begin{tabular}{p{4.0cm}p{2.2cm}rrrp{6.5cm}}
\toprule
Firm & Mark & year & $KL_\text{pros}$ & $\Delta KL$ & First-filing g/s excerpt \\
\midrule
INA ACQUISITION CORP. &  & 2001 & 4.99 & -0.13 & Liners and replacement pipes of polymeric materials, including laminated tubular liners of thermosetting resins and resin absorbent materials with polymeric coatings and plastic liners applied against \\
INA Acquisition Corp. &  & 2001 & 5.59 & +0.03 & Non-metal liners and replacement pipes for repair of existing pipelines and passageways of substantially rigid thermoplastic material \\
Wolverine Holding Company, LLC & PRECISION ENGINEERED PRODUCTS & 2002 & 4.44 & -0.03 & Metal hose stations for the food, dairy, beverage, pharmaceutical, health and beauty, and processing industries; metal tubing; metal pipe fittings; metal gaskets for the food, dairy, beverage, pharmac \\
WAL-MART STORES, INC. & EXPERT GARDENER & 2003 & 5.03 & -0.09 & Wheel barrows Garden apparel, namely, aprons, hats, boots Lawn sprinklers; gardening gloves Plant food; lawn and plant fertilizer for domestic use; potting soil; compost; spike-shaped fertilizer for h \\
Wal-Mart Stores, Inc & AUTO DRIVE & 2017 & 3.55 & -0.01 & Automobile tire inflator sealers; flushing agents for automobile radiators All purpose cleaning preparations with deodorizing properties; automobile carpet and upholstery cleaning preparations; carbur \\
Wal-Mart Stores, Inc. & M & 1998 & 2.43 & -0.08 & Acoustical insulation barrier panels, Acoustical insulation for buildings, Acoustical panels for buildings, Acrylic resin sheeting for use in the manufacture of laminated glass, Adhesive bands for sea \\
Wal-mart Stores, Inc. &  & 2011 & 4.21 & -0.04 & On-line retail store services featuring a wide variety of consumer goods of others; Retail department store services; Retail store services featuring a wide variety of consumer goods of others Air bru \\
EXXON MOBIL CORPORATION & LABEL-LYTE & 1991 & 6.02 & -0.15 & polypropylene film for use in the further manufacture of label facings \\
Exxon Mobil Corporation & XX & 1998 & 3.79 & -0.03 & Synthetic rubber for use in the manufacture of tires, tubes, vehicle body sealing, belts, hoses, seals, [ glass run channels, ] mechanical goods, [ electrical cable jacketing, roof sheeting, garden ho \\
ConocoPhillips Company & CARBON SOLUTIONS FROM CONOCOPH & 2003 & 4.52 & -0.37 & Carbon fibers, not for textile use, namely, carbon fibers and carbon fiber composites based on a wide range of matrices, namely, polymeric, asphaltic, cementitious, metallic, ceramic and carbonaceous, \\
\bottomrule
\end{tabular}
\clearpage

\subsection*{Class 018: Leather \& Luggage}
\label{app:deep-018}

\begin{tabular}{ll}
Clean filings (1990--2021): & 147,807 \\
Mean $KL_\text{pros}$ / $KL_\text{retr}$ / $\Delta KL$: & 4.32 / 4.31 / +0.008 nats \\
5y survival rate: & 43.4\% \\
Share with SEC-listed owner: & 1.98\% \\
\end{tabular}

\begin{center}
\includegraphics[width=0.7\linewidth]{results/per_industry_deep/class018_scatter.png}
\end{center}

\textbf{Trending tokens (top-decile $\Delta KL$)} --- with intra-class 5y survival rate.

\begin{tabular}{lrrr}
\toprule
token & lift & $n$ filings & intra-class survival \\
\midrule
\texttt{restraining} & 5.48$\times$ & 4,738 & 31\% \\
\texttt{locking} & 5.42$\times$ & 4,104 & 28\% \\
\texttt{tie-out} & 5.11$\times$ & 800 & 41\% \\
\texttt{pigs} & 4.34$\times$ & 583 & 17\% \\
\texttt{attaching} & 4.20$\times$ & 1,993 & 35\% \\
\texttt{bellybands} & 4.20$\times$ & 693 & 24\% \\
\texttt{coin-sized} & 4.11$\times$ & 158 & 20\% \\
\texttt{disposing} & 3.91$\times$ & 2,114 & 40\% \\
\texttt{stakes} & 3.82$\times$ & 1,104 & 40\% \\
\texttt{dogs} & 3.47$\times$ & 4,421 & 35\% \\
\bottomrule
\end{tabular}

\textbf{Ending tokens (bottom-decile $\Delta KL$)} --- intra-class 5y survival rate.

\begin{tabular}{lrrr}
\toprule
token & lift & $n$ filings & intra-class survival \\
\midrule
\texttt{totebags} & 2.84$\times$ & 444 & 65\% \\
\texttt{windshirts} & 2.72$\times$ & 224 & 71\% \\
\texttt{blemish} & 2.71$\times$ & 155 & 60\% \\
\texttt{soaks} & 2.43$\times$ & 152 & 55\% \\
\texttt{brief-case} & 2.39$\times$ & 209 & 64\% \\
\texttt{checker} & 2.31$\times$ & 290 & 61\% \\
\texttt{ointment} & 2.27$\times$ & 181 & 65\% \\
\texttt{bronzing} & 2.25$\times$ & 169 & 70\% \\
\texttt{fanny} & 2.16$\times$ & 9,633 & 66\% \\
\texttt{ready-to-eat} & 2.16$\times$ & 190 & 55\% \\
\bottomrule
\end{tabular}

\textbf{Flux-neutral most-frequent tokens ($|\Delta KL|\leq 0.02$)}.

\begin{tabular}{lrrr}
\toprule
token & $n$ in stratum & $n$ in class & intra-class survival \\
\midrule
\texttt{bags} & 13,893 & 116,954 & 42\% \\
\texttt{wallets} & 8,227 & 69,301 & 40\% \\
\texttt{leather} & 7,898 & 65,667 & 35\% \\
\texttt{clothing} & 7,829 & 60,849 & 39\% \\
\texttt{handbags} & 7,826 & 66,332 & 38\% \\
\texttt{backpacks} & 7,351 & 62,116 & 41\% \\
\texttt{cases} & 7,294 & 58,149 & 39\% \\
\texttt{purses} & 6,662 & 55,490 & 40\% \\
\texttt{sports} & 6,256 & 47,630 & 36\% \\
\texttt{shirts} & 6,029 & 43,166 & 43\% \\
\bottomrule
\end{tabular}

\textbf{Top-10 SEC-listed firms with debut in 1990--2021} (ranked by max observed revenue, debut filing g/s shown).

\begin{tabular}{p{4.0cm}p{2.2cm}rrrp{6.5cm}}
\toprule
Firm & Mark & year & $KL_\text{pros}$ & $\Delta KL$ & First-filing g/s excerpt \\
\midrule
AGE International, Inc. & BLANTON'S & 2021 & 4.84 & +0.03 & wall mounted coat rack; wooden display stands for bottle stoppers, bottles and tasting glasses; wall art made of wood reclaimed from bourbon barrels face coverings being masks for protection against v \\
ALTO CO., LTD. & BELLAVIENTE & 2021 & 3.86 & +0.02 & Winter gloves of leather; neckties; scarves; headwear; belts, namely, clothing belts made from imitation leather, and clothing belts of textile; clothing, namely, casual shirts, dress shirts, pants, s \\
ATK Entertainment, Inc. & AALAM TUCK KHALED & 2020 & 2.09 & +0.01 & Clothing, namely, shirts, t-shirts, pants, dresses, skirts, jeans, slacks, belts, hats, caps being headwear, gloves, ear muffs, headwear, neckwear, sweaters, fleece pullovers, jogging suits, warm-up s \\
Kabushiki Kaisha Sunrise & G-SAVIOUR & 1996 & 2.91 & +0.05 & bags, namely, all purpose athletic bags, all purpose sport bags, animal game bags, athletic bags, barrel bags, beach bags, book bags, Boston bags, carry-on bags, clutch bags, diaper bags, duffel bags, \\
Lollipop AB & LOLLIPOP & 1996 & 3.90 & -0.02 & stationery, paper materials for packing paper and cardboard boxes, paper and plastic bags for merchandise packaging, photo and memory albums and scrapbooks, calendars and diaries, bookmarkers braid, h \\
Omagine, Inc. & OMAGINE & 2008 & 2.26 & -0.18 & Athletic bags; Baby backpacks; All-purpose athletic bags; Beach bags; Book bags; Diaper bags; Duffel bags; Gym bags; Coin purses; Hand bags; Purses; Tote bags; Leather, mesh and textile shopping bags; \\
Plantation Holdings Group, Inc. & THE PLANTATION SHOP & 2004 & 4.92 & +0.00 & Gourmet cookies, coffee, wafers, biscuits, snack nut mix consisting primarily of chocolate covered nuts and pralines, candy, toffee and cheese straws Jewelry Retail store services and retail shop at h \\
WAL-MART STORES, INC. & DW & 2003 & 3.92 & -0.06 & Jewelry handbags; wallets; purses Bath and body products, namely, bath, body, baby and talcum powders, bath, body and skin soaps in liquid solid or gel form, bath beads, crystals, foams, gels, lotions \\
Wal-Mart Stores, Inc & DAILY'S & 2000 & 3.09 & -0.09 & Cooking appliances, namely, electric Dutch ovens, food steamers, electric skillets, fryers, microwaves, toasters, breadmakers, popcorn poppers, waffle irons, electric coffee pots, coffee makers; elect \\
Wal-Mart Stores, Inc. & CATCH ONE & 1990 & 4.67 & +0.36 & denim and cloth backpacks and bookbags women's clothing and outerwear, namely blouses, tops, shorts, jackets, coats, sweaters, and scarves \\
\bottomrule
\end{tabular}
\clearpage

\subsection*{Class 019: Building Materials}
\label{app:deep-019}

\begin{tabular}{ll}
Clean filings (1990--2021): & 70,015 \\
Mean $KL_\text{pros}$ / $KL_\text{retr}$ / $\Delta KL$: & 5.08 / 5.10 / -0.015 nats \\
5y survival rate: & 45.6\% \\
Share with SEC-listed owner: & 3.45\% \\
\end{tabular}

\begin{center}
\includegraphics[width=0.7\linewidth]{results/per_industry_deep/class019_scatter.png}
\end{center}

\textbf{Trending tokens (top-decile $\Delta KL$)} --- with intra-class 5y survival rate.

\begin{tabular}{lrrr}
\toprule
token & lift & $n$ filings & intra-class survival \\
\midrule
\texttt{tournaments} & 5.43$\times$ & 70 & 59\% \\
\texttt{cologne} & 4.59$\times$ & 74 & 66\% \\
\texttt{luxury} & 4.37$\times$ & 247 & 25\% \\
\texttt{fundraising} & 4.34$\times$ & 99 & 66\% \\
\texttt{unmounted} & 4.30$\times$ & 79 & 72\% \\
\texttt{reference} & 4.12$\times$ & 80 & 54\% \\
\texttt{camps} & 4.09$\times$ & 83 & 63\% \\
\texttt{comic} & 3.95$\times$ & 86 & 60\% \\
\texttt{wrappers} & 3.92$\times$ & 79 & 58\% \\
\texttt{breakfast} & 3.86$\times$ & 88 & 48\% \\
\bottomrule
\end{tabular}

\textbf{Ending tokens (bottom-decile $\Delta KL$)} --- intra-class 5y survival rate.

\begin{tabular}{lrrr}
\toprule
token & lift & $n$ filings & intra-class survival \\
\midrule
\texttt{j-channels} & 4.74$\times$ & 76 & 43\% \\
\texttt{prehung} & 4.10$\times$ & 78 & 47\% \\
\texttt{log} & 3.17$\times$ & 278 & 50\% \\
\texttt{fabric-covered} & 3.04$\times$ & 79 & 61\% \\
\texttt{entranceways} & 2.87$\times$ & 143 & 34\% \\
\texttt{solariums} & 2.84$\times$ & 74 & 50\% \\
\texttt{castable} & 2.83$\times$ & 300 & 55\% \\
\texttt{oriented} & 2.67$\times$ & 412 & 49\% \\
\texttt{starter} & 2.67$\times$ & 225 & 41\% \\
\texttt{prefinished} & 2.54$\times$ & 114 & 44\% \\
\bottomrule
\end{tabular}

\textbf{Flux-neutral most-frequent tokens ($|\Delta KL|\leq 0.02$)}.

\begin{tabular}{lrrr}
\toprule
token & $n$ in stratum & $n$ in class & intra-class survival \\
\midrule
\texttt{building} & 3,753 & 24,574 & 39\% \\
\texttt{non-metal} & 3,588 & 25,190 & 42\% \\
\texttt{materials} & 3,418 & 22,246 & 39\% \\
\texttt{metal} & 3,143 & 19,399 & 36\% \\
\texttt{use} & 2,837 & 18,354 & 39\% \\
\texttt{construction} & 2,431 & 15,130 & 38\% \\
\texttt{panels} & 2,216 & 14,011 & 38\% \\
\texttt{plastic} & 1,793 & 10,581 & 38\% \\
\texttt{tiles} & 1,756 & 10,563 & 36\% \\
\texttt{wood} & 1,747 & 12,081 & 40\% \\
\bottomrule
\end{tabular}

\textbf{Top-10 SEC-listed firms with debut in 1990--2021} (ranked by max observed revenue, debut filing g/s shown).

\begin{tabular}{p{4.0cm}p{2.2cm}rrrp{6.5cm}}
\toprule
Firm & Mark & year & $KL_\text{pros}$ & $\Delta KL$ & First-filing g/s excerpt \\
\midrule
Cellular Concrete Technologies, Inc. & PRECIOUS PAVERS & 2014 & 6.70 & -0.07 & Light weight cement pavers, namely, stones and pet memorials \\
INA Acquisition Corp. & INTEGRA & 2010 & 5.96 & -0.17 & Nonmetal tubular liners for the lining of sewer systems for the repair or reconstruction of existing sanitary sewer and storm sewer passageways by or on behalf of municipalities, cities, states, and o \\
Next Generation Energy LLC & SUNGALOW & 2007 & 4.92 & -0.01 & On-line retail store services featuring housing and storage structures designed and constructed to enhance solar energy efficiency, namely, mobile homes, transportable homes, prefabricated homes, modu \\
RS Technologies Inc. & POWERON & 2021 & 5.84 & +0.36 & Non-metal utility poles; non-metal utility pole cross arms; non-metal covers for utility poles and utility pole cross arms Technical support services, namely, consulting relating to the use of polyure \\
SPA ACQUISITION CORP. & GAZEBO ORIGINALS & 1998 & 6.17 & +0.08 & Non-metal gazebos sold in prefabricated and kit form \\
WAL-MART STORES, INC. & PALM VALLEY & 2004 & 4.72 & -0.18 & Metal fence panels, posts and stays; metal garden fence panels; metal patio fences; hand-operated metal garden hose reels and cabinets Hand-operated non-metal garden hose reels and cabinets Non-metal  \\
Wal-Mart Stores, Inc & MAINSTAYS BABY & 2016 & 4.32 & -0.03 & Wall and floor art made of natural stone and stone tile Decorative wall art for the home, namely, framed artworks of paper, printed graphics and wall decals Fabric window coverings and treatments, nam \\
Wal-Mart Stores, Inc. & MEMBERS MARK & 1998 & 2.55 & -0.14 & Alabaster, Asphalt, Asphalt composition paving, Asphalt patching compound made of a tar base, Asphalt roof coatings, Asphalt roofing paper, Asphalted felts, Balustrades, Beams of wood and concrete, Bi \\
EXXON MOBIL CORPORATION & EXXONMOBIL & 1999 & 5.54 & +0.17 & asphalts periodical publications gasoline; diesel fuel; aviation fuels; lubricating oils and greases for industrial, marine, aviation, agricultural, automotive and other vehicular applications; heatin \\
Exxon Mobil Corporation & MULTIPAVE & 2002 & 7.14 & +0.35 & Asphalt, pitch and bitumen \\
\bottomrule
\end{tabular}
\clearpage

\subsection*{Class 020: Furniture}
\label{app:deep-020}

\begin{tabular}{ll}
Clean filings (1990--2021): & 165,306 \\
Mean $KL_\text{pros}$ / $KL_\text{retr}$ / $\Delta KL$: & 5.16 / 5.09 / +0.075 nats \\
5y survival rate: & 42.3\% \\
Share with SEC-listed owner: & 2.76\% \\
\end{tabular}

\begin{center}
\includegraphics[width=0.7\linewidth]{results/per_industry_deep/class020_scatter.png}
\end{center}

\textbf{Trending tokens (top-decile $\Delta KL$)} --- with intra-class 5y survival rate.

\begin{tabular}{lrrr}
\toprule
token & lift & $n$ filings & intra-class survival \\
\midrule
\texttt{icons} & 6.03$\times$ & 635 & 26\% \\
\texttt{translucent} & 5.63$\times$ & 671 & 26\% \\
\texttt{slumber} & 5.40$\times$ & 632 & 33\% \\
\texttt{themed} & 5.13$\times$ & 752 & 30\% \\
\texttt{attacks} & 5.02$\times$ & 247 & 30\% \\
\texttt{dangerous} & 4.73$\times$ & 262 & 31\% \\
\texttt{situations} & 4.52$\times$ & 281 & 35\% \\
\texttt{defense} & 4.46$\times$ & 280 & 34\% \\
\texttt{pet-actuated} & 4.23$\times$ & 499 & 21\% \\
\texttt{polypropylene} & 4.07$\times$ & 936 & 28\% \\
\bottomrule
\end{tabular}

\textbf{Ending tokens (bottom-decile $\Delta KL$)} --- intra-class 5y survival rate.

\begin{tabular}{lrrr}
\toprule
token & lift & $n$ filings & intra-class survival \\
\midrule
\texttt{item} & 5.21$\times$ & 614 & 53\% \\
\texttt{non-leather} & 4.47$\times$ & 1,120 & 70\% \\
\texttt{replica} & 3.49$\times$ & 212 & 55\% \\
\texttt{pleated} & 3.00$\times$ & 270 & 43\% \\
\texttt{giftware} & 2.96$\times$ & 213 & 60\% \\
\texttt{souvenir} & 2.90$\times$ & 651 & 60\% \\
\texttt{stationery-type} & 2.89$\times$ & 273 & 64\% \\
\texttt{juvenile} & 2.72$\times$ & 298 & 61\% \\
\texttt{license} & 2.70$\times$ & 2,102 & 63\% \\
\texttt{novelty} & 2.68$\times$ & 5,663 & 54\% \\
\bottomrule
\end{tabular}

\textbf{Flux-neutral most-frequent tokens ($|\Delta KL|\leq 0.02$)}.

\begin{tabular}{lrrr}
\toprule
token & $n$ in stratum & $n$ in class & intra-class survival \\
\midrule
\texttt{furniture} & 5,888 & 69,240 & 34\% \\
\texttt{plastic} & 5,080 & 50,581 & 39\% \\
\texttt{use} & 4,713 & 39,994 & 38\% \\
\texttt{metal} & 4,509 & 45,072 & 35\% \\
\texttt{non-metal} & 3,297 & 32,698 & 38\% \\
\texttt{frames} & 3,109 & 35,284 & 35\% \\
\texttt{storage} & 2,980 & 28,304 & 33\% \\
\texttt{pillows} & 2,907 & 35,588 & 39\% \\
\texttt{containers} & 2,793 & 22,028 & 38\% \\
\texttt{chairs} & 2,596 & 30,226 & 33\% \\
\bottomrule
\end{tabular}

\textbf{Top-10 SEC-listed firms with debut in 1990--2021} (ranked by max observed revenue, debut filing g/s shown).

\begin{tabular}{p{4.0cm}p{2.2cm}rrrp{6.5cm}}
\toprule
Firm & Mark & year & $KL_\text{pros}$ & $\Delta KL$ & First-filing g/s excerpt \\
\midrule
AGE International, Inc. & BLANTON'S & 2021 & 4.67 & -0.06 & wall mounted coat rack; wooden display stands for bottle stoppers, bottles and tasting glasses; wall art made of wood reclaimed from bourbon barrels face coverings being masks for protection against v \\
ALTO CO., LTD. & BELLAVIENTE & 2021 & 4.52 & -0.00 & Winter gloves of leather; neckties; scarves; headwear; belts, namely, clothing belts made from imitation leather, and clothing belts of textile; clothing, namely, casual shirts, dress shirts, pants, s \\
ATK Entertainment, Inc. & AALAM TUCK KHALED & 2020 & 2.29 & -0.11 & Clothing, namely, shirts, t-shirts, pants, dresses, skirts, jeans, slacks, belts, hats, caps being headwear, gloves, ear muffs, headwear, neckwear, sweaters, fleece pullovers, jogging suits, warm-up s \\
American Retail Corporation & JUST DREAM MATTRESS GALLERY & 2013 & 4.89 & -0.05 & Bed covers; Bed pads; Mattress covers; Mattress pads; Pillow cases; Pillow covers; Sheet sets Beds; Box springs; Mattresses \\
Next Generation Energy LLC & ZILLA & 2009 & 5.78 & -0.03 & universal modular mounting system, made primarily of non-metal, the major components of which comprise rails, brackets, a four piece leak proof flashing and standoffs, which supports, reinforces and s \\
WAL-MART STORES, INC. & PINE HOLLOW & 2003 & 6.82 & -0.12 & patio furniture \\
WALMART INC. & SWITCHART & 2014 & 5.53 & +0.17 & Magnetic picture frames \\
Wal-Mart Stores, Inc & DAILY'S & 2000 & 2.75 & -0.02 & Cooking appliances, namely, electric Dutch ovens, food steamers, electric skillets, fryers, microwaves, toasters, breadmakers, popcorn poppers, waffle irons, electric coffee pots, coffee makers; elect \\
Wal-Mart Stores, Inc. & AMERICAN RESIDENCE & 1992 & 4.60 & -0.14 & ceramic toothbrush holders, ceramic soap dishes, ceramic cups, and trash cans rugs, carpets, and wallpaper toss pillows, plastic shower curtain hooks, shower curtain rods, and window coverings; namely \\
Wal-mart Stores, Inc. &  & 2011 & 4.52 & +0.23 & Beds for household pets; Crate covers for pets; Dog kennels; Inflatable pet beds; Kennels with carriers; Pet crates; Pet cushions; Pet furniture; Pet ramp; Pillows for household pets; Portable beds fo \\
\bottomrule
\end{tabular}
\clearpage

\subsection*{Class 021: Household Utensils}
\label{app:deep-021}

\begin{tabular}{ll}
Clean filings (1990--2021): & 219,996 \\
Mean $KL_\text{pros}$ / $KL_\text{retr}$ / $\Delta KL$: & 4.85 / 4.78 / +0.067 nats \\
5y survival rate: & 40.8\% \\
Share with SEC-listed owner: & 2.23\% \\
\end{tabular}

\begin{center}
\includegraphics[width=0.7\linewidth]{results/per_industry_deep/class021_scatter.png}
\end{center}

\textbf{Trending tokens (top-decile $\Delta KL$)} --- with intra-class 5y survival rate.

\begin{tabular}{lrrr}
\toprule
token & lift & $n$ filings & intra-class survival \\
\midrule
\texttt{connects} & 4.55$\times$ & 899 & 28\% \\
\texttt{built} & 4.24$\times$ & 768 & 30\% \\
\texttt{vacuums} & 3.94$\times$ & 553 & 38\% \\
\texttt{repellers} & 3.89$\times$ & 1,803 & 37\% \\
\texttt{elevating} & 3.60$\times$ & 994 & 28\% \\
\texttt{releasing} & 3.60$\times$ & 3,119 & 41\% \\
\texttt{animal-activated} & 3.57$\times$ & 4,097 & 36\% \\
\texttt{feeder} & 3.56$\times$ & 3,215 & 41\% \\
\texttt{therein} & 3.44$\times$ & 3,237 & 41\% \\
\texttt{non-mechanized} & 3.42$\times$ & 1,536 & 31\% \\
\bottomrule
\end{tabular}

\textbf{Ending tokens (bottom-decile $\Delta KL$)} --- intra-class 5y survival rate.

\begin{tabular}{lrrr}
\toprule
token & lift & $n$ filings & intra-class survival \\
\midrule
\texttt{robe} & 4.61$\times$ & 421 & 59\% \\
\texttt{candlestick} & 3.72$\times$ & 247 & 62\% \\
\texttt{non-leather} & 3.70$\times$ & 535 & 64\% \\
\texttt{saute} & 3.46$\times$ & 416 & 54\% \\
\texttt{graduate} & 3.33$\times$ & 372 & 47\% \\
\texttt{checker} & 3.29$\times$ & 283 & 69\% \\
\texttt{lavatory} & 3.28$\times$ & 265 & 51\% \\
\texttt{undergraduate} & 3.19$\times$ & 273 & 51\% \\
\texttt{yo-yos} & 3.16$\times$ & 525 & 59\% \\
\texttt{stir} & 3.13$\times$ & 284 & 58\% \\
\bottomrule
\end{tabular}

\textbf{Flux-neutral most-frequent tokens ($|\Delta KL|\leq 0.02$)}.

\begin{tabular}{lrrr}
\toprule
token & $n$ in stratum & $n$ in class & intra-class survival \\
\midrule
\texttt{use} & 8,685 & 88,182 & 33\% \\
\texttt{household} & 7,561 & 78,836 & 30\% \\
\texttt{containers} & 6,865 & 70,291 & 33\% \\
\texttt{cups} & 6,036 & 62,298 & 37\% \\
\texttt{brushes} & 5,624 & 62,369 & 30\% \\
\texttt{kitchen} & 5,260 & 54,504 & 29\% \\
\texttt{plastic} & 4,995 & 48,673 & 40\% \\
\texttt{mugs} & 4,824 & 48,389 & 44\% \\
\texttt{holders} & 4,782 & 46,572 & 36\% \\
\texttt{paper} & 4,611 & 45,762 & 38\% \\
\bottomrule
\end{tabular}

\textbf{Top-10 SEC-listed firms with debut in 1990--2021} (ranked by max observed revenue, debut filing g/s shown).

\begin{tabular}{p{4.0cm}p{2.2cm}rrrp{6.5cm}}
\toprule
Firm & Mark & year & $KL_\text{pros}$ & $\Delta KL$ & First-filing g/s excerpt \\
\midrule
AGE International, Inc. & BLANTON'S & 2021 & 4.43 & +0.02 & wall mounted coat rack; wooden display stands for bottle stoppers, bottles and tasting glasses; wall art made of wood reclaimed from bourbon barrels face coverings being masks for protection against v \\
ATK Entertainment, Inc. & AALAM KHALED & 2020 & 2.33 & -0.15 & Clothing, namely, shirts, t-shirts, pants, dresses, skirts, jeans, slacks, belts, hats, caps being headwear, gloves, ear muffs, headwear, neckwear, sweaters, fleece pullovers, jogging suits, warm-up s \\
BETTERLIFE PHARMA INC. & BETTERX & 2019 & 4.01 & +0.07 & Printed publications, namely, newsletters, brochures, reports and guides in the field of cannabis bath additives, namely, bath beads, bath bombs, bath crystals and bubble bath; bath herbs; bath oils;  \\
BetterLife Pharma Inc. & LIVE BETTERX & 2020 & 3.56 & -0.19 & Downloadable electronic publications, namely, newsletters, brochures, reports and guides in the field of cannabis Retail store and wholesale store services featuring industrial hemp food and body care \\
MIM Construction AB & CAREY & 2016 & 5.58 & -0.08 & storage containers adapted for use in vehicle boots; vehicle partitions for pets; fitted liners for the cargo area of vehicles; boot tidies, namely, storage bins especially adapted to fit in cars; dog \\
Next Generation Energy LLC & ZILLA & 2009 & 6.28 & -0.21 & universal modular mounting system, made primarily of non-metal, the major components of which comprise rails, brackets, a four piece leak proof flashing and standoffs, which supports, reinforces and s \\
SLA Marketing Group LLC & CB & 2014 & 4.78 & -0.12 & Clothing, namely, scarves, dresses, tops, t-shirts, shirts, jackets, pants, hats Personal lifestyle consulting services; personal fashion consulting services Cosmetic brushes Various colored paper dia \\
WAL-MART STORES, INC. & EXPERT GARDENER & 2003 & 5.07 & -0.11 & Wheel barrows Garden apparel, namely, aprons, hats, boots Lawn sprinklers; gardening gloves Plant food; lawn and plant fertilizer for domestic use; potting soil; compost; spike-shaped fertilizer for h \\
Wal-Mart Stores, Inc & DAILY'S & 2000 & 2.59 & -0.03 & Cooking appliances, namely, electric Dutch ovens, food steamers, electric skillets, fryers, microwaves, toasters, breadmakers, popcorn poppers, waffle irons, electric coffee pots, coffee makers; elect \\
Wal-Mart Stores, Inc. & AMERICAN RESIDENCE & 1992 & 4.69 & -0.15 & ceramic toothbrush holders, ceramic soap dishes, ceramic cups, and trash cans rugs, carpets, and wallpaper toss pillows, plastic shower curtain hooks, shower curtain rods, and window coverings; namely \\
\bottomrule
\end{tabular}
\clearpage

\subsection*{Class 022: Ropes, Tents \& Textiles (raw)}
\label{app:deep-022}

\begin{tabular}{ll}
Clean filings (1990--2021): & 24,515 \\
Mean $KL_\text{pros}$ / $KL_\text{retr}$ / $\Delta KL$: & 4.70 / 4.59 / +0.111 nats \\
5y survival rate: & 40.1\% \\
Share with SEC-listed owner: & 2.31\% \\
\end{tabular}

\begin{center}
\includegraphics[width=0.7\linewidth]{results/per_industry_deep/class022_scatter.png}
\end{center}

\textbf{Trending tokens (top-decile $\Delta KL$)} --- with intra-class 5y survival rate.

\begin{tabular}{lrrr}
\toprule
token & lift & $n$ filings & intra-class survival \\
\midrule
\texttt{infestation} & 4.23$\times$ & 104 & 25\% \\
\texttt{minimize} & 4.21$\times$ & 102 & 27\% \\
\texttt{lice} & 4.17$\times$ & 103 & 27\% \\
\texttt{plastic-covered} & 3.50$\times$ & 306 & 19\% \\
\texttt{belongings} & 3.46$\times$ & 130 & 28\% \\
\texttt{jammed} & 3.40$\times$ & 94 & 30\% \\
\texttt{tensile} & 3.34$\times$ & 826 & 21\% \\
\texttt{stuck} & 3.33$\times$ & 96 & 30\% \\
\texttt{soiled} & 3.32$\times$ & 364 & 21\% \\
\texttt{one-way} & 3.30$\times$ & 191 & 20\% \\
\bottomrule
\end{tabular}

\textbf{Ending tokens (bottom-decile $\Delta KL$)} --- intra-class 5y survival rate.

\begin{tabular}{lrrr}
\toprule
token & lift & $n$ filings & intra-class survival \\
\midrule
\texttt{hairpieces} & 5.33$\times$ & 60 & 53\% \\
\texttt{baler} & 5.19$\times$ & 77 & 45\% \\
\texttt{percent} & 4.61$\times$ & 52 & 77\% \\
\texttt{thc} & 4.44$\times$ & 54 & 78\% \\
\texttt{delta-} & 4.28$\times$ & 56 & 79\% \\
\texttt{cannabis} & 4.12$\times$ & 51 & 84\% \\
\texttt{basis} & 3.55$\times$ & 76 & 71\% \\
\texttt{tetrahydrocannabinol} & 3.33$\times$ & 51 & 78\% \\
\texttt{concentration} & 3.29$\times$ & 88 & 58\% \\
\texttt{deflectors} & 3.20$\times$ & 50 & 44\% \\
\bottomrule
\end{tabular}

\textbf{Flux-neutral most-frequent tokens ($|\Delta KL|\leq 0.02$)}.

\begin{tabular}{lrrr}
\toprule
token & $n$ in stratum & $n$ in class & intra-class survival \\
\midrule
\texttt{use} & 1,031 & 9,878 & 36\% \\
\texttt{bags} & 1,025 & 10,999 & 36\% \\
\texttt{textile} & 795 & 9,060 & 30\% \\
\texttt{plastic} & 762 & 6,723 & 34\% \\
\texttt{materials} & 665 & 7,806 & 29\% \\
\texttt{storage} & 598 & 6,511 & 33\% \\
\texttt{tents} & 590 & 7,779 & 30\% \\
\texttt{metal} & 584 & 5,284 & 34\% \\
\texttt{covers} & 518 & 4,493 & 34\% \\
\texttt{nature} & 487 & 3,655 & 33\% \\
\bottomrule
\end{tabular}

\textbf{Top-10 SEC-listed firms with debut in 1990--2021} (ranked by max observed revenue, debut filing g/s shown).

\begin{tabular}{p{4.0cm}p{2.2cm}rrrp{6.5cm}}
\toprule
Firm & Mark & year & $KL_\text{pros}$ & $\Delta KL$ & First-filing g/s excerpt \\
\midrule
WAL-MART STORES, INC. & FUSION 3-D & 2004 & 4.57 & -0.10 & Clothing and footwear for hunting, namely, pants, shirts, bib overalls, jackets, coats, ponchos, gloves, vests, balaclavas, neck gaiters, scarves, bandanas, hats, boots, shoes, slippers, rainwear, nam \\
Wal-Mart Stores, Inc & HOMETRENDS & 2016 & 3.36 & -0.06 & Storage racks; clothes hangers; packaging containers of plastic; plastic storage containers for commercial or industrial use; plastic boxes; non-metal shoe racks; wardrobes; furniture; mirrors; shelvi \\
Wal-Mart Stores, Inc. & AUTO FX & 1994 & 5.39 & +0.27 & automobile accessories, namely portable cooler bags trunk organizers and spare tire covers pumps for inflating tires net bags for storage \\
APPLE INC. & IPOD & 2005 & 5.24 & +0.15 & Lanyards for portable and handheld digital electronic devices for recording, organizing, transmitting, manipulating, and reviewing text, data, audio, image, and video files [ full line of holders, str \\
Apple Inc. &  & 2007 & 4.61 & +0.05 & Lanyards for portable and handheld digital electronic devices for recording, organizing, transmitting, manipulating, and reviewing text, data, audio, image, and video files Housewares, namely, [ decor \\
McKesson Corporation & MCKESSON & 2005 & 3.60 & +0.06 & [ Cadaver bags; ] bags for medical linen and laundry Sitz baths House mark for a full line of medical and surgical products, equipment and supplies, namely patient apparel consisting of bibs, gowns, c \\
GENERAL MOTORS LLC &  & 2008 & 7.30 & +0.35 & Lanyards for holding badges and passes \\
General Motors LLC & GM & 2021 & 1.65 & -0.07 & Custom manufacture of fuel cells; custom manufacture of motor vehicles and motor vehicle parts Automobile lease financing services; financial information provided by electronic means in the field of m \\
Comcast Corporation & XFINITY & 2021 & 3.10 & +0.11 & Birdhouses; Non-metal hooks; Travel pillows Beach towels; Blanket throws; Table cloths not of paper; Vinyl banners Balloons; Building games; Checkers games; Dice games; Dominoes; Golf balls; Ring game \\
Humana Inc. & FREEWHEELIN & 2009 & 5.83 & +0.12 & lanyards for holding keys or badges; bands for wrapping or binding; bands for wrapping pants legs when riding bikes or bicycles \\
\bottomrule
\end{tabular}
\clearpage

\subsection*{Class 024: Textiles \& Bedding}
\label{app:deep-024}

\begin{tabular}{ll}
Clean filings (1990--2021): & 97,945 \\
Mean $KL_\text{pros}$ / $KL_\text{retr}$ / $\Delta KL$: & 4.25 / 4.21 / +0.042 nats \\
5y survival rate: & 42.9\% \\
Share with SEC-listed owner: & 2.23\% \\
\end{tabular}

\begin{center}
\includegraphics[width=0.7\linewidth]{results/per_industry_deep/class024_scatter.png}
\end{center}

\textbf{Trending tokens (top-decile $\Delta KL$)} --- with intra-class 5y survival rate.

\begin{tabular}{lrrr}
\toprule
token & lift & $n$ filings & intra-class survival \\
\midrule
\texttt{cityscapes} & 4.37$\times$ & 176 & 11\% \\
\texttt{background} & 4.00$\times$ & 275 & 20\% \\
\texttt{bundle} & 3.53$\times$ & 3,700 & 29\% \\
\texttt{pad} & 3.18$\times$ & 4,319 & 32\% \\
\texttt{scenes} & 3.12$\times$ & 250 & 25\% \\
\texttt{swaddling} & 3.05$\times$ & 4,885 & 29\% \\
\texttt{scenery} & 2.95$\times$ & 359 & 17\% \\
\texttt{home-health} & 2.77$\times$ & 1,610 & 30\% \\
\texttt{bio-polymer} & 2.77$\times$ & 1,608 & 30\% \\
\texttt{backdrops} & 2.76$\times$ & 1,377 & 22\% \\
\bottomrule
\end{tabular}

\textbf{Ending tokens (bottom-decile $\Delta KL$)} --- intra-class 5y survival rate.

\begin{tabular}{lrrr}
\toprule
token & lift & $n$ filings & intra-class survival \\
\midrule
\texttt{fingertip} & 3.21$\times$ & 112 & 59\% \\
\texttt{bedskirts} & 3.15$\times$ & 203 & 59\% \\
\texttt{ensembles} & 3.12$\times$ & 208 & 77\% \\
\texttt{solicitation} & 3.01$\times$ & 103 & 59\% \\
\texttt{manufacturers} & 3.00$\times$ & 170 & 56\% \\
\texttt{directed} & 2.88$\times$ & 139 & 65\% \\
\texttt{comforter} & 2.86$\times$ & 1,189 & 63\% \\
\texttt{napery} & 2.85$\times$ & 193 & 41\% \\
\texttt{encasements} & 2.74$\times$ & 208 & 64\% \\
\texttt{mist} & 2.48$\times$ & 113 & 58\% \\
\bottomrule
\end{tabular}

\textbf{Flux-neutral most-frequent tokens ($|\Delta KL|\leq 0.02$)}.

\begin{tabular}{lrrr}
\toprule
token & $n$ in stratum & $n$ in class & intra-class survival \\
\midrule
\texttt{textile} & 4,941 & 49,505 & 34\% \\
\texttt{bed} & 4,795 & 50,531 & 38\% \\
\texttt{towels} & 4,508 & 45,030 & 39\% \\
\texttt{covers} & 4,222 & 45,164 & 35\% \\
\texttt{blankets} & 3,963 & 42,694 & 38\% \\
\texttt{use} & 3,758 & 32,565 & 38\% \\
\texttt{clothing} & 3,293 & 27,312 & 40\% \\
\texttt{sheets} & 3,266 & 35,674 & 40\% \\
\texttt{bags} & 3,139 & 28,143 & 37\% \\
\texttt{table} & 3,139 & 30,422 & 35\% \\
\bottomrule
\end{tabular}

\textbf{Top-10 SEC-listed firms with debut in 1990--2021} (ranked by max observed revenue, debut filing g/s shown).

\begin{tabular}{p{4.0cm}p{2.2cm}rrrp{6.5cm}}
\toprule
Firm & Mark & year & $KL_\text{pros}$ & $\Delta KL$ & First-filing g/s excerpt \\
\midrule
ATK Entertainment, Inc. & AALAM & 2020 & 2.00 & -0.02 & Clothing, namely, shirts, t-shirts, pants, dresses, skirts, jeans, slacks, belts, hats, caps being headwear, gloves, ear muffs, headwear, neckwear, sweaters, fleece pullovers, jogging suits, warm-up s \\
American Retail Corporation & JUST DREAM MATTRESS GALLERY & 2013 & 4.67 & +0.18 & Bed covers; Bed pads; Mattress covers; Mattress pads; Pillow cases; Pillow covers; Sheet sets Beds; Box springs; Mattresses \\
Empirical Ventures, LLC & HOLLYBERRY \& COMPANY & 2014 & 6.58 & +0.34 & Christmas tree skirts; Christmas tree stand covers [Table cloth of textile] \\
Kabushiki Kaisha Sunrise & G-SAVIOUR & 1996 & 3.92 & -0.16 & textiles, namely, napkins, placemats, handkerchiefs, towels, bed covers, and table covers not of paper \\
WAL-MART STORES, INC. & NOBO NO BOUNDARIES & 2003 & 6.08 & -0.58 & Household linens for bed \\
Wal-Mart Stores, Inc & DAILY'S & 2000 & 2.56 & -0.05 & Cooking appliances, namely, electric Dutch ovens, food steamers, electric skillets, fryers, microwaves, toasters, breadmakers, popcorn poppers, waffle irons, electric coffee pots, coffee makers; elect \\
Wal-Mart Stores, Inc. & AMERICAN RESIDENCE & 1992 & 3.50 & +0.10 & shams, comforters, plastic and cloth shower curtains, bed sheets, dust ruffles, table cloths not of paper, window coverings; namely, curtains, and fabric; namely, cotton and polycotton \\
Wal-mart Stores, Inc. & BACKYARD GR LL & 2011 & 5.44 & +0.15 & Barbecue smokers; Barbecues and grills; Ceramic briquettes for use in barbecue grills; Charcoal grills; Gas grills; Lighters primarily for lighting grills, fireplaces and candles; Roasters powered by  \\
McKesson Corporation & MCKESSON & 2005 & 4.11 & +0.10 & [ Cadaver bags; ] bags for medical linen and laundry Sitz baths House mark for a full line of medical and surgical products, equipment and supplies, namely patient apparel consisting of bibs, gowns, c \\
MICROSOFT CORPORATION & MICROSOFT & 1991 & 5.94 & -0.17 & towels and blankets \\
\bottomrule
\end{tabular}
\clearpage

\subsection*{Class 025: Clothing, Footwear \& Headwear}
\label{app:deep-025}

\begin{tabular}{ll}
Clean filings (1990--2021): & 668,593 \\
Mean $KL_\text{pros}$ / $KL_\text{retr}$ / $\Delta KL$: & 4.36 / 4.40 / -0.041 nats \\
5y survival rate: & 51.6\% \\
Share with SEC-listed owner: & 1.39\% \\
\end{tabular}

\begin{center}
\includegraphics[width=0.7\linewidth]{results/per_industry_deep/class025_scatter.png}
\end{center}

\textbf{Trending tokens (top-decile $\Delta KL$)} --- with intra-class 5y survival rate.

\begin{tabular}{lrrr}
\toprule
token & lift & $n$ filings & intra-class survival \\
\midrule
\texttt{mental} & 4.74$\times$ & 741 & 45\% \\
\texttt{cannabis} & 4.73$\times$ & 842 & 65\% \\
\texttt{speaking} & 3.88$\times$ & 902 & 45\% \\
\texttt{facilitating} & 3.81$\times$ & 679 & 50\% \\
\texttt{webinars} & 3.70$\times$ & 671 & 40\% \\
\texttt{brand} & 3.69$\times$ & 1,251 & 68\% \\
\texttt{podcasts} & 3.67$\times$ & 3,407 & 46\% \\
\texttt{leadership} & 3.67$\times$ & 935 & 43\% \\
\texttt{wellness} & 3.59$\times$ & 2,449 & 52\% \\
\texttt{blogs} & 3.57$\times$ & 3,274 & 46\% \\
\bottomrule
\end{tabular}

\textbf{Ending tokens (bottom-decile $\Delta KL$)} --- intra-class 5y survival rate.

\begin{tabular}{lrrr}
\toprule
token & lift & $n$ filings & intra-class survival \\
\midrule
\texttt{mens} & 3.85$\times$ & 2,339 & 67\% \\
\texttt{womens} & 3.76$\times$ & 1,983 & 67\% \\
\texttt{childrens} & 3.27$\times$ & 1,354 & 68\% \\
\texttt{girl} & 3.18$\times$ & 1,482 & 55\% \\
\texttt{boys} & 2.87$\times$ & 3,811 & 54\% \\
\texttt{enhancing} & 2.66$\times$ & 767 & 62\% \\
\texttt{intimate} & 2.64$\times$ & 2,463 & 56\% \\
\texttt{wearing} & 2.54$\times$ & 10,771 & 55\% \\
\texttt{girls} & 2.37$\times$ & 4,983 & 55\% \\
\texttt{cheerleading} & 2.32$\times$ & 1,342 & 58\% \\
\bottomrule
\end{tabular}

\textbf{Flux-neutral most-frequent tokens ($|\Delta KL|\leq 0.02$)}.

\begin{tabular}{lrrr}
\toprule
token & $n$ in stratum & $n$ in class & intra-class survival \\
\midrule
\texttt{shirts} & 47,595 & 452,457 & 51\% \\
\texttt{clothing} & 40,201 & 403,438 & 51\% \\
\texttt{hats} & 39,722 & 358,104 & 52\% \\
\texttt{jackets} & 39,691 & 370,707 & 51\% \\
\texttt{pants} & 38,278 & 356,666 & 52\% \\
\texttt{t-shirts} & 32,486 & 325,880 & 53\% \\
\texttt{caps} & 28,253 & 242,240 & 51\% \\
\texttt{shorts} & 25,313 & 260,271 & 54\% \\
\texttt{footwear} & 24,482 & 205,166 & 50\% \\
\texttt{tops} & 21,513 & 209,252 & 51\% \\
\bottomrule
\end{tabular}

\textbf{Top-10 SEC-listed firms with debut in 1990--2021} (ranked by max observed revenue, debut filing g/s shown).

\begin{tabular}{p{4.0cm}p{2.2cm}rrrp{6.5cm}}
\toprule
Firm & Mark & year & $KL_\text{pros}$ & $\Delta KL$ & First-filing g/s excerpt \\
\midrule
AGE International, Inc. & BLANTON'S & 2021 & 4.93 & +0.01 & wall mounted coat rack; wooden display stands for bottle stoppers, bottles and tasting glasses; wall art made of wood reclaimed from bourbon barrels face coverings being masks for protection against v \\
ALTO CO., LTD. & BELLAVIENTE & 2021 & 4.26 & +0.06 & Winter gloves of leather; neckties; scarves; headwear; belts, namely, clothing belts made from imitation leather, and clothing belts of textile; clothing, namely, casual shirts, dress shirts, pants, s \\
ATK Entertainment, Inc. & ASAHD TUCK KHALED & 2018 & 3.54 & +0.04 & Clothing, namely, shirts, pants, hats, caps, headwear, socks \\
Artix Entertainment, LLC & ADVENTUREQUEST & 2008 & 3.16 & -0.23 & Clothing, namely, hats, caps, golf shirts, t-shirts, sweaters, sweatshirts, jackets, socks, pajamas, underwear and Halloween costumes \\
Atlis Motor Vehicles, Inc. & AMV & 2021 & 4.72 & -0.15 & Clothing, namely, hats; baseball caps; beanie hats; t-shirts; baseball t-shirts; polo shirts; tank tops; jersey t-shirt; hoodies; and sweatshirts \\
BETTERLIFE PHARMA INC. & MOVE BETTERX & 2019 & 4.25 & +0.10 & Printed publications, namely, newsletters, brochures, reports and guides in the field of cannabis bath additives, namely, bath beads, bath bombs, bath crystals and bubble bath; bath herbs; bath oils;  \\
BLH ACQUISITION CO., LLC & DRINK POURN & 2017 & 5.24 & -0.14 & Headwear, namely, caps and hats; Wearable garments and clothing, namely, shirts Bar and restaurant services \\
BetterLife Pharma Inc. & LIVE BETTERX & 2020 & 3.84 & -0.06 & Downloadable electronic publications, namely, newsletters, brochures, reports and guides in the field of cannabis Retail store and wholesale store services featuring industrial hemp food and body care \\
HB Global Enterprises, Inc. & SOUL RECIPROCITY & 2020 & 4.29 & +0.19 & Retail store services featuring subscription boxes containing household items, personal care items, clothing, jewelry and gifts; Retail shops featuring gifts, clothing and jewelry Educational and ente \\
Image International, Inc. & NEW DEFINITIONS & 1993 & 3.70 & -0.14 & wearing apparel; namely, boys' suits, pants, shirts, jackets, topcoats, neck wear, sweat shirts, sweat pants and shorts \\
\bottomrule
\end{tabular}
\clearpage

\subsection*{Class 026: Trim, Lace \& Notions}
\label{app:deep-026}

\begin{tabular}{ll}
Clean filings (1990--2021): & 42,323 \\
Mean $KL_\text{pros}$ / $KL_\text{retr}$ / $\Delta KL$: & 4.47 / 4.34 / +0.124 nats \\
5y survival rate: & 40.4\% \\
Share with SEC-listed owner: & 1.76\% \\
\end{tabular}

\begin{center}
\includegraphics[width=0.7\linewidth]{results/per_industry_deep/class026_scatter.png}
\end{center}

\textbf{Trending tokens (top-decile $\Delta KL$)} --- with intra-class 5y survival rate.

\begin{tabular}{lrrr}
\toprule
token & lift & $n$ filings & intra-class survival \\
\midrule
\texttt{cabbages} & 4.52$\times$ & 84 & 12\% \\
\texttt{lei} & 4.18$\times$ & 146 & 23\% \\
\texttt{topiaries} & 4.03$\times$ & 387 & 16\% \\
\texttt{lemons} & 3.70$\times$ & 119 & 7\% \\
\texttt{orchids} & 3.61$\times$ & 316 & 17\% \\
\texttt{pre-lit} & 3.46$\times$ & 324 & 22\% \\
\texttt{botanicals} & 3.44$\times$ & 697 & 18\% \\
\texttt{ferns} & 3.36$\times$ & 330 & 18\% \\
\texttt{apples} & 3.36$\times$ & 137 & 7\% \\
\texttt{tulips} & 3.31$\times$ & 368 & 14\% \\
\bottomrule
\end{tabular}

\textbf{Ending tokens (bottom-decile $\Delta KL$)} --- intra-class 5y survival rate.

\begin{tabular}{lrrr}
\toprule
token & lift & $n$ filings & intra-class survival \\
\midrule
\texttt{programme} & 5.53$\times$ & 103 & 37\% \\
\texttt{ordered} & 4.87$\times$ & 76 & 28\% \\
\texttt{simulcasting} & 4.65$\times$ & 226 & 37\% \\
\texttt{law} & 4.58$\times$ & 201 & 31\% \\
\texttt{risk} & 4.57$\times$ & 94 & 29\% \\
\texttt{languages} & 4.52$\times$ & 188 & 29\% \\
\texttt{alerts} & 4.32$\times$ & 213 & 31\% \\
\texttt{voices} & 4.11$\times$ & 246 & 37\% \\
\texttt{curricula} & 3.98$\times$ & 138 & 32\% \\
\texttt{scores} & 3.96$\times$ & 182 & 38\% \\
\bottomrule
\end{tabular}

\textbf{Flux-neutral most-frequent tokens ($|\Delta KL|\leq 0.02$)}.

\begin{tabular}{lrrr}
\toprule
token & $n$ in stratum & $n$ in class & intra-class survival \\
\midrule
\texttt{hair} & 1,635 & 25,003 & 35\% \\
\texttt{clothing} & 1,394 & 15,603 & 40\% \\
\texttt{pins} & 1,092 & 12,971 & 37\% \\
\texttt{accessories} & 1,015 & 14,563 & 39\% \\
\texttt{ribbons} & 958 & 11,145 & 31\% \\
\texttt{bands} & 944 & 11,719 & 33\% \\
\texttt{buttons} & 819 & 9,739 & 38\% \\
\texttt{bags} & 800 & 6,509 & 44\% \\
\texttt{clips} & 794 & 10,251 & 37\% \\
\texttt{caps} & 785 & 9,151 & 35\% \\
\bottomrule
\end{tabular}

\textbf{Top-10 SEC-listed firms with debut in 1990--2021} (ranked by max observed revenue, debut filing g/s shown).

\begin{tabular}{p{4.0cm}p{2.2cm}rrrp{6.5cm}}
\toprule
Firm & Mark & year & $KL_\text{pros}$ & $\Delta KL$ & First-filing g/s excerpt \\
\midrule
ATK Entertainment, Inc. & AALAM KHALED & 2020 & 1.95 & -0.01 & Clothing, namely, shirts, t-shirts, pants, dresses, skirts, jeans, slacks, belts, hats, caps being headwear, gloves, ear muffs, headwear, neckwear, sweaters, fleece pullovers, jogging suits, warm-up s \\
Lollipop AB & LOLLIPOP & 1996 & 3.54 & +0.00 & stationery, paper materials for packing paper and cardboard boxes, paper and plastic bags for merchandise packaging, photo and memory albums and scrapbooks, calendars and diaries, bookmarkers braid, h \\
WAL-MART STORES, INC. & GIRL CONNECTION & 2003 & 5.04 & +0.15 & Hair ornaments Handbags Clothing, namely, hats, shoes and belts \\
Wal-Mart Stores, Inc & DAILY'S & 2000 & 2.68 & -0.04 & Cooking appliances, namely, electric Dutch ovens, food steamers, electric skillets, fryers, microwaves, toasters, breadmakers, popcorn poppers, waffle irons, electric coffee pots, coffee makers; elect \\
Wal-Mart Stores, Inc. & COLOR TREND BASICS & 1994 & 5.00 & -0.21 & purses belts hair ornaments \\
MCKESSON CORPORATION & FOOD GUARD & 2001 & 5.69 & +0.01 & [hair nets] disposable latex and vinyl gloves for use in the food service industry [protective clothing for use in the food service area, namely, bouffant caps, beard covers, aprons, sleeve covers, an \\
McKesson Corporation & SUNMARK & 2003 & 3.64 & +0.04 & humidifiers and vaporizers House mark for a line of personal care products, namely, tweezers, nail clippers, nail and cuticle scissors, nail files, razors and razor blades and eyelash curlers House ma \\
Microsoft Corporation &  & 2015 & 1.85 & +0.02 & [ Household or kitchen utensils and ] containers [, namely, rolling pins, spatulas, whisks, and containers; chinaware, namely, plates and ] mugs; beverage glassware; porcelain, pottery and earthenware \\
GENERAL MOTORS COMPANY & H2 & 2001 & 4.67 & +0.05 & Belt buckles for clothing not made of precious metal, embroidered emblems, hat pins, zipper pulls \\
GENERAL MOTORS CORPORATION &  & 1991 & 4.45 & +0.37 & belt buckles not of precious metal \\
\bottomrule
\end{tabular}
\clearpage

\subsection*{Class 027: Carpets \& Wall Hangings}
\label{app:deep-027}

\begin{tabular}{ll}
Clean filings (1990--2021): & 26,727 \\
Mean $KL_\text{pros}$ / $KL_\text{retr}$ / $\Delta KL$: & 4.55 / 4.49 / +0.053 nats \\
5y survival rate: & 40.9\% \\
Share with SEC-listed owner: & 3.14\% \\
\end{tabular}

\begin{center}
\includegraphics[width=0.7\linewidth]{results/per_industry_deep/class027_scatter.png}
\end{center}

\textbf{Trending tokens (top-decile $\Delta KL$)} --- with intra-class 5y survival rate.

\begin{tabular}{lrrr}
\toprule
token & lift & $n$ filings & intra-class survival \\
\midrule
\texttt{logos} & 3.68$\times$ & 133 & 28\% \\
\texttt{anti-breakage} & 3.68$\times$ & 76 & 28\% \\
\texttt{mushiro} & 3.55$\times$ & 124 & 21\% \\
\texttt{voices} & 3.40$\times$ & 50 & 14\% \\
\texttt{hana-mushiro} & 3.36$\times$ & 143 & 24\% \\
\texttt{slopes} & 3.32$\times$ & 698 & 16\% \\
\texttt{resemble} & 3.27$\times$ & 52 & 2\% \\
\texttt{slippage} & 3.27$\times$ & 912 & 30\% \\
\texttt{fire-resistant} & 3.15$\times$ & 746 & 16\% \\
\texttt{multiplayer} & 3.15$\times$ & 54 & 0\% \\
\bottomrule
\end{tabular}

\textbf{Ending tokens (bottom-decile $\Delta KL$)} --- intra-class 5y survival rate.

\begin{tabular}{lrrr}
\toprule
token & lift & $n$ filings & intra-class survival \\
\midrule
\texttt{wallcovering} & 6.26$\times$ & 254 & 30\% \\
\texttt{sample} & 4.51$\times$ & 206 & 43\% \\
\texttt{voices} & 4.20$\times$ & 50 & 14\% \\
\texttt{smooth} & 4.11$\times$ & 56 & 30\% \\
\texttt{resemble} & 4.04$\times$ & 52 & 2\% \\
\texttt{underarm} & 3.96$\times$ & 53 & 13\% \\
\texttt{generated} & 3.89$\times$ & 54 & 4\% \\
\texttt{run} & 3.86$\times$ & 57 & 2\% \\
\texttt{gravity} & 3.82$\times$ & 55 & 13\% \\
\texttt{multiplayer} & 3.70$\times$ & 54 & 0\% \\
\bottomrule
\end{tabular}

\textbf{Flux-neutral most-frequent tokens ($|\Delta KL|\leq 0.02$)}.

\begin{tabular}{lrrr}
\toprule
token & $n$ in stratum & $n$ in class & intra-class survival \\
\midrule
\texttt{mats} & 1,236 & 16,348 & 37\% \\
\texttt{floor} & 1,041 & 14,362 & 35\% \\
\texttt{rugs} & 871 & 11,505 & 37\% \\
\texttt{use} & 807 & 9,061 & 36\% \\
\texttt{textile} & 746 & 8,724 & 34\% \\
\texttt{wall} & 723 & 9,083 & 35\% \\
\texttt{carpets} & 717 & 10,155 & 31\% \\
\texttt{coverings} & 717 & 10,148 & 33\% \\
\texttt{bath} & 672 & 7,736 & 35\% \\
\texttt{plastic} & 630 & 6,553 & 37\% \\
\bottomrule
\end{tabular}

\textbf{Top-10 SEC-listed firms with debut in 1990--2021} (ranked by max observed revenue, debut filing g/s shown).

\begin{tabular}{p{4.0cm}p{2.2cm}rrrp{6.5cm}}
\toprule
Firm & Mark & year & $KL_\text{pros}$ & $\Delta KL$ & First-filing g/s excerpt \\
\midrule
WAL-MART STORES, INC. & HOMETRENDS & 2003 & 6.49 & -0.57 & wall borders \\
Wal-Mart Stores, Inc & DAILY'S & 2000 & 2.53 & -0.01 & Cooking appliances, namely, electric Dutch ovens, food steamers, electric skillets, fryers, microwaves, toasters, breadmakers, popcorn poppers, waffle irons, electric coffee pots, coffee makers; elect \\
Wal-Mart Stores, Inc. & AMERICAN RESIDENCE & 1992 & 4.10 & -0.05 & ceramic toothbrush holders, ceramic soap dishes, ceramic cups, and trash cans rugs, carpets, and wallpaper toss pillows, plastic shower curtain hooks, shower curtain rods, and window coverings; namely \\
GENERAL MOTORS LLC & SATURN & 1992 & 5.55 & +0.58 & carpeting and floor mats for motor vehicles \\
General Motors LLC & Z-SPEC & 2010 & 5.40 & -0.23 & Automotive aftermarket accessories, namely, pedal kits consisting of a gas pedal and a brake pedal, grills for vehicles, spoilers for vehicles, wheels for vehicles, badges for vehicles, automotive int \\
FORD MOTOR COMPANY & MERCURY & 1993 & 4.81 & +0.02 & vehicle floor mats, [ door mats, ] carpet \\
Ford Motor Company & FORD OUTFITTERS & 1999 & 4.96 & -0.09 & floor mats for vehicles and rubber mats \\
LOWE'S COMPANIES, INC. & FASTLOCK & 2001 & 5.75 & +0.25 & Floor covering material consisting of fiberboard or other wood-based material with a laminated decorative surface \\
TESLA, INC. & T & 2009 & 4.49 & -0.15 & [ printed maintenance manuals for others for use in maintaining and repairing automobiles ] book bags; school bags [ mats ] Brake calipers for land vehicles; anti-skid chains, roof-racks, shock absorb \\
Tesla Motors, Inc. & TESLA T & 2009 & 2.40 & -0.10 & vehicle axles; brake pads and brake linings, all for land vehicles; brake calipers, braking systems, bumper bars, clutches, differential gears, drive gears, drive shafts, gear change selectors, transm \\
\bottomrule
\end{tabular}
\clearpage

\subsection*{Class 028: Toys, Games \& Sporting Goods}
\label{app:deep-028}

\begin{tabular}{ll}
Clean filings (1990--2021): & 325,521 \\
Mean $KL_\text{pros}$ / $KL_\text{retr}$ / $\Delta KL$: & 4.94 / 4.94 / -0.006 nats \\
5y survival rate: & 47.1\% \\
Share with SEC-listed owner: & 3.83\% \\
\end{tabular}

\begin{center}
\includegraphics[width=0.7\linewidth]{results/per_industry_deep/class028_scatter.png}
\end{center}

\textbf{Trending tokens (top-decile $\Delta KL$)} --- with intra-class 5y survival rate.

\begin{tabular}{lrrr}
\toprule
token & lift & $n$ filings & intra-class survival \\
\midrule
\texttt{scoring} & 4.69$\times$ & 1,464 & 87\% \\
\texttt{loose-leaf} & 3.79$\times$ & 1,771 & 76\% \\
\texttt{itself} & 3.48$\times$ & 2,536 & 23\% \\
\texttt{securely} & 3.48$\times$ & 2,534 & 23\% \\
\texttt{trimmer} & 3.47$\times$ & 2,552 & 23\% \\
\texttt{palm} & 3.27$\times$ & 1,212 & 19\% \\
\texttt{twisting} & 3.15$\times$ & 907 & 28\% \\
\texttt{grippers} & 3.14$\times$ & 1,896 & 24\% \\
\texttt{developers} & 3.10$\times$ & 984 & 22\% \\
\texttt{noisemaker} & 3.08$\times$ & 1,676 & 18\% \\
\bottomrule
\end{tabular}

\textbf{Ending tokens (bottom-decile $\Delta KL$)} --- intra-class 5y survival rate.

\begin{tabular}{lrrr}
\toprule
token & lift & $n$ filings & intra-class survival \\
\midrule
\texttt{ride-in} & 4.16$\times$ & 358 & 93\% \\
\texttt{scoring} & 3.33$\times$ & 1,464 & 87\% \\
\texttt{squirt} & 3.12$\times$ & 528 & 88\% \\
\texttt{convertible} & 2.78$\times$ & 644 & 43\% \\
\texttt{loose-leaf} & 2.73$\times$ & 1,771 & 76\% \\
\texttt{headcovers} & 2.64$\times$ & 356 & 72\% \\
\texttt{windup} & 2.64$\times$ & 478 & 78\% \\
\texttt{molding} & 2.61$\times$ & 959 & 69\% \\
\texttt{chambers} & 2.59$\times$ & 425 & 48\% \\
\texttt{compound} & 2.54$\times$ & 965 & 73\% \\
\bottomrule
\end{tabular}

\textbf{Flux-neutral most-frequent tokens ($|\Delta KL|\leq 0.02$)}.

\begin{tabular}{lrrr}
\toprule
token & $n$ in stratum & $n$ in class & intra-class survival \\
\midrule
\texttt{games} & 9,844 & 104,843 & 46\% \\
\texttt{toys} & 9,796 & 122,074 & 44\% \\
\texttt{toy} & 8,750 & 109,103 & 42\% \\
\texttt{equipment} & 7,002 & 86,517 & 44\% \\
\texttt{game} & 6,629 & 72,060 & 47\% \\
\texttt{use} & 6,397 & 73,369 & 41\% \\
\texttt{playing} & 5,820 & 60,296 & 47\% \\
\texttt{accessories} & 5,017 & 58,390 & 52\% \\
\texttt{electronic} & 4,977 & 52,730 & 48\% \\
\texttt{balls} & 4,824 & 58,268 & 43\% \\
\bottomrule
\end{tabular}

\textbf{Top-10 SEC-listed firms with debut in 1990--2021} (ranked by max observed revenue, debut filing g/s shown).

\begin{tabular}{p{4.0cm}p{2.2cm}rrrp{6.5cm}}
\toprule
Firm & Mark & year & $KL_\text{pros}$ & $\Delta KL$ & First-filing g/s excerpt \\
\midrule
AGE International, Inc. & BLANTON'S & 2021 & 5.39 & +0.04 & wall mounted coat rack; wooden display stands for bottle stoppers, bottles and tasting glasses; wall art made of wood reclaimed from bourbon barrels face coverings being masks for protection against v \\
ATK Entertainment, Inc. & AALAM & 2020 & 2.24 & -0.12 & Clothing, namely, shirts, t-shirts, pants, dresses, skirts, jeans, slacks, belts, hats, caps being headwear, gloves, ear muffs, headwear, neckwear, sweaters, fleece pullovers, jogging suits, warm-up s \\
Artix Entertainment, LLC & ADVENTUREQUEST & 2008 & 3.85 & -0.19 & Toys, namely, toy figures, dolls, stuffed plush toys, board games and puzzles \\
BIOVIE LLC & CELAFORME & 2021 & 5.32 & -0.36 & Balance boards for improving strength, toning, conditioning, balance, and proprioception; Manually-operated exercise equipment \\
Empirical Ventures, LLC & HOLLYBERRY \& COMPANY & 2014 & 6.60 & +0.10 & Christmas tree skirts; Christmas tree stand covers [Table cloth of textile] \\
KABUSHIKI KAISHA SUNRISE & MOBILE SUIT GUNDAM SEED ECLIPS & 2021 & 3.19 & +0.11 & Toys, namely, toy figures, toy robots and stuffed toys; Toy construction sets; Video game machines; Arcade video game machines; Battery-powered computer game with LCD screen; Game controllers for comp \\
Kabushiki Kaisha Sunrise & G-SAVIOUR & 1996 & 3.63 & -0.21 & television games and software for television games, namely, interactive video games of virtual reality comprised of computer hardware and software; automatic video games in the nature of amusement mac \\
MIM Construction AB & MIMSAFE & 2016 & 5.15 & -0.04 & Cages for household pets Storage containers adapted for use in vehicle boots; vehicle partitions for cargo storage; fitted liners for the cargo area of vehicles; dog guards for use in vehicles; safety \\
Omagine, Inc. & OMAGINE & 2008 & 2.15 & -0.12 & Address books; Almanacs; Appliques in the form of decals; Appointment books; Art prints; Arts and craft paint kits; Autograph books; Baby books; Ball point pens; Baseball cards; Binders; Bookends; Boo \\
Sunrise Incorporated & CRUSH GEAR & 2002 & 2.75 & -0.07 & Accessories for dolls clothes, Action figures, Bubble making wand and solution sets, Doll accessories, Doll cases, Doll clothing, Doll furniture, Doll house furnishings, Dolls; Kites, Musical toys, Pa \\
\bottomrule
\end{tabular}
\clearpage

\subsection*{Class 029: Meat, Dairy \& Prepared Foods}
\label{app:deep-029}

\begin{tabular}{ll}
Clean filings (1990--2021): & 134,683 \\
Mean $KL_\text{pros}$ / $KL_\text{retr}$ / $\Delta KL$: & 4.63 / 4.65 / -0.020 nats \\
5y survival rate: & 51.8\% \\
Share with SEC-listed owner: & 3.42\% \\
\end{tabular}

\begin{center}
\includegraphics[width=0.7\linewidth]{results/per_industry_deep/class029_scatter.png}
\end{center}

\textbf{Trending tokens (top-decile $\Delta KL$)} --- with intra-class 5y survival rate.

\begin{tabular}{lrrr}
\toprule
token & lift & $n$ filings & intra-class survival \\
\midrule
\texttt{relieving} & 5.46$\times$ & 141 & 59\% \\
\texttt{mood} & 5.40$\times$ & 150 & 61\% \\
\texttt{transdermal} & 5.13$\times$ & 185 & 73\% \\
\texttt{depression} & 5.06$\times$ & 170 & 60\% \\
\texttt{salves} & 4.82$\times$ & 303 & 79\% \\
\texttt{relaxation} & 4.55$\times$ & 191 & 60\% \\
\texttt{fatigue} & 4.54$\times$ & 196 & 51\% \\
\texttt{cannabidiol} & 4.52$\times$ & 168 & 89\% \\
\texttt{anxiety} & 4.48$\times$ & 203 & 62\% \\
\texttt{emollients} & 4.40$\times$ & 150 & 66\% \\
\bottomrule
\end{tabular}

\textbf{Ending tokens (bottom-decile $\Delta KL$)} --- intra-class 5y survival rate.

\begin{tabular}{lrrr}
\toprule
token & lift & $n$ filings & intra-class survival \\
\midrule
\texttt{least} & 4.19$\times$ & 172 & 42\% \\
\texttt{franks} & 4.04$\times$ & 166 & 48\% \\
\texttt{arranged} & 3.48$\times$ & 187 & 37\% \\
\texttt{pasteurized} & 3.27$\times$ & 394 & 51\% \\
\texttt{hams} & 3.08$\times$ & 523 & 44\% \\
\texttt{designs} & 3.04$\times$ & 250 & 35\% \\
\texttt{roasts} & 2.96$\times$ & 186 & 47\% \\
\texttt{item} & 2.94$\times$ & 163 & 61\% \\
\texttt{percent} & 2.91$\times$ & 354 & 73\% \\
\texttt{delta-} & 2.89$\times$ & 408 & 75\% \\
\bottomrule
\end{tabular}

\textbf{Flux-neutral most-frequent tokens ($|\Delta KL|\leq 0.02$)}.

\begin{tabular}{lrrr}
\toprule
token & $n$ in stratum & $n$ in class & intra-class survival \\
\midrule
\texttt{food} & 6,117 & 43,355 & 50\% \\
\texttt{vegetables} & 5,426 & 39,964 & 49\% \\
\texttt{fruit} & 5,230 & 35,249 & 50\% \\
\texttt{processed} & 5,133 & 40,733 & 50\% \\
\texttt{meat} & 4,745 & 37,723 & 49\% \\
\texttt{frozen} & 4,680 & 36,038 & 49\% \\
\texttt{dried} & 4,594 & 29,899 & 45\% \\
\texttt{snack} & 4,475 & 28,644 & 52\% \\
\texttt{fruits} & 4,160 & 28,277 & 47\% \\
\texttt{products} & 4,142 & 32,292 & 47\% \\
\bottomrule
\end{tabular}

\textbf{Top-10 SEC-listed firms with debut in 1990--2021} (ranked by max observed revenue, debut filing g/s shown).

\begin{tabular}{p{4.0cm}p{2.2cm}rrrp{6.5cm}}
\toprule
Firm & Mark & year & $KL_\text{pros}$ & $\Delta KL$ & First-filing g/s excerpt \\
\midrule
AGE International, Inc. & BLANTON'S & 2021 & 4.68 & -0.06 & wall mounted coat rack; wooden display stands for bottle stoppers, bottles and tasting glasses; wall art made of wood reclaimed from bourbon barrels face coverings being masks for protection against v \\
ATK Entertainment, Inc. & AALAM KHALED & 2020 & 3.09 & -0.01 & Clothing, namely, shirts, t-shirts, pants, dresses, skirts, jeans, slacks, belts, hats, caps being headwear, gloves, ear muffs, headwear, neckwear, sweaters, fleece pullovers, jogging suits, warm-up s \\
BETTERLIFE PHARMA INC. & BETTERX & 2019 & 2.92 & -0.36 & Printed publications, namely, newsletters, brochures, reports and guides in the field of cannabis bath additives, namely, bath beads, bath bombs, bath crystals and bubble bath; bath herbs; bath oils;  \\
Better For You Wellness, Inc. & BETTER FOR YOU WELLNESS & 2021 & 3.85 & -0.03 & Protein bars, namely, protein supplements formed and packaged as bars Advisory and consultation services in the fields of agronomy, agriculture, and horticulture; agronomic consulting services; agrono \\
BetterLife Pharma Inc. & LIVE BETTERX & 2020 & 2.92 & -0.26 & Downloadable electronic publications, namely, newsletters, brochures, reports and guides in the field of cannabis Retail store and wholesale store services featuring industrial hemp food and body care \\
Golden Global Limited & O MY YOGURT & 2010 & 4.46 & +0.03 & Drinks based on yogurt; Yogurt; Yogurt drinks; Yogurt-based beverages \\
WAL-MART STORES, INC. & SAM'S CHOICE & 1991 & 4.08 & +0.19 & [ processed meats, processed edible seeds and fruit preserves ] breakfast and processed cereals; candy; [ condiments; namely, mustards ] fruit drinks \\
Wal-Mart Stores, Inc. & GREAT VALUE & 1993 & 4.27 & +0.12 & processed fish and poultry[; processed edible seeds;] and soups and soup mixes \\
McKesson Corporation & INSTAPRO & 2006 & 5.99 & +0.05 & Protein for use as a food additive \\
Ford Motor Company & CARACTERS POWERED BY FORD & 2009 & 5.02 & -0.04 & Yogurt, fruit-based snack food, fruit leathers \\
\bottomrule
\end{tabular}
\clearpage

\subsection*{Class 030: Coffee, Tea, Bakery \& Spices}
\label{app:deep-030}

\begin{tabular}{ll}
Clean filings (1990--2021): & 205,442 \\
Mean $KL_\text{pros}$ / $KL_\text{retr}$ / $\Delta KL$: & 4.90 / 4.91 / -0.011 nats \\
5y survival rate: & 51.3\% \\
Share with SEC-listed owner: & 2.72\% \\
\end{tabular}

\begin{center}
\includegraphics[width=0.7\linewidth]{results/per_industry_deep/class030_scatter.png}
\end{center}

\textbf{Trending tokens (top-decile $\Delta KL$)} --- with intra-class 5y survival rate.

\begin{tabular}{lrrr}
\toprule
token & lift & $n$ filings & intra-class survival \\
\midrule
\texttt{cannabidiol} & 4.00$\times$ & 285 & 91\% \\
\texttt{transdermal} & 3.95$\times$ & 205 & 71\% \\
\texttt{cartomizers} & 3.72$\times$ & 234 & 63\% \\
\texttt{salves} & 3.58$\times$ & 427 & 78\% \\
\texttt{anxiety} & 3.35$\times$ & 257 & 65\% \\
\texttt{well-being} & 3.26$\times$ & 353 & 58\% \\
\texttt{refill} & 3.19$\times$ & 1,314 & 58\% \\
\texttt{smokeless} & 3.17$\times$ & 347 & 69\% \\
\texttt{relaxation} & 3.12$\times$ & 260 & 62\% \\
\texttt{fatigue} & 3.07$\times$ & 215 & 49\% \\
\bottomrule
\end{tabular}

\textbf{Ending tokens (bottom-decile $\Delta KL$)} --- intra-class 5y survival rate.

\begin{tabular}{lrrr}
\toprule
token & lift & $n$ filings & intra-class survival \\
\midrule
\texttt{e-liquid} & 4.87$\times$ & 1,024 & 60\% \\
\texttt{refill} & 4.33$\times$ & 1,314 & 58\% \\
\texttt{propylene} & 4.14$\times$ & 249 & 67\% \\
\texttt{cartomizers} & 4.06$\times$ & 234 & 63\% \\
\texttt{glycol} & 3.97$\times$ & 257 & 66\% \\
\texttt{least} & 3.60$\times$ & 361 & 48\% \\
\texttt{designs} & 3.38$\times$ & 314 & 34\% \\
\texttt{glycerin} & 3.23$\times$ & 288 & 66\% \\
\texttt{vaporizer} & 3.22$\times$ & 351 & 74\% \\
\texttt{smokeless} & 3.17$\times$ & 347 & 69\% \\
\bottomrule
\end{tabular}

\textbf{Flux-neutral most-frequent tokens ($|\Delta KL|\leq 0.02$)}.

\begin{tabular}{lrrr}
\toprule
token & $n$ in stratum & $n$ in class & intra-class survival \\
\midrule
\texttt{food} & 8,009 & 56,718 & 50\% \\
\texttt{beverages} & 6,030 & 42,857 & 48\% \\
\texttt{coffee} & 5,960 & 44,886 & 47\% \\
\texttt{tea} & 5,617 & 39,383 & 46\% \\
\texttt{fruit} & 5,609 & 37,759 & 49\% \\
\texttt{snack} & 5,589 & 34,248 & 51\% \\
\texttt{chocolate} & 5,310 & 35,356 & 46\% \\
\texttt{foods} & 5,196 & 32,925 & 51\% \\
\texttt{frozen} & 5,145 & 37,251 & 51\% \\
\texttt{products} & 5,048 & 35,916 & 48\% \\
\bottomrule
\end{tabular}

\textbf{Top-10 SEC-listed firms with debut in 1990--2021} (ranked by max observed revenue, debut filing g/s shown).

\begin{tabular}{p{4.0cm}p{2.2cm}rrrp{6.5cm}}
\toprule
Firm & Mark & year & $KL_\text{pros}$ & $\Delta KL$ & First-filing g/s excerpt \\
\midrule
AGE International, Inc. & BLANTON'S & 2021 & 4.51 & -0.01 & wall mounted coat rack; wooden display stands for bottle stoppers, bottles and tasting glasses; wall art made of wood reclaimed from bourbon barrels face coverings being masks for protection against v \\
ATK Entertainment, Inc. & AALAM & 2020 & 2.99 & +0.01 & Clothing, namely, shirts, t-shirts, pants, dresses, skirts, jeans, slacks, belts, hats, caps being headwear, gloves, ear muffs, headwear, neckwear, sweaters, fleece pullovers, jogging suits, warm-up s \\
BETTERLIFE PHARMA INC. & SLEEP BETTERX & 2019 & 3.05 & -0.25 & Printed publications, namely, newsletters, brochures, reports and guides in the field of cannabis bath additives, namely, bath beads, bath bombs, bath crystals and bubble bath; bath herbs; bath oils;  \\
Better For You Wellness, Inc. & BETTER FOR YOU WELLNESS & 2021 & 3.95 & +0.01 & Protein bars, namely, protein supplements formed and packaged as bars Advisory and consultation services in the fields of agronomy, agriculture, and horticulture; agronomic consulting services; agrono \\
BetterLife Pharma Inc. & LIVE BETTERX & 2020 & 2.98 & -0.16 & Downloadable electronic publications, namely, newsletters, brochures, reports and guides in the field of cannabis Retail store and wholesale store services featuring industrial hemp food and body care \\
Golden Global Limited & O MY YOGURT & 2010 & 6.36 & +0.19 & Frozen yogurt; Frozen yogurt confections \\
Kabushiki Kaisha Sunrise & GUNDAM & 1996 & 4.21 & -0.38 & candies, bubble gum, chewing gum and frozen confections \\
Plantation Holdings Group, Inc. & THE PLANTATION SHOP & 2004 & 4.26 & -0.06 & Gourmet cookies, coffee, wafers, biscuits, snack nut mix consisting primarily of chocolate covered nuts and pralines, candy, toffee and cheese straws Jewelry Retail store services and retail shop at h \\
Sunrise Incorporated & OUTLAW STAR & 2002 & 4.20 & -0.02 & Candy, chocolate, bubble gum and chewing gum, cookies, ice cream, sorbets and confections, namely, frozen confections and confectioners sugar \\
WAL-MART STORES, INC. & SAM'S CHOICE & 1991 & 4.41 & +0.23 & [ processed meats, processed edible seeds and fruit preserves ] breakfast and processed cereals; candy; [ condiments; namely, mustards ] fruit drinks \\
\bottomrule
\end{tabular}
\clearpage

\subsection*{Class 031: Agriculture \& Live Animals}
\label{app:deep-031}

\begin{tabular}{ll}
Clean filings (1990--2021): & 64,135 \\
Mean $KL_\text{pros}$ / $KL_\text{retr}$ / $\Delta KL$: & 4.92 / 4.92 / -0.004 nats \\
5y survival rate: & 49.8\% \\
Share with SEC-listed owner: & 2.21\% \\
\end{tabular}

\begin{center}
\includegraphics[width=0.7\linewidth]{results/per_industry_deep/class031_scatter.png}
\end{center}

\textbf{Trending tokens (top-decile $\Delta KL$)} --- with intra-class 5y survival rate.

\begin{tabular}{lrrr}
\toprule
token & lift & $n$ filings & intra-class survival \\
\midrule
\texttt{opioid} & 8.48$\times$ & 79 & 65\% \\
\texttt{comestibles} & 8.12$\times$ & 80 & 62\% \\
\texttt{epilepsy} & 7.74$\times$ & 93 & 66\% \\
\texttt{addiction} & 7.39$\times$ & 88 & 60\% \\
\texttt{mood} & 6.49$\times$ & 111 & 62\% \\
\texttt{depression} & 6.03$\times$ & 126 & 59\% \\
\texttt{pre-baked} & 5.78$\times$ & 147 & 52\% \\
\texttt{fatigue} & 5.68$\times$ & 118 & 57\% \\
\texttt{glosses} & 5.60$\times$ & 91 & 53\% \\
\texttt{relieving} & 5.57$\times$ & 122 & 69\% \\
\bottomrule
\end{tabular}

\textbf{Ending tokens (bottom-decile $\Delta KL$)} --- intra-class 5y survival rate.

\begin{tabular}{lrrr}
\toprule
token & lift & $n$ filings & intra-class survival \\
\midrule
\texttt{deciduous} & 4.72$\times$ & 106 & 45\% \\
\texttt{percent} & 2.41$\times$ & 555 & 72\% \\
\texttt{nectarines} & 2.38$\times$ & 181 & 47\% \\
\texttt{delta-} & 2.35$\times$ & 639 & 73\% \\
\texttt{sod} & 2.33$\times$ & 476 & 53\% \\
\texttt{less} & 2.31$\times$ & 91 & 79\% \\
\texttt{tetrahydrocannabinol} & 2.28$\times$ & 522 & 74\% \\
\texttt{amounts} & 2.27$\times$ & 75 & 79\% \\
\texttt{wholesalers} & 2.25$\times$ & 71 & 70\% \\
\texttt{concentration} & 2.25$\times$ & 632 & 72\% \\
\bottomrule
\end{tabular}

\textbf{Flux-neutral most-frequent tokens ($|\Delta KL|\leq 0.02$)}.

\begin{tabular}{lrrr}
\toprule
token & $n$ in stratum & $n$ in class & intra-class survival \\
\midrule
\texttt{fresh} & 1,900 & 19,830 & 47\% \\
\texttt{vegetables} & 1,512 & 13,668 & 47\% \\
\texttt{food} & 1,480 & 18,269 & 48\% \\
\texttt{seeds} & 1,280 & 12,362 & 46\% \\
\texttt{fruit} & 1,224 & 10,604 & 48\% \\
\texttt{fruits} & 1,209 & 10,998 & 47\% \\
\texttt{live} & 1,164 & 13,107 & 42\% \\
\texttt{plants} & 1,163 & 12,612 & 45\% \\
\texttt{products} & 1,007 & 9,554 & 46\% \\
\texttt{use} & 960 & 9,176 & 46\% \\
\bottomrule
\end{tabular}

\textbf{Top-10 SEC-listed firms with debut in 1990--2021} (ranked by max observed revenue, debut filing g/s shown).

\begin{tabular}{p{4.0cm}p{2.2cm}rrrp{6.5cm}}
\toprule
Firm & Mark & year & $KL_\text{pros}$ & $\Delta KL$ & First-filing g/s excerpt \\
\midrule
WAL-MART STORES, INC. & GARDEN CLASSICS & 1991 & 6.75 & -0.16 & bagged and potted plants \\
Wal-Mart Stores, Inc & WAY TO CELEBRATE & 2017 & 4.61 & +0.14 & Easter blow mold bunnies, namely, unlit molded plastic toy Easter bunny figurines; Easter egg coloring dye kits Easter blow mold bunnies, namely, pre-lit illuminated molded Easter bunny figurines of p \\
Wal-Mart Stores, Inc. & CAT FOOD & 1991 & 4.51 & +0.09 & cat food \\
Comcast Corporation & COMCAST & 2021 & 3.42 & -0.05 & Candy; Coffee; Cookies; Chocolate candies; Chocolate confections; Chocolate covered pretzels; Chocolates and chocolate-based ready to eat candies and snacks; Kettle corn; Lollipops; Hot chocolate mixe \\
General Electric Company & GE & 2005 & 6.71 & +0.22 & Fluid filters, liquid filtering separators, and gas-liquid filtering separators; apparatus in the nature of filters for removing solids, liquids, and vapors from gases \\
PFIZER INC. & FLAVASURE & 1998 & 5.35 & -0.05 & Meat, meat extracts, preserved meat, and sausages Live animals \\
ARCHER DANIELS MIDLAND COMPANY & SUPERMARKET TO THE WORLD & 1990 & 3.47 & +0.40 & barges, vans and railway cars cosmetic preparations - namely face powder and body powder gin and grain alcohol natural agriculture products; namely, raw peanuts, unprocessed malt, dog and cat food, oa \\
ARCHER-DANIELS-MIDLAND COMPANY & MOORTEC & 1990 & 4.40 & +0.79 & [ animal feed for livestock, poultry and game animals; namely, feed supplements, medicated feeds and mineral feeds; worm expeller preparations; insecticide, pesticide, and rodenticide for domestic and \\
Archer Daniels Midland Company & TRYPTOSINE & 1995 & 4.16 & -0.39 & animal feed premix \\
Archer-Daniels-Midland Company & ADM & 2001 & 3.13 & +0.19 & Barges, vans, trucks and railway cars Staple foods, namely, corn germ, corn syrups, flour, corn starch, cornmeal, grits, spaghetti, sugar, noodles, and macaroni flours, meals and grits; processed semo \\
\bottomrule
\end{tabular}
\clearpage

\subsection*{Class 032: Beer \& Soft Drinks}
\label{app:deep-032}

\begin{tabular}{ll}
Clean filings (1990--2021): & 100,951 \\
Mean $KL_\text{pros}$ / $KL_\text{retr}$ / $\Delta KL$: & 4.18 / 4.20 / -0.014 nats \\
5y survival rate: & 56.6\% \\
Share with SEC-listed owner: & 3.11\% \\
\end{tabular}

\begin{center}
\includegraphics[width=0.7\linewidth]{results/per_industry_deep/class032_scatter.png}
\end{center}

\textbf{Trending tokens (top-decile $\Delta KL$)} --- with intra-class 5y survival rate.

\begin{tabular}{lrrr}
\toprule
token & lift & $n$ filings & intra-class survival \\
\midrule
\texttt{smokeless} & 4.11$\times$ & 124 & 73\% \\
\texttt{sugar-based} & 4.00$\times$ & 130 & 47\% \\
\texttt{cannabidiol} & 3.86$\times$ & 272 & 86\% \\
\texttt{mood} & 3.85$\times$ & 161 & 66\% \\
\texttt{transdermal} & 3.75$\times$ & 192 & 74\% \\
\texttt{relieving} & 3.73$\times$ & 150 & 57\% \\
\texttt{depression} & 3.63$\times$ & 157 & 71\% \\
\texttt{e-liquid} & 3.63$\times$ & 135 & 83\% \\
\texttt{glosses} & 3.62$\times$ & 163 & 56\% \\
\texttt{emollients} & 3.58$\times$ & 173 & 68\% \\
\bottomrule
\end{tabular}

\textbf{Ending tokens (bottom-decile $\Delta KL$)} --- intra-class 5y survival rate.

\begin{tabular}{lrrr}
\toprule
token & lift & $n$ filings & intra-class survival \\
\midrule
\texttt{bock} & 4.66$\times$ & 277 & 67\% \\
\texttt{perishable} & 4.59$\times$ & 146 & 86\% \\
\texttt{lager} & 3.54$\times$ & 7,562 & 56\% \\
\texttt{stout} & 3.33$\times$ & 5,990 & 57\% \\
\texttt{porter} & 3.32$\times$ & 6,431 & 57\% \\
\texttt{ale} & 3.25$\times$ & 10,702 & 55\% \\
\texttt{pasteurized} & 3.22$\times$ & 214 & 76\% \\
\texttt{shandy} & 3.09$\times$ & 3,031 & 58\% \\
\texttt{shelf-stable} & 3.06$\times$ & 121 & 85\% \\
\texttt{craft} & 3.00$\times$ & 2,030 & 49\% \\
\bottomrule
\end{tabular}

\textbf{Flux-neutral most-frequent tokens ($|\Delta KL|\leq 0.02$)}.

\begin{tabular}{lrrr}
\toprule
token & $n$ in stratum & $n$ in class & intra-class survival \\
\midrule
\texttt{beverages} & 5,855 & 53,054 & 55\% \\
\texttt{fruit} & 5,207 & 46,913 & 53\% \\
\texttt{drinks} & 5,032 & 51,807 & 55\% \\
\texttt{non-alcoholic} & 4,669 & 40,713 & 53\% \\
\texttt{juices} & 3,853 & 33,108 & 52\% \\
\texttt{water} & 3,412 & 30,730 & 57\% \\
\texttt{making} & 2,663 & 22,491 & 47\% \\
\texttt{soft} & 2,548 & 26,662 & 54\% \\
\texttt{mineral} & 2,525 & 20,232 & 46\% \\
\texttt{juice} & 2,373 & 19,839 & 54\% \\
\bottomrule
\end{tabular}

\textbf{Top-10 SEC-listed firms with debut in 1990--2021} (ranked by max observed revenue, debut filing g/s shown).

\begin{tabular}{p{4.0cm}p{2.2cm}rrrp{6.5cm}}
\toprule
Firm & Mark & year & $KL_\text{pros}$ & $\Delta KL$ & First-filing g/s excerpt \\
\midrule
AGE International, Inc. & BLANTON'S & 2021 & 4.31 & -0.03 & wall mounted coat rack; wooden display stands for bottle stoppers, bottles and tasting glasses; wall art made of wood reclaimed from bourbon barrels face coverings being masks for protection against v \\
ATK Entertainment, Inc. & AALAM TUCK KHALED & 2020 & 2.71 & +0.01 & Clothing, namely, shirts, t-shirts, pants, dresses, skirts, jeans, slacks, belts, hats, caps being headwear, gloves, ear muffs, headwear, neckwear, sweaters, fleece pullovers, jogging suits, warm-up s \\
BETTERLIFE PHARMA INC. & MOVE BETTERX & 2019 & 2.83 & -0.32 & Printed publications, namely, newsletters, brochures, reports and guides in the field of cannabis bath additives, namely, bath beads, bath bombs, bath crystals and bubble bath; bath herbs; bath oils;  \\
Better For You Wellness, Inc. & BETTER FOR YOU WELLNESS & 2021 & 3.73 & +0.02 & Protein bars, namely, protein supplements formed and packaged as bars Advisory and consultation services in the fields of agronomy, agriculture, and horticulture; agronomic consulting services; agrono \\
BetterLife Pharma Inc. & LIVE BETTERX & 2020 & 2.74 & -0.23 & Downloadable electronic publications, namely, newsletters, brochures, reports and guides in the field of cannabis Retail store and wholesale store services featuring industrial hemp food and body care \\
WAL-MART STORES, INC. & SAM'S CHOICE & 1991 & 4.52 & +0.16 & [ processed meats, processed edible seeds and fruit preserves ] breakfast and processed cereals; candy; [ condiments; namely, mustards ] fruit drinks \\
Wal-Mart Stores, Inc. & CONTE & 1992 & 3.20 & +0.13 & soft drinks \\
Exxon Mobil Corporation & CRYO & 2007 & 3.78 & -0.04 & Tea based cold beverages, coffee based cold beverages Fruit flavored beverages, nonalcoholic carbonated beverages, sports beverages, fruit flavored syrups, nonalcoholic frozen carbonated beverages, no \\
MCKESSON CORPORATION & THE WATERSTOP & 1990 & 5.41 & +1.26 & bottled drinking water \\
Cardinal Health, Inc. & CARDINAL BERGEN HEALTH & 1997 & 4.97 & +0.50 & computer software for inventory management, and preparation and administration of advertising and ordering, for drug stores, pharmacies, hospitals, clinics, alternate care centers, and other healthcar \\
\bottomrule
\end{tabular}
\clearpage

\subsection*{Class 033: Wine \& Spirits}
\label{app:deep-033}

\begin{tabular}{ll}
Clean filings (1990--2021): & 68,232 \\
Mean $KL_\text{pros}$ / $KL_\text{retr}$ / $\Delta KL$: & 4.09 / 4.06 / +0.038 nats \\
5y survival rate: & 53.1\% \\
Share with SEC-listed owner: & 0.90\% \\
\end{tabular}

\begin{center}
\includegraphics[width=0.7\linewidth]{results/per_industry_deep/class033_scatter.png}
\end{center}

\textbf{Trending tokens (top-decile $\Delta KL$)} --- with intra-class 5y survival rate.

\begin{tabular}{lrrr}
\toprule
token & lift & $n$ filings & intra-class survival \\
\midrule
\texttt{blue} & 5.76$\times$ & 2,203 & 53\% \\
\texttt{agave} & 5.19$\times$ & 3,090 & 55\% \\
\texttt{weber} & 4.65$\times$ & 71 & 54\% \\
\texttt{tequilana} & 4.65$\times$ & 71 & 55\% \\
\texttt{mexico} & 3.92$\times$ & 153 & 65\% \\
\texttt{region} & 3.53$\times$ & 541 & 31\% \\
\texttt{certain} & 3.52$\times$ & 88 & 5\% \\
\texttt{seltzers} & 3.45$\times$ & 116 & 49\% \\
\texttt{provence} & 3.42$\times$ & 76 & 17\% \\
\texttt{accordance} & 3.39$\times$ & 882 & 33\% \\
\bottomrule
\end{tabular}

\textbf{Ending tokens (bottom-decile $\Delta KL$)} --- intra-class 5y survival rate.

\begin{tabular}{lrrr}
\toprule
token & lift & $n$ filings & intra-class survival \\
\midrule
\texttt{viticultural} & 4.33$\times$ & 134 & 50\% \\
\texttt{napa} & 4.32$\times$ & 148 & 47\% \\
\texttt{valley} & 3.93$\times$ & 168 & 51\% \\
\texttt{called} & 3.59$\times$ & 145 & 17\% \\
\texttt{american} & 3.56$\times$ & 177 & 50\% \\
\texttt{county} & 3.45$\times$ & 84 & 55\% \\
\texttt{advertised} & 3.38$\times$ & 449 & 57\% \\
\texttt{labeled} & 3.38$\times$ & 450 & 58\% \\
\texttt{guaranteed} & 3.30$\times$ & 103 & 11\% \\
\texttt{grown} & 3.20$\times$ & 531 & 56\% \\
\bottomrule
\end{tabular}

\textbf{Flux-neutral most-frequent tokens ($|\Delta KL|\leq 0.02$)}.

\begin{tabular}{lrrr}
\toprule
token & $n$ in stratum & $n$ in class & intra-class survival \\
\midrule
\texttt{alcoholic} & 3,789 & 47,399 & 52\% \\
\texttt{beverages} & 3,485 & 43,642 & 51\% \\
\texttt{except} & 2,408 & 29,046 & 48\% \\
\texttt{beers} & 2,036 & 24,003 & 50\% \\
\texttt{wine} & 1,591 & 17,805 & 52\% \\
\texttt{spirits} & 1,583 & 20,106 & 51\% \\
\texttt{wines} & 1,271 & 16,851 & 44\% \\
\texttt{fruit} & 1,211 & 11,721 & 46\% \\
\texttt{beer} & 1,113 & 12,622 & 49\% \\
\texttt{distilled} & 1,047 & 16,141 & 54\% \\
\bottomrule
\end{tabular}

\textbf{Top-10 SEC-listed firms with debut in 1990--2021} (ranked by max observed revenue, debut filing g/s shown).

\begin{tabular}{p{4.0cm}p{2.2cm}rrrp{6.5cm}}
\toprule
Firm & Mark & year & $KL_\text{pros}$ & $\Delta KL$ & First-filing g/s excerpt \\
\midrule
AGE INTERNATIONAL, INC. & BLANTON'S & 1990 & 4.35 & +0.47 & alcoholic beverages; namely, bourbon whiskey \\
AGE International, Inc. & BLANTON'S & 2021 & 3.77 & -0.02 & wall mounted coat rack; wooden display stands for bottle stoppers, bottles and tasting glasses; wall art made of wood reclaimed from bourbon barrels face coverings being masks for protection against v \\
WAL-MART STORES, INC. & CURIOUSLY J & 2017 & 1.81 & +0.00 & Dyes for use in the manufacture of textiles and fabrics; food coloring; paints; finger paint; inks, toners, and dyes for use with or on textiles and fabrics; edible food coloring Alcoholic cocktail mi \\
Wal-Mart Stores, Inc. & MEMBERS MARK & 1998 & 3.94 & -0.33 & Aperitif wines, Aperitifs with a wine base, Cooking wine, wine, wine cocktail, wine coolers and wine punches \\
ARCHER DANIELS MIDLAND COMPANY & SUPERMARKET TO THE WORLD & 1990 & 4.25 & +0.41 & barges, vans and railway cars cosmetic preparations - namely face powder and body powder gin and grain alcohol natural agriculture products; namely, raw peanuts, unprocessed malt, dog and cat food, oa \\
Archer Daniels Midland Company & AMERICA'S CORN TRADITION & 1994 & 3.60 & +0.22 & alcoholic beverages, namely vodka and distilled spirits \\
PEPSICO, INC. & ROCKARITA & 2017 & 4.01 & -0.07 & Brewed malt-based alcoholic beverages in the nature of beers; flavored malt beverages, namely, brewed malt-based beers Alcoholic beverages produced from a brewed malt base with natural flavors; flavor \\
Tesla, Inc. & TESLAQUILA & 2018 & 5.15 & +0.65 & Distilled agave liquor; Distilled blue agave liquor \\
SYSCO CORPORATION & SYSCO & 1998 & 2.42 & -0.21 & Electric light bulbs, heat lamp bulbs and flood lamps; plumbing supplies, namely, urinal screens; grease and shortening filters for use in range hoods over stoves, fryers, and other cooking equipment  \\
Sysco Corporation & IMPERIAL DYNASTY & 1994 & 3.26 & -0.39 & food colors for domestic use fresh fruits, namely green bell peppers, jalapeno peppers, pasilla chili peppers, purple bell peppers, red bell peppers, serrano peppers, pineapple, bananas, quince, banan \\
\bottomrule
\end{tabular}
\clearpage

\subsection*{Class 034: Tobacco \& Smokers' Articles}
\label{app:deep-034}

\begin{tabular}{ll}
Clean filings (1990--2021): & 36,678 \\
Mean $KL_\text{pros}$ / $KL_\text{retr}$ / $\Delta KL$: & 3.66 / 3.50 / +0.163 nats \\
5y survival rate: & 55.4\% \\
Share with SEC-listed owner: & 1.42\% \\
\end{tabular}

\begin{center}
\includegraphics[width=0.7\linewidth]{results/per_industry_deep/class034_scatter.png}
\end{center}

\textbf{Trending tokens (top-decile $\Delta KL$)} --- with intra-class 5y survival rate.

\begin{tabular}{lrrr}
\toprule
token & lift & $n$ filings & intra-class survival \\
\midrule
\texttt{comestibles} & 8.37$\times$ & 80 & 64\% \\
\texttt{opioid} & 7.64$\times$ & 89 & 67\% \\
\texttt{chronic} & 7.10$\times$ & 62 & 55\% \\
\texttt{epilepsy} & 6.67$\times$ & 108 & 70\% \\
\texttt{depression} & 6.07$\times$ & 122 & 75\% \\
\texttt{addiction} & 5.95$\times$ & 116 & 57\% \\
\texttt{fatigue} & 5.90$\times$ & 122 & 64\% \\
\texttt{enhancement} & 5.62$\times$ & 121 & 71\% \\
\texttt{melts} & 5.38$\times$ & 78 & 74\% \\
\texttt{mood} & 5.35$\times$ & 129 & 75\% \\
\bottomrule
\end{tabular}

\textbf{Ending tokens (bottom-decile $\Delta KL$)} --- intra-class 5y survival rate.

\begin{tabular}{lrrr}
\toprule
token & lift & $n$ filings & intra-class survival \\
\midrule
\texttt{ride} & 7.36$\times$ & 53 & 19\% \\
\texttt{suitcase} & 6.73$\times$ & 52 & 19\% \\
\texttt{cabarets} & 6.49$\times$ & 77 & 23\% \\
\texttt{museum} & 6.45$\times$ & 76 & 20\% \\
\texttt{velvet} & 6.36$\times$ & 55 & 7\% \\
\texttt{blackboards} & 6.35$\times$ & 63 & 19\% \\
\texttt{nameplates} & 6.33$\times$ & 90 & 33\% \\
\texttt{half-boots} & 6.32$\times$ & 57 & 25\% \\
\texttt{detachable} & 6.30$\times$ & 54 & 17\% \\
\texttt{terracotta} & 6.29$\times$ & 97 & 20\% \\
\bottomrule
\end{tabular}

\textbf{Flux-neutral most-frequent tokens ($|\Delta KL|\leq 0.02$)}.

\begin{tabular}{lrrr}
\toprule
token & $n$ in stratum & $n$ in class & intra-class survival \\
\midrule
\texttt{cigarette} & 1,110 & 21,549 & 52\% \\
\texttt{tobacco} & 956 & 18,056 & 53\% \\
\texttt{cigarettes} & 950 & 19,333 & 54\% \\
\texttt{electronic} & 844 & 17,025 & 52\% \\
\texttt{smokers} & 689 & 12,430 & 49\% \\
\texttt{lighters} & 620 & 11,907 & 53\% \\
\texttt{smoking} & 592 & 11,616 & 58\% \\
\texttt{cases} & 565 & 9,079 & 48\% \\
\texttt{cigars} & 556 & 10,468 & 56\% \\
\texttt{liquid} & 550 & 10,534 & 51\% \\
\bottomrule
\end{tabular}

\textbf{Top-10 SEC-listed firms with debut in 1990--2021} (ranked by max observed revenue, debut filing g/s shown).

\begin{tabular}{p{4.0cm}p{2.2cm}rrrp{6.5cm}}
\toprule
Firm & Mark & year & $KL_\text{pros}$ & $\Delta KL$ & First-filing g/s excerpt \\
\midrule
Breathe eCig Corp. & BREATHE ECIGS & 2015 & 4.31 & -0.22 & Electronic cigarettes for use as an alternative to traditional cigarettes \\
JETBLACK LTD. & SIERRA & 2015 & 3.93 & -0.07 & Smokeless cigar vaporizer pipes; Smokeless cigarette vaporizer pipe \\
Rich Cigars, Inc. & RICH CIGARS & 2015 & 4.04 & -0.06 & Cigars, cigar wraps, humidors, cigar boxes, cigar cases, ash trays, match boxes, cigar holders, cigar tubes, lighters for tobacco products, match containers, and cigar cutters \\
WAL-MART STORES, INC. & CURIOUSLY J & 2017 & 2.01 & -0.13 & Dyes for use in the manufacture of textiles and fabrics; food coloring; paints; finger paint; inks, toners, and dyes for use with or on textiles and fabrics; edible food coloring Alcoholic cocktail mi \\
Wal-Mart Stores, Inc. & M & 1998 & 2.56 & -0.16 & Ashtrays not of precious metal, Chewing tobacco, Cigar cutters, Cigarette cases not of precious metal, Cigarette holders not of precious metal, Cigarette lighters not for land vehicles, Cigarette ligh \\
GENERAL MOTORS CORPORATION & CADILLAC & 1990 & 4.95 & +0.54 & cigarette lighters and matches \\
GENERAL MOTORS LLC & HUMMER & 2008 & 4.84 & +0.29 & Lighters for cigars and cigarettes \\
General Motors LLC &  & 2013 & 4.30 & +0.13 & Lighters for cigars and cigarettes not for land vehicles, cigar holders \\
FORD MOTOR COMPANY & AEROSTAR & 1993 & 2.78 & +0.12 & cigarette lighters not of precious metal, cigarette cases not of precious metal, matches and ashtrays not of precious metal \\
Ford Motor Company & THINK & 2000 & 2.73 & -0.02 & smokers' articles, namely, cigarette and cigar lighters not of precious metal, ashtrays not of precious metal, matches, cigarette cases not of precious metal and cigar cutters not of precious metal \\
\bottomrule
\end{tabular}
\clearpage

\subsection*{Class 035: Advertising \& Retail}
\label{app:deep-035}

\begin{tabular}{ll}
Clean filings (1990--2021): & 996,559 \\
Mean $KL_\text{pros}$ / $KL_\text{retr}$ / $\Delta KL$: & 5.38 / 5.40 / -0.019 nats \\
5y survival rate: & 47.0\% \\
Share with SEC-listed owner: & 3.65\% \\
\end{tabular}

\begin{center}
\includegraphics[width=0.7\linewidth]{results/per_industry_deep/class035_scatter.png}
\end{center}

\textbf{Trending tokens (top-decile $\Delta KL$)} --- with intra-class 5y survival rate.

\begin{tabular}{lrrr}
\toprule
token & lift & $n$ filings & intra-class survival \\
\midrule
\texttt{blockchain} & 5.73$\times$ & 1,717 & 53\% \\
\texttt{worlds} & 5.66$\times$ & 1,167 & 63\% \\
\texttt{cryptocurrency} & 5.28$\times$ & 1,395 & 53\% \\
\texttt{delta-} & 4.16$\times$ & 1,621 & 61\% \\
\texttt{tetrahydrocannabinol} & 4.06$\times$ & 1,376 & 62\% \\
\texttt{thc} & 4.01$\times$ & 1,719 & 61\% \\
\texttt{percent} & 3.88$\times$ & 1,491 & 59\% \\
\texttt{cbd} & 3.75$\times$ & 1,793 & 68\% \\
\texttt{hemp} & 3.65$\times$ & 2,978 & 61\% \\
\texttt{currencies} & 3.54$\times$ & 1,023 & 44\% \\
\bottomrule
\end{tabular}

\textbf{Ending tokens (bottom-decile $\Delta KL$)} --- intra-class 5y survival rate.

\begin{tabular}{lrrr}
\toprule
token & lift & $n$ filings & intra-class survival \\
\midrule
\texttt{chamber} & 3.27$\times$ & 2,109 & 37\% \\
\texttt{price-comparison} & 3.03$\times$ & 3,355 & 52\% \\
\texttt{page} & 2.56$\times$ & 3,999 & 50\% \\
\texttt{earned} & 2.45$\times$ & 1,100 & 61\% \\
\texttt{accessed} & 2.45$\times$ & 6,596 & 57\% \\
\texttt{placing} & 2.42$\times$ & 13,179 & 50\% \\
\texttt{dealership} & 2.40$\times$ & 2,143 & 48\% \\
\texttt{billboard} & 2.35$\times$ & 2,537 & 54\% \\
\texttt{rebates} & 2.32$\times$ & 8,195 & 51\% \\
\texttt{county} & 2.31$\times$ & 1,138 & 35\% \\
\bottomrule
\end{tabular}

\textbf{Flux-neutral most-frequent tokens ($|\Delta KL|\leq 0.02$)}.

\begin{tabular}{lrrr}
\toprule
token & $n$ in stratum & $n$ in class & intra-class survival \\
\midrule
\texttt{featuring} & 49,481 & 364,651 & 44\% \\
\texttt{field} & 46,826 & 317,057 & 42\% \\
\texttt{business} & 46,425 & 320,210 & 42\% \\
\texttt{retail} & 43,302 & 328,625 & 46\% \\
\texttt{management} & 38,743 & 249,195 & 42\% \\
\texttt{information} & 37,357 & 258,493 & 45\% \\
\texttt{others} & 35,624 & 273,004 & 45\% \\
\texttt{store} & 35,602 & 259,918 & 44\% \\
\texttt{on-line} & 29,783 & 228,501 & 46\% \\
\texttt{online} & 28,460 & 190,528 & 44\% \\
\bottomrule
\end{tabular}

\textbf{Top-10 SEC-listed firms with debut in 1990--2021} (ranked by max observed revenue, debut filing g/s shown).

\begin{tabular}{p{4.0cm}p{2.2cm}rrrp{6.5cm}}
\toprule
Firm & Mark & year & $KL_\text{pros}$ & $\Delta KL$ & First-filing g/s excerpt \\
\midrule
AGE International, Inc. & BLANTON'S & 2021 & 5.51 & -0.05 & wall mounted coat rack; wooden display stands for bottle stoppers, bottles and tasting glasses; wall art made of wood reclaimed from bourbon barrels face coverings being masks for protection against v \\
ALTO CO., LTD. & ALTA & 2012 & 5.08 & +0.14 & Optical goods, namely, optical reflectors, optical lenses; actinometers; electrical terminal blocks; electric light dimmers regulators; ballasts for fluorescent lamps; connections for electric lines;  \\
ANCHORAGE INTERNATIONAL, INC. & NETEQUIPMENT & 1999 & 4.78 & -0.04 & cooperative advertising and marketing via a global computer network and print media \\
ATK Entertainment, Inc. & AALAM KHALED & 2020 & 3.01 & -0.01 & Clothing, namely, shirts, t-shirts, pants, dresses, skirts, jeans, slacks, belts, hats, caps being headwear, gloves, ear muffs, headwear, neckwear, sweaters, fleece pullovers, jogging suits, warm-up s \\
ATLAS TECHNOLOGY GROUP, INC. & IGE & 2006 & 5.51 & +0.01 & Entertainment services, namely, providing on-line information via the Internet about video games; providing on-line news and information via the Internet on video games; providing on-line databases fo \\
American Retail Corporation & WATSON'S & 2008 & 5.20 & -0.10 & Retail store services and on-line retail store services featuring swimming pools and accessories therefor, hot tubs, spas, saunas and accessories therefor, fireplaces, fireplace inserts, woodburning s \\
BETTERLIFE PHARMA INC. & LOVE BETTERX & 2019 & 4.43 & +0.12 & Printed publications, namely, newsletters, brochures, reports and guides in the field of cannabis bath additives, namely, bath beads, bath bombs, bath crystals and bubble bath; bath herbs; bath oils;  \\
BLH ACQUISITION CO., LLC & BE PAID TO PARTY LOUIE NATION  & 2019 & 5.89 & -0.06 & Providing incentive award programs through issuance and processing of loyalty points for purchase of a company's goods and services \\
BetterLife Pharma Inc. & LIVE BETTERX & 2020 & 3.97 & -0.00 & Downloadable electronic publications, namely, newsletters, brochures, reports and guides in the field of cannabis Retail store and wholesale store services featuring industrial hemp food and body care \\
Canna, LLC & CANNA ADVISORS DESIGN| BUILD|  & 2016 & 5.69 & +0.07 & Business development consulting services \\
\bottomrule
\end{tabular}
\clearpage

\subsection*{Class 036: Insurance \& Financial}
\label{app:deep-036}

\begin{tabular}{ll}
Clean filings (1990--2021): & 399,896 \\
Mean $KL_\text{pros}$ / $KL_\text{retr}$ / $\Delta KL$: & 4.92 / 4.96 / -0.042 nats \\
5y survival rate: & 47.3\% \\
Share with SEC-listed owner: & 6.69\% \\
\end{tabular}

\begin{center}
\includegraphics[width=0.7\linewidth]{results/per_industry_deep/class036_scatter.png}
\end{center}

\textbf{Trending tokens (top-decile $\Delta KL$)} --- with intra-class 5y survival rate.

\begin{tabular}{lrrr}
\toprule
token & lift & $n$ filings & intra-class survival \\
\midrule
\texttt{decentralized} & 7.76$\times$ & 817 & 61\% \\
\texttt{blockchain-based} & 7.49$\times$ & 929 & 62\% \\
\texttt{cryptographic} & 7.18$\times$ & 890 & 62\% \\
\texttt{token} & 7.11$\times$ & 908 & 57\% \\
\texttt{blockchain} & 6.78$\times$ & 2,507 & 59\% \\
\texttt{crypto} & 6.39$\times$ & 638 & 57\% \\
\texttt{cryptocurrency} & 6.10$\times$ & 3,954 & 59\% \\
\texttt{transferrable} & 5.90$\times$ & 402 & 68\% \\
\texttt{cryptocurrencies} & 5.87$\times$ & 828 & 57\% \\
\texttt{validating} & 5.64$\times$ & 456 & 52\% \\
\bottomrule
\end{tabular}

\textbf{Ending tokens (bottom-decile $\Delta KL$)} --- intra-class 5y survival rate.

\begin{tabular}{lrrr}
\toprule
token & lift & $n$ filings & intra-class survival \\
\midrule
\texttt{guaranty} & 4.10$\times$ & 459 & 45\% \\
\texttt{calling} & 3.19$\times$ & 3,618 & 49\% \\
\texttt{annuity} & 2.67$\times$ & 5,246 & 47\% \\
\texttt{timeshare} & 2.40$\times$ & 1,151 & 60\% \\
\texttt{universal} & 2.39$\times$ & 615 & 46\% \\
\texttt{inland} & 2.36$\times$ & 445 & 49\% \\
\texttt{malpractice} & 2.30$\times$ & 556 & 53\% \\
\texttt{underwriting} & 2.26$\times$ & 42,087 & 44\% \\
\texttt{fidelity} & 2.25$\times$ & 636 & 46\% \\
\texttt{casualty} & 2.24$\times$ & 10,709 & 46\% \\
\bottomrule
\end{tabular}

\textbf{Flux-neutral most-frequent tokens ($|\Delta KL|\leq 0.02$)}.

\begin{tabular}{lrrr}
\toprule
token & $n$ in stratum & $n$ in class & intra-class survival \\
\midrule
\texttt{financial} & 19,537 & 172,088 & 47\% \\
\texttt{management} & 16,051 & 130,568 & 45\% \\
\texttt{information} & 12,538 & 102,892 & 45\% \\
\texttt{field} & 12,115 & 105,707 & 44\% \\
\texttt{investment} & 10,544 & 92,365 & 47\% \\
\texttt{estate} & 8,787 & 88,964 & 46\% \\
\texttt{real} & 8,386 & 85,253 & 46\% \\
\texttt{others} & 8,316 & 70,650 & 45\% \\
\texttt{business} & 8,123 & 58,442 & 43\% \\
\texttt{electronic} & 7,133 & 56,743 & 48\% \\
\bottomrule
\end{tabular}

\textbf{Top-10 SEC-listed firms with debut in 1990--2021} (ranked by max observed revenue, debut filing g/s shown).

\begin{tabular}{p{4.0cm}p{2.2cm}rrrp{6.5cm}}
\toprule
Firm & Mark & year & $KL_\text{pros}$ & $\Delta KL$ & First-filing g/s excerpt \\
\midrule
AmeriTrust Group, Inc. & AMERITRUST & 2018 & 4.60 & -0.23 & Insurance services, namely, administration, brokerage, and underwriting in the fields of commercial property, casualty, umbrella, excess and general liability; insurance claims administration and clai \\
American Life Group, Inc. & AMERICAN LIFE GROUP, INC. & 1995 & 4.11 & -0.33 & insurance underwriting services in the fields of life insurance and annuity products \\
BLUE CHIP CAPITAL GROUP, INC. & RAISEWISE & 2020 & 6.01 & -0.30 & Crowdfunding services in the nature of accepting and administering monetary contributions from a group of individuals \\
Blackstar Enterprise Group, Inc. & BLOCKCHAIN FIRST & 2021 & 4.24 & +0.18 & Financial services, namely, providing a securities exchange on an immutable block chain that prevents the disruption of order flow of customer trades; transaction processing and recording, namely, fin \\
China Foundation, The & CHINA FOUNDATION & 2008 & 4.85 & +0.12 & Arranging and conducting educational conferences; organizing conferences in the field of public health and medical science Charitable services, namely, organizing and conducting volunteer programs and \\
Ching Technologies, Inc. & CANAL CIRCLE & 2018 & 2.98 & +0.12 & Computer services for providing on-line, non-downloadable, hosted software applications for credit scoring and data analytics about credit scoring of persons or businesses lacking or having limited ve \\
ENVOY LLC & CHAMBERLIN EDMONDS & 2004 & 6.55 & -0.08 & provider of revenue recovery services, namely, securing government reimbursement for uncompensated medical care expenses \\
INA Acquisition Corp. & INSITUFORM FINANCIAL SERVICES & 2006 & 6.42 & +0.30 & Financial planning; Fiscal assessments of funding in the field of repair and reconstruction of existing sewer pipes, water pipes, gas mains and industrial pipelines; Financial forecasting needs of fun \\
Infinity Real Estate LLC & INFINITY GROUP & 2019 & 3.15 & -0.05 & Assessment and management of real estate; Commercial and residential real estate agency services; Financial planning and investment advisory services; Investment advice; Land acquisition, namely, real \\
Megas Corporation, The & MEGAS & 1998 & 5.50 & -0.30 & Development and management of associations, and development and marketing of benefits for association members marketing and administration of insurance programs for association members \\
\bottomrule
\end{tabular}
\clearpage

\subsection*{Class 037: Construction \& Repair}
\label{app:deep-037}

\begin{tabular}{ll}
Clean filings (1990--2021): & 178,756 \\
Mean $KL_\text{pros}$ / $KL_\text{retr}$ / $\Delta KL$: & 5.18 / 5.18 / +0.001 nats \\
5y survival rate: & 42.5\% \\
Share with SEC-listed owner: & 4.84\% \\
\end{tabular}

\begin{center}
\includegraphics[width=0.7\linewidth]{results/per_industry_deep/class037_scatter.png}
\end{center}

\textbf{Trending tokens (top-decile $\Delta KL$)} --- with intra-class 5y survival rate.

\begin{tabular}{lrrr}
\toprule
token & lift & $n$ filings & intra-class survival \\
\midrule
\texttt{charging} & 3.45$\times$ & 2,326 & 38\% \\
\texttt{recharging} & 3.10$\times$ & 678 & 40\% \\
\texttt{rootering} & 2.96$\times$ & 412 & 35\% \\
\texttt{microorganisms} & 2.83$\times$ & 205 & 40\% \\
\texttt{suitable} & 2.58$\times$ & 329 & 35\% \\
\texttt{spread} & 2.48$\times$ & 222 & 32\% \\
\texttt{waiting} & 2.09$\times$ & 187 & 30\% \\
\texttt{renewable} & 2.07$\times$ & 2,872 & 37\% \\
\texttt{solar-powered} & 2.03$\times$ & 276 & 34\% \\
\texttt{swimming-pool} & 2.00$\times$ & 200 & 28\% \\
\bottomrule
\end{tabular}

\textbf{Ending tokens (bottom-decile $\Delta KL$)} --- intra-class 5y survival rate.

\begin{tabular}{lrrr}
\toprule
token & lift & $n$ filings & intra-class survival \\
\midrule
\texttt{arcades} & 4.08$\times$ & 201 & 70\% \\
\texttt{lounges} & 3.81$\times$ & 194 & 63\% \\
\texttt{casinos} & 3.76$\times$ & 250 & 67\% \\
\texttt{geologic} & 3.66$\times$ & 227 & 44\% \\
\texttt{dialing} & 3.64$\times$ & 198 & 41\% \\
\texttt{convention} & 3.61$\times$ & 524 & 57\% \\
\texttt{dinner} & 3.59$\times$ & 217 & 57\% \\
\texttt{formations} & 3.54$\times$ & 268 & 43\% \\
\texttt{galleries} & 3.35$\times$ & 200 & 66\% \\
\texttt{housewares} & 3.33$\times$ & 183 & 67\% \\
\bottomrule
\end{tabular}

\textbf{Flux-neutral most-frequent tokens ($|\Delta KL|\leq 0.02$)}.

\begin{tabular}{lrrr}
\toprule
token & $n$ in stratum & $n$ in class & intra-class survival \\
\midrule
\texttt{repair} & 13,046 & 85,741 & 39\% \\
\texttt{maintenance} & 11,787 & 76,986 & 39\% \\
\texttt{installation} & 10,028 & 66,270 & 38\% \\
\texttt{construction} & 8,381 & 53,612 & 40\% \\
\texttt{field} & 7,737 & 46,645 & 37\% \\
\texttt{systems} & 6,480 & 40,587 & 37\% \\
\texttt{equipment} & 6,363 & 38,385 & 37\% \\
\texttt{management} & 6,213 & 37,024 & 40\% \\
\texttt{building} & 5,852 & 33,623 & 38\% \\
\texttt{development} & 5,798 & 36,545 & 40\% \\
\bottomrule
\end{tabular}

\textbf{Top-10 SEC-listed firms with debut in 1990--2021} (ranked by max observed revenue, debut filing g/s shown).

\begin{tabular}{p{4.0cm}p{2.2cm}rrrp{6.5cm}}
\toprule
Firm & Mark & year & $KL_\text{pros}$ & $\Delta KL$ & First-filing g/s excerpt \\
\midrule
DNA PRECIOUS METALS INC. & DNA PRECIOUS METALS & 2013 & 7.40 & -0.12 & mining extraction services mining exploration services \\
DONGFANG CITY HOLDING GROUP COMPANY LIMITED & DONGFANGCITY & 2021 & 4.50 & +0.00 & Real estate development; Real estate development and construction of commercial, residential and hotel property; Real estate site selection; Building construction supervision services for real estate  \\
Gen Serv & GENSERV & 2020 & 7.85 & +0.06 & We are a world wide Generator Service company in the Emergency Response company. We work for the DOD and all companies hotels Etc \\
INA ACQUISITION CORP. & INSITUFORM TECHNOLOGIES & 1993 & 6.62 & -0.45 & repair or reconstruction of pipes such as sewer pipes or other conduits by inserting liners therewithin \\
INA Acquisition Corp. & CLEAN WATER FOR THE WORLD & 2006 & 6.23 & +0.11 & Repair and reconstruction of existing sewer pipes, water pipes, gas mains and industrial pipelines by installing liners within the existing conduit and reinstatement of existing services \\
Infinity Real Estate LLC &  & 2019 & 4.61 & +0.06 & Real estate development Assessment and management of real estate; Commercial and residential real estate agency services; Financial planning and investment advisory services; Investment advice; Land a \\
Integrated Technologies, Inc. & INTEGRATED TECHNOLOGIES, INC. & 2016 & 3.88 & -0.07 & Consulting services in the field of factory construction; consulting services in the field of renovation planning; construction of buildings Engineering services; design of buildings; engineering desi \\
KX TECHNOLOGIES LLC & THE SCIENCE BEHIND BETTER WATE & 2007 & 5.00 & -0.14 & [ Installation, repair and maintenance of water treatment equipment ] Water treatment equipment and parts thereof, namely, water softening and conditioning units, water filters, reverse-osmosis system \\
Next Generation Energy LLC & SUNGALOW & 2007 & 5.59 & +0.25 & On-line retail store services featuring housing and storage structures designed and constructed to enhance solar energy efficiency, namely, mobile homes, transportable homes, prefabricated homes, modu \\
Omagine, Inc. & IMAGINATION. INSPIRATION. ILLU & 2008 & 4.15 & -0.04 & Real estate services, namely, listing, leasing and financing of residential and commercial properties; Property management services, namely, management of residential and commercial properties Broadca \\
\bottomrule
\end{tabular}
\clearpage

\subsection*{Class 038: Telecommunications}
\label{app:deep-038}

\begin{tabular}{ll}
Clean filings (1990--2021): & 170,826 \\
Mean $KL_\text{pros}$ / $KL_\text{retr}$ / $\Delta KL$: & 4.32 / 4.37 / -0.050 nats \\
5y survival rate: & 51.5\% \\
Share with SEC-listed owner: & 7.36\% \\
\end{tabular}

\begin{center}
\includegraphics[width=0.7\linewidth]{results/per_industry_deep/class038_scatter.png}
\end{center}

\textbf{Trending tokens (top-decile $\Delta KL$)} --- with intra-class 5y survival rate.

\begin{tabular}{lrrr}
\toprule
token & lift & $n$ filings & intra-class survival \\
\midrule
\texttt{blockchain-based} & 5.53$\times$ & 206 & 44\% \\
\texttt{blockchain} & 4.71$\times$ & 567 & 44\% \\
\texttt{cryptocurrency} & 4.66$\times$ & 496 & 42\% \\
\texttt{decentralized} & 4.32$\times$ & 192 & 40\% \\
\texttt{currencies} & 3.44$\times$ & 488 & 41\% \\
\texttt{token} & 3.30$\times$ & 376 & 44\% \\
\texttt{cryptographic} & 2.86$\times$ & 357 & 41\% \\
\texttt{ledger} & 2.65$\times$ & 339 & 36\% \\
\texttt{wallet} & 2.45$\times$ & 469 & 44\% \\
\texttt{avatars} & 2.41$\times$ & 191 & 64\% \\
\bottomrule
\end{tabular}

\textbf{Ending tokens (bottom-decile $\Delta KL$)} --- intra-class 5y survival rate.

\begin{tabular}{lrrr}
\toprule
token & lift & $n$ filings & intra-class survival \\
\midrule
\texttt{pager} & 3.64$\times$ & 744 & 50\% \\
\texttt{intrastate} & 3.17$\times$ & 287 & 56\% \\
\texttt{long} & 2.96$\times$ & 10,732 & 49\% \\
\texttt{interstate} & 2.85$\times$ & 382 & 54\% \\
\texttt{dot} & 2.78$\times$ & 230 & 36\% \\
\texttt{microfiche} & 2.75$\times$ & 233 & 36\% \\
\texttt{dial-up} & 2.71$\times$ & 1,145 & 47\% \\
\texttt{distance} & 2.71$\times$ & 11,827 & 47\% \\
\texttt{prepaid} & 2.69$\times$ & 3,601 & 48\% \\
\texttt{calling} & 2.67$\times$ & 6,096 & 53\% \\
\bottomrule
\end{tabular}

\textbf{Flux-neutral most-frequent tokens ($|\Delta KL|\leq 0.02$)}.

\begin{tabular}{lrrr}
\toprule
token & $n$ in stratum & $n$ in class & intra-class survival \\
\midrule
\texttt{computer} & 12,603 & 98,267 & 48\% \\
\texttt{transmission} & 11,558 & 97,943 & 47\% \\
\texttt{electronic} & 11,439 & 89,931 & 47\% \\
\texttt{information} & 10,920 & 82,632 & 47\% \\
\texttt{data} & 10,279 & 81,950 & 46\% \\
\texttt{internet} & 10,056 & 74,844 & 46\% \\
\texttt{software} & 9,004 & 57,208 & 44\% \\
\texttt{network} & 8,604 & 68,598 & 49\% \\
\texttt{field} & 8,558 & 56,975 & 42\% \\
\texttt{networks} & 8,336 & 59,701 & 44\% \\
\bottomrule
\end{tabular}

\textbf{Top-10 SEC-listed firms with debut in 1990--2021} (ranked by max observed revenue, debut filing g/s shown).

\begin{tabular}{p{4.0cm}p{2.2cm}rrrp{6.5cm}}
\toprule
Firm & Mark & year & $KL_\text{pros}$ & $\Delta KL$ & First-filing g/s excerpt \\
\midrule
Bubblr Limited & C J CITIZENS JOURNALIST & 2019 & 2.95 & +0.09 & Television and radio broadcasting services; audio and video broadcasting services, namely, broadcasting of news in the fields of current affairs, sports and topics of general interest over the interne \\
HB Global Enterprises, Inc. & SOUL RECIPROCITY NETWORK & 2021 & 5.13 & +0.03 & Broadcasting of programmes via the internet; Streaming of audio, video and audiovisual material on the Internet; Video broadcasting \\
IOLITE TECHNOLOGIES CORPORATION & SOCIAL MOSAIC & 2011 & 2.88 & +0.04 & Providing access to on-line computer databases and on-line searchable databases in the field of classifieds, commercial listings and advertisements of services and products of others, social networkin \\
LYYNKS INC. & LYYNKS & 2012 & 2.81 & -0.05 & Audio and video broadcasting services over the Internet; Broadcasting of video and audio programming over the Internet; Broadcasting services and provision of telecommunication access to video and aud \\
Ludwig Enterprises, Inc. & THE ONE & 2008 & 6.41 & -0.67 & Television and radio broadcasting and rebroadcasting services Equipment for receiving broadcast radio transmissions, namely, radios, tuners, decoders, demodulators, receivers and antennas \\
Magine Holding AB & HEADWEB & 2010 & 3.53 & +0.03 & Printed matter, namely, newsletters about movies, cinematic films, the entertainment and showbiz industry Broadcasting and streaming of cinematographic films and programs through a computer network; v \\
Omagine, Inc. & OMAGINE & 2008 & 4.62 & -0.21 & Broadcasting of radio and television programs \\
R4 TECHNOLOGIES, LLC. & THUMBPASS & 2010 & 4.77 & -0.05 & Computer software for controlling and managing access to vending machines Providing user access to a global computer network to access vending machine accounts Providing an on-line portal that provide \\
RK International, Inc. & RK LIVE-TV & 1997 & 4.84 & -0.24 & providing a telecommunications service and software support for the Internet related network computer boxes for the television marketing and selling Internet related network computer boxes for the tel \\
TECTONA Ltd. & BORED CO. & 2021 & 3.64 & +0.25 & Telecommunications services in the nature of television broadcasting, television transmission, broadcasting, streaming and providing telecommunication access to films Technological services and design \\
\bottomrule
\end{tabular}
\clearpage

\subsection*{Class 039: Transport \& Logistics}
\label{app:deep-039}

\begin{tabular}{ll}
Clean filings (1990--2021): & 120,417 \\
Mean $KL_\text{pros}$ / $KL_\text{retr}$ / $\Delta KL$: & 5.03 / 5.05 / -0.020 nats \\
5y survival rate: & 45.4\% \\
Share with SEC-listed owner: & 6.90\% \\
\end{tabular}

\begin{center}
\includegraphics[width=0.7\linewidth]{results/per_industry_deep/class039_scatter.png}
\end{center}

\textbf{Trending tokens (top-decile $\Delta KL$)} --- with intra-class 5y survival rate.

\begin{tabular}{lrrr}
\toprule
token & lift & $n$ filings & intra-class survival \\
\midrule
\texttt{dumpsters} & 4.54$\times$ & 427 & 35\% \\
\texttt{payloads} & 3.50$\times$ & 143 & 56\% \\
\texttt{roll-off} & 3.43$\times$ & 169 & 34\% \\
\texttt{renewable} & 3.21$\times$ & 1,456 & 37\% \\
\texttt{vicinity} & 3.15$\times$ & 197 & 57\% \\
\texttt{debris} & 3.04$\times$ & 963 & 35\% \\
\texttt{biofuels} & 2.86$\times$ & 241 & 41\% \\
\texttt{recharging} & 2.85$\times$ & 267 & 34\% \\
\texttt{regenerative} & 2.85$\times$ & 123 & 20\% \\
\texttt{on-street} & 2.83$\times$ & 212 & 47\% \\
\bottomrule
\end{tabular}

\textbf{Ending tokens (bottom-decile $\Delta KL$)} --- intra-class 5y survival rate.

\begin{tabular}{lrrr}
\toprule
token & lift & $n$ filings & intra-class survival \\
\midrule
\texttt{business-to-business} & 3.35$\times$ & 164 & 23\% \\
\texttt{accompanied} & 2.53$\times$ & 154 & 58\% \\
\texttt{detailing} & 2.32$\times$ & 250 & 30\% \\
\texttt{tangible} & 2.19$\times$ & 151 & 23\% \\
\texttt{pvc} & 2.05$\times$ & 132 & 29\% \\
\texttt{occasion} & 1.99$\times$ & 241 & 50\% \\
\texttt{dances} & 1.94$\times$ & 263 & 57\% \\
\texttt{microcomputers} & 1.87$\times$ & 123 & 42\% \\
\texttt{reconditioning} & 1.87$\times$ & 284 & 18\% \\
\texttt{emerging} & 1.84$\times$ & 196 & 36\% \\
\bottomrule
\end{tabular}

\textbf{Flux-neutral most-frequent tokens ($|\Delta KL|\leq 0.02$)}.

\begin{tabular}{lrrr}
\toprule
token & $n$ in stratum & $n$ in class & intra-class survival \\
\midrule
\texttt{transportation} & 5,604 & 51,832 & 43\% \\
\texttt{information} & 4,698 & 36,418 & 42\% \\
\texttt{others} & 4,567 & 34,665 & 41\% \\
\texttt{storage} & 4,466 & 34,843 & 41\% \\
\texttt{field} & 4,075 & 29,124 & 40\% \\
\texttt{transport} & 3,666 & 25,720 & 35\% \\
\texttt{delivery} & 3,637 & 29,995 & 43\% \\
\texttt{travel} & 3,563 & 33,822 & 42\% \\
\texttt{management} & 3,403 & 22,987 & 39\% \\
\texttt{computer} & 3,081 & 23,387 & 43\% \\
\bottomrule
\end{tabular}

\textbf{Top-10 SEC-listed firms with debut in 1990--2021} (ranked by max observed revenue, debut filing g/s shown).

\begin{tabular}{p{4.0cm}p{2.2cm}rrrp{6.5cm}}
\toprule
Firm & Mark & year & $KL_\text{pros}$ & $\Delta KL$ & First-filing g/s excerpt \\
\midrule
Atlis Motor Vehicles, Inc. & AMV & 2021 & 4.91 & -0.05 & Subscription-based leasing services of vehicles; vehicle delivery services, namely, delivering vehicles to members of a subscription-based service; providing an online searchable database via the inte \\
Baltia Air Lines, Inc. & BALTIA & 1990 & 6.30 & +0.41 & airline services; namely, transportation of people and property \\
Odenza & LET US TAKE YOU THERE & 2015 & 4.47 & +0.03 & Travel agency services, namely, making reservations and bookings for transportation \\
Omagine, Inc. & BASED ON A TRUE STORY & 2008 & 3.62 & -0.12 & Travel agency services, namely, making reservations and booking for temporary lodging and dining; Hotel and resort hotel services; Resort lodging services; Temporary accommodation services; Restaurant \\
SunVesta Holding AG & SUNVESTA & 2005 & 4.87 & +0.05 & Services for providing food and drink in cafes, cafeterias, restaurants and bars; hotel, motel, boarding house, restaurant, cafe, cafeteria and bar services; providing temporary guest accommodation; r \\
WAL-MART STORES, INC. & CURIO & 2012 & 4.02 & -0.06 & Providing forums, electronic bulletin boards, RSS feeds and list servers for transmission of messages among computer users and subscribers concerning general merchandise and general consumer goods, pr \\
Wal-Mart Stores, Inc &  & 2017 & 5.04 & +0.16 & Shipping and delivery of goods \\
Wal-Mart Stores, Inc. & SAM'S TRAVEL CLUB & 1995 & 4.73 & -0.29 & travel club and information services \\
Amazon.com, Inc. & EAMAZON.COM & 2000 & 4.29 & -0.02 & transport; travel arrangement; travel clubs; providing an interactive database featuring an automatically updating travel planner; delivery of goods; delivery services; rental services; distribution o \\
EXXON MOBIL CORPORATION & DENBURY & 2012 & 6.21 & -0.01 & Transportation by pipeline Extraction of oil, gas, carbon dioxide and helium; Laying, constructing, maintaining and operating pipelines; Well drilling for oil, gas, helium and carbon dioxide Explorati \\
\bottomrule
\end{tabular}
\clearpage

\subsection*{Class 040: Materials Treatment}
\label{app:deep-040}

\begin{tabular}{ll}
Clean filings (1990--2021): & 94,015 \\
Mean $KL_\text{pros}$ / $KL_\text{retr}$ / $\Delta KL$: & 5.37 / 5.37 / -0.004 nats \\
5y survival rate: & 44.6\% \\
Share with SEC-listed owner: & 5.75\% \\
\end{tabular}

\begin{center}
\includegraphics[width=0.7\linewidth]{results/per_industry_deep/class040_scatter.png}
\end{center}

\textbf{Trending tokens (top-decile $\Delta KL$)} --- with intra-class 5y survival rate.

\begin{tabular}{lrrr}
\toprule
token & lift & $n$ filings & intra-class survival \\
\midrule
\texttt{vectors} & 5.64$\times$ & 101 & 53\% \\
\texttt{biomanufacturing} & 5.00$\times$ & 168 & 39\% \\
\texttt{cbd} & 4.62$\times$ & 147 & 76\% \\
\texttt{therapies} & 4.62$\times$ & 171 & 40\% \\
\texttt{delta-} & 4.62$\times$ & 210 & 66\% \\
\texttt{thc} & 4.51$\times$ & 204 & 67\% \\
\texttt{tetrahydrocannabinol} & 4.49$\times$ & 185 & 68\% \\
\texttt{autoimmune} & 4.38$\times$ & 130 & 32\% \\
\texttt{percent} & 4.37$\times$ & 192 & 65\% \\
\texttt{gene} & 4.27$\times$ & 295 & 44\% \\
\bottomrule
\end{tabular}

\textbf{Ending tokens (bottom-decile $\Delta KL$)} --- intra-class 5y survival rate.

\begin{tabular}{lrrr}
\toprule
token & lift & $n$ filings & intra-class survival \\
\midrule
\texttt{snapshots} & 5.67$\times$ & 104 & 50\% \\
\texttt{taproom} & 4.39$\times$ & 123 & 64\% \\
\texttt{portraits} & 3.56$\times$ & 191 & 42\% \\
\texttt{brewing} & 3.49$\times$ & 1,037 & 54\% \\
\texttt{brewed} & 3.36$\times$ & 143 & 62\% \\
\texttt{digitizing} & 3.31$\times$ & 956 & 45\% \\
\texttt{retouching} & 3.18$\times$ & 1,036 & 44\% \\
\texttt{lager} & 3.07$\times$ & 114 & 67\% \\
\texttt{tee-shirt} & 3.05$\times$ & 95 & 82\% \\
\texttt{beer} & 2.97$\times$ & 1,318 & 52\% \\
\bottomrule
\end{tabular}

\textbf{Flux-neutral most-frequent tokens ($|\Delta KL|\leq 0.02$)}.

\begin{tabular}{lrrr}
\toprule
token & $n$ in stratum & $n$ in class & intra-class survival \\
\midrule
\texttt{others} & 4,275 & 37,163 & 41\% \\
\texttt{custom} & 3,709 & 32,029 & 41\% \\
\texttt{field} & 3,591 & 25,985 & 38\% \\
\texttt{use} & 2,907 & 21,005 & 38\% \\
\texttt{manufacture} & 2,835 & 24,294 & 40\% \\
\texttt{treatment} & 2,699 & 19,339 & 35\% \\
\texttt{design} & 2,559 & 17,922 & 35\% \\
\texttt{products} & 2,532 & 18,350 & 38\% \\
\texttt{materials} & 2,426 & 16,692 & 33\% \\
\texttt{manufacturing} & 2,289 & 17,273 & 39\% \\
\bottomrule
\end{tabular}

\textbf{Top-10 SEC-listed firms with debut in 1990--2021} (ranked by max observed revenue, debut filing g/s shown).

\begin{tabular}{p{4.0cm}p{2.2cm}rrrp{6.5cm}}
\toprule
Firm & Mark & year & $KL_\text{pros}$ & $\Delta KL$ & First-filing g/s excerpt \\
\midrule
ECO Integrated Technologies, Inc. & ECO INTEGRATED TECHNOLOGIES & 2015 & 5.63 & -0.20 & Environmental remediation services, namely, treatment of solid waste and water \\
Fox Integrated Technologies, Inc. & FOXCAM & 1999 & 6.03 & -0.16 & Recycling services, namely, monitoring to allow verification of recycling of electronic components as an integral part of the recycling services \\
Miza Pharmaceuticals, Inc. & OPTOPICS & 1999 & 4.70 & +0.03 & eye care products, namely, eye drops, eye washes, artificial tears, contact lens solutions and generic pharmaceutical drug solutions for use in the treatment of eye related diseases and disorders; and \\
New Dimension LLC & N NEW DIMENSION & 2007 & 5.90 & -0.07 & digital printing services \\
Next Generation Energy LLC & NEXT GENERATION ENERGY & 2006 & 5.42 & +0.27 & Develop new technology for others in the field of renewable energy related products and resources Contract manufacturing in the field of renewable energy related products \\
RS Technologies Inc. & POWERON & 2021 & 5.98 & +0.14 & Non-metal utility poles; non-metal utility pole cross arms; non-metal covers for utility poles and utility pole cross arms Technical support services, namely, consulting relating to the use of polyure \\
WAL-MART STORES, INC. & WAL MART DIGITAL PHOTO CENTER & 2004 & 6.60 & -0.08 & online photo processing services \\
WALMART INC. & ART.COM & 1999 & 4.73 & -0.52 & custom framing services offered via computer networks and global communications networks; computerized on-line custom art framing services, consulting and advisory services on choosing art framing mat \\
Wal-Mart Stores, Inc. & DEVELOP YOUR BUSINESS & 2003 & 7.34 & -0.55 & film developing services \\
EXXON MOBIL CORPORATION & DENBURY & 2012 & 5.64 & -0.07 & Transportation by pipeline Extraction of oil, gas, carbon dioxide and helium; Laying, constructing, maintaining and operating pipelines; Well drilling for oil, gas, helium and carbon dioxide Explorati \\
\bottomrule
\end{tabular}
\clearpage

\subsection*{Class 041: Education \& Entertainment}
\label{app:deep-041}

\begin{tabular}{ll}
Clean filings (1990--2021): & 920,449 \\
Mean $KL_\text{pros}$ / $KL_\text{retr}$ / $\Delta KL$: & 5.18 / 5.19 / -0.006 nats \\
5y survival rate: & 47.0\% \\
Share with SEC-listed owner: & 3.06\% \\
\end{tabular}

\begin{center}
\includegraphics[width=0.7\linewidth]{results/per_industry_deep/class041_scatter.png}
\end{center}

\textbf{Trending tokens (top-decile $\Delta KL$)} --- with intra-class 5y survival rate.

\begin{tabular}{lrrr}
\toprule
token & lift & $n$ filings & intra-class survival \\
\midrule
\texttt{blockchain} & 6.41$\times$ & 1,106 & 53\% \\
\texttt{cryptocurrency} & 6.08$\times$ & 957 & 50\% \\
\texttt{worlds} & 3.63$\times$ & 1,126 & 52\% \\
\texttt{wager} & 3.33$\times$ & 1,327 & 42\% \\
\texttt{accept} & 3.26$\times$ & 1,382 & 41\% \\
\texttt{tokens} & 3.21$\times$ & 2,330 & 43\% \\
\texttt{in-company} & 3.17$\times$ & 2,402 & 42\% \\
\texttt{currency} & 3.05$\times$ & 2,174 & 45\% \\
\texttt{keynote} & 2.88$\times$ & 3,167 & 39\% \\
\texttt{vlogs} & 2.83$\times$ & 2,230 & 32\% \\
\bottomrule
\end{tabular}

\textbf{Ending tokens (bottom-decile $\Delta KL$)} --- intra-class 5y survival rate.

\begin{tabular}{lrrr}
\toprule
token & lift & $n$ filings & intra-class survival \\
\midrule
\texttt{instrumental} & 3.76$\times$ & 2,579 & 51\% \\
\texttt{pageants} & 3.63$\times$ & 2,304 & 57\% \\
\texttt{live-action} & 2.75$\times$ & 2,579 & 67\% \\
\texttt{batting} & 2.66$\times$ & 999 & 41\% \\
\texttt{skating} & 2.63$\times$ & 2,840 & 47\% \\
\texttt{band} & 2.59$\times$ & 21,828 & 49\% \\
\texttt{on-going} & 2.53$\times$ & 20,697 & 61\% \\
\texttt{highlight} & 2.53$\times$ & 924 & 54\% \\
\texttt{nightshirts} & 2.53$\times$ & 1,021 & 67\% \\
\texttt{hockey} & 2.52$\times$ & 5,857 & 49\% \\
\bottomrule
\end{tabular}

\textbf{Flux-neutral most-frequent tokens ($|\Delta KL|\leq 0.02$)}.

\begin{tabular}{lrrr}
\toprule
token & $n$ in stratum & $n$ in class & intra-class survival \\
\midrule
\texttt{field} & 60,199 & 419,101 & 42\% \\
\texttt{entertainment} & 49,261 & 407,659 & 50\% \\
\texttt{educational} & 44,270 & 309,765 & 42\% \\
\texttt{featuring} & 36,816 & 255,080 & 45\% \\
\texttt{information} & 36,171 & 230,417 & 45\% \\
\texttt{nature} & 36,114 & 267,431 & 45\% \\
\texttt{conducting} & 35,086 & 242,144 & 41\% \\
\texttt{training} & 30,092 & 202,842 & 40\% \\
\texttt{seminars} & 28,345 & 188,062 & 40\% \\
\texttt{education} & 26,981 & 176,329 & 42\% \\
\bottomrule
\end{tabular}

\textbf{Top-10 SEC-listed firms with debut in 1990--2021} (ranked by max observed revenue, debut filing g/s shown).

\begin{tabular}{p{4.0cm}p{2.2cm}rrrp{6.5cm}}
\toprule
Firm & Mark & year & $KL_\text{pros}$ & $\Delta KL$ & First-filing g/s excerpt \\
\midrule
ATK Entertainment, Inc. & AALAM TUCK KHALED & 2020 & 3.05 & -0.08 & Clothing, namely, shirts, t-shirts, pants, dresses, skirts, jeans, slacks, belts, hats, caps being headwear, gloves, ear muffs, headwear, neckwear, sweaters, fleece pullovers, jogging suits, warm-up s \\
ATLAS TECHNOLOGY GROUP, INC. & IGE & 2006 & 5.41 & +0.01 & Entertainment services, namely, providing on-line information via the Internet about video games; providing on-line news and information via the Internet on video games; providing on-line databases fo \\
Artix Entertainment, LLC & BATTLEON & 2006 & 4.70 & +0.09 & Entertainment services, namely, providing an on-line computer game \\
BAT Entertainment LLC & MISS INTERNATIONAL UNIVERSE & 2018 & 7.44 & -0.23 & Entertainment in the nature of beauty pageants \\
BETTERLIFE PHARMA INC. & BETTER+ & 2019 & 4.81 & +0.06 & Printed publications, namely, newsletters, brochures, reports and guides in the field of cannabis bath additives, namely, bath beads, bath bombs, bath crystals and bubble bath; bath herbs; bath oils;  \\
BLH ACQUISITION CO., LLC & ROCK THE TOT & 2012 & 5.66 & +0.08 & Entertainment services, namely, arranging and conducting of eating competitions \\
BetterLife Pharma Inc. & LIVE BETTERX & 2020 & 4.46 & -0.10 & Downloadable electronic publications, namely, newsletters, brochures, reports and guides in the field of cannabis Retail store and wholesale store services featuring industrial hemp food and body care \\
Bubblr Limited & C J CITIZENS JOURNALIST & 2019 & 3.48 & +0.05 & Television and radio broadcasting services; audio and video broadcasting services, namely, broadcasting of news in the fields of current affairs, sports and topics of general interest over the interne \\
China Foundation, The & CHINA FOUNDATION & 2008 & 5.31 & -0.05 & Arranging and conducting educational conferences; organizing conferences in the field of public health and medical science Charitable services, namely, organizing and conducting volunteer programs and \\
Clickstream Corporation & THINK QUIK, PLAY QUIK, WINQUIK & 2020 & 3.94 & -0.02 & Entertainment services, namely, providing temporary use of non-downloadable interactive trivia games; conducting live entertainment in the nature of games featuring trivia questions and answers via th \\
\bottomrule
\end{tabular}
\clearpage

\subsection*{Class 042: Scientific \& Tech Services}
\label{app:deep-042}

\begin{tabular}{ll}
Clean filings (1990--2021): & 743,414 \\
Mean $KL_\text{pros}$ / $KL_\text{retr}$ / $\Delta KL$: & 5.17 / 5.20 / -0.022 nats \\
5y survival rate: & 46.1\% \\
Share with SEC-listed owner: & 5.61\% \\
\end{tabular}

\begin{center}
\includegraphics[width=0.7\linewidth]{results/per_industry_deep/class042_scatter.png}
\end{center}

\textbf{Trending tokens (top-decile $\Delta KL$)} --- with intra-class 5y survival rate.

\begin{tabular}{lrrr}
\toprule
token & lift & $n$ filings & intra-class survival \\
\midrule
\texttt{non-fungible} & 8.30$\times$ & 855 & 65\% \\
\texttt{blockchain-based} & 5.53$\times$ & 1,341 & 57\% \\
\texttt{cryptocurrency} & 4.88$\times$ & 3,183 & 54\% \\
\texttt{blockchain} & 4.42$\times$ & 3,877 & 53\% \\
\texttt{decentralized} & 4.38$\times$ & 1,252 & 55\% \\
\texttt{token} & 4.28$\times$ & 1,178 & 51\% \\
\texttt{ledger} & 3.80$\times$ & 1,398 & 52\% \\
\texttt{cryptographic} & 3.78$\times$ & 1,424 & 53\% \\
\texttt{wallet} & 3.75$\times$ & 1,618 & 51\% \\
\texttt{tokens} & 3.72$\times$ & 3,652 & 48\% \\
\bottomrule
\end{tabular}

\textbf{Ending tokens (bottom-decile $\Delta KL$)} --- intra-class 5y survival rate.

\begin{tabular}{lrrr}
\toprule
token & lift & $n$ filings & intra-class survival \\
\midrule
\texttt{carry-out} & 7.92$\times$ & 1,135 & 52\% \\
\texttt{sandwiches} & 6.41$\times$ & 1,234 & 52\% \\
\texttt{restaurant} & 6.33$\times$ & 18,456 & 55\% \\
\texttt{pizza} & 6.31$\times$ & 932 & 50\% \\
\texttt{motel} & 6.10$\times$ & 1,201 & 49\% \\
\texttt{bakery} & 5.93$\times$ & 1,397 & 46\% \\
\texttt{resort} & 5.67$\times$ & 2,116 & 56\% \\
\texttt{nightclub} & 5.64$\times$ & 948 & 66\% \\
\texttt{cafe} & 5.38$\times$ & 1,076 & 59\% \\
\texttt{lounge} & 5.27$\times$ & 1,595 & 54\% \\
\bottomrule
\end{tabular}

\textbf{Flux-neutral most-frequent tokens ($|\Delta KL|\leq 0.02$)}.

\begin{tabular}{lrrr}
\toprule
token & $n$ in stratum & $n$ in class & intra-class survival \\
\midrule
\texttt{software} & 46,113 & 355,301 & 43\% \\
\texttt{computer} & 41,094 & 345,425 & 46\% \\
\texttt{field} & 35,666 & 273,520 & 40\% \\
\texttt{information} & 31,871 & 244,099 & 45\% \\
\texttt{use} & 30,111 & 228,642 & 41\% \\
\texttt{design} & 27,704 & 202,311 & 38\% \\
\texttt{development} & 27,340 & 187,141 & 39\% \\
\texttt{featuring} & 26,369 & 209,628 & 41\% \\
\texttt{management} & 26,332 & 181,534 & 40\% \\
\texttt{others} & 26,085 & 208,844 & 42\% \\
\bottomrule
\end{tabular}

\textbf{Top-10 SEC-listed firms with debut in 1990--2021} (ranked by max observed revenue, debut filing g/s shown).

\begin{tabular}{p{4.0cm}p{2.2cm}rrrp{6.5cm}}
\toprule
Firm & Mark & year & $KL_\text{pros}$ & $\Delta KL$ & First-filing g/s excerpt \\
\midrule
AmeriTrust Group, Inc. & SIMPLE.FAST.DIRECT. & 2020 & 5.58 & -0.09 & Insurance services, namely, administration, brokerage, and underwriting for workers' compensation; insurance claims administration and claims processing services; insurance agency services, namely, in \\
Artelo Biosciences, Inc. & ARTELO BIOSCIENCES & 2019 & 7.77 & +0.08 & Research and development of pharmaceuticals; namely, therapeutic agents that modulate the endocannabinoid system \\
BLH ACQUISITION CO., LLC & BAR LOUIE & 1993 & 5.82 & -0.62 & tavern and restaurant services \\
Blackstar Enterprise Group, Inc. & INTERNET DIGITAL OFFERING & 2021 & 5.37 & +0.08 & Downloadable financial services software for managing and hosting indications of interest for potential initial or secondary future offerings on an immutable block chain Financial services, namely, pr \\
Bubblr Limited & C J CITIZENS JOURNALIST & 2019 & 3.31 & -0.02 & Television and radio broadcasting services; audio and video broadcasting services, namely, broadcasting of news in the fields of current affairs, sports and topics of general interest over the interne \\
CHATCHING INC. & USE IT, OWN IT & 2011 & 3.62 & +0.12 & Internet based social networking and social introduction services; On-line social networking services; Online social networking services provided through a community website; Social introduction agenc \\
CellCyte Genetics, Inc. & CELLCYTE & 2005 & 6.63 & -0.24 & Research and development of therapeutic and diagnostic products for others related to a wide variety of diseases, disorders and conditions \\
ChatAnd, Inc. & CHAT\& & 2011 & 4.95 & +0.17 & Providing temporary use of online, non-downloadable computer software for web browsing and information sharing in the field of sales and conversion; providing temporary use of online, non-downloadable \\
Ching Technologies, Inc. & CHING WALLET & 2018 & 4.01 & +0.46 & Software for use in electronically trading, storing, sending, receiving, accepting and transmitting digital currency, and managing digital currency payment and exchange transactions; services in the f \\
CohBar, Inc. & MITO+ & 2021 & 6.10 & -0.05 & Research and development of novel pharmaceutical products, namely, research and development of novel peptides for use as pharmaceuticals \\
\bottomrule
\end{tabular}
\clearpage

\subsection*{Class 043: Restaurants \& Hotels}
\label{app:deep-043}

\begin{tabular}{ll}
Clean filings (1990--2021): & 136,699 \\
Mean $KL_\text{pros}$ / $KL_\text{retr}$ / $\Delta KL$: & 4.87 / 4.85 / +0.014 nats \\
5y survival rate: & 48.4\% \\
Share with SEC-listed owner: & 2.80\% \\
\end{tabular}

\begin{center}
\includegraphics[width=0.7\linewidth]{results/per_industry_deep/class043_scatter.png}
\end{center}

\textbf{Trending tokens (top-decile $\Delta KL$)} --- with intra-class 5y survival rate.

\begin{tabular}{lrrr}
\toprule
token & lift & $n$ filings & intra-class survival \\
\midrule
\texttt{truck} & 3.11$\times$ & 4,705 & 44\% \\
\texttt{vegan} & 2.85$\times$ & 652 & 45\% \\
\texttt{specified} & 2.37$\times$ & 156 & 76\% \\
\texttt{contest} & 2.15$\times$ & 149 & 81\% \\
\texttt{ccrcs} & 2.14$\times$ & 299 & 33\% \\
\texttt{proportioned} & 2.06$\times$ & 608 & 52\% \\
\texttt{plant-based} & 2.02$\times$ & 213 & 49\% \\
\texttt{properly} & 2.02$\times$ & 615 & 53\% \\
\texttt{brewpub} & 1.99$\times$ & 649 & 45\% \\
\texttt{facilitation} & 1.98$\times$ & 162 & 72\% \\
\bottomrule
\end{tabular}

\textbf{Ending tokens (bottom-decile $\Delta KL$)} --- intra-class 5y survival rate.

\begin{tabular}{lrrr}
\toprule
token & lift & $n$ filings & intra-class survival \\
\midrule
\texttt{portions} & 4.35$\times$ & 170 & 78\% \\
\texttt{mixed-use} & 3.88$\times$ & 209 & 78\% \\
\texttt{puppet} & 3.50$\times$ & 177 & 59\% \\
\texttt{bistros} & 3.16$\times$ & 256 & 54\% \\
\texttt{contest} & 3.15$\times$ & 149 & 81\% \\
\texttt{drugs} & 3.01$\times$ & 219 & 69\% \\
\texttt{specified} & 3.01$\times$ & 156 & 76\% \\
\texttt{preferred} & 2.93$\times$ & 954 & 66\% \\
\texttt{facilitation} & 2.78$\times$ & 162 & 72\% \\
\texttt{carry-out} & 2.75$\times$ & 3,427 & 59\% \\
\bottomrule
\end{tabular}

\textbf{Flux-neutral most-frequent tokens ($|\Delta KL|\leq 0.02$)}.

\begin{tabular}{lrrr}
\toprule
token & $n$ in stratum & $n$ in class & intra-class survival \\
\midrule
\texttt{restaurant} & 9,526 & 76,814 & 51\% \\
\texttt{food} & 7,062 & 55,402 & 46\% \\
\texttt{bar} & 5,865 & 41,470 & 51\% \\
\texttt{featuring} & 4,987 & 36,194 & 46\% \\
\texttt{catering} & 4,843 & 32,880 & 45\% \\
\texttt{beverages} & 3,598 & 23,790 & 45\% \\
\texttt{information} & 3,558 & 21,479 & 41\% \\
\texttt{temporary} & 3,541 & 23,817 & 40\% \\
\texttt{restaurants} & 3,511 & 25,081 & 49\% \\
\texttt{hotel} & 3,344 & 22,958 & 50\% \\
\bottomrule
\end{tabular}

\textbf{Top-10 SEC-listed firms with debut in 1990--2021} (ranked by max observed revenue, debut filing g/s shown).

\begin{tabular}{p{4.0cm}p{2.2cm}rrrp{6.5cm}}
\toprule
Firm & Mark & year & $KL_\text{pros}$ & $\Delta KL$ & First-filing g/s excerpt \\
\midrule
BLH ACQUISITION CO., LLC & DRINK POURN & 2017 & 6.28 & -0.29 & Headwear, namely, caps and hats; Wearable garments and clothing, namely, shirts Bar and restaurant services \\
Infinity Real Estate LLC &  & 2020 & 4.50 & -0.05 & Hotel services Real estate development Assessment and management of real estate; Commercial and residential real estate agency services; Financial planning and investment advisory services; Investment \\
Omagine, Inc. & OMAGINE & 2008 & 3.03 & -0.08 & Travel agency services, namely, making reservations and booking for temporary lodging and dining; Hotel and resort hotel services; Resort lodging services; Temporary accommodation services; Restaurant \\
SunVesta Holding AG & SUNVESTA & 2005 & 4.11 & +0.19 & Services for providing food and drink in cafes, cafeterias, restaurants and bars; hotel, motel, boarding house, restaurant, cafe, cafeteria and bar services; providing temporary guest accommodation; r \\
TGI SOLAR POWER GROUP INC. & ADVENT & 2021 & 4.36 & -0.13 & Hotel and restaurant services; Providing hotel accommodation; Resort hotel services; Restaurant and bar services; Restaurant and café services; Restaurant and hotel services \\
Wal-Mart Stores, Inc & YUMPRINT & 2017 & 4.74 & +0.01 & Providing a website featuring information about recipes, ingredients and cooking information; Providing an online computer database to consumers featuring information about recipes, ingredients and co \\
Wal-Mart Stores, Inc. & KID CONNECTION & 2003 & 6.62 & -0.49 & restaurant services retail candy stores services \\
Exxon Mobil Corporation & BENGAL TRADERS GOURMET COFFEES & 2004 & 6.04 & +0.01 & Coffee Preparation and serving of coffee for consumption on or off the premises \\
Apple Inc. &  & 2012 & 3.99 & +0.12 & providing consumer product information for the purpose of selecting general consumer merchandise to meet the consumer's specifications; advisory and consultancy services relating to all the aforesaid  \\
UNITEDHEALTH GROUP INCORPORATED & CARE24 & 2002 & 5.21 & +0.37 & Referral services in the field of financial information, employee assistance and work life programs, legal information, dependent care information, health care and mental health information Providing  \\
\bottomrule
\end{tabular}
\clearpage

\subsection*{Class 044: Medical \& Personal Care}
\label{app:deep-044}

\begin{tabular}{ll}
Clean filings (1990--2021): & 188,794 \\
Mean $KL_\text{pros}$ / $KL_\text{retr}$ / $\Delta KL$: & 4.96 / 4.95 / +0.005 nats \\
5y survival rate: & 45.8\% \\
Share with SEC-listed owner: & 5.56\% \\
\end{tabular}

\begin{center}
\includegraphics[width=0.7\linewidth]{results/per_industry_deep/class044_scatter.png}
\end{center}

\textbf{Trending tokens (top-decile $\Delta KL$)} --- with intra-class 5y survival rate.

\begin{tabular}{lrrr}
\toprule
token & lift & $n$ filings & intra-class survival \\
\midrule
\texttt{candidates} & 3.21$\times$ & 218 & 40\% \\
\texttt{inside} & 3.20$\times$ & 222 & 40\% \\
\texttt{hydration} & 3.19$\times$ & 700 & 42\% \\
\texttt{adding} & 3.10$\times$ & 203 & 36\% \\
\texttt{cell-based} & 3.02$\times$ & 212 & 53\% \\
\texttt{analyzed} & 2.94$\times$ & 204 & 33\% \\
\texttt{semiconductors} & 2.86$\times$ & 203 & 25\% \\
\texttt{intravenous} & 2.69$\times$ & 1,143 & 43\% \\
\texttt{abnormalities} & 2.61$\times$ & 280 & 60\% \\
\texttt{delta-} & 2.59$\times$ & 405 & 69\% \\
\bottomrule
\end{tabular}

\textbf{Ending tokens (bottom-decile $\Delta KL$)} --- intra-class 5y survival rate.

\begin{tabular}{lrrr}
\toprule
token & lift & $n$ filings & intra-class survival \\
\midrule
\texttt{winery} & 5.86$\times$ & 341 & 64\% \\
\texttt{vineyard} & 5.76$\times$ & 481 & 59\% \\
\texttt{grapes} & 4.99$\times$ & 597 & 58\% \\
\texttt{condominiums} & 4.90$\times$ & 208 & 74\% \\
\texttt{stud} & 4.36$\times$ & 551 & 46\% \\
\texttt{lounges} & 4.35$\times$ & 209 & 66\% \\
\texttt{regimes} & 4.15$\times$ & 207 & 63\% \\
\texttt{percent} & 3.96$\times$ & 338 & 67\% \\
\texttt{shrub} & 3.93$\times$ & 252 & 40\% \\
\texttt{vacation} & 3.84$\times$ & 268 & 58\% \\
\bottomrule
\end{tabular}

\textbf{Flux-neutral most-frequent tokens ($|\Delta KL|\leq 0.02$)}.

\begin{tabular}{lrrr}
\toprule
token & $n$ in stratum & $n$ in class & intra-class survival \\
\midrule
\texttt{health} & 10,877 & 79,196 & 45\% \\
\texttt{medical} & 10,847 & 79,310 & 44\% \\
\texttt{information} & 10,386 & 70,993 & 44\% \\
\texttt{care} & 10,206 & 71,534 & 45\% \\
\texttt{field} & 9,573 & 64,713 & 42\% \\
\texttt{treatment} & 6,025 & 40,540 & 42\% \\
\texttt{featuring} & 5,423 & 37,942 & 44\% \\
\texttt{wellness} & 4,779 & 33,138 & 45\% \\
\texttt{use} & 4,369 & 28,832 & 43\% \\
\texttt{purposes} & 4,284 & 28,697 & 40\% \\
\bottomrule
\end{tabular}

\textbf{Top-10 SEC-listed firms with debut in 1990--2021} (ranked by max observed revenue, debut filing g/s shown).

\begin{tabular}{p{4.0cm}p{2.2cm}rrrp{6.5cm}}
\toprule
Firm & Mark & year & $KL_\text{pros}$ & $\Delta KL$ & First-filing g/s excerpt \\
\midrule
BETTERLIFE PHARMA INC. & LOVE BETTERX & 2019 & 3.51 & +0.02 & Printed publications, namely, newsletters, brochures, reports and guides in the field of cannabis bath additives, namely, bath beads, bath bombs, bath crystals and bubble bath; bath herbs; bath oils;  \\
Better For You Wellness, Inc. & BETTER FOR YOU WELLNESS & 2021 & 4.88 & -0.19 & Protein bars, namely, protein supplements formed and packaged as bars Advisory and consultation services in the fields of agronomy, agriculture, and horticulture; agronomic consulting services; agrono \\
BetterLife Pharma Inc. & LIVE BETTERX & 2020 & 3.34 & -0.15 & Downloadable electronic publications, namely, newsletters, brochures, reports and guides in the field of cannabis Retail store and wholesale store services featuring industrial hemp food and body care \\
Bone Biologics, Inc. & BONE BIOLOGICS & 2006 & 7.32 & -0.01 & provider of orthobiologics, bone growth factors and scaffolds, bone graft substitutes,cartilage growth factors and scaffolds,spinal disc growth factors and scaffolds, articular cartilage growth factor \\
Brookfield Resources, LLC & BROOKFIELD RESOURCES, LLC BR & 2004 & 5.17 & -0.14 & Consulting services in medical communications for the pharmaceutical and related industries \\
Event Cardio Group & EVENT CARDIO & 2016 & 5.40 & +0.09 & Medical device, namely, a cardiac and cardiac event monitor having a sensor and capability to transmit cardiac and other physiological data of patients remotely Medical Services, namely, remote monito \\
IPR Pharmaceuticals, Inc. & RISING ABOVE THE STATIN QUO & 2002 & 4.77 & +0.69 & Educational services, namely, conducting classes, conferences, workshops, seminars and symposia in the field of treatment of cardiovascular diseases, and cardiovascular health; and providing the above \\
Omagine, Inc. & IMAGINATION. INSPIRATION. ILLU & 2008 & 3.55 & -0.13 & Real estate services, namely, listing, leasing and financing of residential and commercial properties; Property management services, namely, management of residential and commercial properties Broadca \\
R2 Technologies, Inc. & CRYOAESTHETICS & 2020 & 4.58 & -0.13 & Medical and cosmetic body care services using cryotherapy \\
Sharps Technology, Inc. & SHARPS PROVENSA & 2018 & 5.12 & -0.01 & Online retail store and online retail services through direct solicitation by distributors directed to end-users featuring syringes and needles for medical use and for injection, disposable syringes a \\
\bottomrule
\end{tabular}
\clearpage

\subsection*{Class 045: Legal \& Personal Services}
\label{app:deep-045}

\begin{tabular}{ll}
Clean filings (1990--2021): & 126,420 \\
Mean $KL_\text{pros}$ / $KL_\text{retr}$ / $\Delta KL$: & 4.96 / 4.99 / -0.022 nats \\
5y survival rate: & 45.2\% \\
Share with SEC-listed owner: & 3.47\% \\
\end{tabular}

\begin{center}
\includegraphics[width=0.7\linewidth]{results/per_industry_deep/class045_scatter.png}
\end{center}

\textbf{Trending tokens (top-decile $\Delta KL$)} --- with intra-class 5y survival rate.

\begin{tabular}{lrrr}
\toprule
token & lift & $n$ filings & intra-class survival \\
\midrule
\texttt{authenticated} & 4.60$\times$ & 126 & 55\% \\
\texttt{mediumship} & 4.58$\times$ & 144 & 27\% \\
\texttt{blockchain} & 4.07$\times$ & 381 & 52\% \\
\texttt{worlds} & 3.79$\times$ & 153 & 44\% \\
\texttt{cryptocurrency} & 3.70$\times$ & 305 & 49\% \\
\texttt{token} & 3.68$\times$ & 174 & 45\% \\
\texttt{user-created} & 3.65$\times$ & 178 & 56\% \\
\texttt{currencies} & 3.38$\times$ & 207 & 41\% \\
\texttt{tarot} & 3.36$\times$ & 393 & 34\% \\
\texttt{augmented} & 3.26$\times$ & 368 & 45\% \\
\bottomrule
\end{tabular}

\textbf{Ending tokens (bottom-decile $\Delta KL$)} --- intra-class 5y survival rate.

\begin{tabular}{lrrr}
\toprule
token & lift & $n$ filings & intra-class survival \\
\midrule
\texttt{localities} & 6.04$\times$ & 139 & 84\% \\
\texttt{retired} & 5.42$\times$ & 155 & 83\% \\
\texttt{practicing} & 5.03$\times$ & 163 & 71\% \\
\texttt{slogans} & 4.70$\times$ & 1,014 & 89\% \\
\texttt{coordinated} & 4.21$\times$ & 197 & 73\% \\
\texttt{cartoon} & 4.06$\times$ & 1,323 & 82\% \\
\texttt{occasion} & 4.03$\times$ & 513 & 50\% \\
\texttt{valuable} & 3.87$\times$ & 289 & 69\% \\
\texttt{evangelical} & 3.86$\times$ & 339 & 55\% \\
\texttt{acquaintances} & 3.85$\times$ & 187 & 50\% \\
\bottomrule
\end{tabular}

\textbf{Flux-neutral most-frequent tokens ($|\Delta KL|\leq 0.02$)}.

\begin{tabular}{lrrr}
\toprule
token & $n$ in stratum & $n$ in class & intra-class survival \\
\midrule
\texttt{field} & 8,452 & 55,961 & 40\% \\
\texttt{information} & 8,089 & 54,797 & 42\% \\
\texttt{computer} & 4,972 & 33,473 & 44\% \\
\texttt{featuring} & 4,962 & 32,041 & 43\% \\
\texttt{personal} & 4,778 & 35,288 & 43\% \\
\texttt{management} & 4,663 & 26,810 & 38\% \\
\texttt{online} & 4,600 & 30,225 & 45\% \\
\texttt{others} & 4,562 & 29,050 & 41\% \\
\texttt{social} & 4,526 & 34,800 & 48\% \\
\texttt{nature} & 4,259 & 25,152 & 38\% \\
\bottomrule
\end{tabular}

\textbf{Top-10 SEC-listed firms with debut in 1990--2021} (ranked by max observed revenue, debut filing g/s shown).

\begin{tabular}{p{4.0cm}p{2.2cm}rrrp{6.5cm}}
\toprule
Firm & Mark & year & $KL_\text{pros}$ & $\Delta KL$ & First-filing g/s excerpt \\
\midrule
CHATCHING INC. & USE IT, OWN IT & 2011 & 3.52 & +0.16 & Internet based social networking and social introduction services; On-line social networking services; Online social networking services provided through a community website; Social introduction agenc \\
Global Smart Energy, LLC & SMART GRID SPRINGBOARD SERVICE & 2010 & 5.03 & -0.13 & Consulting services to utilities and electric power clients, namely, conducting energy audits of commercial and institutional facilities for the purpose of improving energy efficiency in the field of  \\
HB Global Enterprises, Inc. & SOUL RECIPROCITY & 2020 & 4.03 & +0.10 & Retail store services featuring subscription boxes containing household items, personal care items, clothing, jewelry and gifts; Retail shops featuring gifts, clothing and jewelry Educational and ente \\
IOLITE TECHNOLOGIES CORPORATION & SOCIAL MOSAIC & 2011 & 2.84 & +0.07 & Providing access to on-line computer databases and on-line searchable databases in the field of classifieds, commercial listings and advertisements of services and products of others, social networkin \\
Mirata Pharmaceuticals, LLC & MIRATA PHARMACEUTICALS & 2020 & 5.76 & -0.09 & licensing of intellectual property pharmaceutical products, namely, gene therapies \\
SLA Marketing Group LLC & CB & 2014 & 5.42 & -0.01 & Clothing, namely, scarves, dresses, tops, t-shirts, shirts, jackets, pants, hats Personal lifestyle consulting services; personal fashion consulting services Cosmetic brushes Various colored paper dia \\
WAL-MART STORES, INC. & JIG & 2011 & 3.51 & +0.06 & computer software, computer software development tools, computer software application interfaces and application program interfaces (APIs), all for operating search engines and managing online service \\
Wal-Mart Stores, Inc. & WALMART TOGO & 2011 & 6.27 & +0.18 & Personal shopping for others Retail and on-line grocery store services featuring home delivery service; Retail department store services \\
Exxon Mobil Corporation & EXXONMOBIL TRAVEL CLUB & 2009 & 6.03 & -0.05 & Reimbursement of emergency travel expenses, of ambulance service expenses, and of unexpected transport expenses Automobile club services, namely, providing members with protection against auto theft b \\
APPLE INC. &  & 2004 & 5.14 & -0.14 & (Based on Use in Commerce) providing a selection of on-line electronic greeting cards; [ (Based on 44(e)) computer dating services; providing a selection on-line electronic greeting cards ] \\
\bottomrule
\end{tabular}
\clearpage


\end{document}
